\documentclass[
    a4paper,
    11pt,
    bibliography=totoc,
    listof=totoc,      
    headinclude=true,
    cleardoublepage=empty,
]{scrbook}

%------------------------------------------------------------------------------
% PDF/A-2u compliance (fonts + metadata)
%------------------------------------------------------------------------------

\usepackage[T1]{fontenc}
\usepackage[utf8]{inputenc} % required for pdfLaTeX
\usepackage{lmodern}
\usepackage[final]{microtype} % improves spacing, reduces pagination shifts

\usepackage[a-2u]{pdfx} % loads hyperref internally in a PDF/A-safe way

%------------------------------------------------------------------------------
% Page layout (KOMA-Script)
%------------------------------------------------------------------------------
% Use KOMA-Script typearea for stable margins
\KOMAoptions{footinclude=false} % Fusszeile wird nicht zu Satzspiegel gezaehlt
\KOMAoptions{headsepline=true} % Trennlinie zwischen Kopfzeile und Text
\KOMAoptions{DIV=12} % beeinflusst Satzspiegel
\KOMAoptions{BCOR=8mm} % Bindekorrektur
\KOMAoptions{numbers=noenddot} % removes dot after chapter/section.. number

%------------------------------------------------------------------------------
% Packages
%------------------------------------------------------------------------------

% Mathematics
\usepackage{amsmath, amssymb, amsthm, mathtools, bm, dsfont, physics, mathrsfs}

% Algorithms
\usepackage{algorithm}
\usepackage{algpseudocode}

% Languages
\usepackage[ngerman, english]{babel}

% Graphics and floats
\usepackage{graphicx, booktabs, tabularx, multirow, leftindex, lscape}
\newcolumntype{C}{>{\centering\arraybackslash}X}  % a centered column type 
\usepackage[format=plain,labelfont=bf,labelsep=colon]{caption}
\usepackage{subcaption}
\renewcommand\thesubfigure{\alph{subfigure}} % enumeration of subfigures

% Code listings
\usepackage{listings}
\usepackage{xcolor}

% Text utilities
\usepackage{csquotes}    % quotation marks according to language
\usepackage{comment}
\usepackage{marginnote}
\usepackage{setspace}

% PGFPlots
\usepackage{pgfplots}
\pgfplotsset{compat=1.18}
% uncomment to compile faster, NOT compliant with pdf/a-2u
%\usepgfplotslibrary{external} 
%\tikzexternalize

%\usepackage{psfrag} % do not know if used for anything?

% Compatibility with third-party packages
\usepackage{scrhack}


%------------------------------------------------------------------------------
% Backref, Hyperref and Cleveref (last)
%------------------------------------------------------------------------------

% enable back-references in bibliography
\usepackage[hyperpageref]{backref} 

% configure hyperref loaded with pdfx package
\hypersetup{
    colorlinks=true,
    linkcolor=blue,
    citecolor=blue,
    urlcolor=blue,
    linktoc=page,
}

% Zum Druck verwende schwarze Links!
%\hypersetup{colorlinks=true, linkcolor=black, citecolor=black, urlcolor=black, linktoc=page}

% uncomment if pdfx is NOT loaded
%% hyperref settings 
%\usepackage[
%    unicode,
%    colorlinks=true,
%    linkcolor=blue,
%    citecolor=blue,
%    pagebackref=true,
%    linktoc=page
%]{hyperref}

% Zum Druck verwende schwarze Links!
%\usepackage[unicode, colorlinks=true, linkcolor=black, citecolor=black, urlcolor=black, pagebackref=true, linktoc=page]{hyperref} 

% load cleverref as last package
\usepackage[capitalise, noabbrev]{cleveref}

%------------------------------------------------------------------------------
% Beautiful Chapter Headings 
%------------------------------------------------------------------------------

% Horizontal rule under chapter title
\renewcommand*{\chapterheadstartvskip}{\vspace*{1cm}}
\renewcommand*{\chapterheadendvskip}{%
  \vspace*{1em}%
  \hrule height 0.4pt%
  \vspace*{2em}%
}

%------------------------------------------------------------------------------
% Draft writing & macros (optional, can remove in final version)
%------------------------------------------------------------------------------

% Todos in floating environments
%\renewcommand{\marginpar}{\marginnote}

% Todonotes workaround with TikZ
%\newcommand{\newtodo}[2][]{\tikzexternaldisable\todo[#1]{#2}\tikzexternalenable}

%------------------------------------------------------------------------------
% Bibliography configuration
%------------------------------------------------------------------------------

\newcommand{\printbibliography}{%
    \bibliography{masterLibrary.bib}%
}
\bibliographystyle{abbrv}

% Backref formatting
\renewcommand*{\backref}[1]{}
\renewcommand*{\backrefalt}[4]{[{%
    \ifcase #1 Not cited.%
          \or Cited on page~#2.%
          \else Cited on pages #2.%
    \fi%
}]}

%------------------------------------------------------------------------------
% Theorem environments
%------------------------------------------------------------------------------

\theoremstyle{plain}
\newtheorem{theorem}{Theorem}[chapter]
\newtheorem{lemma}[theorem]{Lemma}
\newtheorem{corollary}[theorem]{Corollary}
\newtheorem{proposition}[theorem]{Proposition}

\newtheorem{assumption}{Assumption}
\renewcommand{\theassumption}{\Alph{assumption}}

\theoremstyle{definition}
\newtheorem{definition}[theorem]{Definition}
\newtheorem{example}[theorem]{Example}
\newtheorem{examples}[theorem]{Examples}

\theoremstyle{remark}
\newtheorem{remark}[theorem]{Remark}
\newtheorem{remarks}[theorem]{Remarks}

%------------------------------------------------------------------------------
% Math macros
%------------------------------------------------------------------------------

\newcommand{\SetR}{\mathbb R}
\newcommand{\SetN}{\mathbb N}
\newcommand{\SetZ}{\mathbb Z}
\newcommand{\E}{\mathbb E}
\newcommand{\PM}{\mathbb P}
\newcommand{\QM}{\mathbb P}
\newcommand{\Var}{\mathbb V}
\newcommand{\BorelSA}{\mathfrak B}

% Signature
\newcommand{\X}[2]{\mathbb X_{#1}^{#2}}

% RL components
\newcommand{\SpaceS}{{\cal S}}
\newcommand{\SpaceA}{{\cal A}}
\newcommand{\SpaceO}{{\cal O}}
\newcommand{\SpaceH}{{\cal H}}
\newcommand{\KernelP}{{\cal P}}
\newcommand{\KernelR}{{\cal R}}
\newcommand{\mfh}{\mathfrak{h}}
\newcommand{\mfo}{\mathfrak{o}}
\newcommand{\mfa}{\mathfrak{a}}

\DeclareMathOperator*{\argmin}{arg\,min}
\DeclareMathOperator*{\argmax}{arg\,max}

%------------------------------------------------------------------------------
% Enumerate style
%------------------------------------------------------------------------------

\renewcommand\theenumi{(\roman{enumi})}	
\renewcommand\labelenumi{\theenumi}
\renewcommand\theenumii{(\alph{enumii})}
\renewcommand\labelenumii{\theenumii}

%------------------------------------------------------------------------------
% Code listing style
%------------------------------------------------------------------------------

%\definecolor{codegreen}{rgb}{0,0.6,0}
%\definecolor{codegray}{rgb}{0.5,0.5,0.5}
%\definecolor{codepurple}{rgb}{0.58,0,0.82}
%\definecolor{backcolor}{rgb}{0.95,0.95,0.92}

%\lstdefinestyle{mystyle}{backgroundcolor=\color{backcolor}, commentstyle=\color{codegreen}, keywordstyle=\color{magenta}, numberstyle=\tiny\color{codegray}, stringstyle=\color{codepurple}, basicstyle=\ttfamily\small, breakatwhitespace=false, breaklines=true, captionpos=b, keepspaces=true, numbers=left, numbersep=5pt, showspaces=false, showstringspaces=false, showtabs=false, tabsize=2}

%\lstset{style=mystyle}

%------------------------------------------------------------------------------
% PGFPlots settings
%------------------------------------------------------------------------------

\pgfplotsset{width=0.6\textwidth,compat=1.18}

% Externalize figures for faster compilation, NOT compliant with pdf/a2-u
%\usepgfplotslibrary{external}
%\tikzexternalize

% Treat period spacing normally
\frenchspacing


%%%%%%%%%%%%%%%%%%%%%%%%%%%%%%%%%%%%%%%%%%%%%%%%%%%%%%%%%%%%%%%%%%%%%%%%%%%%%%%
%%%%%%%%%%%%%%%%%%%%%%%%%%%%%%%%%%%%%%%%%%%%%%%%%%%%%%%%%%%%%%%%%%%%%%%%%%%%%%%
% Begin document
%%%%%%%%%%%%%%%%%%%%%%%%%%%%%%%%%%%%%%%%%%%%%%%%%%%%%%%%%%%%%%%%%%%%%%%%%%%%%%%
%%%%%%%%%%%%%%%%%%%%%%%%%%%%%%%%%%%%%%%%%%%%%%%%%%%%%%%%%%%%%%%%%%%%%%%%%%%%%%%

\begin{document}

%------------------------------------------------------------------------------
% Title page DEUTSCH V1
%------------------------------------------------------------------------------

\selectlanguage{ngerman}
\pagenumbering{roman}

\newcommand{\svskip}{\vskip 0.1cm}
\newcommand{\mvskip}{\vskip 0.5cm}
\newcommand{\lvskip}{\vskip 1cm}

\begin{titlepage}
  %\selectlanguage{ngerman}
  %\vspace*{-2cm}
  \begin{center}
    \includegraphics[width=0.63\textwidth]{graphics/Logo__FAM__name.png}
    \hskip 1em
    \includegraphics[width=0.16\textwidth]{graphics/UCBC_logo.png}
    \includegraphics[width=0.16\textwidth]{graphics/IEOR_logo.jpg}
    
    \lvskip 
    {\LARGE D\Large ~I~P~L~O~M~A~R~B~E~I~T}
    \vskip 8mm
    {\huge\bfseries Signature-$Q$-Learning: A Non-Markovian Reinforcement Learning Approach \\}
    \lvskip 
    \large 
    Zur Erlangung des akademischen Grades 
    \mvskip 
    {\Large\bfseries Diplom-Ingenieur} 
    \mvskip 
    im Rahmen des Studiums 
    \mvskip 
    {\Large\bfseries Finanz- und Versicherungsmathematik}
    \svskip 
    {am Institut für Stochastik und Wirtschaftsmathematik\\ 
    der Fakultät für Mathematik und Geoinformation\\ 
    TU Wien}
    \mvskip  
    eingereicht von 
    \mvskip
    {\Large\bfseries Felix Stubenvoll-Pretschuh, BSc. BSc.}
    \svskip
    Matrikelnummer {01114103} 
    \lvskip
    Betreuerin: 
    \svskip 
    {Univ.-Prof. Dr. Dipl.-Ing. Christa Cuchiero}\\
    Institut für Statistik und Operations Research, Universität Wien
    \mvskip
    In Zusammenarbeit mit: 
    \svskip
    {Prof. Xin Guo, Coleman Fung Chair in Financial Modeling}\\
    Department of Industrial Engineering and Operations Research \\
    University of California, Berkeley
  \end{center}    
  \vfill    
  \small
    Wien, 17. Dezember 2025 %\today
    %
    \hfill
    % 
    \begin{minipage}[t]{4.5cm}
    \centering
    \underline{\hspace*{4.5cm}}\\
    \small Verfasser
    \end{minipage}
    %
    \hspace{0.6cm}
    \begin{minipage}[t]{4.5cm}
    \centering
    \underline{\hspace*{4.5cm}}\\
    \small Betreuerin
    \end{minipage}  
    \vspace*{-15mm}
\end{titlepage}


\cleardoublepage

%------------------------------------------------------------------------------
% Title page ENGLISH V1
%------------------------------------------------------------------------------

\selectlanguage{english}
\pagenumbering{roman}

\renewcommand{\svskip}{\vskip 0.1cm}
\renewcommand{\mvskip}{\vskip 0.5cm}
\renewcommand{\lvskip}{\vskip 1cm}

\begin{titlepage}
  %\vspace*{-2cm}
  \begin{center}
    \includegraphics[width=0.63\textwidth]{graphics/Logo__FAM__name.png}
    \hskip 1em
    \includegraphics[width=0.16\textwidth]{graphics/UCBC_logo.png}
    \includegraphics[width=0.16\textwidth]{graphics/IEOR_logo.jpg}
    
    \lvskip 
    {\LARGE M\Large ~A~S~T~E~R ~ T~H~E~S~I~S}
    \vskip 8mm
    {\huge\bfseries Signature-$Q$-Learning: A Non-Markovian Reinforcement Learning Approach \\}
    \lvskip 
    \large 
    Submitted for the degree of   
    \mvskip 
    {\Large\bfseries Master of Science}
    \mvskip 
    within the master's program
    \mvskip 
    {\Large\bfseries Financial and Actuarial Mathematics}
    \svskip 
    Institute of Statistics and Mathematical Methods in Economics\\
    Faculty of Mathematics and Geoinformation\\
    Vienna University of Technology\\
    \mvskip  
    by 
    \mvskip
    {\Large\bfseries Felix Stubenvoll-Pretschuh, BSc. BSc.}
    \svskip
    Matriculation number {01114103}
    \lvskip
    Supervisor: 
    \svskip 
    {Univ.-Prof. Dr. Dipl.-Ing. Christa Cuchiero}\\
    Institute for Statistics and Operations Research, University of Vienna
    \mvskip
    In cooperation with: 
    \svskip
    {Prof. Xin Guo, Coleman Fung Chair in Financial Modeling}\\
    Department of Industrial Engineering and Operations Research \\
    University of California, Berkeley
  \end{center}    
  \vfill    
  \small
    Vienna, December 17, 2025 %\today
    %
    \hfill
    % 
    \begin{minipage}[t]{4.5cm}
    \centering
    \underline{\hspace*{4.5cm}}\\
    \small Author
    \end{minipage}
    %
    \hspace{0.6cm}
    \begin{minipage}[t]{4.5cm}
    \centering
    \underline{\hspace*{4.5cm}}\\
    \small Supervisor
    \end{minipage}  
    \vspace*{-15mm}
\end{titlepage}

\cleardoublepage

%------------------------------------------------------------------------------
% Title page DEUTSCH V2
%------------------------------------------------------------------------------

\begin{comment}

\selectlanguage{ngerman}
\pagenumbering{roman}

\renewcommand{\svskip}{\vskip 0.2cm}
\renewcommand{\mvskip}{\vskip 0.6cm}
\renewcommand{\lvskip}{\vskip 1.1cm}

\begin{titlepage}
  %\vspace*{-2cm}
  \begin{center}
    \includegraphics[width=0.63\textwidth]{graphics/Logo__FAM__name.png}
    \hskip 1em
    \includegraphics[width=0.16\textwidth]{graphics/UCBC_logo.png}
    \includegraphics[width=0.16\textwidth]{graphics/IEOR_logo.jpg}
    
    \lvskip 
    {\LARGE D\Large ~I~P~L~O~M~A~R~B~E~I~T}
    \vskip 8mm
    {\huge\bfseries Signature-$Q$-Learning: A Non-Markovian Reinforcement Learning Approach \\} 
    \lvskip
    \large 
    Zur Erlangung des akademischen Grades
    \mvskip   
    {\Large\bfseries Diplom-Ingenieur} 
    \mvskip
    im Rahmen des Studiums
    \mvskip
    {\Large\bfseries Finanz- und Versicherungsmathematik}
    \svskip
    {am Institut für Stochastik und Wirtschaftsmathematik\\ 
    der Fakultät für Mathematik und Geoinformation\\ 
    TU Wien}
    \mvskip
    eingereicht von
    \mvskip
    {\Large\bfseries Felix Stubenvoll-Pretschuh, BSc. BSc.}
    \svskip
    Matrikelnummer {01114103}
  \end{center}    
  \vfill    
    \begin{tabbing}
    \hspace*{44mm} \= \kill
    Betreuerin:
        \> Univ.-Prof. Dr. Dipl.-Ing. Christa Cuchiero, \\
        \> Institut für Statistik und Operations Research, Universität Wien \\[2ex]
    In Zusammenarbeit mit:
        \> Prof. Xin Guo, Coleman Fung Chair Financial Modeling,\\
        \> Department of Industrial Engineering and Operations Research, \\
        \> University of California, Berkeley
    \end{tabbing}
    \lvskip
    \small
    %
    Wien, \today
    %
    \hfill
    % 
    \begin{minipage}[t]{4.5cm}
    \centering
    \underline{\hspace*{4.5cm}}\\
    \small Verfasser
    \end{minipage}
    %
    \hspace{0.6cm}
    \begin{minipage}[t]{4.5cm}
    \centering
    \underline{\hspace*{4.5cm}}\\
    \small Betreuerin
    \end{minipage}  
    \vspace*{-15mm}
\end{titlepage}

\cleardoublepage

\end{comment}



%------------------------------------------------------------------------------
% Title page ENGLISH V2
%------------------------------------------------------------------------------

\begin{comment}

\selectlanguage{english}
\pagenumbering{roman}

\renewcommand{\svskip}{\vskip 0.2cm}
\renewcommand{\mvskip}{\vskip 0.6cm}
\renewcommand{\lvskip}{\vskip 1.1cm}

\begin{titlepage}
  %\vspace*{-2cm}
  \begin{center}
    \includegraphics[width=0.63\textwidth]{graphics/Logo__FAM__name.png}
    \hskip 1em
    \includegraphics[width=0.16\textwidth]{graphics/UCBC_logo.png}
    \includegraphics[width=0.16\textwidth]{graphics/IEOR_logo.jpg}
    
    \lvskip 
    {\LARGE M\Large ~A~S~T~E~R ~ T~H~E~S~I~S}
    \vskip 8mm
    {\huge\bfseries Signature-$Q$-Learning: A Non-Markovian Reinforcement Learning Approach \\} 
    \lvskip
    \large 
    Submitted for the degree of   
    \mvskip   
    {\Large\bfseries Master of Science}
    \mvskip
    within the master's program
    \mvskip
    {\Large\bfseries Financial and Actuarial Mathematics}
    \svskip
    Institute of Statistics and Mathematical Methods in Economics\\
    Faculty of Mathematics and Geoinformation\\
    Vienna University of Technology\\
    \mvskip
    by
    \mvskip
    {\Large\bfseries Felix Stubenvoll-Pretschuh, BSc. BSc.}
    \svskip
    Matriculation number {01114103}
  \end{center}    
  \vfill    
    \begin{tabbing}
    \hspace*{44mm} \= \kill
    Supervision:
        \> Univ.-Prof. Dr. Dipl.-Ing. Christa Cuchiero, \\
        \> Institute of Statistics and Operations Research, University of Vienna \\[2ex]
    In cooperation with:
        \> Prof. Xin Guo, Coleman Fung Chair Financial Modeling,\\
        \> Department of Industrial Engineering and Operations Research, \\
        \> University of California, Berkeley
    \end{tabbing}
    \lvskip
    \small
    %
    Vienna, \today
    %
    \hfill
    % 
    \begin{minipage}[t]{4.5cm}
    \centering
    \underline{\hspace*{4.5cm}}\\
    \small Author
    \end{minipage}
    %
    \hspace{0.6cm}
    \begin{minipage}[t]{4.5cm}
    \centering
    \underline{\hspace*{4.5cm}}\\
    \small Supervisor
    \end{minipage}  
    \vspace*{-15mm}
\end{titlepage}

\cleardoublepage

\end{comment}



%------------------------------------------------------------------------------
% Abstract german and english
%------------------------------------------------------------------------------

\chapter*{Kurzfassung}
\selectlanguage{ngerman}

Die meisten Methoden des Reinforcement Learning (bestärkendes Lernen) gehen davon aus, dass die zugrunde liegende Umgebung durch einen Markov-Entscheidungsprozess (MDP, Markov Decision Process) beschrieben werden kann, in dem der aktuelle Zustand des Prozesses ausreicht, um eine optimale Strategie zu erlernen. Viele reale Probleme weisen jedoch historienabhängige Dynamiken auf, bei denen Zustandsübergänge und Belohnungen von vergangenen Zuständen und Aktionen abhängen können, sodass MDP-basierte Lernmethoden sich als unzureichend erweisen. Diese Arbeit widmet sich dieser Problematik und entwickelt einen Rahmen für Reinforcement Learning in nicht-Markov'schen Umgebungen. Dazu werden historienbasierte Entscheidungsprozesse (HDP, History-based Decision Process) als formales Modell eingeführt sowie der Signature-Q-Learning-Algorithmus vorgestellt, der Pfadsignaturen nutzt, um optimale Strategien in solchen Umgebungen zu approximieren.

Der Hauptteil dieser Arbeit gliedert sich in drei Teile. Zunächst werden in Kapitel \labelcref{chapter:HistoryRLEnvironment} historienbasierte Entscheidungsprozesse eingeführt, die den Rahmen der klassischen MDPs erweitern, um Umgebungen mit historienabhängigen Dynamiken zu modellieren. HDPs bieten einen nichtstationären Rahmen mit endlichem Zeithorizont, in dem die Strategie des Agenten nicht nur auf dem aktuellen Zustand, sondern auf der gesamten beobachteten Zustands- und Aktionshistorie beruht. In diesem Modell entwickeln wir ein dynamisches Programmierprinzip für die optimalen Wertfunktionen, in Analogie zu den Bellman-Gleichungen für MDPs, charakterisieren optimale Strategien als solche, die diesem Prinzip genügen, und beweisen deren Existenz.

Zweitens wird in Kapitel \labelcref{chapter:SignatureQLearning} der Signature-Q-Learning-Algorithmus vorgestellt, eine neue Reinforcement-Learning-Methode für historienabhängige Umgebungen. Dieser basiert auf den universellen Approximationseigenschaften von Pfadsignaturen, einem Konzept, das in der Theorie der rauen Pfade (rough paths) eine wichtige Rolle spielt. Der Algorithmus approximiert die optimale Q-Funktion als lineares Funktional der Pfadsignaturen von beobachteten Zustands-Aktions-Historien und ermöglicht so eine effiziente Repräsentation von zeitlichen Abhängigkeiten. Im Gegensatz zum klassischen Q-Learning in MDPs, das auf einzelnen Zustands-Aktions-Paaren basiert, nutzt Signature-Q-Learning die gesamte Historie und bleibt dabei rechnerisch effizient, indem es gekürzte Pfadsignaturen und eine inkrementelle Aktualisierung auf Basis der Chen-Identität verwendet.

Drittens evaluieren wir in Kapitel \labelcref{chapter:Experiments} den Algorithmus anhand zweier numerischer Experimente. Das Erste stellt eine nur teilweise beobachtbare Umgebung dar und zeigt, dass Signature-Q-Learning vergangene Zustände erfolgreich integriert, um das Problem zu lösen. Das Zweite wendet den Algorithmus auf das Problem der Liquidation einer Aktienposition unter Nutzung des ABIDES-Marktsimulators an, in dem historienabhängige Preiseffekte eine wesentliche Rolle spielen. In beiden Fällen lernt Signature-Q-Learning erfolgreich annähernd optimale Strategien, trotz der nicht-markovschen Natur der Umgebungen. Wir analysieren die Auswirkungen der Länge der Signaturekürzung, diskutieren die Überschätzung von Q-Werten und vergleichen die gelernte Strategie mit einer annähernd optimalen Basisstrategie.

Diese Arbeit leistet einen Beitrag zur Weiterentwicklung des Reinforcement Learning, indem sie historienbasierte Entscheidungsprozesse als  Modellierungsrahmen einführt und das Lernen in historienabhängigen Umgebungen mit der Approximation von Pfadfunktionalen durch die Signatur von Pfaden verbindet. Der vorgeschlagene Algorithmus ist eine modellfreie, online und off-policy Methode und stellt damit ein vielseitiges Werkzeug zur Lösung nicht-markovscher Probleme in verschiedensten Anwendungsbereichen dar. Zukünftige Forschungsschwerpunkte umfassen die Konvergenzanalyse der Methode, eine Verbesserung der Stichprobeneffizienz sowie die Erweiterung des Ansatzes auf kontinuierliche Aktionsräume.

\clearpage

\chapter*{Abstract}
\selectlanguage{english}

Most reinforcement learning (RL) methods assume that the environment follows Markov Decision Process (MDP) dynamics, where the current state is sufficient for learning an optimal policy. However, many real-world problems exhibit history-dependent characteristics, in which transitions and rewards depend on past states and actions, rendering MDP-based learning methods inadequate. This thesis addresses this challenge, with the aim of developing a reinforcement learning framework for non-Markovian environments by introducing History-Based Decision Processes (HDPs) as a model and proposing Signature-Q-Learning, an algorithm that leverages path signatures to approximate optimal policies in history-dependent settings.

This thesis's main body is structured into three parts. First, we introduce History-based Decision Processes in \cref{chapter:HistoryRLEnvironment}, extending the MDP framework to account for history-dependent dynamics. HDPs provide a finite-horizon, non-stationary modeling framework in which an agent’s decision-making is based on observation–action histories rather than individual states. We establish a dynamic programming principle for optimal value functions in HDPs, analogous to the Bellman equations for MDPs. We characterize optimal policies as those that maximize value functions satisfying this principle and prove their existence.

Second, in \cref{chapter:SignatureQLearning}, we propose Signature-Q-Learning, a history-dependent reinforcement learning algorithm based on the universal approximation properties of path signatures, a notion that plays a prominent role in rough path theory. The algorithm formulates the optimal Q-function as a linear functional of the signature of an observation-action history, providing an efficient and expressive representation of temporal dependencies. Unlike traditional Q-learning in MDPs, which operates on state-action pairs, Signature-Q-learning encodes the entire history while maintaining computational efficiency through signature truncation and incremental updates based on Chen’s identity.

Third, in \cref{chapter:Experiments}, we implement and evaluate the algorithm through two distinct numerical experiments. The first experiment modifies the Mountain Car problem to a partially observable setting, demonstrating that Signature-Q-Learning effectively integrates past observations to solve the problem. The second experiment applies the algorithm to a financial execution problem using the ABIDES market simulator, where history-dependent price impacts play a crucial role. In both cases, Signature-Q-Learning successfully learns policies that maximize rewards despite the non-Markovian nature of the environments. We analyze the effects of signature truncation, discuss the overestimation of Q-values, and compare performance against a baseline policy.

This work contributes to reinforcement learning by providing History-based Decision Processes as a novel modeling framework and bridging learning in history-dependent domains with function approximation via path signatures. The proposed method is model-free, online, and off-policy, making it a versatile tool for solving non-Markovian RL problems across various domains. Future research directions include a convergence analysis, improving the method's sample efficiency, and extending it to continuous action spaces.

%------------------------------------------------------------------------------
% Acknowledgments
%------------------------------------------------------------------------------

\chapter*{Acknowledgments}

First of all, I would like to express my deepest gratitude to my supervisor, Christa Cuchiero, for her continuous support, guidance, and genuine curiosity in my project. Her invaluable advice whenever I encountered difficulties, together with the freedom she granted me in shaping the direction of my thesis, greatly enriched my research experience. I am especially thankful for Prof. Cuchiero's support in enabling my six-month stay at the University of California, Berkeley, which broadened both my research and perspective, as well as for the opportunities to present my work within and beyond her group, through which I learned a great deal about collaboration and to reflect on my own research.

I would like to extend my sincere gratitude to my host at Berkeley, Xin Guo. From the very beginning, she welcomed me as a full member of her research group, was open to my ideas, and offered her support through thoughtful discussions and valuable input. My time in her department exposed me to many fascinating and interconnected aspects of financial mathematical research and the many interesting people I met, making my stay in the Bay Area truly inspiring and unforgettable.

I am grateful to the Austrian Marshall Plan Foundation for providing financial support for my stay at Berkeley through their scholarship. Their ongoing efforts to foster transatlantic collaboration first inspired me to pursue a research stay in the US, and I hope many more researchers will have the opportunity to benefit from their support.

Furthermore, I would like to thank Sandra Trenovatz, who always offered a timely, knowledgeable, and helpful hand with all organizational matters during my studies.

I am incredibly thankful to my colleagues at TU Wien, especially Hozan, Dominik, Katrin and David, who were with me until the very end. At the start of my studies, I could not have imagined spending so much time together, but every hour shared, both in and outside of university, was truly worthwhile and made this journey a great experience. I hope this friendship continues indefinitely.

A particular thank you goes to my family, my parents, who trusted my decision to pursue further studies and provided their unconditional love and support, both emotionally and financially, throughout this time. Thank you for always encouraging me, even though you had little idea of what I had been working on over the years. Without you, I would not be where I am today.

Finally, my deepest and most heartfelt thank you to you, Katharina. You have been with me all these years. You shared the most enthusiastic and exuberant times with me. You were there when times got rough, when I doubted my decision to pursue this study. You helped me grow, live, and appreciate the good in every aspect of life. I cannot express the amount of gratitude and love I feel towards you. Your sincerity, kindness, and joyfulness continue to inspire me every day. Without you, I could not have done this. I am with you. Always.

Thank you!
    
\clearpage

%------------------------------------------------------------------------------
% Statement of originality 
%------------------------------------------------------------------------------

\chapter*{Statement of originality}

\vspace*{2cm}

I hereby declare that I have authored the present master thesis independently and did not use any sources other than specified. I have not yet submitted the work to any other examining authority in the same or comparable form. It has not been published yet.

\vspace*{3cm}


\noindent
Vienna, December 17, 2025 %\today
%
\hfill 
%
\begin{minipage}[t]{5cm}
\centering
\underline{\hspace*{5cm}}\\
\small Felix Stubenvoll-Pretschuh
\end{minipage}

\cleardoublepage


%------------------------------------------------------------------------------
% Table of contents
%------------------------------------------------------------------------------

%\onehalfspacing % to increase spacing in toc
\tableofcontents
%\singlespacing
\clearpage
\pagenumbering{arabic}


%%%%%%%%%%%%%%%%%%%%%%%%%%%%%%%%%%%%%%%%%%%%%%%%%%%%%%%%%%%%%%%%%%%%%%%%%%%%%%%
\chapter{Introduction} 
\label{chapter:Introduction}
%%%%%%%%%%%%%%%%%%%%%%%%%%%%%%%%%%%%%%%%%%%%%%%%%%%%%%%%%%%%%%%%%%%%%%%%%%%%%%%


% First paragraph: A short general introduction to reinforcement learning
Reinforcement learning (RL) is a machine learning paradigm designed to address 
sequential decision-making problems through goal-oriented learning. More 
precisely, a RL problem involves an agent interacting with an environment to 
learn optimal sequential decisions that lead to the achievement of a specific 
goal. This learning is typically based on feedback from repeated interactions, 
and occurs without explicit supervision. Many RL approaches do not even require 
prior knowledge of the environment's dynamics, however, they may be learned 
and used to improve decision-making. RL algorithms have achieved notable 
success across various domains, including problems in finance such as optimal 
execution, portfolio management, and option pricing (see 
\cite{hambly2022_rarl} for a recent survey),
surpassing human champions in complex board games such as 
Chess~\cite{silver2018_grlatmcsts} and Go~\cite{silver2017_mgwhk}, as well as 
achieving human-level performance in video games 
\cite{mnih2015_hctdrl,vinyals2019_glsumrl}. Furthermore, the RL variant 
\textit{reinforcement learning from human feedback} was used to fine-tune 
OpenAI's language model \textit{InstructGPT}~\cite{ouyang2022_tlmfiwhf}, and 
its widely successful sibling model \textit{ChatGPT}~\cite{openai2022_chatgpt}.

% Second paragraph: Description of Markov Decision Process and its components as the usual modeling framework in reinforcement learning
Reinforcement learning employs an abstract framework to model the interaction 
between a learning agent and its environment, characterized by states (or 
observations), actions, and rewards. This framework is typically formalized as 
a Markov Decision Process (MDP) \cite{puterman2005_mdpdsdp}, defined over a 
set of states $\SpaceS$, where transitions between states are influenced by actions chosen 
from a set $\SpaceA$. Interactions occur at discrete time steps and we consider a finite 
decision-making horizon $T\in\SetN$. At each step $t\in\{0,1,\dots,T\}$, the 
process is in some state $s_t\in\SpaceS$. Upon taking an action 
$a_t\in\SpaceA$, the process transitions to a new state 
$s_{t+1}\in\SpaceS$ and emits a reward $R_t(s_t,a_t)\in\SetR$, representing 
the immediate payoff associated with the state-action pair $(s_t,a_t)$. The 
transition dynamics are governed by transition functions or kernels $P_t$, 
mapping pairs $(s,a)\in\SpaceS\times\SpaceA$ to probability distributions over 
$\SpaceS$ for each step $t\in\{0,1,\dots,T-1\}$. More specifically, for finite state 
and action spaces, the state transition probability at step $t$ is given by
\begin{equation}
\label{eq:IntroductionMarkovProperty}   
\begin{split}
    P_t(s_t,a_t,s') &= \PM[S_{t+1}=s'\mid S_t=s_t, A_t=a_t] \\
    &=\PM[S_{t+1}=s'\mid S_0=s_0, A_0=a_0, \dots, S_t=s_t, A_t=a_t] 
\end{split}    
\end{equation}
for any $s', s_0,s_1, \dots, s_t\in\SpaceS$ and $a_0,a_1,\dots,a_t\in\SpaceA$, 
where $\{S_u, A_u\}_{u=0,\dots,t}$ and $S_{t+1}$ denote the process' random 
states and random actions, respectively. \cref{eq:IntroductionMarkovProperty} 
is known as the eponymous \textit{Markov property}, which states that the next 
state's distribution at $t+1$ depends solely on the current state and action 
values, and one may discard any information prior to time $t$. A policy is a tuple $\pi = \{\pi_t\}
_{t=0}^{T-1}$ which governs action selection, where $\pi_t$ specifies a rule, either deterministic or stochastic, for selecting actions at step $t$. The goal is to identify a policy that maximizes the 
expected sum of discounted rewards over the decision-making horizon. 
Mathematically, this is equivalent to finding the optimal value functions 
defined as
\begin{equation}
\label{eq:MDPValueFunctionIntroduction}
    V^*_t(s) = \sup_\pi \E\Bigl[
        \sum_{u=t}^{T-1} 
            \gamma^{u-t} R_u(S_u,A_u) + \gamma^{T-t}R_T(S_T) 
        \ \Big|\ S_t=s
    \Bigr], \quad \forall s\in\SpaceS,
\end{equation}
for each $t\in\{0,1\dots,T-1\}$, and as $V_T(s):= R_T(s)$ for $t=T$, with the terminal reward function $R_T: 
\SpaceS\to\SetR$, where the 
discount factor $\gamma \in [0,1]$ weights future rewards. A policy that 
achieves the supremum in \eqref{eq:MDPValueFunctionIntroduction} is called \textit{optimal}. Due to property
\eqref{eq:IntroductionMarkovProperty}, optimal polices in MDPs are functions
of the current state $s_t$ at each $t \in \{0,1, 
\dots,T-1\}$, ignoring prior states visited and actions taken. This Markovian 
nature of optimal policies is a key advantage of MDPs, contributing to their 
popularity as modeling framework in reinforcement learning.

% Third paragraph: Comments on the solvability of Markov Decision Processes via dynamic programming
If $\SpaceS$ and $\SpaceA$ are finite, MDPs can be solved effectively using 
dynamic programming methods, such as backward induction for finite-horizon 
problems \cite[Section~4.5]{puterman2005_mdpdsdp}, and value or policy 
iteration for the infinite-horizon cases \cite[Sections~6.3~and~6.4]%
{puterman2005_mdpdsdp}. These methods can be generalized to continuous state 
and action spaces under certain regularity conditions on the transition 
kernels $P_t$ and reward functions $R_t$, ensuring the measurability of 
policies and the existence of expectations, compare \cite{hinderer2016_do}. 
However, their computational cost grows exponentially with the dimension of 
state and action spaces, a phenomenon known as the \textit{curse of 
dimensionality}. Additionally, these methods require complete knowledge of the 
environment, including transition kernels and reward functions, which is 
usually unavailable in practice.

% Fourth paragraph: Q-learning as reinforcement learning method for learning an optimal policy in MDPs when transition dynamics are unknown
Reinforcement learning algorithms circumvent the need for explicit knowledge 
of the environment by learning policies through sampling transitions from it. A prominent example is the Q-learning algorithm, introduced by 
Watkins \cite{watkins1989_lfdr,watkins1992_q}. The algorithm keeps estimates of 
the optimal action-values or Q-values $Q^*_t(s,a)\in\SetR$ for all pairs $(s,a)$, which represent the expected cumulative discounted reward over the remaining decision horizon from taking action 
$a\in\SpaceA$ in state $s\in\SpaceS$ at time $t$, and acting optimally 
thereafter. It builds on the fact that an optimal policy can be 
derived from the optimal Q-functions $Q_t^*:\SpaceS\times\SpaceA\to\SetR$ by selecting an action that maximizes $Q^*_t(s,\cdot)$ at each step for all $s\in\SpaceS$. Given that the process is in state $s_t$ at step $t$, action $a_t$ is obtained by
\begin{equation}
\label{eq:IntroductionMDPMaximizingPolicy}
    a_t = \pi_t^*(s_t) \in \argmax_{a\in\SpaceA} Q^*_t(s_t,a),
\end{equation}
implying that the knowledge of the functions $Q^*_t$ is equivalent to 
an optimal policy. A dynamic programming relation for the optimal Q-functions 
$Q_t^*$ allows for iterative updating of the estimates towards the true optimal Q-values $Q^*_t(s,a)$ by sampling transition tuples $(s_t,a_t,r_t,s_{t+1})$ from 
the environment. This procedure is a form of \textit{temporal-difference 
learning} \cite[Chapter~6]{sutton2018_rli}, as it updates estimates using the 
immediately obtained rewards together with estimates at the subsequent step as the 
target, a process called \textit{bootstrapping}. When $\SpaceS$ and $\SpaceA$ 
are finite and small the estimates of the values $Q^*_t(s,a)$ can be stored in 
a $\abs{\SpaceS}\times\abs{\SpaceA}$ matrix and the method is referred to as 
\textit{tabular} Q-learning. For large or continuous spaces function 
approximation techniques are required, where one considers a class of 
parametrized functions $\{Q^\theta_t\}_\theta$ with finite-dimensional 
parameter vector $\theta\in\SetR^d$ for some $d\in\SetN$, and learns the 
parameter values that best approximate $Q_t^*$ through gradient-based methods.

% Fifth paragraph: Non-Markovian problems and the problem of learning in them with reinforcement learning methods built on Markovian assumptions
The ability of most reinforcement learning algorithms, including Q-learning, to 
learn an optimal policy depends critically on the assumption that the 
environment adheres to the Markov property. When this assumption is violated, 
learned policies that are functions of the current state may be suboptimal. 
This raises a fundamental question: How can reinforcement learning address 
non-Markovian environments where transitions and rewards depend on the history 
of states and actions? Such dynamics arise in various contexts. In finance, for 
instance, models based on fractional Brownian motion lead to non-Markovian 
price dynamics. In an optimal execution setting, trading 
impacts prices in a history-dependent manner. In memory tasks, for example considered in \cite{heess2015_mcwrnn}, explicit information has to be remembered over a number of steps, and temporal dependencies, such as achieving 
sub-goals in a specific order, also exemplify history-dependent problems.%
\footnote{
    \cite{heess2015_mcwrnn} considers a memory task of controlling a reacher 
    where the target's position is only available at $t=0$. Temporal 
    dependencies are, for example, present in learning to follow a recipe or 
    assembly manual, where the correct order of steps is important and 
    transition dynamics depend on the past.
} Furthermore, %in partially observable domains,
in situations where observations provide only 
partial information about states, the agent is required to infer knowledge 
about the true state from observation histories. As an example consider a robot 
navigating a maze with limited visibility, relying on memory of past moves to 
determine its location. Such scenarios necessitate modeling frameworks 
accommodating history dependence and corresponding learning methods.

% Sixth paragraph: Thesis objective in one sentence and description of the history-dependent modeling framework proposed in this thesis
This thesis has two primary objectives: (1) to introduce \textit{history-based 
decision processes} (HDPs) as a general modeling framework for 
history-dependent dynamics, and (2) to develop \textit{Signature-Q-learning}, a 
novel reinforcement learning algorithm leveraging \textit{path signatures} as 
a core feature for handling history-dependence. HDPs provide a finite-horizon, 
non-stationary modeling framework for reinforcement learning. The agent 
again interacts with the 
environment at steps $t\in\{0,1,\dots,T\}$, now receiving observations 
$o_t\in\SpaceO$ without assumptions about whether the observation space 
$\SpaceO$ represents true states or partial observations thereof. Transition 
dynamics may depend arbitrarily on the full history of observations and 
actions up to the current step,
\[
    h_t := (o_0,a_0,o_1,a_1,\dots,a_{t-1},o_t).
\]
Assuming countable $\SpaceO$, the transition probabilities at step $t$ are 
governed by a transition kernel $P_t$, now a function of $h_t$, and given as
\begin{equation}
\label{eq:IntroductionHDPTransition}    
    \PM[O_{t+1}=o \mid H_t=h_t,A_t=a] = P_t(h_t,a,o),
\end{equation}
where $H_t$ denotes the random history, $A_t$ the random action, and $O_{t+1}$ 
the random observation at steps $t$ and $t+1$ respectively. Similarly, the 
rewards $R_t(h_t,a_t)$ depend on the history-action pair $(h_t,a_t)$. 
\Cref{eq:IntroductionHDPTransition} implies that optimal value functions in this 
framework that are too functions of histories, and we establish a dynamic 
programming procedure for them, analogous to the Markovian case. This provides 
the basis for characterizing optimal policies as maximizing policies in 
\cref{thm:HDPOptimalValueIteration}, which select actions maximizing the sum of 
the immediate reward and the expected value over the remaining episode at each 
single decision step. We prove the existence of such maximizing and thereby optimal 
policies in HDPs in \cref{thm:HDPOptimalPolicyExistence}. 

% Seventh paragraph: Description of the Signature-Q-Learning learning algorithm for learning in the history-dependent framework and based Q-learning in Markov Decision Processes
The Signature-Q-learning algorithm builds upon the optimal Q-functions in HDPs. With some abuse of notation, they are again
denoted as $Q^*_t$ for $t\in\{0,1,\dots,T\}$. However, in the HDP setting, optimal Q-functions are now 
functions of history-action pairs, such that $Q^*_t(h_t, a_t)$ provides the 
expected cumulative reward from taking action $a_t\in\SpaceA$ in history $h_t$, 
and acting optimally thereafter. We show that, analogously to their Markovian
counterparts, knowledge of the functions $Q_t^*$ yields an optimal policy by 
maximizing over actions at each step in the spirit of 
\eqref{eq:IntroductionMDPMaximizingPolicy}, now with histories $h_t $ instead of 
states $s_t$ as input. From the respective procedure for optimal 
value-functions, we show that the optimal Q-functions also satisfy a dynamic 
programming relation, \textit{the optimal Q-iteration for HDPs}, which allows 
us to learn estimates through sampling histories by interaction with the 
environment.

% Eighth paragraph: Information about path signatures as the main ingredient of the proposed Signature-Q-learning algorithm for function approximation
A core component of Signature-Q-learning is the \textit{signature of a path}, 
a mathematical concept originating in the work of Chen 
\cite{chen1957_ipgigbhf,chen1977_ipi} and foundational to rough path theory 
introduced by Lyons \cite{lyons1998_dedrs,lyons2007_dedrp}. The signature 
exhibits some important 
properties, specific references are provided in \cref{chapter:Signature:section:BoundedVariation}. For a path including time as one of its components, its signature at time $T\in(0,\infty)$ uniquely 
determines the path's trajectories over the interval $[0,T]$, providing a 
means of encoding the path's information. Chen's identity 
states an algebraic structure, which facilitates incremental updates of the 
signature via tensor operations, whenever new path information arrives. Linear functionals of the signature admit a universal approximation property, which establishes the 
time-augmented paths' signatures as a linear regression basis for continuous (with respect to some variation metric on path spaces)
path functionals.
These properties have driven the successful development of signature-based 
methods in mathematical finance 
\cite{akyildirim2022_asmmad,bayer2021_osws,cuchiero2023_smtc,
    kalsi2020_oewrps,lyons2020_nphed,min2021_sdfpmfgwcn} 
and machine learning 
\cite{xie2018_lscwfcrnohctr,lyons1998_dedrs,li2019_sgrusfclwpsfttm,
    min2021_sdfpmfgwcn}.  

% Ninth paragraph: The specific application of the path signature in Signature-Q-learning and history-dependent framework

In Signature-Q-learning, path signatures enable compact representation of 
histories. Since enumerating histories and keeping an estimate for each pair 
$(h_t,a)$ is computationally unfeasible, the method represents the optimal 
Q-functions as linear functionals of the signature of observation-action 
histories and learns the functional's parameters. For a history $h_t$, we 
denote its time-augmented history, which includes the time-steps $t$ as an 
additional component, by $\hat{h}_t$. Any pair $(\hat{h}_t,a_t)$ is then given 
as a stream
\[
    \hat{h}_t^{a_t} := (u, o_u, a_u)_{u=0}^t.
\]
The linear interpolation of $\hat{h}_t^{a_t}$ constitutes a continuous path in 
$\SetR^{d+2}$, with $d\in\SetN$ the dimension of observation space $\SpaceO$. 
Let its signature be denoted by $S_t(\hat{h}_t^{a_t})$. It provides a feature 
vector for approximating the optimal Q-functions by parametrized maps
\begin{equation}
\label{eq:IntroductionQParametrization}    
    Q(h_t,a_t\mid\ell) 
    %:= \langle \ell, S_t(\hat{h}_t^{a_t}) \rangle,
    := \ell\bigl(S_t(\hat{h}_t^{a_t}) \bigr),
\end{equation}
where $\ell$ is a linear functional acting on the 
signature $S_t(\hat{h}_t^{a_t})$. We show that, under some regularity conditions, the universal approximation property guarantees that the optimal Q-functions can be approximated arbitrarily well by such functionals. Signature-Q-Learning learns the parameters of a linear functional $\ell$, which provides a sufficiently good estimation. The algorithm's learning mechanism is based on the optimal Q-iteration in HDPs. We highlight that the approximation argument together with the learning procedure yields a \textit{single} functional of 
the form \eqref{eq:IntroductionQParametrization} that approximates the optimal 
Q-functions at \textit{all} steps $t\in\{0,1,\dots,T\}$, instead of a separate 
functional at each step, ensuring efficient use of samples. Furthermore, Chen's 
identity allows for computationally inexpensive updating of signatures with 
sequentially arriving data, aligning naturally with reinforcement learning's 
iterative nature. The proposed method constitutes a model-free, online, and 
off-policy RL algorithm, and offers a versatile tool for learning in 
history-dependent environments.

%------------------------------------------------------------------------------
\paragraph{Related Work}
\label{chapter:Introduction:par:RelatedWork}

Using path signatures to approximate Q-functions in an RL setting is partly
motivated by the work of Bonnier et al. \cite{bonnier2019_dst}, who propose a model
for processing sequential data, with the signature acting as a layer
within deep neural networks. Experiments include a partially observable 
RL problem similar to that considered in 
\cref{chapter:Experiments:section:MountainCar}. The method differs from the 
one presented in this thesis as it relies on universal approximation 
properties for (deep) neural networks, resulting in a non-linear functional of 
the signature as Q-function approximation. To the best of our knowledge, apart 
from the experiment in \cite{bonnier2019_dst}, the signature has not been 
explicitly incorporated into an approximation procedure in an RL setting.

A substantial body of literature exists on Markov Decision Processes. We 
highlight the two primary references for the present work. The monograph 
of Puterman \cite{puterman2005_mdpdsdp} provides a unified and 
comprehensive treatment of MDPs and their theoretical and computational 
solvability, whereas Hinderer \cite{hinderer2016_do} concentrates on 
finite-horizon MDPs and develops several frameworks ranging from models with 
finite state and action spaces to those on arbitrary measurable spaces and even 
random environments. In its mathematical treatment, the 
latter refrains from advanced measure theory, emphasizing the structure of 
solutions and the construction of specific models in various applications.

Approaches to address non-Markovian RL include expanding the state 
representation and introducing partial observability. In the former, one 
defines a new state variable $\tilde{s}_t := (s_{t-k+1},
s_{t-k+2}, \dots, s_{t-1}, 
s_t)$, that includes the states over the last $k$ steps at each time
point as input to the learning algorithm. This is equivalent to implicitly 
assuming that the environment follows a \textit{$k$-order} MDP, where the 
distribution of the next state, given a chosen action, depends on the last $k$ 
states. If this assumption is true, the problem becomes Markovian with respect 
to $\tilde{s}_t$. This is, for example, the case in the work of Mnih et al. 
\cite{mnih2015_hctdrl}, who introduce Deep-Q-Networks, in which $\tilde{s}_t$ 
contains four frames of a computer game. In other cases, however, it remains 
unclear whether this approach indeed induces an MDP; see, for instance, Sadighian 
\cite{sadighian2019_drlcmm,sadighian2020_edrlfcmm}. The latter approach models 
the problem as a \textit{partially observable Markov Decision Process} (POMDP), 
following the classical formulation by Kaelbling et al. \cite{kaelbling1998_paposd}, 
which is a latent MDP for which the decision maker only receives partial 
observations of the true Markovian states, and must infer more information 
about them from the history of observations. While POMDPs are of theoretical 
interest on their own, learning in such an environment proved to be notoriously 
hard; see, for example, Liu et al. \cite{liu2022_wporls}. 

Note that the history-based environment model proposed in the present thesis 
incorporates both mentioned approaches. If the environment follows a POMDP, 
it naturally uses the history of available observations and actions to create
learning policies. If the environment follows a $k$-order MDP, the
observations in the HDP setting represent the true state, and in incorporating
the entire history prior to the current step, we include the necessary $k$
past steps on which the transition dynamics depend.

Another line of work on general non-Markovian RL is the Feature-RL framework proposed by Hutter \cite{hutter2009_frlpum,hutter2009_frlpsm,hutter2016_esabmdp}. 
It aims at learning a map $\phi$ that aggregates histories to a finite set of 
states, $s_t = \phi(h_t)$, such that the aggregated process is (approximately) 
a small MDP, which can be solved with classical methods such as policy and 
value iteration. Learning $\phi$ crucially depends on the class of functions 
to choose from, which dictates the possible MDP representation, the search 
method applied within this class, and the cost functional, which defines what 
a good MDP representation looks like. Typical function classes in this setting 
are tree-based, for example, suffix trees. Although Feature-RL is inherently 
model-based, Daswani and Hutter \cite{daswani2013_qhrl} propose a model-free 
variant, in which the learning objective for $\phi$ is the direct optimization 
of rewards. The authors propose a value-based cost function\footnote{
    More precisely, a regularized, pathwise squared Bellman error for 
    Q-functions.
}, which intertwines the selection of the feature map $\phi$ and Q-learning, 
such that the optimization for $\phi$ compares different maps based on 
the resulting Q-learning performance in the corresponding aggregated process. 
The process induced by $\phi$ no longer has to be a strict MDP, and the 
approach extends to linear function approximation. There are some 
similarities between \cite{daswani2013_qhrl} and the method in this thesis. 
Both consider general RL starting from histories; however, they differ in that 
we do not assume a finite observation space, and consider an a priori fixed 
feature map $\phi$, specifically the signature transform. Furthermore, the resulting 
range of our feature map is still uncountable and does not aim at inducing any 
Markovian structure. Both works rely on linear function approximation, with 
the difference that we approximate optimal Q-functions defined directly in the 
history-based setting, as opposed to Q-functions for some finite-state 
aggregated process.

Learning in partially observable domains related to the present thesis is 
addressed in \cite{hausknecht2015_drqpom,heess2015_mcwrnn}. Both works utilize 
recurrent neural networks as function approximators to memorize information 
through time. Heess et al. \cite{heess2015_mcwrnn} extend two existing 
actor–critic algorithms, where the critic is a parametrized estimate of the 
Q-function whose gradient is used to update a parametrized policy. Policies and 
Q-functions are considered as functions of histories $h_t$, and are 
parametrized as recurrent and long-short-term memory networks. Note that if the 
policy in Algorithm 1 in \cite{heess2015_mcwrnn} were to select actions according 
to the maximal action value provided by the current Q-function estimate, the 
algorithm would amount to a Q-learning variant, closely related to 
Signature-Q-Learning in \cref{chapter:SignatureQLearning}. Hausknecht and Stone 
\cite{hausknecht2015_drqpom} build on the Deep-Q-Networks introduced by Mnih 
et al. \cite{mnih2015_hctdrl}, extending them through the addition of a recurrent 
layer to encode past information, whereby the network serves as a parametrized 
Q-function.

Finally, we provide some references to Q-learning in MDPs, as it is the 
starting point for Signature-Q-Learning. Under suitable conditions on the sampling 
policy and learning rates, tabular Q-learning was shown to converge in the 
stationary case by Tsitsiklis \cite{tsitsiklis1994_asaq} and by Jaakkola et al. 
\cite{jaakkola1993_csidpa}. Convergence with function approximation is less 
established and may fail due to the combination of off-policy sampling, 
bootstrapping, and approximation; see \cite[Section~11.3]{sutton2018_rli} and 
\cite{hasselt2018_drldt}. 
Convergence with linear function approximation was established by Szepesvári
\cite{szepesvari2004_iq} under restrictions on the features, and by Melo
\cite{melo2008_arlwfa} under restrictive assumptions on the sampling policy.
Maei \cite{maei2010_tolcwfa} relaxes some of these assumptions and proposes an 
algorithm that is guaranteed to converge to a local minimum.
Furthermore, Bhatnagar et al. \cite{bhatnagar2016_mqwlfa} and Carvalho et al.
\cite{carvalho2020_cvqwlfa} establish convergence for two time-scale variants of 
Q-learning with linear function approximation. Convergence results for Q-learning 
with non-linear function approximation remain scarce, although numerical 
performance is often strong; examples include Neural-fitted Q-learning
by Riedmiller \cite{riedmiller2005_nfifewdenrlm} and, most notably, the success 
of Deep-Q-Networks by Mnih et al. \cite{mnih2015_hctdrl}, where the Q-function 
is represented by a deep neural network.

We note that most work on Q-learning, including the references in the
preceding paragraph, addresses infinite-horizon, stationary
MDPs, where $T=\infty$ and the transition dynamics are time-homogeneous. In
such cases, the optimal Q-function is time-independent, so that a single function
$Q^*$ describes the expected cumulative reward given a pair $(s,a)$ at any time
$t\in\mathbb N$, assuming optimal behaviour thereafter. Tabular Q-learning
for finite-horizon, non-stationary MDPs is treated by Garcia and Ndiaye 
\cite{garcia1998_lrarlaf}, and recently, Veness and Papadimitriou 
\cite{vp2024_eemsgwfhq} established its convergence via stochastic approximation 
techniques. Moreover, stationary, infinite-horizon MDPs are often used as modeling 
assumptions in RL problems, either explicitly or implicitly through the application 
of RL methods built on such assumptions, even when the true problem is 
finite-horizon and potentially non-stationary. The work of Majeed and Hutter 
\cite{majeed2018_qcndp} explains why Q-learning often still performs well beyond 
the MDP setting by proving that convergence extends to certain classes of 
non-Markovian environments. Drawing on ideas from Feature-RL, the authors show that
state-uniformity of the optimal Q-function is both necessary and sufficient for 
convergence. This means that whenever two histories are mapped to the same aggregated state via a feature map $\phi$, their optimal Q-values must
agree. This substantially strengthens the theoretical foundations of Q-learning, 
showing that it may remain a sound learning mechanism even in non-stationary and 
non-Markovian domains, which are common in practical applications.


%------------------------------------------------------------------------------
\paragraph{Thesis Outline}
\label{chapter:Introduction:par:ThesisOutline}

% chapter signatures
In \cref{chapter:Signature} the signature of paths of bounded variation is 
defined. We state its properties which make it a natural choice for the 
application as function approximator in our proposed reinforcement learning 
method, most importantly its universal approximation property given in 
\cref{thm:UniversalApproximation}, and provide details on the signature's
computation for discrete data streams.

% chapter reinforcement learning and MDP
\Cref{chapter:RLAndMDP} provides a thorough introduction to reinforcement 
learning, its components and MDPs as the usual modeling framework, which will 
serve as the basis for the generalization to the history-dependent set-up. 
We define in mathematical detail a finite-horizon, non-stationary MDP with 
Borel state space and finite action space and make precise the notion 
of optimality, with the help of optimal value and Q-functions. 
\Cref{eq:MDPMaximizingPolicy} characterizes optimal policies as 
maximizing policies and the existence of such in the considered set-up is 
proved in \cref{thm:MDPOptimalPolicyExistence}. Furthermore, we present the Q-%
learning method for MDPs based on the Bellman equations \eqref{eq:MDPBellmanQ}, 
including a detailed procedure for the tabular case in 
\cref{alg:MDPQLearning}, and a description on how to use function approximation 
if the state space is large or infinite. The chapter concludes with remarks on 
the shortcomings of using MDPs as the model for RL problems, serving as 
motivation for the introduction of a more general framework.

% chapter history dependent decision process
\Cref{chapter:HistoryRLEnvironment} constitutes the main theoretical part of 
the thesis. We introduce History-dependent decision processes (HDPs) as a 
generic model for RL problems without any Markovian assumptions on transition 
dynamics. Generalizing the exposition in the previous chapter a 
detailed description of its components is given, most notably transition 
dynamics depending arbitrarily on observation-action histories. This then 
leads to the definition of optimal value functions 
(\cref{def:HDPOptimalValueFunctions}), now in terms of histories instead of 
single states. Analogously to MDPs, \cref{thm:HDPOptimalValueIteration} 
characterizes optimal policies as maximizing policies, and further provides us 
with the \textit{optimal value iteration} in HDPs, a dynamic programming 
procedure from which an optimal policy could be obtained, if the HDPs 
components are known and computable. The chapter concludes by proofing the 
existence of maximizing policies in \cref{thm:HDPOptimalPolicyExistence}.

% chapter Signature-Q-learning
\Cref{chapter:SignatureQLearning} introduces Signature-Q-Learning as learning 
method for history-dependent decision 
processes, and as a generalization of Q-learning in MDPs. To this end, optimal 
Q-functions in HDPs are defined as functions of observation-action histories, 
and we detail how their knowledge is sufficient to obtain an optimal policy 
$\pi^*$. \Cref{eq:HDPOptimalQIteration} states the optimal Q-iteration in 
HDPs, which serves as the algorithm's basis. Building on the path signature's 
properties collected in \cref{chapter:Signature}, most notably the universal 
approximation property, we parametrize optimal Q-functions as a linear 
functional of the histories' signature and detail the process of learning the 
functional's parameters by sampling histories from the environment. A detailed 
algorithm for Signature-Q-Learning is provided in \cref{alg:SignatureQLearning}, 
and aspects regarding its implementation are discussed.

% chapter experiments
Finally, in \cref{chapter:Experiments} we show Signature-Q-Learning's success 
in solving reinforcement learning problems in two experiments. In each 
experiment, we construct the environment with a description of its elements, 
give details on the implementation and the training procedure. The 
first experiment is a partially observable modification of the classic 
Mountain Car problem from the RL library \texttt{Gymnasium} 
\cite{towers2023_gv0}. The algorithm's performance is discussed by comparing 
different configurations for parameters distinct to the Signature-Q-learning, 
namely the signature's truncation order and the basepoint used for its 
computation. The environment offers a clear terminal state that must be 
reached to count an episode as a success, and while the method learns a 
strategy that reaches this states, it overestimates Q-values. We discuss this 
problem, including a hypothesis that identifies the 
different episode lengths as the reason. The second experiment constitutes a 
simple execution problem. It is implemented as custom environment based on the 
market simulator ABIDES \cite{byrd2020_athmms,amrouni2022_aemdesafm}. We 
compare the policies obtained from Signature-Q-Learning to a baseline policy 
that serves as proxy for optimality. Albeit the algorithm does not achieve 
full optimality, it converges towards the baseline policy, and to the true 
optimal Q-values, with the issues observed in the first experiment not being 
present in the second. The code for both experiments is available on the 
\href{https://github.com/phaelicks}{author's github account.}\footnote{
    \url{https://github.com/phaelicks}
} 

% chapter stochastic kernel preliminaries
Additionally, \cref{chapter:PreliminariesStochasticKernels}, presents key results 
on stochastic transition kernels and provides a theoretical foundation for the 
decision processes discussed in earlier chapters. In particular, we revisit 
the construction of discrete-time stochastic processes governed by a specified 
collection of transition kernels. This framework underpins the modeling of 
reinforcement learning problems, where the evolution of the environment is 
determined step-by-step by these kernels.

\begin{comment}
%------------------------------------------------------------------------------
\paragraph{Limitations}
\label{chapter:Introduction:par:Limitations}

We want to highlight four limitations of this work and the proposed method.
\begin{itemize}
    \item The environment is defined in terms of observations, such that it naturally applies when either true states are observed or partial observations thereof. Optimal policies, however, are only defined with respect of those observations. That is if there are some underlying hidden Markovian dynamics, optimal policies as defined in \cref{def:HDPOptimalPolicy} may not be optimal with respect to the hidden true states.

    Observability assumptions: if observations are only partial, the policy derived is only optimal with respect to the information contained in observations. If there are hidden Markovian dynamics, the results give no indication about whether some policy is optimal for the hidden MDP.

    \item Both experiments are considered successes because they 
    show the algorithm exhibiting the desired property of maximizing rewards 
    and learning a policy that is approximately optimal. However, limitations 
    persist, and further improvement are necessary.

    \item No convergence result for the proposed reinforcement learning method. One approach would be to use stochastic approximation results where one update step in the stochastic approximation procedure would happen over episodes, in a similar fashion to \cite{vp2024_eemsgwfhq}.

    \item In the application of the universal approximation theorem for linear functionals of path signatures to the HDP Q-functions, continuity of the optimal Q-function is assumed in order for the universal application property to be applicable. 

\end{itemize}
\end{comment}

%------------------------------------------------------------------------------
\paragraph{Notation}
\label{chapter:Introduction:par:Notation}

The set of all probability measures on a measurable space $(\Omega, 
\mathcal{F})$ will be denoted $\text{Prob}(\Omega,\mathcal F)$ or simply 
$\text{Prob}(\Omega)$ if the corresponding $\sigma$-algebra is clear from the 
context. For a real-valued random variable $X$ on $(\Omega, \mathcal{F})$ with
distribution $P \in \text{Prob}(\Omega)$ we write $X \sim P$, while a 
realization of $X$ will usually be written as lower case letter $x = X(\omega) 
\in \SetR$ for some $\omega\in \Omega$. Given a second measurable space 
$(\Omega', \mathcal{F}')$ the product-$\sigma$-algebra on $\Omega\times\Omega'$ 
is denoted $\mathcal{F} \otimes\mathcal{F}'$. With some notational overload, 
the symbol $\otimes$ will also be used to define product kernels, and as the 
tensor product, however, the respective meaning will always be clear from the 
objects involved. 
%
Given a topological space $(\mathcal{X},
\mathcal{T})$, we will denote the Borel $\sigma$-algebra on $\mathcal{X}$ 
induced by the topology $\mathcal{T}$ as $\mathcal{B}_\mathcal{X}$.
%
Furthermore, for any $T\in\SetN$ define 
\[
    [T]:= \{0,1,\dots,T\}
\]
and the extended real numbers as $\overline{\SetR} := \SetR\cup\{-\infty, 
+\infty\}$. Finally, for a function $f:\Omega\to\overline{\SetR}$ we denote its positive and negative absolute part by
\[
    f^+:= \max\{0,f\} \quad \text{and} \quad f^-:= \max\{0,-f\}.
\]

%%%%%%%%%%%%%%%%%%%%%%%%%%%%%%%%%%%%%%%%%%%%%%%%%%%%%%%%%%%%%%%%%%%%%%%%%%%%%%%
\chapter{The Signature of a Path} 
\label{chapter:Signature}
%%%%%%%%%%%%%%%%%%%%%%%%%%%%%%%%%%%%%%%%%%%%%%%%%%%%%%%%%%%%%%%%%%%%%%%%%%%%%%%


This sections describes the main ingredient of the proposed reinforcement 
learning method, namely the \textit{signature of a path}. We restrict 
ourselves to continuous paths of bounded $1$-variation, however, we highlight
that the definition of the signature and most of its properties can be 
extended to a much broader class of less regular paths, namely (weakly) 
geometric rough paths, see \cite{lyons2007_dedrp} for a thorough introduction 
to the theory.\footnote{
    It is possible to define the signature also for non-geometric rough paths, however some properties change, such as Chen's identity and integration by parts as provided by the shuffle product, compare~\cite{hairer2015_gvnrp}.
}
We start with some basic notion before giving the precise definition of the
signature, followed by an exposition of some of its algebraic and analytic
properties that make it suitable for the application in a reinforcement
learning setting. We especially emphasise a version of the universal
approximation theorem in \cref{thm:UniversalApproximation}, which enables the
signature's use to approximate linear functions on path spaces. Furthermore,
we define the signature of a discrete stream of data and demonstrate its
computation in practice. The exposition's structure loosely follows 
\cite{bonnier2019_dst,lyons2023_smml} while its contents and notation draw from
\cite{cuchiero2023_smtc,bayer2021_osws}. 
%
%------------------------------------------------------------------------------
\section{The Signature of Paths of Bounded Variation}
\label{chapter:Signature:section:BoundedVariation}
%------------------------------------------------------------------------------

We start by defining the space in which the signature takes its values. 
For $d\in\SetN$ and $n\in\SetN_0$ the $n$-th tensor product of 
$\SetR^d$ is given by
\[
    (\SetR^d)^{\otimes 0} := \SetR, 
    \quad 
    (\SetR^d)^{\otimes n} := \underbrace{
        \SetR^d \otimes \dots \otimes \SetR^d
            }_{
        n \text{ times}
        },
\]
with $\otimes$ the tensor product.

\begin{definition}[Extended tensor algebra]
\label{def:ExtendedTensorAlgebra}
    Let $d\in\SetN$. The extended tensor algebra over $\SetR^d$ is defined as
    \[
        T((\SetR^d)) := \prod_{n\geq 0} (\SetR^d)^{\otimes n}
        = \{\bm{a} := (a_0,a_1,\dots, a_n,\dots) 
            \mid a_n\in(\SetR^d)^{\otimes n}, n\in \SetN_0\}.
    \]
    We equip $T((\SetR^d))$ with an addition, a multiplication and scalar 
    multiplication, defined for $\bm{a}, \bm{b} \in T((\SetR^d))$ and 
    $\lambda\in\SetR$ by
    \begin{gather*}
        \bm{a} + \bm{b} := (a_0+b_0, a_1+b_1, \dots, a_n+b_n,\dots) 
            \in T((\SetR^d)), \\
        \bm{a} \otimes \bm{b} := (c_0, c_1, \dots, c_n, \dots ) 
            \in T((\SetR^d)),        
    \end{gather*}
    where
    \[
        c_n := \sum_{k=0}^n (a_k \otimes b_{n-k}) \in (\SetR^d)^{\otimes n},
    \]
    and $\lambda \bm{a} := (\lambda a_0, \lambda a_1, \dots, \lambda a_n, 
    \dots)$. Then, $(T((\SetR^d)), +, \cdot, \otimes)$ forms a non-commutative 
    algebra over $\SetR$ with unit $\bm{1} = (1,0,0, \dots)$.
\end{definition}

Intuitively, $T((\SetR^d))$ contains all sequences of tensors with of 
increasing order with components in $\SetR$. We further define the space 
of all finite sequences of tensors as well as the sub-set of all sequences of
fixed length.

\begin{definition}[(Truncated) tensor algebra]
\label{def:TensorAlgebra}
    Let $d\in\SetN$. The \textit{truncated tensor algebra} of order 
    $N\in\SetN$ over $\SetR^d$ is defined as
    \[
        T^{(N)}(\SetR^d) := \{\bm{a} \in T((\SetR^d))
            \mid \forall n\geq N: a_n = 0\},        
    \]
    with dimension\footnote{
        Strictly speaking, $T^{(N)}(\SetR^d)\subset T((\SetR^d))$ is 
        infinite-dimensional, however it can readily be identified with the 
        $d^N$-dimensional space $\bigoplus_{n=0}^N(\SetR^d)^{\otimes n}$.
    } 
    $d_N := \sum_{k=0}^n d^k = (d^{N+1} - 1) / (d-1)$. The \textit{tensor 
    algebra} over $\SetR^d$ is defined as
    \[
        T(\SetR^d) := \bigcup_{N\in\SetN} T^{(N)}(\SetR^d).
    \]
    Both, the truncated and the tensor algebra, equipped with addition, 
    multiplication and scalar multiplication as in
    \cref{def:ExtendedTensorAlgebra} form a a non-commutative, unitary algebra
    over $\SetR$.    
\end{definition}

For $n\in\SetN$, and a multiindex $I=(i_1, i_2\dots,i_n)$ with entries 
$i_j\in\{1,2,\dots,d\}$ for $1\leq j \leq n$, we set $\abs{I}:= n$. Further, 
let $E = \{e_1,\dots, e_d\}$ be the canonical basis of $\SetR^d$. For each multiindex $I$ with $\abs{I} = n \geq 1$ we set
\[
    e_I := e_{i_1}\otimes e_{i_2}\otimes \dots \otimes e_{i_n}.
\]
Observe that the set $\{e_I \mid \abs{I} = n\}$ is the canonical orthonormal 
basis of $(\SetR^d)^{\otimes n}$ for each $n\in\SetN$. If we denote the basis 
element of $(\SetR^d)^{\otimes 0} = \SetR$ by $e_\emptyset$, and set 
$\abs{\emptyset} := 0$ for the empty index $I=\emptyset$, each element 
$\bm{a}\in T((\SetR^d))$ can be written as the formal sum
\[
    \bm{a} = \sum_{n\geq0}\sum_{\abs{I}=n} a_Ie_I,
\]
with $a_I\in\SetR$ for each multiindex $I$.

An element $\bm{a}\in T((\SetR^d))$ is invertible if and only if
$a_0\neq0$ \cite{lyons2007_dedrp}. In particular, we consider the subset 
$T_1((\SetR^d)) := \{\bm{a}\in T((\SetR^d)\mid a_0 = 1\}$. For 
$\bm{b}\in T_1((\SetR^d))$, its inverse is given as
\[
    \bm{b}^{-1} := \sum_{n\geq 0} (\bm{1} - \bm{b})^{\otimes n},
\]
which is well defined since for each $n\geq 0$, due to $(\bm 1-\bm{b})_0 = 0$,
only finitely many terms of the sum contain non-zero tensors of degree $n$.
Note that $(T_1((\SetR^d)), \otimes)$ forms a group, and the signature
of a path defined below takes its values in this group.

Next, we describe linear functionals acting on elements $\bm{a}\in
T((\SetR^d))$. Given a multiindex $I$ and basis element $e_I$ we define the
corresponding linear functional acting on $\bm{a}$ by $\langle e_I,\bm{a}
\rangle := a_I$. These basic functionals are extended to linear functionals
$\langle \ell, \cdot \rangle: T((\SetR^d)) \to \SetR$ for any element $\ell = 
\sum_{\abs{I}\geq 0} \ell_I e_I \in T(\SetR^d)$ with $\ell_I \in\SetR$, that 
act on $\bm{a}$ by 
%$\ell\in T(\SetR^d)$ we define the linear functional 
%$\langle \ell, \cdot \rangle: T((\SetR^d)) \to \SetR$ by
\[
    \langle \ell, \bm{a} \rangle
    := \sum_{\abs{I} \geq 0} \ell_I \langle e_I, \bm{a} \rangle
    = \sum_{\abs{I}\geq 0} \ell_I a_I.
\]
Observe that the sum is well defined since $\ell_I=0$ for all but finitely
many $I$. With a slight abuse of notation we will refer to 
$\ell\in T(\SetR^d)$ as the corresponding linear functional $ \langle \ell,
\cdot \rangle$ as well as the family of coefficients 
$\{\ell_I\}_{\abs{I}\geq n}$. For practical purposes, we will consider linear 
functionals of fixed length $N\in\SetN$ in \cref{chapter:SignatureQLearning}, 
such that $\ell_I = 0$ for all multiindices with $\abs{I}\geq N$. In 
this case $\ell \in T^{(N)}(\SetR^d)$, and we will implicitly identify $\ell$
with a finite-dimensional vector in $\SetR^{d_N}$, where $d_N$ is the 
dimension  of $T^{(N)}(\SetR^d)$, when considering practical implementations.

A path $X:[0,T]\to\SetR^d$ is said to be of finite 1-variation on the interval 
$[s,t] \subset [0,T]$, if its total variation is finite,
\begin{equation}
\label{eq:1-VariationNorm}
    \norm{X}_{1,[s,t]} := 
    \sup_{\mathcal{D} \subset [s,t]} \sum_{t_i\in\mathcal{D}} 
    \norm{X_{t_i} - X_{t_{i-1}}}_1
        < \infty,
\end{equation}
where the supremum is taken over all finite partitions $\mathcal{D}=\{t_i\}_i$
of $[s,t]$, and $\norm{\cdot}_1$ denotes the 1-norm on $\SetR^d$. We denote 
the space of all continuous paths of finite 1-variation on $[0,T]$ as 
$\mathcal{V}_1([0,T], \SetR^d)$ and the subset of all
such paths starting at the origin by $\mathcal{V}_1^0([0,T], \SetR^d) := \{X
\in \mathcal{V}_1 ([0,T], \SetR^d) \mid X_0=0 \}$. For two paths 
$X,Y\in\mathcal{V}_1([0,T],\SetR^d)$ we define their 1-variation distance over 
an interval $[s,t] \subset [0,T]$ as
\begin{equation}
\label{eq:1-VariationDistance}
    d_{1,[s,t]} (X,Y) := \norm{X-Y}_{1,[s,t]}.
\end{equation}
Note, that the space $\mathcal{V}_1^0([0,T], \SetR^d)$ equipped with 
$d_{1,[0,T]}$, becomes a metric space.

We are now ready to define the signature of a path of finite variation.

\begin{definition}%[Signature]
\label{def:Signature}
    Let $X\in\mathcal{V}_1([0,T], \SetR^d)$. We define the \textit{signature} 
    of $X$ over the interval $[s,t]\subset[0,T]$ as  
    \begin{equation}
    \label{eq:Signature}
        S_{s,t}(X) 
        := (1, S^1_{s,t}(X), \dots, S^n_{s,t}(X), \dots) \in T_1((\SetR^d)),
    \end{equation}
    with
    \begin{equation*}
        S^n_{s,t}(X) 
        := \int_{s<u_1<\dots<u_n<t} \dd{X_{u_1}}\otimes\dots
            \otimes \dd{X_{u_n}} \in (\SetR^d)^{\otimes n}.
    \end{equation*}
    Further, we define the \textit{truncated signature up to order}
    $N\in\SetN$ of $X$ as
    \begin{equation}
        S^{(N)}_{s,t}(X) 
        := (1, S^1_{s,t}(X), \dots, S^N_{s,t}(X)) \in T^{(N)}(\SetR^d).
    \end{equation}
    For the case $s=0$ we write $S_t(X) := S_{s,t}(X)$ and if we only
    refer to the signature of $X$, we implicitly mean the signature
    over the whole interval, i.e. $S_{T}(X)$. 
\end{definition}

The $n$-th component of the signature of $X$ is an element in the $n$-th 
tensor product of $\SetR^d$. The individual coordinate values of those tensors 
are given by the following.

\begin{definition}%[Coordinate iterated integrals]
\label{def:CoordinateInteratedIntegrals}
    For a multiindex $I=(i_1,\dots,i_n) \in\{1,\dots,d\}^n$, $n\in\SetN$,
    and $X\in\mathcal{V}_1([0,T], \SetR^d)$, we define the \textit{coordinate
    iterated integral} of $X$ over $[s,t]\subset [0,T]$ as
    \begin{equation}
        S_{s,t}^I(X) := \big\langle e_I, S_{s,t}(X) \big\rangle
        = \int_{s<u_1<\dots<u_n<t} \dd{X^{(i_1)}_{u_1}} \dots
            \dd{X^{(i_n)}_{u_n}}.
    \end{equation}
\end{definition}

\begin{remark}
    \begin{enumerate}
        \item Since we consider paths of finite 1-variation, the integrals
        in \cref{def:Signature} are understood in the sense of the Stieltjes
        integral. We have noted that the signature can be defined for less regular paths \cite{lyons2007_dedrp}. For example, if $X$ is a 
        $\SetR^d$-valued semimartingale, then the integrals can be understood
        as either the Itô integral or the Stratonovich integral. We refer to
        \cite{cuchiero2023_smtc} for the definition of the signature for (paths
        of) continuous semimartingales in the context of finance and to
        \cite{cuchiero2025_uatcfcplsm, friz2017_grilrplf} for the more general
        setting of c\`adl\`ag paths.
        
        \item For the truncated signature up to order $N$ seen as vector 
        in $\SetR^{d_N}$, its coordinate iterated integrals are ordered according to ordering of multiindices,
        \[
            ((1),(2),\dots,(d),
            (1,1),(1,2),\dots,(1,d),(2,1),\dots,(d,d),
            (1,1,1),\dots, 
            (\overbrace{d,\dots,d}^{N-\text{times}})).
        \]
        \item Note the difference between $S^n_{s,t}(X), 
            S^{(N)}_{s,t}(X)$ and $S^I_{s,t}(X)$ for $n,N\in\SetN$ and a 
            mutliindex $I$. The first denotes the $n$-th component of the 
            signature of $X$ as element in $(\SetR^d)^{\otimes n}$, the
            second one is the truncated signature up to order $N$ in $T^{(N)}
            (\SetR^d)$, while the third is the real-valued coordinate
            iterated integral for multiindex $I$.
    \end{enumerate}
\end{remark}

\begin{definition}
\label{def:TimeAugmentedPath}
    For $X\in\mathcal{V}_1([0,T],\SetR^d)$ we define its corresponding
    \textit{time augmented} path as $\widehat{X}_t := (t, X_t)$, which is a
    $\SetR^{d+1}$-valued path of finite variation. The $0$-th component 
    corresponds to the time component, that is $\widehat X^{(0)}_t = t$.
\end{definition}

\begin{remark}
    One may alternatively define a time-augmented path $\widehat{X}_t$ with
    the first component given as the normalized time $t/T$, the remaining time
    $T-t$, as well as its normalized version $(T-t)/T$. The following
    \cref{lem:Uniqueness} regarding the uniqueness of the signature still
    holds true for the time-augmentation given in any of those ways.
\end{remark}

\begin{example}
\label{exa:SigExamples}
    \begin{enumerate}
        \item \label{exa:SigExamples:item:TimeAugmentedPath}
        Consider $\widehat{X}_t = (t,X_t)$ for some $X \in
        \mathcal{V}_1([0,T],\SetR)$. Its signature over $[s,t] \subset [0,T]$
        is given as
        \begin{multline*}
            S_{s,t}(\widehat{X})
            = \Bigl(1, t-s, X_t-X_s, 
                \frac{1}{2}(t-s)^2, \int_s^t u \dd{X_u}, \int_s^t X_u \dd{u},
                    \frac{1}{2}(X_t - X_s)^2, \\
                \frac{1}{6}(t-s)^3, \frac{1}{2} \int_s^t (u_3-s)^2 
                    \dd{ X_{u_3}}, \dots \Bigr).   
        \end{multline*}
        In general, for any path $X \in \mathcal{V}_1([0,T],\SetR)$ and 
        $I=(1,1,\dots,1)$  with $\abs{I}=n$ we have
        \[
            S^I_{s,t}(X) = \frac{(X_t-X_s)^n}{n!}.
        \]
        
        \item \label{exa:SigExamples:item:LinearPath}
        Let $X:[0,T]\to\SetR^d$ be a linear path, that is 
        $X_t=X_0+t\frac{X_T-X_0}{T}$ for $X_0,X_T \in \SetR^d$. Since we
        have
        \[
            \dd{X^{(i_j)}_u} = \frac{X_T^{(i_j)}-X_0^{(i_j)}}{T} \dd{u}
            \quad \text{and} \quad
            \int_{0<u_1<\dots<u_n<t} \dd{u_1} \dots \dd{u_n} = \frac{T^n}{n!},
        \]
        for $i_j \in \{1,\dots,d\}$, the coordinate iterated integral of $X$
        for multiindex $I=(i_1,\dots,i_n)$ is
            \begin{equation}      
                S^I_T(X) 
                = \frac{1}{n!}\prod_{j=1}^n X_T^{(i_j)}-X_0^{(i_j)}.
            \end{equation}
        With the map defined as $\SetR^d \ni a \mapsto \exp(a) := 
        \sum_{n\geq0} \frac{1}{n!} a^{\otimes n} \in T((\SetR^d))$, with 
        $a^{\times 0} := 1$, the signature of $X$ takes the simple form 
        \begin{equation}
            S_T(X) = \exp(X_T - X_0).
        \end{equation}
    \end{enumerate}
\end{example}

The signature exhibits various key properties that make it useful in the
reinforcement learning setting considered in this thesis. The first property, 
which is due to \cite{chen1957_ipgigbhf, chen1977_ipi}, states that 
concatenation of paths corresponds to tensor multiplication of the paths'
signatures. 

\begin{definition}
    For $0\leq u \leq T$ let $X:[0,u]\to\SetR^d$ and $Y:[u,T]\to\SetR^d$ be
    two continuous paths. Their concatenation is the continuous path
    $X\star Y:[0,T]\to\SetR^d$ defined as 
    \begin{equation}
        (X\star Y)_t :=
        \begin{cases}
            X_t, & \text{if } t\in [0,u), \\
            X_u + Y_t - Y_u, & \text{if } t\in [u,T].
        \end{cases}
    \end{equation}
\end{definition}

\begin{lemma}[Chen's identity]
\label{lem:ChensIdendity}
    For $u\in [0,T]$, and two paths $X \in \mathcal{V}_1([0,u],\SetR^d)$ and 
    $Y \in \mathcal{V}_1([u,T],\SetR^d)$, it holds that
    \begin{equation}
    \label{eq:ChensIdendity}
        S_T(X\star Y) = S_{0,u}(X) \otimes S_{u,T}(Y).
    \end{equation}
\end{lemma}
For the proof we refer to \cite{lyons2007_dedrp}. 
\Cref{lem:ChensIdendity} has some implications. First, it allows for easy 
updating of the signature over an extended interval whenever new information 
about a path arrives, a property that is crucial in a reinforcement learning 
setting. Second, a path's signature over the whole interval can be pieced 
together by considering the path on subintervals, which indeed is the basis 
for explicitly computing the signature of a stream of data, see 
\cref{chapter:Signature:section:StreamComputation}. Furthermore, Chen's 
identity states that the range of the map $X\mapsto S_T(X)$ is closed under 
multiplication, and since $S_T(X)_0 = 1$, it turns out that it is even a group.

\begin{proposition}[Proposition 2.14 in \cite{lyons2007_dedrp}]
\label{prop:InverseSignature}
    For $X\in\mathcal{V}_1([0,T],\SetR^d)$, let $\overleftarrow{X}$ be the
    path run backwards, that is $\overleftarrow{X}_t:=X_{T-t}$ for
    $t\in[0,T]$. Then, for any $[s,t]\subset [0,T]$ it holds that $S_{s,t}
    (\overleftarrow{X}) = S_{s,t}(X)^{-1}$ and the range of the map $X\mapsto 
    S_{s,t}(X)$ with the multiplication $\otimes$ on $T((\SetR^d))$ is a 
    subgroup in $T((\SetR^d))$.
\end{proposition}

\Cref{prop:InverseSignature} together with Chen's identity provides a way to 
obtain the signature of a shortened path by tensor multiplication with the
inverse of the signature of the path over a subinterval. More precisely for 
$0 \leq t \leq T$ we have
\[
    S_{t,T}(X) = S_{0,t}(X)^{-1} \otimes S_{0,T}(X).
\]

The next property essentially states that the signature captures the order of
events along a path, while ignoring the ``speed'' at which those events take 
place.

\begin{proposition}[Invariance under time reparametrization]
\label{prop:ReparametrizationInvariance}
    Let $X\in\mathcal{V}_1([0,T],\SetR^d)$ and $\varphi:[0,T]\to[0,T]$ be 
    non-decreasing and surjective. Then, $S_T(X) = S_T(X\circ \varphi)$.
\end{proposition}
This property can be an advantage in handling paths as input data in 
applications since identifying paths with different parametrization 
constitutes a substantial dimensionality reduction. If, however, the time 
parametrization is in fact important, as it is often the case with financial 
data, one can simply consider the time augmented path $\widehat{X}$ from 
\cref{def:TimeAugmentedPath}. Adding a time component is also a sufficient 
condition for the uniqueness of the signature.

\begin{lemma}[Uniqueness]
\label{lem:Uniqueness}
    Let $X, Y \in \mathcal{V}^0_1([0,T],\SetR^d)$ and consider their time 
    augmented paths $\widehat{X}_t$ and $\widehat{Y}_t$. Then, 
    $S_T(\widehat{X}) = S_T(\widehat{Y})$ if, and only if $X_t=Y_t$ for each 
    $t\in[0,T]$.
\end{lemma}

In general, the signature of $X\in\mathcal{V}_1([0,T],\SetR^d)$ and of the 
translated path $X+c_1$ for some $c_1\in\SetR^d$, agree, and the same holds 
true for $\widehat{X}+c_2$ and $\widehat X$ for any $c_2\in\SetR^{d+1}$, 
however one can always consider $X_t-X_0$. In \cite{hambly2010_uspbvrpg}, the 
authors showed that the signature $S_T(X)$ of $X$ without time-component, 
determines $X$ up to so called tree-like equivalences.%
\footnote{
    One can think of a tree-like path as one that retraces itself. For the 
    exact definition in the bounded variation case see \cite[Definition
    1.2]{hambly2010_uspbvrpg}.
} 
\Cref{lem:Uniqueness} is a direct implication of this result, since adding 
time as strictly monotone component eliminates any tree-like sections of a 
path. The result was later extended to the more general rough path setting in
\cite{boedihardjo2016_srpu}.

The next property is the main ingredient for the signature's application in 
the approximation procedure in our reinforcement learning setting. It states 
that continuous functionals on paths can be approximated arbitrarily well by 
linear functionals on the paths' signature, uniformly in time. For a path 
$X\in\mathcal{V}_1([0,T],\SetR^d)$ we set $X_{[0,t]} := (X_s)_{s\in[0,t]}$ 
for its trajectory over the interval $[0,t]$, and analogue for its 
time-augmentation $\widehat{X}_{[0,t]} := (s,X_s)_{s\in[0,t]}$. Define the 
space of all time-augmented paths starting at the origin until $t\in[0,T]$ as
\[
    \Omega^0_t 
    := \{\widehat{X}_{[0,t]} 
        \mid X \in\mathcal{V}^0_1([0,T],\SetR^d)\}.
\]
which is a metric space equipped with the 1-variation distance $d_{1,[0,t]}$ .
Given $\widehat{X}\in\Omega^0_t$, we define the extended path $\widehat{X}^t$ 
on $[0,T]$ as
\begin{equation}
\label{eq:ExtendedPath}
    \widehat{X}^t_s := \begin{cases}
        \widehat{X}_s, & \text{if } s\in[0,t]\\
        (s, X_t), & \text{if } s\in(t,T].
    \end{cases}    
\end{equation}

\begin{definition}
\label{def:StoppedPathSpace}
    We call $\Lambda_T := \bigcup_{t=0}^T \Omega^0_t$ the \textit{space of 
    stopped paths} and equip it with the metric
    \begin{equation}
    \label{eq:StoppedPathsMetric}
        d_{\Lambda_T}(\widehat{X}_{[0,t]}, \widehat{Y}_{[0,s]}) 
        := d_{1,[0,t]}(\widehat{X}_{[0,t]}, \widehat{Y}^s_{[0,t]}) + \abs{t-s},
    \end{equation}
    for any $\widehat{X}, \widehat{Y} \in \Lambda_T$, where we assume w.l.o.g. 
    $s\leq t$, and $d_{1,[0,t]}$ is the 1-variation distance 
    \eqref{eq:1-VariationDistance} $[0,t]$, and $\widehat{Y}^s$ 
    is the extended path defined in \eqref{eq:ExtendedPath}.
\end{definition}

\begin{theorem}[Universal approximation] 
\label{thm:UniversalApproximation}
    Let $K\subset \Omega^0_T$ be compact and $F:\Lambda_T \to \SetR$ be a
    continuous path functional. Then, for any $\epsilon>0$ there exists a
    linear functional $\ell\in T(\SetR^{d+1})$
    such that
    \begin{equation}
    \label{eq:UniversalNonLinearity}
        \sup_{\widehat X \in K} \sup_{t\in[0,T]} \abs{
            F\bigl(\widehat X_{[0,t]}\bigr) 
            - \bigl\langle \ell, S_t(\widehat{X}) \bigr\rangle
        } < \epsilon.
    \end{equation}
\end{theorem}

The theorem is proved in a more general form in \cite{kalsi2020_oewrps} and 
\cite{bayer2021_osws}. The case of continuous paths of bounded
1-variation follows as special case.
    
%------------------------------------------------------------------------------
\section{The Signature of a Stream of Data and its Computation} 
\label{chapter:Signature:section:StreamComputation}
%------------------------------------------------------------------------------

So far we have been considering the signature of continuous paths on some
interval $[0,T]$ for $T\in\SetR$. In practice, however, one usually is given a
set of discrete data points. By embedding such a stream into a linear path, we
can define the signature of a stream of discrete data, and we illustrate
how it can be computed explicitly.

Following \cite{bonnier2019_dst}, we define the set of streams of data in 
$\SetR^d$ as
\[
    \mathcal{S}(\SetR^d) 
    := \{\mathrm x= (x_0, x_1, \dots, x_n) \mid x_i \in \SetR^d, n\in\SetN\}.
\]
For $\mathrm x = (x_0,\dots,x_n)\in\mathcal{S}(\SetR^d)$, its length is given 
by $\abs{\mathrm x} = n+1$. We assume that there is a natural time series 
structure and the elements of $\mathcal{S}(\SetR^d)$ consist of sequential 
data points. If the data does not represent a time series, the methods in this 
section still apply, but one first has to decide on the order in 
which the points are visited. We embed a stream $\mathrm x$ into a continuous 
path with values in $\SetR^d$ by linear interpolation. For a stream $\mathrm x 
\in \mathcal{S}(\SetR^d)$ with $\abs{\mathrm x}=n+1$, fix a time $T>0$, and 
choose $t_0, t_1, \dots, t_n \in[0,1]$ such that
\[
    0=t_0 < t_1 < \dots < t_n = 1.
\]
If the stream's indices already correspond to time points the canonical choice 
would be $T=n$. We define the path $X:[0,T]\to\SetR^d$ as the linear 
interpolation of 
$\mathrm x$,
\[
    X_t = \begin{cases}
        x_i, & \text{if } t\in\{t_0, \dots, t_n\}, \\
        x_{i} + (t-t_i) \displaystyle \frac{x_{i+1} - x_i}{t_{i+1}-t_i}, &
            \text{if } t\in (t_i, t_{i+1}) \text{ for } i\in \{0,\dots,n-1\}.
    \end{cases}
\]

\begin{definition}
\label{def:SignatureOfStream}
    Let $\mathrm x\in\mathcal{S}(\SetR^d)$ be a $d$-dimensional stream of data 
    with length $\abs{\mathrm x}=n+1$, and let $X:[0,T]\to\SetR$ be its linear 
    interpolation. For $i,j\in\{0, 1, \dots, n\}$ with $i\leq j$, we define 
    the signature of $\mathrm x$ from $i$ to $j$ as
    \begin{equation}
    \label{eq:SignatureStream}
        S_{i,j}(\mathrm x) := S_{t_i,t_j}(X).
    \end{equation}
    The truncated signature up to order $N\in\SetN$ of $\mathrm x$ is defined
    analogously via $X$. If $i=0$ we write $S_j(\mathrm x) = 
    S_{0,j}(\mathrm x)$ for any $j\in\{0,\dots,n\}$.
\end{definition}

\Cref{prop:ReparametrizationInvariance} ensures that $S_{i,j}(\mathrm x)$
is well defined, as it does not depend on the specific choice of the partition
points $t_i$. 
\Cref{exa:SigExamples}\ref{exa:SigExamples:item:TimeAugmentedPath} 
illustrates that the signature of a linear path takes a simple form as the 
exponential in $T((\SetR^d))$ of the total increment of the path. The 
continuous interpolation $X$ of some $\mathrm x \in 
\mathcal{S}(\SetR^d)$ is just a concatenation of linear paths. Together with 
Chen's identity \eqref{eq:ChensIdendity} we are provided with a way of 
explicitly computing $S_{t_i,t_j}(X)$.

\begin{proposition}
\label{prop:SignatureComputation}
    Let $\mathrm x\in\mathcal{S}(\SetR^d)$ with length $\abs{\mathrm x}
    =n+1$. Then, for any $i,j\in\{0,1,\dots,n\}$ with $i<j$,  
    \[
        S_{i,j}(\mathrm x) = S_{t_i,t_j}(X) 
        = \exp(x_{i+1}-x_i)\otimes\exp(x_{i+2}-x_{i+1})\otimes\dots
            \otimes\exp(x_j-x_{j-1}), 
    \]
    where for $a\in\SetR$,
    \[
        \exp(a):=\sum_{n\geq0} \frac{1}{n!} a^{\otimes n} \in T((\SetR^d)).
    \]
\end{proposition}
\begin{proof}
    W.l.o.g. fix $i=0$ and $j=n$. For any two consecutive points 
    $x_k,x_{k+1}$, the path segment $X_{[t_k,t_{k+1}]}$ is 
    linear, thus $S_{t_k,t_{k+1}}(X) = \exp(x_{k+1}-x_{k})$. 
    Since $X_{[0,T]}$ is the concatenation of all linear segments 
    $\bm{\mathrm x}_{[t_k,t_{k+1}]}$ for $k=0,\dots,n-1$, iterative 
    application of Chen's identity yields the result.
\end{proof}
\begin{remark}
    \cref{prop:SignatureComputation} illustrates that computing the signature
    for an incoming stream of data can be done in an efficient way. Assume
    that for $\mathrm{x}\in\mathcal{S}(\SetR^d)$ with length $\abs{\mathrm x}
    =n+1$ we have computed $S_{n}(\mathrm x)$, and additional data $\mathrm x' 
    \in\SetR^d$ with $\abs{\mathrm x'}=m+1$ arrives. Further, assume we have 
    appended $x_n$ to the beginning of $\mathrm x'$ such that $x_n = x'_0$. 
    Instead of concatenating $\mathrm x$ and $\mathrm x'$, and computing the 
    signature of the concatenation containing $n+m+1$ elements, one only needs 
    to compute $S_m(\mathrm x')$ and tensor multiply it with $S_n(\mathrm x)$ 
    to obtain the signature of the full stream. This property 
    proves to be of particular advantage in a reinforcement learning setting 
    since there, new information arrives step-wise at each transition. Its 
    applicability will be seen in \cref{alg:SignatureQLearning}.
\end{remark}

There are multiple Python libraries that provide functionalities for computing 
signatures such as the \texttt{iisignature}
\cite{reizenstein2020_a1ileciss}, \texttt{esig} \cite{lyons2002_esig}, which 
was originally developed by Terry Lyons et al., as well as the
\texttt{Signatory} \cite{kidger2021_sdcsltb}, with the latter used in this 
work. Further computational details are provided for the experiments in 
\cref{chapter:Experiments:section:MountainCar:sub:TrainingDetails,%
chapter:Experiments:section:Execution:sub:TrainingDetails}.

%%%%%%%%%%%%%%%%%%%%%%%%%%%%%%%%%%%%%%%%%%%%%%%%%%%%%%%%%%%%%%%%%%%%%%%%%%%%%%%
\chapter{Reinforcement Learning and Markov Decision Processes} 
\label{chapter:RLAndMDP}
%%%%%%%%%%%%%%%%%%%%%%%%%%%%%%%%%%%%%%%%%%%%%%%%%%%%%%%%%%%%%%%%%%%%%%%%%%%%%%%

%------------------------------------------------------------------------------
\section{Reinforcement Learning}
\label{chapter:RLAndMDP:section:RL}
%------------------------------------------------------------------------------

Reinforcement learning \cite{sutton2018_rli} is a machine learning paradigm 
where an algorithm, commonly referred to as an `agent', learns to make a 
series of decisions by interacting with an evolving environment and adjusting 
its behaviour based on numeric feedback signals in the form of rewards or 
penalties. The goal is to maximize its cumulative performance over time by 
learning which actions yield the highest subsequent rewards, given a 
particular state or history of states the underlying environment was in. 
Environments can take various forms, ranging from the physical world when 
training robots to simulated environments like game simulators or replayed 
sequences of real-life events such as financial time series. 
%
A typical reinforcement learning setting is depicted in \cref{fig:RLNutshell}. 
At step $t$, the agent receives the state from the environment and a
reward associated with the transition from the previous step's state. At step 
$t$, the agent observes that the environment is in some state $S_t$. Based on 
this information and, potentially, information from previous steps, the agent 
calculates an action $A_t$, which it sends to the environment. Subsequently, 
the environment transitions to a new state $S_{t+1}$ and emits a reward 
$R_{t+1}$ associated with this transition. Both are sent back to the agent, who
internally adapts its action selection mechanism based on the reward received, 
and the cycle repeats.
%
\begin{figure}[htpb!]
    \centering
    \includegraphics[width=0.6\textwidth]{graphics/RL-in-a-nutshell.png}
    \caption[Reinforcement learning in a nutshell]{%
        Reinforcement learning in a nutshell. \cite[Figure~3.1]{sutton2018_rli}
    }
    \label{fig:RLNutshell}
\end{figure}
%
This concept derives from how humans learn by trial and error. As we try out 
new behaviour and, in this way, interact with our environment, we 
receive feedback in various forms, such as gaining points in a game or
experiencing physical sensations like pain or pleasure. Based on this feedback,
we adjust our behaviour accordingly to achieve the desired outcome in the
future.

Reinforcement learning differs from supervised and unsupervised learning, 
constituting a third machine learning paradigm \cite[Section~1.1]%
{sutton2018_rli}. In supervised learning, a model is trained on a labelled
data set of correct input-output examples, with the aim of generalizing the 
model's responses to input not included in the training set. While we can 
retrain a model on new data, it is usually impossible to obtain correct 
examples in interactive, evolving settings, and thus, in reinforcement 
learning, an agent must learn without relying on some ground truth provided by 
an external supervisor. Although reinforcement learning does not rely on 
labelled data, it differs from the unsupervised learning paradigm in terms of 
its objective. While the latter aims to find structure in unlabelled data, 
reinforcement learning agents try to maximize a reward signal in a dynamic 
environment.

The usual formal model of the agent-environment framework depicted in
\cref{fig:RLNutshell} is that of a \textit{Markov Decision Process},  
presented in the following \cref{chapter:RLAndMDP:section:MDP}. However, this 
work introduces a \textit{history-dependent decision process} as an 
alternative model in \cref{chapter:HistoryRLEnvironment:section:HDP}. Aside 
from the agent-environment model, further essential aspects of reinforcement 
learning are the representation of states, the value function, the policy, the 
reward design, and the trade-off between exploration and exploitation.

\textit{State representation} is crucial in reinforcement learning. States 
encapsulate any information available to the agent about the environment at 
certain times and serve as input to the agent's decision-making rule or 
policy. From a mathematical point of view, the state is often assumed to be 
the fully observable true state of the environment, as is the case in the 
Markov Decision Process framework in \cref{chapter:RLAndMDP:section:MDP}. 
Since, in practice, this is mostly not achievable, one can drop this 
assumption. Letting the agent only have access to a partial observation of the 
true environment state leads to a partially observable Markov decision 
process; see also \cref{chapter:RLAndMDP:section:IssueWithMDP}. In the 
history-dependent framework proposed in 
\cref{chapter:HistoryRLEnvironment:section:HDP}, we work with observations, 
not distinguishing whether these represent fully or partially observable 
states.
 
The \textit{reward signal} defines the objective in a reinforcement learning 
problem. At each step, the environment provides the agent with a single 
numeric value as a reward. By defining which state-action-state combinations 
are favourable or unfavourable, the reward signal serves as the primary basis 
for altering the agent's behaviour and guiding it towards maximizing the 
cumulative reward received. The function that produces the signal is generally 
stochastic and unknown to the agent and is usually part of the problem 
formulation. Thus, the agent's ability to successfully solve a problem hinges 
on the reward signal it provides.

While the reward signal indicates what is good in an immediate sense, the
\textit{value function} determines what is good in the long run. The value of 
a state represents its long-term desirability and is given as the expected 
total reward achievable in the future, taking into account states and actions 
that are likely to follow in the future. Decisions are made based on value 
estimates from observed states and actions chosen in those states. The 
agent chooses the actions it expects to lead to states with the 
highest value. Thus, efficiently estimating values is one of the most 
critical tasks in reinforcement learning. Value functions originally stem from 
the problem of optimal control of a dynamically evolving (stochastic) system. 
Its solution, using value functions and dynamic programming principles such as 
Bellman equations, serves as the foundation for many methods in reinforcement 
learning.

A \textit{policy} defines the agent's decision-making rules and is a mapping 
from perceived states or sequences of states to actions. It is at the core of 
each reinforcement learning method and fully determines an agent's behaviour. 
A policy can be stochastic in that it outputs a probability distribution over 
actions. Policies can take on various forms. In value-based methods, the agent 
learns a value function from which a policy is derived by selecting actions 
that lead to states with the highest value. Such methods include the famous 
Q-learning \cite{watkins1992_q} or SARSA \cite[Section~6.4]{sutton2018_rli}. 
In policy-based methods, the agent directly learns a parametrized mapping from 
states to actions instead of a value function. Examples include REINFORCE 
\cite{williams1992_ssgacrl}, Trust Region Policy Optimization 
\cite{schulman2015_trpo}, and state-of-the-art Proximal Policy Optimization 
\cite{schulman2017_ppoa}. Actor-critic methods are a mixture of learning  
a value function and a policy. The value function serves as a critic and 
provides a gradient direction for the policy update, while the policy as actor 
decides which action to choose next \cite{konda1999_aa}. In this work, we 
propose a value-based method, which is an adaptation of the Q-learning 
algorithm for history-dependent environments.

Another key aspect and inherent challenge in reinforcement learning, distinct
from other learning approaches, is navigating the \textit{trade-off between
exploration and exploitation}. To maximize rewards, an agent must prioritize
actions previously found to be effective. However, discovering such actions
requires exploring alternatives that have not yet been selected. Thus, the 
agent must balance exploiting known strategies for immediate reward and
exploring new possibilities to improve future decision-making. It may achieve 
this by following a strategy of gradually favouring actions known to be good
ones while maintaining a level of exploration. Especially in stochastic
environments, repeated trials are essential to estimate the expected rewards 
of different actions accurately. The exploration-exploitation trade-off 
remains a complex issue, not present in other learning paradigms.

Reinforcement learning has diverse applications across various domains, 
ranging from robotics and gaming to finance and healthcare. Classic 
reinforcement learning methods have been instrumental in solving problems such 
as game playing, automated control systems, and scheduling tasks. In recent 
years, integrating reinforcement learning with deep learning has led to 
significant advancements in domains like computer vision, natural language 
processing, and autonomous vehicles. Deep reinforcement learning techniques, 
leveraging the power of neural networks, have demonstrated remarkable success 
in complex tasks such as image recognition, language understanding, and 
sequential decision-making. This connection of reinforcement learning with 
deep learning has opened up new possibilities for solving intricate problems 
that were previously challenging to address using traditional methods alone.

%From training agents to play video games at superhuman levels to optimizing resource allocation in dynamic environments, the combined approach of reinforcement learning and deep learning continues to expand the boundaries of what is achievable in artificial intelligence.


%------------------------------------------------------------------------------
\section{Markov Decision Processes}
\label{chapter:RLAndMDP:section:MDP}
%------------------------------------------------------------------------------

The usual mathematical framework for modeling a reinforcement learning problem
is that of a Markov Decision Process (MDP). We give an introduction to this
framework, characterize optimal policies, and illustrate the Q-learning method 
in MDPs. In \cref{chapter:HistoryRLEnvironment}, we will build on this 
exposition, generalize it to a history-dependent decision process, and adapt 
the Q-learning method to this history-dependent framework. Various modeling
assumptions can be taken when defining an MDP. In view of the Q-learning
method illustrated in the following section, we introduce an MDP with 
finite-time horizon, finite action space, and Borel state space. For a 
thorough introduction to MDPs we refer to the monograph 
\cite{puterman2005_mdpdsdp} for finite and countable state and action spaces, 
and to \cite{hinderer2016_do} for the case where the involved spaces are 
allowed to be more general measurable spaces. If not mentioned otherwise, this 
exposition builds on Sections 2 and 4 in \cite{puterman2005_mdpdsdp}, as well 
as \cite{szepesvari2010_arl}, and Chapter 16 in \cite{hinderer2016_do}, where 
the latter reference guarantees the existence of an optimal deterministic 
Markovian policy in the considered setup with general Borel state space.

A discrete-time, finite time-horizon, Markov Decision Process is given as the tuple
\begin{equation}
\label{eq:MDPTuple}
    (\SpaceS, \SpaceA, \KernelP, \KernelR, \gamma, T),
\end{equation}
where its components are defined as follows. 
% state space
The \textit{state space} $\SpaceS$ is the set of possible environment states. 
It is assumed to be a separable complete metric space $(\SpaceS, d_\SpaceS)$, 
endowed with its Borel $\sigma$-algebra $\mathcal{B}_{\SpaceS}$.
% action space
The \textit{action space} $\SpaceA$ is a finite, non-empty set of actions with 
cardinality $\abs{\SpaceA}\in\SetN$.\footnote{
    More generally, the action space too can be assumed to be a separable,
    complete metric space without any restriction on cardinality, cf.
    \cite{hinderer2016_do}. However, dealing with general, continuous action
    spaces is more involved and in view of the computational method proposed in
    this work, not necessary.
    %More generally the state space and action space can both be assumed  to be Polish (topological) spaces or Borel spaces. While in the case of the state space this would not change much in this works exposition, dealing with a general, continuous action space is more involved and in view of the computational method proposed in this work, not necessary.
} 
We may assume without loss of generality $\SpaceA = \{1,2,\dots,
\abs{\SpaceA}\}$, endow $\SpaceA$ with the discrete metric, and denote 
the $\sigma$-algebra on $\SpaceA$ as $\mathcal{B}_\SpaceA=2^{\SpaceA}$. If 
both $\SpaceS$ and $\SpaceA$ were finite, the tuple \eqref{eq:MDPTuple} is 
referred to as \textit{finite} Markov Decision Process.
% time horizon
Moreover, 
$T\in\SetN$ is the finite time-horizon at which the process terminates, if no other 
termination criterion was met before.
% transition kernels
The collection of \textit{state transition
kernels} is denoted by 
$\KernelP = (P_t)_{t\in[T-1]}$. 
For each $t\in[T-1]$, the state transition kernel $P_t$ is a
stochastic kernel from $\SpaceS \times \SpaceA$ to $\SpaceS$, such that for 
any state-action pair $(s,a) \in \SpaceS \times \SpaceA$ and measurable set
$S \in \mathcal{B}_\SpaceS$, $P_t(S\,|\,s,a)$ returns the probability that the
next state is in $S$, provided the current state is $s$ and action $a$ is 
selected (more on this in \cref{eq:MDPMarkovNature} below). The probability 
measure on $\SpaceS$ corresponding to the pair $(s,a)$ is denoted $P_t(\,
\cdot\,|s,a)$.
% reward functions
Further, $\KernelR = (R_t)_{t\in[T]}$ is a collection of real valued
\textit{reward functions}. For each $t\in[T-1]$ the reward function
$R_t:\SpaceS \times \SpaceA \to \SetR$ is measurable with respect to the
product $\sigma$-algebra $\mathcal{B}_\SpaceS\otimes\mathcal{B}_\SpaceA$, and 
for $t=T$, the \textit{terminal reward function} $R_T:\SpaceS \to \SetR$ is 
$\mathcal{B}_\SpaceS$-measurable. We assume that rewards are bounded by some  
$0<C<\infty$, such that $\abs{R_t(s,a)} \leq C$ for each $t\in[T-1]$ and any 
pair $(s,a)\in\SpaceS\times\SpaceA$, and $\abs{R_T(s)} \leq C$ for any 
$s\in\SpaceS$.
% discount factor
Finally, $\gamma\in[0,1]$ is a discount factor, which weights future rewards.

An agent interacts with a given MDP in consecutive episodes at discrete steps. 
At the beginning of episode $k\in\SetN$, the process is in some state 
$s_0\in\SpaceS$ sampled from a fixed initial distribution $\mu \in 
\text{Prob}(\SpaceS)$. Then, at each step $t\in[T-1]$ the agent observes the 
process' state $s_t$, and chooses an action $a_t\in\SpaceA$ based on some 
decision rule at $t$. As a result, the it receives a reward $R_t(s_t, a_t)$ 
and the process transitions to a new state $s_{t+1}$, which is sampled from 
the probability distribution $P_t(\cdot\,|\,s_t,a_t)$ on $\SpaceS$. The last
action is taken at step $T-1$ after which the process reaches a terminal
state $s_T$, and the agent receives the additional terminal reward $R_T(s_T)$
before the cycle starts anew in episode $k+1$. 
% state-action history
This behaviour produces a state-action history\footnote{
    More specifically it produces an state-action-reward history, however we
    exclude the reward in \eqref{eq:MDPHistory} since in practice the reward
    part $r_{t+1}$ can usually be derived from the transition tuple 
    $(s_t,a_t,s_{t+1})$.
}
% a history in MDP
\begin{equation}
    \label{eq:MDPHistory}
    h_T = (s_0, a_0, s_1, a_1, \dots, s_{T-1}, a_{T-1}, s_T) 
        \in (\SpaceS \times \SpaceA)^{T} \times \SpaceS.
\end{equation}

% history spaces in MDP
\begin{definition}
    For a given MDP \eqref{eq:MDPTuple} and any $t\in[T]$, the space of
    histories up to time $t$ is defined as 
    \[
        \SpaceH_t := (\SpaceS \times \SpaceA)^{t} \times \SpaceS, 
    \]
    with the convention $\SpaceH_0=\SpaceS$ and elements $h_t = (s_0,a_0, s_1, a_1, \dots, s_{t-1}, a_{t-1}, s_t) \in \SpaceH_t$. Each $\SpaceH_t$ is equipped with the 
    Borel $\sigma$-algebra\footnote{
        Since $\SpaceS$ is a separable metric space, the Borel $\sigma$-algebra on $\SpaceH_t$ induced by the product topology on $\SpaceH_t$ coincides with the product $\sigma$-algebra of the $\sigma$-algebras $\mathcal{B}_\SpaceS$ and $\mathcal{B}_\SpaceA$. 
    } 
    $
    \mathcal{B}_{\SpaceH_t} 
        = \bigotimes_{i=0}^t(
            \mathcal{B}_\SpaceS \otimes \mathcal{B}_\SpaceA 
        ) \otimes \mathcal{B}_\SpaceS.
    $
\end{definition}

% a remark on rewards
\begin{remark}
\label{rem:MDPRewardsRemark}
    The rewards may also depend on the environment's state at the next
    step, such that they are determined by measurable reward functions 
    $\tilde R_t:\SpaceS\times\SpaceA\times\SpaceS \to \SetR$ for $t\in[T-1]$, 
    where $\tilde R_t(s,a,s')$ returns the reward received when the state of 
    environment at time $t$ is $s$, action $a$ is chosen and the environment 
    transitions to state $s'$ at time $t+1$. In this case, we define the 
    \textit{expected immediate reward function} at $t\in[T-1]$ by
    \begin{equation}
    \label{eq:MDPExpectedImmediateReward}
        R_t(s,a) 
        %:= \E_{P_t(\dd{s'}|s,a,)}[R_t(s,a,S')]
        := \int_\SpaceS \tilde R_t(s,a,s') P_t(\dd{s'}|s,a),
        \quad \text{for all }(s,a)\in \SpaceS\times\SpaceA.
    \end{equation}
    Since $\tilde R_t$ is measurable, so is $R_t$ for any $t\in[T-1]$ 
    due to \cref{lem:KernelsAuxiliaryMeasurableExpectation}, and the 
    subsequent results in this section remain true if $R_t$ is defined by 
    \eqref{eq:MDPExpectedImmediateReward}. From a reinforcement learning 
    perspective, it is irrelevant if the reward sample $r_t\in\SetR$ that the 
    agent receives at time $t$, depends only on $(s_t,a_t)$ or, additionally 
    on the next state $s_{t+1}$, because we only require, the value $r_t$ to 
    be fixed before the agent selects the next action $a_{t+1}$, and that it 
    does not depend on its choice.
\end{remark}


We now formalize the concept of a policy which the agent uses to select
actions.
% policy in MDP
\begin{definition}[Policy in MDPs]
\label{def:MDPPolicy}
    For an MDP, a policy is defined as the family $\pi = (\pi_t)_{t\in[T-1]}$,
    where we call $\pi$ a 
    \begin{itemize}
        \item \textit{randomized policy}, if $\pi_t$ is a stochastic kernel
        from
        $(\SpaceS, \mathcal{B}_{\SpaceS})$ to $(\SpaceA, \mathcal{B}_\SpaceA)$
        for each $t\in[T-1]$,      
        \item \textit{deterministic policy}, if $\pi_t$ is a measurable map 
        $\pi_t:\SpaceS \to \SpaceA$ for each $t\in[T-1]$.
    \end{itemize}
    The sets of all randomized and deterministic policies are denoted as 
    $\Pi^{\text{R}}$ and $\Pi^{\text{D}}$, respectively. Any 
    $\pi\in\Pi^\text{D}$ can be considered as a special case of a randomized 
    policy by setting $\tilde{\pi}_t(\cdot|s) := \delta_{\pi_t(s)}(\cdot)$ 
    for each $s\in\SpaceS$, thus $\Pi^{\text{D}}\subset\Pi^{\text{R}}$, and 
    $\pi\in\Pi^{\text{R}}$ always includes the case $\pi\in\Pi^{\text{D}}$.
\end{definition}

Policies in \cref{def:MDPPolicy} are also called \textit{Markovian}
policies (either deterministic or randomized), because $\pi_t$ is a function
of the current state $s_t$ only, and we will see that those are sufficient for 
optimality in an MDP. In general, policies may be history-dependent (compare  
\cref{def:HDPPolicy} below).

A fixed policy $\pi\subset \Pi^\text{R}$ together with an initial distribution
$\mu\in\text{Prob}(\SpaceS)$ induces the canonical probability space
\begin{equation}
\label{eq:MDPCanonicalProbSpace}
    (\SpaceH_T, \mathcal{B}_{\SpaceH_T}, \PM^\pi),
\end{equation}
where the probability measure $\PM^\pi$ is defined as in
\cref{chapter:PreliminariesStochasticKernels}, by the product of the initial 
distribution $\mu$, and the stochastic kernels $\pi = (\pi_t)_{t\in[T-1]}$  and 
$\mathcal{P} = (P_t)_{t\in[T-1]}$, such that it satisfies
\begin{equation}
\label{eq:MDPCanonicalMeasure}
\begin{split}
    \PM^\pi [\dd{h_T}] 
    = &\ P_{T-1}(\dd s_T | s_{T-1}, a_{T-1})\, \pi_{T-1}(\dd a_{T-1}|s_{T-1})\,
    P_{T-2}(\dd s_{T-1} | s_{T-2}, a_{T-2}) \\
    %\pi_{T-2}(\dd a_{T-2}|h_{T-2})
    &\dots P_{0}(\dd s_1 | s_0, a_0)\, \pi_0(\dd a_0|s_0)\, \mu(\dd s_0),
\end{split}    
\end{equation}
where $h_T = (s_0, a_0, s_1,\dots,s_T)\in\SpaceH_T$. Note, that the notation 
$\PM^\pi$ omits the dependence on $\mu$ and $\mathcal{P}$, since 
both are fixed, to highlight, that the measure $\PM^\pi$ is determined by the 
choice of $\pi$.
% 
For any history $h_T \in \SpaceH_T$, define the random states $S_t, t\in[T]$, 
and the random actions $A_t,t\in[T-1]$ as the coordinate maps on
\eqref{eq:MDPCanonicalProbSpace},
\[
    S_t(h_T) := s_t, 
    \quad \text{and} \quad
    A_t(h_T) := a_t.
\]
Note, that $S_t$ and $A_t$ are $\SpaceS$ and $\SpaceA$-valued random variables
with their conditional probability distributions determined by the kernels 
$P_t$ and $\pi_t$, respectively. The process $S_0, A_0, S_1, \dots, \\ A_{T-1}, S_T$ 
is just the canonical process on the probability space 
\eqref{eq:MDPCanonicalProbSpace}. Further, we define the history process by
\[
    H_0 := S_0, 
    \quad \text{and} \quad
    H_t := (S_0,A_0,S_1,\dots,S_t), \text{ for } t\in[T]\setminus\{0\},
\]
where again, $H_t$ is a random variable on \eqref{eq:MDPCanonicalProbSpace}
with values in $\SpaceH_t$. For any $t\in[T]$, history $h_t = (s_0, a_0, s_1,
\dots, s_t) \in\SpaceH_t$, Borel set $B\in\mathcal{B}_\SpaceS$, and action
$a_t\in\SpaceA$, the measure $\PM^\pi$ satisfies
\begin{gather}
    \PM^\pi[S_0 \in B] = \mu(B), \\
    \label{eq:MDPActionProb}
    \PM^\pi[A_t = a_t \,|\, H_t = h_t] 
        = \PM^\pi[A_t = a_t \,|\, S_t = s_t] 
        = \pi_t(a\,|\,s_t), \\
    \label{eq:MDPMarkovNature}
    \PM^\pi[S_{t+1}\in B \,|\, H_t=h_t, A_t=a_t] 
        = \PM^\pi[S_{t+1}\in B \,|\, S_t=s_t, A_t=a_t] 
        = P_t(B\,|\,s_t,a_t).
\end{gather}
\Cref{eq:MDPMarkovNature} is the Markov property, which states that the 
distribution of the next state $S_{t+1}$ under $\PM^\pi$ depends on the 
history past states and actions only through the current values $s_t$ and 
$a_t$.
% Markovian policy
If $\pi\in \Pi^\text{D}$, then for any $t\in[T-1]$, $A_t := \pi_t(S_t)$, the right-hand side in \cref{eq:MDPActionProb} is $\delta_{\pi_t(s_t)}(a_t)$, and \cref{eq:MDPCanonicalMeasure} changes to
\begin{equation*}
\begin{split}    
    \PM^\pi[\dd{h_T}] 
    = &\ P_{T-1}(\dd s_T | s_{T-1}, a_{T-1}) 
        \delta_{\pi_{T-1}(s_{T-1})}(\dd a_{T-1})
            P_{T-2}(\dd s_{T-1} | s_{T-2}, a_{T-2}) \\
    &\dots P_{0}(\dd s_1 | s_0, a_0) 
        \delta_{\pi_0(s_0)}(\dd a_0) 
            \mu(\dd s_0).
\end{split}
\end{equation*}
Additionally, the state process $(S_t)_{t\in[T]}$ is an $\SpaceS$-valued
Markov chain with initial distribution $\mu$ and transition kernels
$P_{\pi_t}(s, ds') := P_t(\dd s'|s,\pi_t(s))$.

\begin{remark}
    In the rest of \cref{chapter:RLAndMDP:section:MDP} and the following 
    \cref{chapter:RLAndMDP:section:QLearningInMDP} we only consider policies 
    $\pi\in\Pi^\text{D}$, because \cref{thm:MDPOptimalPolicyExistence} below 
    shows that the deterministic policies are sufficient to obtain optimality.
\end{remark}
%
Following a policy $\pi$, the agent receives a sequence of rewards, which, 
from it's point of view, is regarded as random. The agent's objective is to 
find a policy that maximizes the expected value of the resulting random reward 
sequence. Value functions allow to make this statement precise. 
% value function in MDP
Fix a policy $\pi$ and $t\in[T]$. The \textit{value function} at time $t$, 
corresponding to $\pi$ is defined as
\begin{equation}
\label{eq:MDPValueFunction}
    V^\pi_t(s) 
    := \E^\pi\Bigl[ 
        \sum_{u=t}^{T-1} \gamma^{u-t}R_u(S_u, A_u) + \gamma^{T-t} R_T(S_T)
        \,\Big|\, S_t = s
    \Bigr],
    \quad 
    \text{ for } s\in\SpaceS,
\end{equation}
and as $V^\pi_T(s) := R_T(s)$ for any $s\in\SpaceS$ if $t=T$, where it is 
understood that $\pi$ is followed from $t$ onwards. The notation $\E^\pi$ 
indicates, that the conditional expectation is taken under $\PM^\pi$. Note, 
that for $t>0$, $V_t^\pi$ is a function of $s$ only, because of the Markovian 
structure of both $P_t$ and $\pi_t$. $V_0^\pi(s)$ gives the \textit{expected 
total discounted reward} over the whole decision making horizon if the agent 
follows policy $\pi$ and the environment starts in state $s$. For any 
$t\in[T]$, we define the \textit{optimal value function} of the MDP 
\eqref{eq:MDPTuple} at time $t$ as
\begin{equation}
\label{eq:MDPOptimalValueFunction}
    V^*_t(s) := \sup_{\pi\in\Pi^R} V^\pi_t(s), 
    \quad \text{ for } s\in\SpaceS,
\end{equation}
and as $V^*_T(s):=R_T(s)$ if $t=T$. The value $V^*_t(s)$ corresponds to the 
maximal attainable expected reward over the remaining decision horizon and is 
also referred to as the optimal value given state $s$ at time $t$. Because of 
the boundedness assumption for reward functions, $V_t^\pi(s)$ exists and 
is bounded by a constant independent of $\pi$, for all $s\in\SpaceS$, and all 
$t\in[T]$, and, thus, the same holds true for holds true for $V^*_t(s)$ for 
any $s\in\SpaceS$ and all $t\in[T]$. Furthermore, since the functions $R_T$, 
$R_t$, and $\pi_t$ are measurable for all $t\in[T]$, and the composition of 
measurable functions is again measurable, $V_t^\pi$ is a measurable function 
from $\SpaceS$ to $\SetR$ for each $t\in[T]$ according to
\cref{lem:KernelsAuxiliaryMeasurableExpectation}. A priori $V^*_t$ may not be
measurable, since the supremum in \eqref{eq:MDPOptimalValueFunction} is taken
over an uncountable family of measurable functions, which need not be
measurable.
%%
However, with our assumption $\abs{\SpaceA}<\infty$, the optimal value
functions $V^*_t$ are indeed measurable, which follows from the fact, that the
supremum in \eqref{eq:MDPOptimalValueFunction} equals the maximum over the 
finite set of all possible action sequences $\{a_i\}_{i=0}^{T-1}$.
%\footnote{Measurability of $V_t^*$ could also be proofed by adapting the \cite[Theorem~16.1.28]{hinderer2016_do}, which provides the measurability of the optimal value function in a stationary MDP, to the non-stationary set-up in Theorem 16.1.28 therein, corresponding to the setting in this work.}
Additionally, $\E^\pi[\abs{V_t^*(S_t)}]<+\infty$ for all $t\in[T]$ and the 
functions $V^*_t$ satisfy the so called \textit{Bellman optimality equations},
\begin{equation}
\label{eq:MDPBellmanValue}
    V^*_t(s) 
    = \max_{a\in\SpaceA} \Bigl\{ 
        R_t(s,a) + \gamma\, \E_{P_{t+1}(\cdot|s,a)}\bigl[V^*_{t+1}(S_{t+1})\bigr] 
    \Bigr\},
\end{equation}
for all $s\in\SpaceS$ and $t\in[T-1]$, with terminal condition 
\begin{equation}
\label{eq:MDPBellmanValueTerminal}
    V^*_T(s) = R_T(s),
\end{equation}
where $\E^\pi_{P_{t+1}(\cdot|s,a)}$ is the conditional expectation of 
$V^*_t(S_{t+1})$ given $S_t=s, A_t=a$, which determined by the transition 
kernel $P_t$, as in \cref{eq:ConditionalExpectationSingle}. Any solution to 
the non-linear system of equations \eqref{eq:MDPBellmanValue} and 
\eqref{eq:MDPBellmanValueTerminal} is a collection of optimal value functions 
for the underlying MDP.\footnote{
    Suppose $(v_t)_{t\in[T]}$ is a solution of 
    \cref{eq:MDPBellmanValue,eq:MDPBellmanValueTerminal}, then 
    $(v_t+c)_{t\in[T]}$ for any $c\in\SetR$ is also a solution. Thus, the 
    optimal value function of a finite-horizon MDP given by \cref{eq:MDPTuple} 
    is not unique.
}
With the notion of value functions, we define optimality of a policy. A 
policy $\pi^*\in\Pi^\text{R}$ is called an \textit{optimal policy} if,
\begin{equation}
\label{eq:MDPOptimalPolicy}
    V_0^{\pi^*}(s) \geq V_0^\pi(s), 
    \quad \text{for all } s\in\SpaceS, \text{ and }\pi\in\Pi^\text{R}.
\end{equation}
The corresponding value functions for an optimal policy $\pi^*$ 
satisfy
\begin{equation}
\label{eq:MDPOptimalPolicyValueFunction}
    V^{\pi^*}_t(s) = V^*_t(s), 
    \quad \text{ for all } s\in\SpaceS, \text{ and } t\in[T].
\end{equation}
Specifically, $V^{\pi*}_0(s) = V^*_0(s)$ for any $s\in\SpaceS$, and $\pi^*$ 
achieves the maximal attainable expected discounted reward over the whole 
decision period. In general, an optimal policy may not exist under all 
modeling assumptions. However, when $\SpaceA$ is finite, an optimal policy 
does indeed exist in the class of deterministic policies, in other word, the 
supremum in \eqref{eq:MDPOptimalValueFunction} is attained for some 
$\pi\in\Pi^\text{D}$.
% Existence of optimal policy
\begin{theorem}
\label{thm:MDPOptimalPolicyExistence}
    For any finite-horizon Markov Decision Process given by \eqref{eq:MDPTuple},
    there exits a deterministic Markovian policy that is optimal.
\end{theorem}
\begin{proof}
    Proposition~4.4.3 in \cite{puterman2005_mdpdsdp} proves the case with 
    finite or countable $\SpaceS$. For general, Borel state space, 
    \cite[Theorem~16.1.6]{hinderer2016_do} shows the existence of an optimal 
    policy for stationary Markov Decision Processes, where $P_t\equiv P$ and 
    $R_t\equiv R$ for all $t\in[T-1]$ for some fixed state transition kernel 
    $P$ and reward function $R$, with the help of the so called Structure 
    Theorem, Theorem~16.1.12 therein. The non-stationary case, follows from 
    the non-stationary version of the Structure Theorem, Theorem 16.1.28 
    therein.
\end{proof}

From the Bellman equations \eqref{eq:MDPBellmanValue} and 
\eqref{eq:MDPBellmanValueTerminal}, and the fact that the value functions of 
an optimal policy satisfy \eqref{eq:MDPOptimalPolicyValueFunction}, it 
is evident that for a policy to be optimal it must select actions that attain 
the maximum on the right hand side of \eqref{eq:MDPBellmanValue} at each 
$t\in[T-1]$, thus,
\begin{equation}
\label{eq:MDPMaximizingPolicy}
    \pi^*_t(s) \in \argmax_{a\in\SpaceA} \Bigl\{ 
        R_t(s,a) + \gamma\, \E_{P_t(\dd{s_{t+1}}|s,a)}\bigl[
            V^*_{t+1}(S_{t+1})
        \bigr] 
    \Bigr\}.
\end{equation}
%
A policy which satisfies \eqref{eq:MDPMaximizingPolicy} for all $t\in[T-1]$
and $s\in\SpaceS$ is called a \textit{maximizing} policy. The measurable 
selection theorem in \cref{appendix:MeasurableSelection}, guarantees the 
measurability of a maximizing policy $\pi^*_t$ at each $t$, thus it is indeed 
a deterministic policy according to \cref{def:MDPPolicy}.

The problem of optimal control of a Markov Decision Process is concerned with 
computing an optimal policy when the transition kernels and rewards functions
are known. This can be done through dynamic programming by recursively solving 
the Bellman equations \eqref{eq:MDPBellmanValue},
\eqref{eq:MDPBellmanValueTerminal} to get optimal value functions, and 
subsequently, choosing for each state a policy that maximizes the right hand 
side of \eqref{eq:MDPBellmanValue}. There exists vast literature on solving 
MDPs, and on the existence of optimal policies with varying modeling 
assumptions, we refer to \cite{puterman2005_mdpdsdp, hinderer2016_do, 
feinberg2002_hmdp} for some monographs on this matter. 
%
In contrast to the optimal control of an MDP, reinforcement learning is
concerned with learning an optimal or approximately policy without the 
knowledge of $\KernelP$ and $\KernelR$, solely from sampled transitions 
$(s_t, a_t, r_t, s_{t+1})$. In this case, the rewards and expectations in 
\eqref{eq:MDPBellmanValue} cannot be computed explicitly. Nevertheless, the 
Bellman optimality equations serve as basis for many RL algorithms, including 
Q-Learning which is presented in the following 
\cref{chapter:RLAndMDP:section:QLearningInMDP}.

%------------------------------------------------------------------------------
\section{Q-Learning in Markov Decision Processes} 
\label{chapter:RLAndMDP:section:QLearningInMDP}
%------------------------------------------------------------------------------

When reward functions $R_t$ and state transition kernels $P_t$ are unknown, 
and the expectation in \eqref{eq:MDPBellmanValue} can not be computed 
explicitly, exact dynamic programming methods are not applicable. In 
Q-Learning, introduced in \cite{watkins1992_q}, the agent interacts with the 
environment modeled as MDP, to collect transition samples $(s_t, a_t, r_t, 
s_{t+1})$. Those are subsequently used in a sample based version of the 
Bellman equations \eqref{eq:MDPBellmanValue}, 
\eqref{eq:MDPBellmanValueTerminal} to learn so called Q-functions from which 
an optimal policy is derived. Note, hat most existing work on Q-Learning 
focuses on the infinite-horizon, stationary case, where transition kernels and 
reward functions do not depend on time. As a result, Q-functions are time 
independent, and optimal policies are stationary. In contrast, this section 
describes Q-Learning for a finite-horizon, non-stationary MDP given by the 
tuple \cref{eq:MDPTuple}, where both, policies and Q-functions retain the 
time-dependence.

For an MDP given by \eqref{eq:MDPTuple}, the \textit{optimal 
$Q$-function} at time $t\in[T-1]$ is defined as
\begin{equation}
\label{eq:MDPQFunctions}
    Q^*_t(s,a) := 
        R_t(s,a) + \gamma\, \E_{P_t(\dd s_{t+1}|s,a)}\bigl[V^*_{t+1}(S_{t+1})\bigr],
        \quad \text{for } (s,a)\in\SpaceS\times\SpaceA,
\end{equation}
and for $t=T$ as $Q^*_T(s,a):=R_T(s)$ for $s\in\SpaceS$, and we refer to the 
collection $(Q^*_t)_{t\in[T]}$ as the optimal Q-functions. The function 
$Q^*_t$ is also known as the optimal \textit{action-value function} at time 
$t$, because it provides the optimal value over the remaining episode if the 
process is in state $s$ at time $t$, and action $a$ is selected. Note, that 
since $R_t$ is measurable, and the map $(s,a)\mapsto \E_{P_t(\cdot|s,a)}
[V^*_{t+1}(S_{t+1})]$ is measurable according to 
\cref{lem:KernelsAuxiliaryMeasurableExpectation}, $(s,a)\mapsto Q^*_t(s,a)$ is 
a measurable map from $\SpaceS\times\SpaceA$ to $\SetR$. Additionally, it is 
bounded.

The optimal value functions and optimal Q-functions are
related in the following way:
\begin{equation}
\label{eq:MDPValueQConnection} 
    V^*_t(s) = \max_{a\in\SpaceA} Q^*_t(s,a), \quad s\in\SpaceS.
\end{equation}
If we plug \eqref{eq:MDPValueQConnection} into \eqref{eq:MDPBellmanValue}, we 
get the \textit{Bellman equations for optimal Q-functions},
\begin{equation}
\label{eq:MDPBellmanQ}
    Q^*_t(s,a) =         
    R_t(s,a) + \gamma\, \E_{P_{t+1}(\cdot|s,a)}\Bigl[
        \max_{\alpha\in\SpaceA} Q^*_{t+1}(S_{t+1},\alpha)
    \Bigr], 
    \quad t\in[T-1],
\end{equation}
with the terminal condition
\begin{equation}
\label{eq:MDPBellmanQTerminal}
	Q^*_T(s,a) = R_T(s),
\end{equation}
for any $(s,a) \in \SpaceS \times \SpaceA$. Since $\SpaceA$ is finite, for any 
$s\in\SpaceS$, the maximum in \eqref{eq:MDPValueQConnection} is attained. 
Thus, according to the measurable selection theorem in 
\cref{appendix:MeasurableSelection}, for any $t\in[T]$, the map $s \mapsto 
\max_{a\in\SpaceA} Q^*_t$ is measurable, and the conditional expectation 
$\E_{P_{t+1}(\cdot|s,a)}$ in \eqref{eq:MDPBellmanQ} is well-defined.
%
Analogous to value functions, any solution of the non-linear system of 
equations \eqref{eq:MDPBellmanQ} and \eqref{eq:MDPBellmanQTerminal}, is a 
collection of optimal action-value functions for the underlying MDP. Since an optimal deterministic policy $\pi^*$ satisfies \eqref{eq:MDPMaximizingPolicy}, it follows directly from the definition of the
Q-functions that 
\begin{equation}
\label{eq:MDPMaximizingPolicyQ}
    \pi^*_t(s) \in \argmax_{a\in\SpaceA} Q^*_t(s,a), 
    \quad \text{for all } s\in\SpaceS, t\in[T-1].
\end{equation}
We call a policy that maximizes $Q^*_t(s,\cdot)$ for each $s\in\SpaceS$ greedy 
with respect to $Q^*_t$. A policy that is greedy with respect to $Q^*_t$ for 
all $t\in[T-1]$ is called \textit{greedy with respect to the collection 
$(Q_t^*)_{t\in[T]}$}. \Cref{eq:MDPMaximizingPolicyQ} implies that knowledge of 
the optimal Q-functions is sufficient for acting optimally in the underlying 
MDP, by following a policy that is greedy with respect to $(Q_t^*)_{t\in[T]}$.

The question of how to find the functions $Q_t^*$ in a computational manner 
without the knowledge $P_t$ and $R_t$ was considered by Watkins in 
\cite{watkins1989_lfdr}. The author proposed the Q-learning algorithm, a 
sample based version of the Bellman equations \eqref{eq:MDPBellmanQ}, where 
observed transition tuples $(s_t, a_t, r_t, s_{t+1})$ are used to update 
estimates of the functions $Q^*_t$. We present the method for a 
finite-horizon, non-stationary MDP \eqref{eq:MDPTuple}, where additionally to $
\abs{\SpaceA}<+\infty$, we assume $\abs{\SpaceS}<+\infty$, a setting that is 
considered in \cite{vp2024_eemsgwfhq}. Then, at each step $t\in[T]$, the 
optimal Q-function $Q^*_t$ can be represented as a $\abs{\SpaceS \times 
\SpaceA}$-matrix, with its entries given as the optimal Q-values
$Q^*_t(s,a)\in\SetR$ for pairs $(s,a)\in\SpaceS\times\SpaceA$. Q-learning 
starts with some arbitrary initial value $Q^0_t(s,a)\in\SetR$ as an estimate 
for $Q^*_t(s,a)$ for each pair $(s,a) \in \SpaceS \times \SpaceA$, and each 
$t\in[T]$. Let $Q^n_t(s,a)$ be the estimate of $Q_t^*(s,a)$ after the $n$-th 
update iteration.
% Q-update description WITH with time index
Assume that, by interaction with the environment, the agent has observed the 
transition sample $(s_t,a_t,r_t,s_{t+1})$ at $t\in[T-1]$, where $s_t$ is the 
state at $t$, $a_t$ the action selected in this state, $r_t=R_t(s_t,a_t)$ is 
the reward the agent received, and $s_{t+1}$ is the new state to which the 
environment transitioned to, sampled from $P_t(\cdot|s_t,a_t)$. Then, for each 
$(s,a)\in\SpaceS\times\SpaceA$, we update the current estimate $Q^n_t(s,a)$ by
\begin{equation}
\label{eq:MDPQUpdate}
    Q^{n+1}_t(s,a) 
    \gets Q^n_t(s,a) + \alpha_{n}\, \Bigl[
        r_t + \gamma \max_{a'\in\SpaceA} Q^{n}_{t+1}(s_{t+1},a') - Q^n_t(s,a)
    \Bigr] \times \mathds{1}_{\{s,a\}}(s_t,a_t),
\end{equation}
where $\alpha_n \in (0,1]$ is the learning rate for update iteration $n+1$. 
For the observed pair $(s_t, a_t)$, it balances the value of the new estimate 
$Q^{n+1}_t(s_t,a_t)$, between the current estimate $Q^n_t(s_t,a_t)$ and the 
target value $r_t + \gamma \max_{a'\in\SpaceA} Q^{n}_{t+1}(s_{t+1},a')$, which 
resembles a sample of the Bellman equation \eqref{eq:MDPBellmanQ}, and is a
supposedly better estimate of the true value $Q^*_t(s_t,a_t)$ due to $r_T$ 
being the actual reward received from the environment. This process of 
learning estimates from other estimates is called bootstrapping. Note that, if 
we have observed the specific pair $(s_t,a_t)$, then for any other pair 
$(s,a)\in\SpaceS\times\SpaceA\setminus\{(s_t,a_t)\}$, we have $Q^{n+1}_t(s,a) 
= Q^{n}_t(s,a)$. If $t=T-1$, then the agent additionally receives terminal 
reward $r_T = R_T(s_T)$, which serves as target to update terminal estimates by
\begin{equation}
\label{eq:MDPQUpdateTerminal}
    Q^{n+1}_T(s,a) 
    \gets Q^n_T(s,a) + \alpha_n\, \bigl[
        r_T - Q^n_T(s,a)\bigr] \times \mathds{1}_{\{s\}}(s_T),
\end{equation}
for all $s\in\SpaceS$, and $a\in\SpaceA$.

Transition samples are obtained by following some fixed exploration policy 
$\pi\in\Pi^\text{R}$. The full Q-learning algorithm for an MDP 
\eqref{eq:MDPTuple}, with finite state and action spaces is provided in
\cref{alg:MDPQLearning}. Using stochastic approximation arguments, the authors
of \cite{vp2024_eemsgwfhq} proved, that $Q^n_t(s,a)$ converges to $Q^*_t(s,a)$ 
for $n\to\infty$ almost surely for each pair $(s,a)\in\SpaceS\times\SpaceA$, 
and all $t\in[T]$, provided the learning rates satisfy the so-called
Robbins-Monro conditions \cite{robbins1951_sam},
\[
    \sum_{n=0}^\infty a_n = \infty, 
    \qquad 
    \sum_{n=0}^\infty a_n^2 < \infty,
\]
and each pair $(s,a)$ is sampled infinitely often at each time $t\in[T-1]$.%
\footnote{
    \cite{vp2024_eemsgwfhq} does not explicitly mention the second condition, 
    however, the used sampling policy implicitly guarantees it. 
}
\cite{watkins1992_q} proved convergence for the case of an infinite-horizon, 
stationary MDP.

\begin{algorithm}
	\caption{Q-learning for finite-horizon, non-stationary, finite MDPs}
	\label{alg:MDPQLearning}
	\begin{algorithmic}[1] 
    % optional parameter [1] means line numbering for every line 
		\Require Total number of iterations $N\in\SetN$, exploration policy 
        $\pi$, learning rates $\{\alpha_n\}_{n=0}^{N-1}$
                
        \State Initialize $Q^0_t(s,a)$ for all $(s,a)\in\SpaceS\times\SpaceA$, 
         and $t\in[T]$.
		\State Sample $s_0\in\SpaceS$ from initial distribution $\mu$.
		% for loop, one for each episode
		\For{iteration $n =0,1,2,\dots,N-1$}
            \State Set $t = n \bmod T$.
            \State Sample action $a_t$ from $\pi_t(\cdot|s_t)$.
            \State Receive reward $r_t$ and new state $s_{t+1}$, 
                after taking $a_t$.
            \State Update estimates $Q^{n}_t(s,a)$ for all $(s,a)$ according to
                \eqref{eq:MDPQUpdate}.
            \If{$t=T-1$}
                \State Receive terminal reward $r_T$.
                \State Update $Q^{n}_T(s,a)$ for all $(s,a)$ according to
                \eqref{eq:MDPQUpdateTerminal}.
            \EndIf
		\EndFor
	\end{algorithmic}
\end{algorithm}

\begin{remark}
    Q-learning is an \textit{off-policy} reinforcement learning method, where 
    the policy that is learned is different from the exploration policy $\pi$
    that generates transition samples. In general, transition samples could be 
    collected before learning estimates takes place. However, the 
    exploration policy must guarantee sufficient \textit{exploration} of the 
    state space to ensure learning of correct estimates and, theoretically, 
    visit each pair $(s, a)$ infinitely often. A popular choice as exploration 
    policy is $\varepsilon$-greedy, with $\varepsilon\in[0,1]$. At each step 
    $t\in[T-1]$, this policy if greedy w.r.t to the current estimate $Q^n_t$ 
    with probability $1-\varepsilon$, and samples an action uniformly from 
    $\SpaceA$ otherwise. To ensure sufficient \textit{exploitation} of 
    estimates already learned, the value of exploration parameter 
    $\varepsilon$ is usually gradually reduced over the training period.
\end{remark}

\Cref{alg:MDPQLearning} hinges on the assumption that $\SpaceS$ is finite.
When the cardinality of $\SpaceS$ is infinite or large, such that storing an
array of size $\abs{\SpaceS}\times\abs{\SpaceA}\times{T}$ becomes infeasible, 
one has to resort to function approximation to estimate optimal Q-values. 
Consider a family of parametrized functions $\mathcal{Q} = 
\{Q_\theta \mid \theta \in \SetR^m\}$ for fixed $m\in\SetN$, with $Q_\theta:
\SpaceS \times \SpaceA \to \SetR$ for any parameter vector $\theta$. For 
$Q_\theta \in \mathcal{Q}$ and a pair $(s,a)\in\SpaceS\times\SpaceA$ we write 
$Q(s,a;\theta) := Q_{\theta}(s,a)$. For each time $t\in[T]$, we approximate 
the optimal Q-function $Q^*_t$ by
\begin{equation}
\label{eq:MDPQApprox}
    Q^*_t(s,a) \approx Q(s,a;\theta_t), 
    \quad \text{for all } (s,a)\in\SpaceS\times\SpaceA,
\end{equation}
with parameter vector $\theta_t \in \SetR^m$, and we replace iteration 
\eqref{eq:MDPQUpdate} by a suitable iteration to find a parameter vector 
$\theta^*_t, t\in[T]$, that yields the best approximation of $Q_t^*$ in 
$\mathcal{Q}$. Analogue to the finite case, Q-learning with function 
approximation starts with some arbitrary initial
parameter vector $\theta^0_t$ for each $t\in[T]$, which provides initial
estimates $Q(s,a;\theta^0_t)$. Let $\theta^n_t$ be parameter vector at step 
$t$, after update iteration $n\in\SetN$.
% Q-update description WIHT time index
Upon observing the transition tuple $(s_t,a_t,r_t,s_{t+1})$ at time $t\in 
[T-1]$, where again $r_t=R_t(s_t,a_t)$, and $s_{t+1}$ is sampled from 
$P_t(\cdot|s_t,a_t)$, we define the target value
\begin{equation}
\label{eq:MDPQApproxTarget}
    y^n_t 
    := r_t 
        + \gamma \max_{a\in\SpaceA} Q(s_{t+1},a;\theta^n_t),
\end{equation}
and update the current approximation given by $\theta^n_t$ towards the target 
value $y_t^n$ via single-step stochastic gradient descent, minimizing the 
squared loss\footnote{
    Other loss functions may be applied, as for example the smooth-$L_1$ loss
    function used in the experiments in \cref{chapter:Experiments}. The loss 
    function's gradient in \eqref{eq:MDPQApproxLoss} changes accordingly.
}
\begin{equation}
\label{eq:MDPQApproxLoss}
    \mathcal{L}(y_t^n, \theta^n_t) 
    = \tfrac{1}{2}\bigl(y^n_t - Q(s_t,a_t;\theta^n_t)\bigr)^2.
\end{equation}
This corresponds to updating the parameter vector $\theta_t^n$ to 
\begin{equation}
\label{eq:MDPQApproxUpdate}
    \theta^{n+1}_t 
    \gets \theta^n_t + \alpha_n
        \Bigl[
            r_t + \gamma \max_{a\in\SpaceA} Q(s_{t+1},a;\theta^n_{t+1})
            - Q(s_t,a_t;\theta^n_t)
        \Bigr]\, \nabla_{\theta}\, Q(s_t,a_t;\theta){\Big|}_{\theta=\theta_t^n}
\end{equation}
with learning rate $\alpha_n\in(0,1]$ in update iteration $n+1$. The algorithm 
for Q-learning with functions approximation is analogue to 
\cref{alg:MDPQLearning}, with the finite-case update \eqref{eq:MDPQUpdate} 
replaced by the parameter vector update \eqref{eq:MDPQApproxUpdate}.

A typical choice for the family $\mathcal{Q}$ is the class of neural networks,
such that $Q(s,a;\theta) = F(s,a;\theta)$ where $F$ is a
multilayer-perceptron parametrized by $\theta=(\bm{W},\bm{b})$ with a
collection of weight matrices $\bm{W}$ and bias vectors $\bm{b}$. Neural Fitted
Q-learning proposed in \cite{riedmiller2005_nfifewdenrlm} applies a shallow 
neural network approximation in an off-line fashion, with pre-collected sets 
of transition samples. In \cite{mnih2015_hctdrl}, the authors introduced Deep 
Q-Network, where a more complex deep neural network is used to estimate
the optimal Q-function (stationary case), together with the use of a separate 
target neural network for action selection, and experience reply to enhance 
learning performance. Their approach demonstrated impressive learning and 
generalization results across a wide range of Atari 2600 computer games 
\cite{bellemare2013_aleepga}, using the same learning network configuration 
and hyperparameter values across games.

Another common approach is to set $\mathcal{Q}$ to be the linear span of some
set of feature maps $\phi_k(s,a):\SpaceS\times\SpaceA\to\SetR$ for $k=1,
\dots,m$, resulting in linear approximations,
\[
    Q(s,a;\theta) 
    = \sum_{k=0}^m \phi_k(s,a)\theta_k
    = \phi(s,a)^\top \theta
\]
with the feature vector $\phi(s,a) = (\phi_1(s,a),\dots,\phi_m(s,a))^\top$, and the 
parameter vector $\theta = (\theta_1, \theta_2, \dots \theta_m)^\top$.
The paper \cite{melo2008_arlwfa} established conditions under which Q-Learning with 
linear function approximation converges, these including that the exploration 
policy induces the state process $(S_t)_{t\geq0}$ to be an ergodic Markov 
chain, and that the feature maps $\phi_k$ are linearly independent. With the
goal of weakening those conditions, \cite{carvalho2020_cvqwlfa} provides a 
convergent two-time scale variant of Q-learning with linear function 
approximation with less stringent assumptions on the exploration policy, as 
well as error bound for the method. Note that the mentioned references on
Q-learning with linear function approximation \cite{melo2008_arlwfa,%
carvalho2020_cvqwlfa}, and on Q-learning with neural network approximation 
\cite{mnih2015_hctdrl,riedmiller2005_nfifewdenrlm}, assumed an 
infinite-horizon, stationary MDP as environment model. Nevertheless, the 
approaches can be readily applied to the finite-horizon, non-stationary case 
by considering separate estimates for each time $t$, however, convergence 
results and sampling efficiency would have to be reassessed. In this work we 
take a similar approach to Q-learning with linear function approximation in 
MDPs, in that we also use linear functions of some feature maps to approximate 
optimal Q-function, albeit in a history-dependent setting. Here, the feature 
maps have to be functions of the observed history $h_t=(s_0, a_0, s_1, \dots, 
a_{t-1}, s_t)$, and are given by the signature of a path, introduced in 
\cref{chapter:Signature}.

%------------------------------------------------------------------------------
\section[The Issue with MDPs as a Modeling Framework]{The Issue with Markov Decision Processes as a Modeling Framework}
\label{chapter:RLAndMDP:section:IssueWithMDP}
%------------------------------------------------------------------------------

Reinforcement learning problems are typically formalized as Markov decision
processes, where states satisfy the Markov property \eqref{eq:MDPMarkovNature},
and the information contained in the current state is sufficient for the agent 
to learn an optimal policy. This formalization however, implicitly assumes 
that the true underlying environment indeed permits to be modeled as such, 
and that the model captures all the relevant information. In many scenarios 
this assumption may not be true, and thus, modeling the environment as MDP is 
problematic. We highlight three cases.

First, the environment might be described better as an $n$-step Markov 
decision process for some $n\in\SetN$. The difference to MDPs defined in 
\cref{chapter:RLAndMDP:section:MDP} is that the distribution of the state 
$S_{t+1}$ depends on states and action at the prior steps $t, t-1, \dots, t-n$,
such that
\[
    \PM[S_{t+1}\in B \,|\,H_t = h_t] 
    = \PM[
        S_{t+1}\in B \,|\, S_{t-k}=s_{t-k}, A_{t-k}=a_{t-k}, 
            \dots S_t = s_t, A_t = a_t 
    ],
\]
for any Borel set $B\in\mathcal{B}_\SpaceS$, with history $h_t = (s_0, a_0, 
s_1,\dots, a_{t-1}, s_{t})$. When the true environment follows an $n$-step 
MDP, one can construct an 1-step MDP by defining new states as 
\[
    \tilde{s}_t := (s_{t-k}, a_{t-k}, \dots, a_{t-1}, s_t).
\]
and corresponding transition kernels $\tilde{P}_t(\dd{\tilde s_{t+1}}|
\tilde s_t, a_t)$.\footnote{
    The new state transition kernels $\tilde P_t$ are determined by the true
    transition kernels $P_t$ which are functions of the last $n$ states and
    action. Since we consider the constructed 1-step MDP in a reinforcement
    learning set-up we do not need  $\tilde P_t$ exactly.
}
The random states $\tilde{S}_{t+1}$ exhibit the Markov
property \eqref{eq:MDPMarkovNature} w.r.t. to $\tilde s_t$ and $a_t$. For 
example, in \cite{mnih2015_hctdrl}, the Deep Q-Network trained to
play Atari 2600 games, receives the last $n=4$ frames as input with selected
actions repeated during those frames, and the model for the games considered
becomes Markovian. Other games, however, would require more than four frames, 
and are still non-Markovian w.r.t to the newly constructed states. This shows 
that, even with the knowledge that an environment is an $n$-step MDP
for \textit{some} $n$, selecting the correct $n$ still poses an issue.

Second, in real-world environments it is seldom possible to determine the full
true state, let alone provide it to the agent, leading to partial 
observability. Such environments are modeled as \textit{partially observable
Markov Decision Processes} (POMDP) \cite{kaelbling1998_paposd}, which are
defined as an MDP given by \eqref{eq:MDPTuple}, together with an additional
measurable space, \textit{observation space}  $\SpaceO$, and a
collection of \textit{observation kernels} $(O_t)_{t\in[T]}$. While the true
states $s\in\SpaceS$ of an POMDP satisfy the Markov property
\eqref{eq:MDPMarkovNature}, this structure is hidden. Assume that at
time $t$ the process is in state $s_t\in\SpaceS$. Instead of state $s_t$, the
agent receives only a partial observation $o_t\in\SpaceO$, which is generated
according to the distribution $O_t(\dd o|s_t)$. Reinforcement learning in
POMDPs is known to be a computationally as well as statistically challenging
task. In the worst case, learning a near-optimal policy may require a sample
size exponential in episode length \cite{krishnamurthy2016_rlwro}.
Although the true underlying environment is an MDP, incomplete state
information prohibits modeling it as such. In terms of Q-learning this means
that in general $Q(o,a;\theta_t)\neq Q(s,a;\theta_t)$
for any Q-function approximation parametrized by $\theta_t$, with $s$ being
the true state and $o$ the observation received. Estimating Q-values as in
\cref{chapter:RLAndMDP:section:QLearningInMDP} simply with states replaced by
observations can be arbitrarily bad. Coming back to the Atari 2600 games as an 
example, with the number of frames 
less than four, the games considered in \cite{mnih2015_hctdrl} are
partially observable environments. Specifically, in the game ``Pong'', the 
agent requires both, the position and velocity of a ball to infer an optimal 
paddle move. However, with only one frame as state, the agent receives just 
the ball's position, with no means to infer its velocity, and can thus not 
determine an optimal paddle move as action. In other cases, computational 
constraints could prevent to store the true state, if it was fully observable, 
and some compact state representation has to be employed which might not 
include all necessary information for the model to be an MDP. This again shows 
the importance of state representation in reinforcement learning as 
discussed at the beginning of this chapter.

Finally, an environment may simply exhibit dependence on the past, in that
the distribution of the state at time $t+1$, and possibly also the reward,
depend on the whole state-action history $h_t$, or parts of it
at specific steps prior to $t+1$. In such environments it is required to
maintain information over many steps and they are also referred to as 
long-term memory tasks. They can be either fully or partially observable and
subsume $n$-step MDPs and POMDPs mentioned above. 
In \cite{heess2015_mcwrnn}, the authors considered reinforcement learning in
such environments using recurrent neural networks to tackle the memory
requirements. One of the problems considered is to control a 3-joint reacher
which must hit a randomly placed target, but the targets position is only
provided in the very first observation. Thus, the probability of hitting a
target state at $t+1$ depends on the observation or state at time $t=0$ and
the agent agent has to memorize the position throughout the whole episode.

In all three cases, the problem's nature does not admit an MDP as model. 
By only considering the current state, important contextual
information and dependencies between previous states and actions are
overlooked and valuable information that can be derived from considering the
agent's history is ignored. This can lead to sub-optimal decision-making as the
agent may not be able to effectively capture and utilize all relevant
information for maximizing long-term rewards. As a result, history dependence 
has to be taken into account in the modeling and learning process. To this 
end, we introduce a history dependent environment in the following
\cref{chapter:HistoryRLEnvironment}, which assumes no Markovian structure in
the first place, and propose a method for learning in such environments in
\cref{chapter:SignatureQLearning}.



%%%%%%%%%%%%%%%%%%%%%%%%%%%%%%%%%%%%%%%%%%%%%%%%%%%%%%%%%%%%%%%%%%%%%%%%%%%%%%%
\chapter[A History-Dependent RL Environment]{A History-Dependent Reinforcement Learning Environment} 
\label{chapter:HistoryRLEnvironment}
%%%%%%%%%%%%%%%%%%%%%%%%%%%%%%%%%%%%%%%%%%%%%%%%%%%%%%%%%%%%%%%%%%%%%%%%%%%%%%%

In this chapter we generalize the exposition outlined in 
\cref{chapter:RLAndMDP:section:IssueWithMDP,%
chapter:RLAndMDP:section:QLearningInMDP} for MDPs to a history-dependent 
reinforcement learning environment. We introduce a history-based decision 
process (HDP) as an alternative modeling framework, describe its components, 
and further define (optimal) value functions and optimality of a policy. No 
Markov assumptions on the dynamics are made, and the distribution of
future observations and rewards may depend on the full history prior to the
current time. This leads to the necessity to adapt policies and optimal value
functions, and define them in terms of histories.

Unless otherwise stated, the definitions and results in this chapter are, to the best of our knowledge, new, but  of course, they build on existing work from the literature. The
definition of HDPs as a general RL framework is inspired by 
\cite{majeed2018_qcndp}, which extends convergence of Q-learning 
from MDPs to a subset of non-MDPs, so called Q-value Uniform Decision Processes, 
and further by \cite{hutter2016_esabmdp}, which is related to Feature-RL 
\cite{hutter2009_frlpum}, and concerned with aggregating histories to states 
such that the aggregated process solves the original problem. Both works are 
discussed in the literature review in \cref{chapter:Introduction:par:RelatedWork}. The approach of 
expressing the conditional expectations that define (optimal) value functions 
as integrals w.r.t. product kernels, and to decompose this product to 
establish recursion relations for said functions, draws upon
\cite[Chapters~9. to 11.]{bertsekas1996_socdc} and 
\cite[Chapter~16]{hinderer2016_do}, where cost respectively value functionals
are a priori introduced as integrals w.r.t. to (products of) stochastic 
kernels. The latter reference further provides the measurable selection 
theorem detailed in \cref{appendix:MeasurableSelection}, foundational for 
proving the existence of a measurable maximizing policy in an HDP, which, as a 
result, yields measurability of the optimal value function and existence of an 
optimal policy.
 
\begin{remark}
    To avoid notational overload, many objects in the history-dependent 
    setting are denoted in the same way as in the Markovian setting in 
    \cref{chapter:RLAndMDP}. In the following, the objects in question 
    refer to the history-dependent case. In the case where an object relates 
    to the Markovian setting, this will be mentioned explicitly.
\end{remark}

%------------------------------------------------------------------------------
\section{History-Based Decision Processes}
\label{chapter:HistoryRLEnvironment:section:HDP}
%------------------------------------------------------------------------------
% environment and spaces
A discrete-time, finite-horizon, \textit{History-based Decision Process} (HDP) 
is defined as the tuple
\begin{equation}
\label{eq:HDPTuple}
    (\SpaceO, \SpaceA, \mathcal{P}, \mathcal{R}, \gamma, T),    
\end{equation}
with \textit{observation space} $\SpaceO\subset\SetR^d$ for $d\in\SetN$, 
finite \textit{action space} $\SpaceA$, finite time horizon $T\in\SetN$, and 
discount factor $\gamma\in[0,1]$. Additionally, $\mathcal{P} =
(P_t)_{t\in[T-1]}$ is a collection of \textit{observation transition kernels}, 
and $\mathcal{R}=(R_t)_{t\in[T]}$ a collection of \textit{reward functions}.
The components are described in detail in the paragraphs below.
% agent 
Analogous to an MDP, the agent interacts with an HDP in 
consecutive episodes. At step $t=0$ of episode $k\in\SetN$, the environment
emits some initial observation $o_0\in\SpaceO$ according to some fixed initial
distribution $\mu\in\text{Prob}(\SpaceO)$. Then, at each step $t\in[T-1]$ the
agent receives some observation of the environment's state $o_t\in\SpaceO$ and
chooses an action $a_t\in\SpaceA$. Upon taking action $a$, the environment
emits a reward  $r_{t}\in\SetR$ and subsequently transitions to a new state
with observation $o_{t+1}\in\SpaceO$. This cycle repeats until at time
$T-1$ the last action is taken, and upon transitioning to a 
terminal observation $o_{T+1}\in\SpaceO$, the agent receives a terminal reward 
$r_T\in\SetR$, and the next episode $k+1$ begins.
% histories
This behaviour produces an observation-action-history\footnote{
    More specifically it produces an observation-action-reward-history, however
    we exclude the reward in \eqref{eq:HDPHistory} since the reward $r_{t+1}$ can
    usually be derived from the tuple $(o_t,a_t,o_{t+1})$.
}, which at time $t\in [T]$ is given as,
\begin{equation}
    \label{eq:HDPHistory}
    h_t := (o_0, a_0, o_1, a_1, \dots, o_{t-1}, a_{t-1}, o_t) 
        \in (\SpaceO \times \SpaceA)^{t} \times \SpaceO.
\end{equation}

%------------------------------------------------------------------------------
% state space, action space
\paragraph{Observation, action and history spaces}
The observation space $\SpaceO$ is given as a compact subset of $\SetR^d$ 
for some $d\in\SetN$, and considered a measurable space with its Borel 
$\sigma$-algebra $\mathcal{B}_\SpaceO$. We use the notion of observations to
describe the environment's condition, as opposed to states, to include
potential partial observability. The action space $\SpaceA$ is again a finite, 
non-empty set of actions with cardinality $\abs{\SpaceA}\in\SetN$. Without 
loss of generality we assume $\SpaceA=\{1,2,\dots,\abs{\SpaceA}\}$, and, where 
needed, we denote its Borel $\sigma$-algebra as $\mathcal{B}_\SpaceA = 
2^\SpaceA$.

\begin{remark}
    In some problems the set of actions to choose from depends on the current 
    observation $o\in\SpaceO$ the environment is in, and is given as subset 
    $\SpaceA_{o}\subset\SpaceA$, such that $\SpaceA = \bigcup_{o\in\SpaceO} 
    \SpaceA_o$, compare
    \cite{ning2020_ddqoe, hinderer2016_do, puterman2005_mdpdsdp}. For
    simplicity, and because the experiments considered in 
    \cref{chapter:Experiments} are modeled such that all actions are 
    selectable at each step, we do not consider observation dependent actions 
    sets. This simplifies notation without a great loss in generality.
    At the same time, one could allow $\SpaceO$ and 
    $\SpaceA$ to depend on time. This case can be included in the formulation 
    \eqref{eq:HDPTuple} by setting $\SpaceO = \bigcup_{t\in[T]}\SpaceO_t$, and 
    analogously for $\SpaceA=\bigcup_{t\in[T]}\SpaceA_t$, and by modifying the 
    transition probabilities accordingly.
\end{remark}

\begin{definition}
\label{def:HistorySpaces}
    For an HDP \eqref{eq:HDPTuple}, the \textit{space of observation-action 
    histories} up to time $t$ is defined as
    \[
        \SpaceH_t := (\SpaceO \times \SpaceA)^t \times \SpaceO,
        \quad \text{for each } t \in [T],
    \]
    with the convention $\SpaceH_0 := \SpaceO$, and elements $h_t\in\SpaceH_t$
    given by \eqref{eq:HDPHistory}. Note, that $\SpaceH_t \subset
    \SetR^{d+t(d+1)}$ for any $t\in[T]$, and we consider it as measurable space
    with its Borel-$\sigma$-algebra
    $
        \mathcal{B}_{\SpaceO_t} 
        = \bigotimes_{i=1}^t (
            \mathcal{B}_\SpaceO \otimes \mathcal{B}_\SpaceA
        ) \otimes \mathcal{B}_\SpaceO
    $.\footnote{
        The Borel $\sigma$-algebra on $\SpaceH_t$ again coincides with the
        product $\sigma$-algebra generated by the $\sigma$-algebras 
        $\mathcal{B}_\SpaceO$ and $\mathcal{B}_\SpaceA$.
    }
\end{definition}

Additionally to the spaces $\SpaceH_t$, their time-augmented versions will be 
important for the proposed reinforcement learning method.

\begin{definition}
\label{def:TimeAugmentedHistorySpaces}
    For $t\in[T]$ and $h_t\in\SpaceH_t$, we define the \textit{time-augmented
    history} $\hat{h}_t$ by
    \begin{equation}
    \label{eq:TimeAugmentedHistory}
        \hat{h}_t  
        := (0, o_0, a_0, 1, o_1, a_1, \dots, t-1, o_{t-1}, a_{t-1}, t, o_t)
        \in\hat\SpaceH_t,
    \end{equation}
    with the time-augmented history space 
    $\hat{\SpaceH}_t := \bigtimes_{u=0}^{t-1} 
    (\{u\}\times\SpaceO\times\SpaceA)\times\{t\}\times\SpaceO$.
\end{definition}

\begin{remark}
    Similar to time-augmented paths of bounded variation in
    \cref{def:TimeAugmentedPath}, the time-augmented history $\hat{h}_t$ may
    be defined with the time-component given by normalized steps $u/T$, 
    remaining time steps $T-u$, or normalized remaining steps $1-u/T$, for 
    $u\in[t]$. From a reinforcement learning perspective, normalizing the 
    inputs to the learning algorithm improves stability, which is why in the 
    experiments in \cref{chapter:Experiments} the time-component will always 
    be normalized.
\end{remark}

Let $t\in[T]$ and consider a history $h_t\in\SpaceH_t$ and its 
time-augmentation $\hat{h}_t\in\hat\SpaceH_t$. We introduce the following 
notation and definitions regarding histories and their spaces:
\begin{enumerate}
    \item For $h_t$ (and $\hat h_t$), the last observation's time index $t$ 
    will be called the (time-augmented) history's length. Note, that $t$ 
    corresponds to the number of actions contained in the (time-augmented) 
    history, while the number of observation is $t+1$.
    
    \item Given time points $s,u\in[t]$ with $s < u$, we define the 
    sub-history from $s$ until $u$ of $h_t$ and $\hat h_t$, respectively, as
    \begin{align}
        h_{s:u} 
        &:= (o_s, a_s, \dots, o_{u-1}, a_{u-1}, o_u) \in \SpaceH_{s:u} \\
        \hat h_{s:u} 
        &:= (s, o_s, a_s, \dots, u-1, o_{u-1}, a_{u-1}, u, o_u)
            \in \hat \SpaceH_{s:u}
    \end{align}
    with the spaces $\SpaceH_{s:u}:=\SpaceH_{u-s}$ and 
    $\hat \SpaceH_{s:u} := \bigtimes_{u=s}^t
    (\{u\}\times\SpaceO\times\SpaceA)\times\{t\}\times\SpaceO$. The length 
    of the sub-histories $h_{s:u}$ and $\hat h_{s:u}$ is given by $u-s$.

    \item Often, we consider the concatenation of $h_t$ with the next 
    selected action $a_t\in\SpaceA$ and the subsequent observation 
    $o_{t+1}\in\SpaceO$, which we denote as $(h_t,a_t,o_t)$, with the 
    understanding that $h_t$ is ``unpacked'', in the sense that its 
    components are inserted into tuple $(h_t,a_t,o_t)$ without parenthesis, 
    yielding a concatenated history of length $t+1$. If  
    $a\in\SpaceA$ and $o\in\SpaceO$ are given without time-index, the 
    concatenation $(h_t,a,o)$ is defined in the same way. Analogously, we 
    denote the concatenation of $\hat h_t$ with $a_t, t+1, o_{t+1}$ as 
    $(\hat h_t, a_t, t+1, o_{t+1})$.

    \item Given a function of histories $f:\SpaceH_t\to\SetR$ for some
    $t\in[T]$, we sometimes want to emphasise certain input components to $f$, 
    by splitting the input history. For example, for any $h_t\in\SpaceA$, we 
    may write $f(h_t) = f(h_{t-1}, a_t, o_{t+1})$ to emphasize the last two 
    history components $a_t$ and $o_{t+1}$, or $f(h_t) = f(h_u, a_u, 
    h_{u+1:t})$ to highlight the history up to, or from some time $0\leq u <t$ 
    as well as action $a_u$. In the same way we may write $f(h_{t-1}, a, o)$ 
    where the input to $f$ is the concatenated history $(h_{t-1},a,o)$ of 
    length $t$ for some $a\in\SpaceA$ and $o\in\SpaceO$.  

    \item The concatenations of $h_t$ and $\hat h_t$ with a single action 
    $a\in\SpaceA$, are denoted $h_t^a$ and $\hat h_t^a$, respectively. They 
    are given in the form of the discrete streams of data
    \[
        h_t^a := (o_u, a_u)_{u\in[t]}
        \quad \text{and} \quad \hat{h}^a_t = (u,o_u, a_u)_{u\in[t]},
        %\in \SpaceH_t \times \SpaceA = (\SpaceO \times \SpaceA)^{t+1}
    \]
    where we set $a_t := a$, and with components $(o_u,a_u)\in\SetR^{d+1}$
    and $(u, o_u,a_u)\in\SetR^{d+1}$, respectively, for $u \in [t]$.    
    
    \item In \cref{chapter:SignatureQLearning} we compute the signature
    of histories $h_t^a$, and their time-augmentations $\hat{h}^a_t$. 
    Since the history as argument already includes the time index, we omit it 
    in the signature, and set
    \[
        S(h^a_t) := S_{0,t}(h^a_t),
        \qquad
        S\bigl(\hat{h}^a_t\bigr) := S_{0,t}\bigl(\hat{h}^a_t\bigr),        
    \]
    where the signature of the discrete streams of data $h_t^a$ and 
    $\hat{h}^a_t$ is defined via the signature of their linear 
    interpolations $H^a_{[0,t]}$ and $\hat{H}^a_{[0,t]}$, as in 
    \cref{def:SignatureOfStream}. In the same way, for the signature of a 
    sub-history, we set $S(h^{a_u}_{s:u}) := S_{s:u}(h^a_t)$. Specifically, 
    we will encounter $S(h^a_{t:t+1})$ when updating the signature of a 
    concatenated history after new data has arrived, which is simply the 
    signature of the linear interpolation between the data point 
    $(t,o_t,a_t)$ at time $t$ and $(t+1,o_{t+1},a)$ at $t+1$.
    
\end{enumerate}

The reason why we consider time augmented histories is twofold. First,
including time-steps in observations creates time-aware agents which, have 
shown superior performance compared to their time-unaware counterparts in 
finite-horizon reinforcement learning tasks in a variety of 
methods \cite{pardo2018_tlrl}. And second, having a time component in 
observation-action histories comes in handy when approximating 
Q-functions by linear functionals of the histories' signatures 
over time in \cref{chapter:SignatureQLearning} by means of
\cref{lem:Uniqueness,thm:UniversalApproximation}.

%------------------------------------------------------------------------------
% transition function
\paragraph{Observation transition kernels}
\label{par:ObservationTransitions}
At the beginning of each episode, the initial observation $o_0$ is 
sampled according to some initial distribution $\mu\in\text{Prob}(\SpaceO)$.
Observations at subsequent steps are emitted according to the observation
transition kernels $\mathcal{P} = (P_t)_{t\in[T-1]}$, defined as follows.
\begin{definition}
\label{def:ObservationTransitionKernel}
    For each $t\in[T-1]$, the \textit{observation transition kernel}
    $\mathcal{P}_t$ is a stochastic kernel from 
    $(\SpaceH_t\times\SpaceA, 
        \mathcal{B}_{\SpaceH_t}\otimes\mathcal{B}_\SpaceA)$
    to $(\SpaceO, \mathcal{B}_\SpaceO)$. 
\end{definition}
The transition kernel $\mathcal{P}$ have the same role as for MDPs in the 
previous chapter, with the difference that here they are functions of 
action-observation histories prior the current time step.
For $t\in[T-1]$, the workings of $\mathcal{P}_t$ are as follows: For any set 
$B\in\mathcal{B}_\SpaceO$, the probability of the next
observation $o_{t+1}$ belonging to $B$, provided that up to $t$ the
environment's history is $h_t\in\SpaceH_t$ and the current action chosen by
the agent is $a_t\in\SpaceA$, is given by $P_t(B\,|\,h_t,a_t)$. Put
differently, $o_{t+1}$ is sampled according to the probability distribution 
$P_t(\,\cdot\,|\, h_t,a_t)$ on the measurable space 
$(\SpaceO, \mathcal{B}_\SpaceO)$.

%------------------------------------------------------------------------------
% reward function
\paragraph{Reward functions}
Analogue to the observation transitions kernels, the rewards are functions of
histories. For $t\in[T-1]$, the \textit{reward function} $R_t:\SpaceH_t \times \SpaceA
\to \SetR$ is measurable w.r.t. $\mathcal{B}_{\SpaceH_t} \otimes 
\mathcal{B}_\SpaceA$. It provides the reward $R_t(h_t, a)$ from choosing
action $a\in\SpaceA$ in history $h_t \in
\SpaceH_t$ at time $t$.\footnote{
    This definition includes the case where the reward depends on history
    $h_t$ only through the last observation $o_t$.
}
Additionally, at terminal time $T$, the terminal
reward is determined by the $\mathcal{B}_{\SpaceH_T}$-measurable
\textit{terminal reward function} $R_T:\SpaceH_T\to \SetR$. We assume 
throughout that there exists a constant
$C\in\SetR_+$, such that $\abs{R_t(h_t,a)} < C$ for all $t\in[T-1]$ and
$(h_t,a)\in\SpaceH_t\times \SpaceA$ as well as $\abs{R_T(h_T)} < C$ for all
$h_T\in\SpaceH_T$. From this assumption it immediately follows that 
\[
    \norm{R_t}_\infty = 
    \sup_{(h_t,a)\in\SpaceH_t\times\SpaceA} \abs{R_t(h_t,a)} \leq C,
\]
as well as $\norm{R_T}_\infty \leq C$.

\begin{remark}
\label{rem:HDPRewardsRemark}
    As in the Markovian case, compare \cref{rem:MDPRewardsRemark}, rewards may
    additionally depend on the environment's observation at the next step, 
    such that the reward at time $t\in[T-1]$ is determined by the measurable 
    reward function $\tilde R_t:\SpaceH_t\times\SpaceA\times\SpaceS\to\SetR$.
    where $\tilde R_t(h_t,a,o)$, gives the reward at time $t$ when the history 
    is $h_t$, action $a$ is selected and the environment transitions to 
    observation $o$ at time $t+1$. In this case we can again define an
    \textit{expected immediate reward function} at $t\in[T-1]$ by
    \begin{equation}
    \label{eq:HDPExpectedImmediateReward}
        R_t(h_t,a) := \int_\SpaceO \tilde R_t(h_t,a,o)\, P_t(\dd o\,|\,h_t,a)
        \quad \text{for all } (h_t,a)\in\SpaceH_t\times\SpaceA.
    \end{equation}
    Since $\tilde R_t$ is measurable, so is $R_t$ for any $t\in[T-1]$ 
    due to \cref{lem:KernelsAuxiliaryMeasurableExpectation}, and the 
    subsequent results in this section remain true if $R_t$ is defined by 
    \eqref{eq:HDPExpectedImmediateReward}. From a reinforcement learning 
    perspective, it is irrelevant if the reward sample $r_t\in\SetR$ that the 
    agent receives at time $t$, depends only on $(h_t,a_t)$ or, additionally 
    on the next observation $o_{t+1}$, since we only require, the value $r_t$
    to be fixed before the agent selects the next action $a_{t+1}$, and that 
    it does not depend on its choice.    
\end{remark}

% alternative reward definition as random variables
\begin{comment}
Similar to the observation transition kernels, we admit the general 
case where rewards are generated by stochastic kernels. For each $t\in[T]$,
the \textit{reward kernel} $\mathcal{R}_{t+1}$ is defined as a stochastic 
kernel from $(\SpaceH_t\times\SpaceA, \mathfrak{B}_{\SpaceH_t}\otimes_\sigma
\mathfrak{B}_\SpaceA)$ to $(\SetR, \mathfrak{B}_\SetR)$, where 
$\mathfrak{B}_\SetR$ is the Borel-$\sigma$-algebra on $\SetR$. Given history
$h_t\in\SpaceH_t$ at time $t$ and that the agent has chosen $a_t\in\SpaceA$ as
next action, the random reward the agent receives upon taking $a_t$ and the
environment transitioning to some $o_{t+1}$ is
\[
    R_{t+1}^{(h_t,a_t)} \sim \mathcal{R}_{t+1}(\,\cdot\,|\,h_t,a_t),
\]
where again $\mathcal{R}_{t+1}(\,\cdot\,|\,h_t,a_t)$ is the probability measure
on $\SetR$ given pair $(h_t,a_t)\in\SpaceH_t\times\SpaceA$. We assume
throughout that there exists a constant $C\in\SetR_+$, such that $\bigl|R_{t+1}^
{(h_t,a_t)}\bigr|<C$ a.s. for all $t\in[T]$ and $(h_t,a_t)\in\SpaceH_t\times
\SpaceA$.

The reward kernels also give rise to the \textit{immediate reward functions}, 
$r_{t+1}:\SpaceH_t\times\SpaceA \to \SetR$. For each $t\in[T]$, the function
$r_{t+1}$ gives the expected reward when choosing action $a_t\in\SpaceA$ given
history $h_t\in\SpaceH_t$,
\[
    r_{t+1}(h_t,a_t) := \E\Bigl[R_{t+1}^{(h_t,a_t)}\Bigr]. 
    %\{for \}\quad t\in[T], (h_t,a_t)\in\SpaceH_t\times\SpaceA.
\]
The assumption that rewards are bounded by $C$, immediately implies 
\[
    \norm{r_{t+1}}_\infty = 
    \sup_{(h_t,a_t)\in\SpaceH_t\times\SpaceA} \abs{r_{t+1}(h_t,a_t)} < C.
\]
In the following, we omit the dependence on $(h_t,a_t)$ and only write 
$R_{t+1}$ for $R_{t+1}^{(h_t,a_t)}$ if it is clear form the context.
\end{comment}

\begin{remark}
    No stationarity or Markov assumptions are made about the observation
    transition kernels $\mathcal{P}$, nor about whether $\SpaceO$ resembles
    partial observations, or the true states itself. This provides that the 
    methods developed in this thesis are readily applicable to partially 
    observable MDPs as environment model or even fully non-Markovian models.
\end{remark}

%------------------------------------------------------------------------------
\paragraph{Policies and the agent's objective}
% policy
The agent's objective is to find a sequence of actions $(a_0,a_1,\dots,a_T)$, 
where at each time $t$ the choice of $a_t$ is dependent on the information 
contained in $h_t$, that maximizes the expected cumulative sum of discounted 
rewards over each episode. Actions are selected according to a policy the 
agent follows.
\begin{definition}
\label{def:HDPPolicy}
    For an HDP \eqref{eq:HDPTuple}, a policy is defined as a family $\pi =
    (\pi_t)_{t\in[T-1]}$, where $\pi$ is a
    \begin{itemize}
        \item \textit{history-dependent, randomized policy}, if $\pi_t$ is a
        stochastic kernel from $(\SpaceH_t,\mathcal{B}_{\SpaceH_t})$ to 
        $(\SpaceA,\mathcal{B}_\SpaceA)$
        
        \item \textit{history-dependent, deterministic policy}, if $\pi_t$ is 
        a measurable map $\pi_t:\SpaceH_t\to\SpaceA$,
    \end{itemize}
    for each $t\in[T-1]$. The sets of all randomized and deterministic 
    history-dependent policies are denoted as $\Pi^{\text{HR}}$ and 
    $\Pi^{\text{HD}}$, respectively. Any $\pi\in\Pi^{\text{HD}}$ can be
    considered as a special case of a randomized policy by setting 
    $\tilde \pi_t(\cdot|h_t) := \delta_{\pi_t(h_t)}(\cdot)$ for each 
    $h_t\in\SpaceH_t$, thus $\Pi^{\text{HD}}\subset\Pi^{\text{HR}}$, and 
    $\pi\in\Pi^\text{HR}$ always includes the case $\pi\in\Pi^\text{HD}$.
\end{definition}
When observation transitions depend on histories, we can not hope for
Markovian optimal policies, and all policies are considered functions of
histories. The existence of at least one such policy is guaranteed by setting
$\pi_t(h_t) = a$ for fixed $a\in\SpaceA$, and all $t\in[T-1]$. Note,
that this is possible because we assumed that the set of possible actions to 
choose from at each step does not depend on the observation the process is in, 
otherwise greater care has to be taken to guarantee the existence of measurable
policies, compare \cite[Chapter~16]{hinderer2016_do}.

A policy $\pi\in\Pi^\text{HR}$ together with a fixed initial probability 
distribution $\mu\in\text{Prob}(\SpaceO)$ induces the canonical probability 
space
\begin{equation}
\label{eq:HDPCanonicalProbSpace}
    (\SpaceH_T, \mathcal{B}_{\SpaceH_T}, \PM^\pi),
\end{equation}
where the measure $\PM^\pi$ is determined as in 
\cref{chapter:PreliminariesStochasticKernels} by the product of the measure 
$\mu$ and the stochastic kernels $(P_t)_{t\in[T-1]}$ and 
$(\pi_t)_{t\in[T-1]}$. We suppress the dependence on the initial distribution 
$\mu$ since it is considered fixed. The measure $\PM^\pi$ satisfies
\begin{equation}
\label{eq:HDPCanonicalMeasure}
\begin{split}
    \PM^\pi [\dd{h_T}] 
    = &\ P_{T-1}(\dd o_T | h_{T-1}, a_{T-1})\, \pi_{T-1}(\dd a_{T-1}|h_{T-1})\,
    P_{T-2}(\dd o_{T-1} | h_{T-2}, a_{T-2}) \\
    %\pi_{T-2}(\dd a_{T-2}|h_{T-2})
    &\dots P_{0}(\dd o_1 | o_0, a_0)\, \pi_0(\dd a_0|o_0)\, \mu(\dd o_0),
\end{split}    
\end{equation}
with $h_T = (o_0, a_0, a_1, \dots, a_{T-1}, o_T)$. We again omit the 
dependence on $\mu$ and $\mathcal{P}$ in $\PM^\pi$, since both are fixed, and 
to highlight that the measure $\PM^\pi$ is determined by the choice of $\pi$.
For any history
$h_t\in\SpaceH_T$, we define the random observations $(O_t)_{t\in[T]}$ and 
random actions $(A_t)_{t\in[T-1]}$ as the coordinate maps on
\eqref{eq:HDPCanonicalProbSpace},
\[
    O_t(h_T) := o_t, \qquad A_t(h_T) := a_t.
\]
The conditional distributions of the random variables $O_t$ and $A_t$ are 
determined by $P_t$ and $\pi_t$, respectively. The process $O_0, A_0, O_1, 
\dots, A_{T-1}, O_T$ is just the canonical process on the probability space
\eqref{eq:HDPCanonicalProbSpace}. Further, we define the history process by
\[
    H_0 := O_0,
    \quad \text{and} \quad
    H_t := (O_0, A_0, O_1, \dots, A_{t-1}, O_t), \text{ for } t\in[T]\setminus\{0\},
\]
where $H_t$ is a $\SpaceH_t$-valued random variable on 
\eqref{eq:HDPCanonicalProbSpace} for each $t\in[T]$. Similar to the Markovian 
case, for any $t\in[T]$, history $h_t\in\SpaceH_t$, Borel set $B\in\mathcal{B}_\SpaceO$, and action $a\in\SpaceA$, the measure $\PM^\pi$ satisfies
\begin{gather}
    \PM^\pi[O_0\in B] = \mu(B), \notag\\
    \label{eq:HDPActionProb}    
    \PM^\pi[A_t = a|H_t = h_t] = \pi(a|h_t), \\
    \PM^\pi[O_{t+1}\in B|H_t=h_t, A_t=a] = P_t(B|h_t,a) \notag.
\end{gather}
The difference to a MDP is apparent in that the conditional 
probabilities given $H_t=h_t$, and potentially $A_t=a$, are functions of the 
full history $h_t$, given by the observation transition and policy kernels.
For a policy $\pi\in\Pi^{\text{HD}}$, the random actions are given as $A_t := 
\pi_t(O_t)$, and \cref{eq:HDPActionProb,eq:HDPCanonicalMeasure} change to
\begin{equation*}
    \PM^\pi[A_t = a|H_t = h_t] = \delta_{\pi(h_t)}(a).
\end{equation*}
and
\begin{equation*}
\begin{split}    
    \PM^\pi [\dd{h_T}] 
    = &\ P_{T-1}(\dd o_T | h_{T-1}, a_{T-1})\, 
        \delta_{\pi_{T-1}(h_{T-1})}(\dd a_{T-1})\,
            P_{T-2}(\dd o_{T-1} | h_{T-2}, a_{T-2}) \\
    &\dots P_{0}(\dd s_1 | o_0, a_0)\, 
        \delta_{\pi_0(o_0)}(\dd a_0)\, 
            \mu(\dd o_0).
\end{split}
\end{equation*}
When $\pi_t$ is deterministic, observation transition kernels only depend on $h_t$ and we define
\[
     P_{\pi_t}(\dd{o} |\,h_t) := P_t(\dd{o} |\, h_t, \pi_t(h_t))
     \quad \text{for } t\in[T-1],
\]
which is a stochastic kernel from $\SpaceH_t$ to $\SpaceO$.

\begin{remark}
For $s,t\in[T]$ with $s<t$, denote by $H_{s:u} := (O_s, A_s, O_{s+1}, \dots, A_{u-1}, O_u)$ the random sub-history from time $s$ until $t$. To ease notation in the following section, we define the stochastic kernels 
\begin{align}
\label{eq:PolicyTransitionKernelProduct}
    (\pi \otimes P)^{(t,T)} 
    & := \bigotimes_{u=t}^{T-1} (\pi_u\otimes P_u), 
    \quad \text{for } t\in[T-1],\\
\label{eq:TransitionKernelPolicyProduct}    
    (P \otimes \pi)^{(t,T)}
    & := P_t \otimes (\pi \otimes P)^{(t+1,T)},
    \quad \text{for } t\in[T-2],
\end{align}
where it follows from \cref{cor:NKernelProduct} that $(\pi \otimes 
P)^{(t,T)}$ is a stochastic kernel from $\SpaceH_t$ to $\SpaceA \times
\SpaceH_{t+1:T}$, and a regular conditional distribution of the random 
sub-history $(A_t, H_{t+1:T})$ for each $t\in[T-1]$. Moreover, $(P \otimes 
\pi)^{(t,T)}$ is a stochastic kernel from $\SpaceH_t\times\SpaceA$ to 
$\SpaceH_{t+1:T}$, and a regular conditional distribution of the random sub-history $H_{t+1:T}$.

\end{remark}

\begin{comment}
\begin{lemma}
\label{lem:MeasureOnHistories}
    Let $P_0\in\text{Prob}(\SpaceO)$ be given and for $t\in[T]_0$, let 
    $\mathcal{P}_t$ be the probability kernel given by
    \cref{def:ObservationTransitionKernel}. Then, for each $t\in[T]$, there exists a
    unique probability measure $\PM_t \in \text{Prob}(\SpaceH_t\times \SpaceA)$, such that
    \begin{multline}
        \PM_t[ O_0\times A_0\times O_1\times \dots 
            \times O_t\times A_{t}] \\
        = \int_{O_0}\int_{A_0}\int_{O_1} \dots \int_{O_t} 
            \pi_t(A_t\,|\,a_0,o_0,a_1,\dots, a_{t-1},o_t) 
                \dd{\mathcal{P}_{t-1}(o_t |\, a_0,o_0,a_1,\dots, a_{t-1})}
                    \\
        \dd{\pi_{t-1}(a_{t-1} |\, a_0,o_0,a_1,\dots,o_{t-1})}
            \cdots \dd{\mathcal{P}_0(o_1 |\, o_0, a_0)}
                \dd{\pi_{0}(a_0 |\,o_0)} \dd{P_0(o_0)}        
    \end{multline} 
\end{lemma}

\begin{proof}
    For the proof we refer to Chapter 10.2, Equation (16) as well as Propositions 7.13 and 7.28 in \cite{bertsekas1996_socdc}.
\end{proof}
\end{comment} 

%------------------------------------------------------------------------------
\section[Value Functions and Optimal Policies in HDPs]{Value Functions and Optimal Policies in History-Based Decision Processes}
\label{chapter:HistoryRLEnvironment:section:ValueQFunctions}
%------------------------------------------------------------------------------

The agent's goal is to find a policy $\pi$ which maximizes the expected total 
discounted reward over an episode. A policy that achieves this goal is called 
an \textit{optimal policy}. As in the MDP case we want to characterize 
optimality of a policy through the notion of value functions. For an HDP, 
those need to be defined in terms of histories.

\begin{definition}[Value functions in HDPs]
\label{def:HDPValueFunctions}
    Let $\pi\in\Pi^\text{HR}$. The value function $V_t^\pi:\SpaceH_t 
    \to \SetR$ at time $t\in[T-1]$ corresponding to $\pi$ in the HDP 
    \eqref{eq:HDPTuple}, is defined for every $h_t\in\SpaceH_t$ by
    \begin{equation}
    \label{eq:HDPValueFunction}
        V^\pi_t(h_t) 
        := \E^\pi\bigg[
            \sum_{u=t}^{T-1} \gamma^{u-t} R_u(H_u, A_u) + \gamma^{T-t} R_T(H_T)
            \,\bigg|\, H_t = h_t
        \bigg]
    \end{equation}
    and for $t=T$ as $V^\pi_T(h_T):=R_T$ for every $h_T\in\SpaceH_T$. It is
    understood that policy $\pi$ is followed from $t$ onwards, and the
    conditional expectation $\E^\pi$ is given on the canonical probability 
    space with probability measure $\PM^\pi$.
\end{definition}

For $t=0$, the quantity $V_0^\pi(o_0)$ provides the expected total discounted 
reward over the whole episode, if the agent follows policy $\pi$, and
the environment starts in observation $o_0\in\SpaceO$. Note, that in this 
setup, with finite time horizon, the discount factor could also be encoded in 
the rewards. Since all reward functions $(R_t)_{t\in[T]}$ are bounded, so is 
$V_t^\pi$ for each policy $\pi\in\Pi^\text{HR}$ and $t\in[T]$. For a fixed 
policy, the corresponding value functions may be computed via the following 
value iteration.

\begin{theorem}[Value iteration in HDPs]
\label{thm:HPDValueIteration}
    Fix a $\pi\in\Pi^\text{HR}$. Then, the corresponding value 
    functions $V_t^\pi$ are $\mathcal{B}_{\SpaceH_t}$-measurable for each 
    $t\in[T]$, and satisfy the recursion
    \begin{equation}
    \label{eq:HDPValueIterationRandomized}
        V^\pi_t(h_t) 
        = \sum_{a\in\SpaceA} \pi_t(a|h_t) \Bigl(
            R_t(h_t, a) + \gamma \, 
            \E_{P_t(\dd{o}|h_t, a)} \bigl[
                V^\pi_{t+1}(h_t, a, O_{t+1})
            \bigr]
        \Bigr), \quad t\in[T-1],        
    \end{equation}
    with the terminal condition $V^\pi_T(h_T) = R_T(h_T)$, and where 
    $\E_{P_t(\dd o|h_t,a)}$ is defined as in 
    \cref{eq:ConditionalExpectationMultipleOne} as the conditional expectation 
    of $V^\pi_{t+1}(H_{t+1}) = V_{t+1}^\pi(H_t,A_t,O_{t+1})$ given $\{H_t=h_t, 
    A_t=a\}$. If $\pi\in\Pi^{\normalfont{\text{HD}}}$ the recursion is
    given as
    \begin{equation}
    \label{eq:HDPValueIterationDeterministic}
        V^\pi_t(h_t) 
        = R_t(h_t, \pi_t(h_t)) + \gamma \, 
        \E_{P_{\pi_t}(\dd{o}|h_t)} \bigl[
            V^\pi_{t+1}(h_t, \pi_t(h_t), O_{t+1})
        \bigr], \quad t\in[T-1],
    \end{equation}
    with the stochastic kernel $P_{\pi_t}(\dd{o}|\,h_t) = P_t(\dd{o}|\,h_t,
    \pi_t(h_t))$.
\end{theorem}
\begin{remark}
    We did not state the value iteration for a given policy in the Markovian case in \cref{chapter:RLAndMDP:section:MDP}, however, we highlight that for an MDP, the value iteration can be stated in the same form as in \cref{eq:HDPValueIterationRandomized,eq:HDPValueIterationDeterministic}, in this case for a fixed Markovian randomized or deterministic policy, respectively. Replace all time-$t$ histories $h_t$ by the state $s_t$, the random observation $O_{t+1}$ by the random state $S_{t+1}$, and let $P_t$ be a Markovian transition kernel, only depending on the pair $(s_t,a)$. Note that in this case, the value function at $t+1$ in the expectation is a function of the current state only and given as $V^\pi_{t+1}(S_{t+1})$.
\end{remark}
\begin{proof}[Proof of \cref{thm:HPDValueIteration}]
    For each $t\in[T-1]$, we define the auxiliary function $f_t:\SpaceH_T \to 
    \SetR$ by
    \begin{equation}
    \label{eq:AuxiliaryRewardFunction}
        f_t(h_T) = f_t(h_{t},a_t,h_{t+1:T}) 
        := R_t(h_t,a_t) + \sum_{u=t+1}^{T-1} 
            \gamma^{u-t} R_u(h_u,a_u) + R_T(h_T)
    \end{equation}
    and for $t=T$ as $f_T(h_T) := R_T(h_T)$ for any history $h_T\in\SpaceH_T$, 
    and where $h_u\in\SpaceH_u$ is the sub-history from $0$ to $u$, and
    $a_u\in\SpaceA$ is the action at time $u$ of the input history $h_T$, for 
    each $u\in\{t, t+1,\dots, T\}$. The function $f_t$ provides the cumulative 
    discounted reward obtained from $t$ onwards, if the action-observation 
    history over the whole episode is $h_T$. Note, that $f_t$ is 
    $\mathcal{B}_{\SpaceH_T}$-measurable for each $t\in[T]$, due to the 
    measurability of the reward functions $R_t$, and that the functions 
    $(f_t)_{t\in[T]}$ satisfy
    \begin{equation}    
    \label{eq:AuxiliaryRewardFunctionRecursion}
        f_t(h_T) = R_t(h_t,a_t) + \gamma\, f_{t+1}(h_T).
    \end{equation}
    Recall, that the stochastic kernel $(\pi\otimes P)^{(t,T)}$ defined in 
    \cref{eq:PolicyTransitionKernelProduct} is a regular conditional 
    distribution of $(A_t,H_{t+1:T})$ given $H_t$. With the help of $f_t$, the 
    kernel $(\pi\otimes P)^{(t,T)}$ and 
    \cref{eq:ConditionalExpectationMultipleTwo} we may express the value 
    function at time $t$ as
    \begin{equation}   
    \label{eq:HDPValueFunctionIntegral}
    \begin{split}
    V^{\pi}_t(h_t) 
    &= \E^\pi[f_t(H_t,A_t,H_{t+1:T}) \,|\, H_t=h_t] \\
    &= \int_{\SpaceA\times\SpaceH_{t+1:T}}
            f_t(h_t, a, h_{t+1:T})\, 
                (\pi\otimes P)^{(t,T)}\bigl(\dd{(a,h_{t+1:T}}) \,|\, h_t\bigr).
    \end{split}
    \end{equation}
    It follows from \cref{lem:KernelsAuxiliaryMeasurableExpectation}, that 
    $V^\pi_t$ is $\mathcal{B}_{\SpaceH_t}$-measurable for each $t\in[T-1]$, 
    and $V^\pi_T$ is $\mathcal{B}_{H_T}$-measurable by definition as $R_T$.
    This means that for $(h_t,a)\in\SpaceH_t\times\SpaceA$, the conditional 
    expectations
    \[
        \E^\pi_{P_t(\dd{o}|h_t,a)}[V^\pi_{t+1}(h_{t},a,O_{t+1})]
        =\E^\pi[V^\pi_{t+1}(h_{t},A_t,O_{t+1})\,|\,H_t=h_t,A_t=a]
    \]
    as in \cref{eq:ConditionalExpectationMultipleOne} is well-defined, and 
    moreover, can be expressed as integral with respect to the kernel 
    $(P\otimes\pi)^{t+1,T}$ defined in 
    \eqref{eq:TransitionKernelPolicyProduct}, which is a regular conditional 
    distribution of $H_{t+1:t}$ given $(H_t,A_t)$.

    For the recursion, at $t=T$, we get by definition the terminal condition. 
    Now let $t\in[T-1]$. We want to apply \cref{thm:KernelFubini} to the 
    measure $(\pi\otimes P)^{(t,T)}(\cdot|h_t)$ on $\SpaceA\times\SpaceH_t$. 
    Fix input $h_t\in\SpaceH_t$, and define the probability measure $\mu_1$ on 
    $\SpaceA$ as $\mu_1 := \pi(\cdot|h_t)$ and the stochastic kernel 
    $\kappa_1$ as
    \begin{align*}
        \kappa_1: 
            \SpaceA \times \mathcal{B}_{\SpaceH_{t+1:T}} &\to [0,1], \\
        (a, F) &\mapsto \kappa_1(a,F) 
            := (P\otimes\pi)^{(t,T)}(F\,|\,h_t,a),
    \end{align*}
    where again $\kappa_1(\dd{h_{t+1:T}}|\,a) :=  \kappa_1(a,
    \dd{h_{t+1:T}})$ for the probability measure on $\SpaceH_{t+1:T}$ for
    action $a\in\SpaceA$. Then, by definition 
   \[
        (\pi\otimes P)^{(t,T)}(\cdot\,|\,h_t) 
        = \mu_1 \otimes \kappa_1,
    \]
    and from \cref{eq:HDPValueFunctionIntegral}, and application of 
    \cref{thm:KernelFubini} to the product $\mu_1\otimes\kappa_1$, we obtain 
    \begin{align}
        V_t^\pi(h_t) 
        &= \E^\pi[f_t(H_T)|H_t=h_t] 
        = \int_{\SpaceA \times \SpaceH_{t+1:T}} f_t(h_t,\cdot, \cdot)
            \dd{(\mu_1\otimes\kappa_1)} 
        \nonumber \\
        &= \sum_{a\in\SpaceA} \int_{\SpaceH_{t+1:T}} 
            f_t(h_t,a,h_{t+1:T}) \, 
                \kappa_1(\dd{h_{t+1:T}} |\, a)\, \mu_1(a) 
        \nonumber \\
        &= \sum_{a\in\SpaceA} \pi_t(a|h_t) \biggl(
            \int_{\SpaceH_{t+1:T}} 
                f_t(h_t,a,h_{t+1:T}) \, 
            \kappa_1(\dd{h_{t+1:T}} |\, a) 
        \biggr)
        \nonumber \\        
        &= \sum_{a\in\SpaceA} \pi_t(a|h_t) \biggl(
            R_t(h_t, a) + \gamma 
            \int_{\SpaceH_{t+1:T}} f_{t+1}(h_t, a, h_{t+1:T}) \,
                \kappa_1(\dd{h_{t+1:T}} |\, a)    
        \biggr) 
        \label{eq:ProofValueIterationFirstStep}
    \end{align}    
    where the second equality follows from inserting the definition of 
    $\mu_1$, and the last equality from the recursive property
    \eqref{eq:AuxiliaryRewardFunctionRecursion} of the function $f_t$ and the
    linearity of the integral. By definition of $\kappa_1$ and 
    $(P\otimes\pi)^{(t,T)}$ we get
    \[
        \kappa_1(\cdot \,|\,a) 
        = \bigl(P_t\otimes (\pi\otimes P)^{(t+1,T)}\bigr)(
            \cdot\,|\, h_t,a),
    \]
    for each $a\in\SpaceA$ and fixed $h_t\in\SpaceH_t$. We want to apply 
    \cref{thm:KernelFubini} again to the probability measure $\kappa(\cdot|a)$ 
    on $\SpaceH_{t+1:T} = \SpaceO\times \SpaceA\times \SpaceH_{t+2:T}$.
    For fixed $(h_t,a)\in\SpaceH_t \times\SpaceA$, define the probability 
    measure $\mu_2$ on $\SpaceO$ as $\mu_2 := P_t(\cdot\,|\, h_t, a_t)$, and 
    the stochastic kernel $\kappa_2$ as
    \begin{align*}
        \kappa_2: 
        \SpaceO \times \mathcal{B}_{\SpaceA\times\SpaceH_{t+2:T}} 
            &\to [0,1], \\
        (o, F) &\mapsto \kappa_2(o,F) 
            := (\pi\otimes P)^{(t+1,T)}(F\,|\,(h_t, a, o))
    \end{align*}
    Then, $\kappa_1(\cdot\,|\,a_t) = \mu_2 \otimes \kappa_2$, and
    application of \cref{thm:KernelFubini} to the product $\mu_2 \otimes 
    \kappa_2$ yields for the integrals in 
    \eqref{eq:ProofValueIterationFirstStep},
    \begin{align}
        \int_{\SpaceH_{t+1:T}} & f_{t+1}(h_t, a, h_{t+1:T}) \,
            \kappa_1(\dd{h_{t+1:T}} |\, a) 
        = \int_{\SpaceH_{t+1:T}} f_{t+1}(h_t, a, \cdot)
            \dd{(\mu_2 \otimes \kappa_2)} 
        \nonumber \\
        &= \int_\SpaceO \int_{\SpaceA\times\SpaceH_{t+2:T}}
                f_{t+1}(h_t,a,o,a',h_{t+2:T}) \,
                    \kappa_2\bigl(\dd{(a',h_{t+2:T}})\,|\,o\bigr) \,
                        \mu_2(\dd{o})
        \nonumber \\                        
        &= \int_\SpaceO V_{t+1}^\pi(h_t,a_t,o) \, \mu_2(\dd{o}) 
        = \E_{P_t(\dd{o}|h_t,a_t)} \bigl[ 
                V_{t+1}^\pi(h_t, a_t , O)
            \bigr],
        \label{eq:ProofValueIterationSecondStep}
    \end{align}
    where the second to last equality follows from expression
    \eqref{eq:HDPValueFunctionIntegral} for the value
    function $V_{t+1}^\pi$ at $t+1$ with input history $(h_t,a,o)$, 
    as integral w.r.t to $\kappa_2(\cdot\,|\,o) = (\pi\otimes P)^{(t+1,T)}
    (\cdot\,|\,(h_t, a,o))$ over $\SpaceA\times\SpaceH_{t+2:T}$ with the 
    auxiliary function $f_{t+1}$, and the last equality follows by definition 
    of $\mu_2=P_t(\cdot|h_t,a)$ and that of the conditional expectation 
    $\E_{P_t(\dd{o}| h_t, a_t)}$ as in 
    \cref{eq:ConditionalExpectationMultipleOne}. Since $h_t\in\SpaceH_t$ was 
    arbitrary, the claim for $\pi\in\Pi^\text{HR}$ follows by plugging 
    \eqref{eq:ProofValueIterationSecondStep} into 
    \eqref{eq:ProofValueIterationFirstStep}.
    
    If $\pi'\in\Pi^{\text{HD}}$, we interpret it as randomized policy by 
    setting $\tilde\pi_t(a|h_t) := \delta_{\pi'_t(h_t)}(a)$ for each 
    $t\in[T-1]$ and any $(h_t,a)\in \SpaceH_t\times\SpaceA$. Substituting 
    $\pi(a|h_t)$ in expression \eqref{eq:HDPValueIterationRandomized} by 
    $\tilde\pi_t(a|h_t)$ yields recursion 
    \eqref{eq:HDPValueIterationDeterministic} for deterministic $\pi'$.    
\end{proof}

%%%%%%%%%%%%%%%%
% value iteration theorem for deterministc policy only
\begin{comment}
\begin{theorem}[Value iteration for deterministic policy]
\label{thm:HPDValueIteration}
    For a fixed policy $\pi\in\Pi^{HD}$, the corresponding value functions
    $V_t^\pi, t\in[T]$ satisfy the recursion
    \begin{equation}
    \label{eq:HDPValueIterationPolicy}
        V^\pi_t(h_t) 
        = R_t(h_t, \pi_t(h_t)) + \gamma \, 
        \E_{P_{\pi_t}(\dd{o}|h_t)} \bigl[
            V^\pi_{t+1}(h_t, \pi_t(h_t), O)
        \bigr], \quad t\in[T-1],
    \end{equation}
    where $O\sim P_{\pi_t(\dd{o}|h_t)}$ and with terminal condition 
    $ V^\pi_T(h_T) = R_T(h_T)$.
\end{theorem}
\begin{proof}
    If $t=T$ we get the terminal condition. Now assume $t\in[T-1]$. Since $\pi_t$
    is given as deterministic function $\pi_t:\SpaceH_t\to\SpaceA$, we
    understand the policy at time $t$ as the stochastic kernel 
    $\tilde\pi_t(h_t,a) := \delta_{\pi_t(h_t)}(a)$ for any $(h_t,a)\in \SpaceH_t
    \times\SpaceA$ at time $t$. For fixed input $h_t\in\SpaceH$, the value
    function $V_t^\pi(h_t)$ is defined as integral w.r.t to the probability
    measure $(\tilde\pi\otimes P)^{(t,T)}(\cdot\,|\,h_t)$ on $\SpaceA \times 
    \SpaceH_{t+1:T}$. With the probability measure $\mu_1 := \delta_{\pi_t(h_t)}$
    on the discrete space $\SpaceA$ and the stochastic kernel $\kappa_1$ as
    \begin{align*}
        \kappa_1: 
        \SpaceA \times \mathcal{B}_{\SpaceH_{t+1:T}} &\to [0,1], \\
        (a, F) &\mapsto \kappa_1(a,F) := (P\otimes\tilde\pi)^{(t,T)}(F\,|\,h_t,a)
    \end{align*}
    where we again denote $\kappa_1(F\,|\,a) :=  \kappa_1(a,F)$, we have 
    \[
        (\tilde\pi\otimes P)^{(t,T)}(\cdot\,|\,h_t) 
        = \mu_1 \otimes \kappa_1
    \]
    and application of \cref{thm:KernelFubini} to the product $\mu_1 \otimes
    \kappa_1$ yields
    \begin{align*}
        V_t^\pi(h_t) 
        &= \E^\pi[f_t(H_T)|H_t=h_t] \\ 
        &= \int_{\SpaceA \times \SpaceH_{t+1:T}} f_t \dd{\mu_1\otimes\kappa_1} \\
        &= \sum_{a\in\SpaceA} \int_{\SpaceH_{t+1:T}} 
            f_t(h_t,a,h_{t+1:T}) \, 
            \kappa_1(\dd{h_{t+1:T}}\,|\, a)\, \mu_1(\{a\}) \\
        &= \int_{\SpaceH_{t+1:T}} f_t(h_t,\pi_t(h_t),h_{t+1:T})) \,
            \kappa_1(\dd{h_{t+1:T}}\, | \, \pi_t(h_t)) \\
        &= R_t(h_t, \pi_t(h_t)) + \gamma 
            \int_{\SpaceH_{t+1:T}} f_{t+1}(h_t,\pi_t(h_t),h_{t+1:T})) \,
                \kappa_1(\dd{h_{t+1:T}}\, | \, \pi_t(h_t))                
    \end{align*}
    where the last equality follows from the recursive property
    \eqref{eq:AuxiliaryRewardFunctionRecursion} of the function $f_t$ and the
    linearity of the integral. Observe that $\kappa_1 = (P\otimes\tilde
    \pi)^{(t,T)} = P_t\otimes (\tilde\pi\otimes P)^{(t+1,T)}$. For fixed
    $h_t\in\SpaceH_t$, with the probability measure $\mu_2 := P_{\pi_t}(\cdot\,|
    \,h_t) = P_t(\cdot\,|\,h_t,\pi_t(h_t))$ on $\SpaceO$ and the stochastic
    kernel $\kappa_2$ given as
    \begin{align*}
        \kappa_2: 
        \SpaceO \times \mathcal{B}_{\SpaceA\times\SpaceH_{t+2:T}} &\to [0,1], \\
        (o, F) &\mapsto \kappa_2(o,F) 
            := (\tilde\pi\otimes P)^{(t+1,T)}(F\,|\,h_t, \pi_t(h_t), o),
    \end{align*}
    we have that $\kappa_1(\cdot\,|\,\pi_t(h_t)) = \mu_2 \otimes \kappa_2$ and
    anew application  of \cref{thm:KernelFubini} yields for the above integral
    \begin{align*}
        \int_{\SpaceH_{t+1:T}} 
            &f_{t+1}(h_t,\pi_t(h_t),h_{t+1:T})) \,
            \kappa(\dd{h_{t+1:T}}\, | \, \pi_t(h_t)) \\                    
        &= \int_\SpaceO \int_{\SpaceA\times\SpaceH_{t+2:T}}
                f_{t+1}(h_t,\pi_t(h_t),o,a,h_{t+2:T}) \,
                    \kappa_2(\dd{(a,h_{t+2:T}})\,|\,o) \,
                        \mu_2(\dd{o}) \\
        &= \int_\SpaceO V_{t+1}^\pi(h_t,\pi_t(h_t),o) \, \mu_2(\dd{o}) 
        = \E_{P_{\pi_t}(\dd{o}|h_t)} \bigl[ 
                V_{t+1}^\pi(h_t, \pi_t(h_t) , O)
            \bigr]
    \end{align*}
    where the second to last equality follow from the definition of the value
    function $V_{t+1}^\pi(h_t,\pi_t(h_t),o)$ as integral w.r.t to 
    $\kappa_2(\cdot\,|\,o) = (\tilde\pi\otimes P)^{(t+1,T)}(\cdot\,|\,h_t,
    \pi_t(h_t),o)$ over $\SpaceA\times\SpaceH_{t+2:T}$ and the last equality
    from the definition of the conditional expectation $\E_{P_{\pi_t}(\dd{o}|
    h_t)}$ and $\mu_2$. The claim follows.
\end{proof}
\end{comment}
%%%%%%%%%%%%%%%%

Analogously to the Markovian case, we define a policy to be optimal if it obtains
the maximal attainable expected total discounted reward over the episode. 
\begin{definition}[Optimal policies in HDPs]
\label{def:HDPOptimalPolicy}
    A policy $\pi^*\in\Pi^\text{HR}$ is called optimal for the HDP 
    \eqref{eq:HDPTuple}, if it satisfies
    \[
        V^{\pi^*}_0(o) \geq V^\pi_0(o),
        \quad \text{for all } \pi\in\Pi^\text{HR}, o\in\SpaceO.
    \]
\end{definition}

The agent's goal is to find such an optimal policy. We will characterize an
optimal policy with the help of optimal value functions.  

\begin{definition}[Optimal value functions in HDPs]
\label{def:HDPOptimalValueFunctions}
    The optimal value function $V^*_t:\SpaceH_t\to\SetR$ at time $t\in[T-1]$ 
    for the HDP \eqref{eq:HDPTuple} is defined as 
    \begin{equation}
    \label{eq:HDPOptimalValueFunction}
        V_t^*(h_t) := \sup_{\pi\in\Pi^\text{HR}} V^\pi_t(h_t),
        \quad \text{for } h_t\in\SpaceH_t,
    \end{equation}
    and at time $t=T$ as $V_T^*(h_T) := R_T(h_T)$ for $h_T\in \SpaceH_T$.  
\end{definition}
The optimal value function at time $t$ gives the maximal attainable value over
the remaining episode, and thus, at $t=0$ the maximum attainable value over the
whole decision period. Due to the boundedness of reward functions, for any 
$t\in[T]$, the value $V_t^\pi$ is bounded uniformly in $\pi$, which implies 
the boundedness of $V^*_t$. 
%Furthermore, since any policy may only select from a finite number of action sequences, the supremum is attained 
In \cref{thm:HDPOptimalValueIteration} we will show, that the
value functions corresponding to an optimal policy $\pi^*$ coincide with the
optimal value functions, and that optimal policies are deterministic and can 
be characterized as maximizing policies. For $\pi\in\Pi^\text{HD}$, we say 
that $\pi_t$ is a \textit{maximizer} at time $t\in[T-1]$ if
\begin{equation}
\label{eq:HDPMaximizerPolicy}    
    \begin{aligned}
        R_t(h_t,\pi^*_t(h_t)) 
        &+ \gamma \, \E_{P_{\pi_t}(\dd{o}|h_t)} \bigl[ 
            V^*_{t+1}(h_t,\pi^*_t(h_t), O_{t+1}) 
        \bigr] \\  
        &= 
        \max_{a\in\SpaceA} \Bigl\{
            R_t(h_t, a) + \gamma \,
            \E_{P_t(\dd{o}|h_t,a)} \bigl[ V^*_{t+1}(h_t, a, O_{t+1}) \bigr]
        \Bigr\},
    \end{aligned}
\end{equation}
whenever the map $V^*_{t+1}$ is $\mathcal{B}_{\SpaceH_{t+1}}$-measurable such 
that the conditional expectation $\E_{P_t(\cdot|h_t,a)}$ is well-defined 
according to \cref{eq:ConditionalExpectationMultipleOne}. A policy
which is a maximizer at each $t\in[T-1]$ is called a \textit{maximizing
policy}. The value functions corresponding to a maximizing policy satisfy the following optimality equations.
\begin{theorem}[Optimal value iteration in HDPs]
\label{thm:HDPOptimalValueIteration}
    Let $\pi^*\in\Pi^\text{\normalfont{HD}}$ be a maximizing policy. Then,
    $V^*_t = V^{\pi^*}_t$ for $t\in[T]$, specifically, $V^{\pi^*}_0 = V^*_0$, 
    which implies $\pi^*$ is an optimal policy. Furthermore, $V_t^*$ is 
    $\mathcal{B}_{\SpaceH_t}$-measurable for each $t\in[T]$, and the 
    \textit{optimal value iteration}
    \begin{equation}
    \label{eq:HDPOptimalValueIteration}
        V^*_t(h_t)
        =
        \max_{a\in\SpaceA} \Bigl\{
            R_t(h_t,a) + \gamma \,
            \E_{P_t(\dd{o}|h_t,a)} \bigl[ V^*_{t+1}(h_t, a, O_{t+1}) \bigr]
        \Bigr\},
        \quad t\in[T-1],
    \end{equation}
    holds for any with terminal condition $V^*_T(h_T) = R_T(h_T)$.
\end{theorem}
\begin{remark}
    The optimality equations \eqref{eq:HDPOptimalValueIteration}, are an 
    analogue of the Bellman optimality equations 
    \eqref{eq:MDPBellmanValue} for an MDP, where the state $s_t$ is 
    defined as the full history up to time $t$. If we replace the time-$t$ 
    history $h_t$ by a state $s_t$, the history-dependent transition 
    kernel $P_t(\cdot | h_t,a)$ by a Markovian kernel $P_t(\cdot |
    s_t,a)$, and the triplet $(h_t,a, O_{t+1})$ in the expectation by the 
    random state $S_{t+1}$, \cref{eq:HDPOptimalValueIteration} becomes the 
    Bellman optimality equation \eqref{eq:MDPBellmanValue} at time $t$, 
    where the respective Markovian optimal value functions are functions 
    of the current state only.
\end{remark}
For the proof of \cref{thm:HDPOptimalValueIteration} we need the following
lemma.
% the following fact is not needed
\begin{comment}
as well as the fact that for $t\in[T-1]$ and a function $g:
\SpaceH_t\times\SpaceA \to \SetR$, due to $\abs{\SpaceA}<\infty$,
\begin{equation}
\label{eq:SupPolicyMaxActions}
    \sup_{\pi\in\Pi} g(h_t,\pi_t(h_t)) 
    = \max_{a\in\SpaceA} g(h_t,a).
\end{equation}
\end{comment}
%
\begin{lemma}
\label{lem:HDPOptimalValueIterationAuxiliary}
    Let $t\in[T-1]$, and assume that $V^*_{t+1}$ is 
    $\mathcal{B}_{\SpaceH_{t+1}}$-measurable. Then, for any $h_t\in\SpaceH_t$,
    \begin{equation}
        V_t^*(h_t) \leq 
        \max_{a\in\SpaceA} \Bigl\{
            R_t(h_t,a) + \gamma \,
            \E_{P_t(\dd{o}|h_t,a)} \bigl[ V^*_{t+1}(h_t, a, O_{t+1}) \bigr]
        \Bigr\}.
    \end{equation}
\end{lemma}
\begin{proof}
    Fix a policy $\pi\in\Pi^\text{HR}$. By definition we have $V^\pi_t\leq
    V^*_t$, which, for $(h_t,a)\in\SpaceH_t\times\SpaceA$ implies 
    \[
        \E_{P_t(\dd{o}|h_t,a)}[V^\pi_t(h_t,a,O_{t+1})]
        \leq
        \E_{P_t(\dd{o}|h_t,a)}[V^*_t(h_t,a,O_{t+1})].        
    \]
    Application of the value iteration in \cref{thm:HPDValueIteration}
    with policy $\pi$ in the first step, and above inequality in the second 
    step yields
    \begin{align*}
        V^\pi_t(h_t) 
        &= \sum_{a\in\SpaceA} \pi_t(a|h_t) \Bigl(
            R_t(h_t, a) + \gamma \,
                \E_{P_{t}(\dd{o}|h_t, a)} \bigl[ 
                    V^\pi_{t+1} (h_t, a, O_{t+1})
                \bigr]
            \Bigr) \\
        &\leq \sum_{a\in\SpaceA} \pi_t(a|h_t) \Bigl(
            R_t(h_t, a) + \gamma \,
                \E_{P_{t}(\dd{o}|h_t, a)} \bigl[ 
                    V^*_{t+1} (h_t, a, O_{t+1})
                \bigr]
            \Bigr) \\   
        &\leq \max_{a\in\SpaceA}\, \Bigl\{
            R_t(h_t, a) + \gamma \,
                \E_{P_{t}(\dd{o}|h_t, a)} \bigl[ 
                    V^*_{t+1} (h_t, a, O_{t+1})
                \bigr]
            \Bigr\},      
        \end{align*}
    where the last inequality follows from the fact that $\pi_t(a|h_t)$ are 
    probabilities and sum to one. If $\pi$ is deterministic, we again 
    interpret it as kernels $\delta_{\pi_t}$, and in the above sums only the 
    term with action $a=\pi_t(h_t)$ remains. Since $\pi$ was arbitrary, the 
    inequality also holds if we take the supremum on the left hand side.
\end{proof}
\begin{proof}[Proof of \cref{thm:HDPOptimalValueIteration}]
    The proof is carried out by backward induction. 
    
    For $t=T$, by definition 
    $V^{\pi^*}_T = R_T = V^*_T$, which yields the terminal condition. By 
    assumption on $R_T$, the optimal value function $V^*_T = R_T$ is 
    $\mathcal{B}_{\SpaceH_T}$-measurable, and as such, the conditional 
    expectation
    \[
        \E^\pi_{P_{T-1}(\dd{o}|h_{T-1},a)}[
            V^*_T(h_{T-1}, a, O_{t+1})
        ]
        = \E^\pi[V^*_T(H_{T-1},A_{T-1},O_T\,|\, H_{t-1}=h_{t-1}, A_{T-1}=a]
    \]
    is well defined for any $(h_{T-1},a)\in\SpaceH_{T-1}\times\SpaceA$. 

    For $t\in[T-1]$, assume that $V^*_{t+1} = V^{\pi^*}_{t+1}$, and that 
    $V^*_{t+1}$ is $\mathcal{B}_{\SpaceH_t}$-measurable.\footnote{
    Measurability of $V^*_{t+1}$ actually follows from the equality $V^*_{t+1}=V^{\pi^*}_{t+1}$ and the measurability of $V^{\pi^*}_{t+1}$.
    }
    Then, the conditional expectation 
    \[
        \E^\pi_{P_{t+1}(\dd{o}|h_{t+1},a)}[V^*_{t+1}(h_{t+1}, a, O_{t+2})]
    \]
    is well-defined according to \cref{eq:ConditionalExpectationMultipleOne}, 
    and application of \cref{lem:HDPOptimalValueIterationAuxiliary}, together 
    with the definition of the maximizing policy $\pi^*$ in 
    \cref{eq:HDPMaximizerPolicy}, the induction hypothesis for $t+1$, and the 
    value iteration \eqref{eq:HDPValueIterationDeterministic} applied to the 
    value functions corresponding to $\pi^*$, yield
    \begin{align*}
        V^*_t(h_t) 
        & \leq \max_{a\in\SpaceA} \Bigl\{ R_t(h_t,a) + \gamma\,
            \E_{P_t(\dd{o}|h_t,a)}\bigl[ V^*_{t+1}(h_t,a,O_{t+1}) \bigr]
        \Bigr\} \\
        & = R_t(h_t,\pi^*_t(h_t)) + \gamma\,
            \E_{P_{\pi^*_t}(\dd{o}|h_t)}\bigl[ 
                V^*_{t+1}(h_t,\pi^*_t(h_t),O_{t+1})
            \bigr]\\
        & = R_t(h_t,\pi^*_t(h_t)) + \gamma\,
            \E_{P_{\pi^*_t}(\dd{o}|h_t)}\bigl[
                V^{\pi^*}_{t+1}(h_t,\pi^*_t(h_t),O_{t+1})
            \bigr] \\
        & = V^{\pi^*}_t(h_t)
        \leq \sup_{\pi\in\Pi} V^{\pi}_t(h_t) = V^*_t(h_t),
    \end{align*}
    which shows the equality $V^*_t=V^{\pi^*}_{t+1}$, and that the optimal 
    value iteration \eqref{eq:HDPOptimalValueIteration} holds for the 
    considered step $t$. Moreover, since the value function $V^{\pi^*}_{t+1}$ 
    for deterministic policy $\pi^*$ is $\mathcal{B}_{\SpaceH_t}$-measurable, 
    so is $V^*_t$, which concludes the induction step. For $t=0$, equality $V^*_0=V_0^{\pi^*}$, shows that the maximizing policy $\pi^*$ is optimal.
\end{proof}

For a maximizing policy $\pi^*$ to be a policy in the sense of
\cref{def:HDPPolicy} ,and for the application of 
\cref{thm:HDPOptimalValueIteration}, at step $t\in[T-1]$, $\pi^*_t$ not only 
has to select an action that maximizes the right hand side of
\eqref{eq:HDPMaximizerPolicy}, but has to do so in a measurable way.

\begin{theorem}[Existence of maximizing policy]
\label{thm:HDPOptimalPolicyExistence}
    For the HDP \eqref{eq:HDPTuple}, there exists a maximizing policy $\pi^*$ 
    \eqref{eq:HDPMaximizerPolicy}.
\end{theorem}
\begin{proof}
    For $t\in[T-1]$, the reward function $R_t$ is by assumption 
    $\mathcal{B}_{\SpaceH_t}\otimes \mathcal{B}_\SpaceA$-measurable. By 
    \cref{lem:KernelsAuxiliaryMeasurableExpectation}, and the expression of 
    the conditional expectation as integral, the same holds true for the map
    \[
        (h_t,a)\mapsto \E_{P_t(\dd{o}|h_t,a)}[V^*_{t+1}(h_t,a,O_{t+1})],
    \]
    and thus, also the for the function $g:\SpaceH_t\times\SpaceA$ defined as
    \[
        g(h_t,a) := 
        R_t(h_t,a) + \gamma\, \E_{P_t(\dd{o}|h_t,a)}[V^*_{t+1}(h_t,a,O)].
    \]
    It follows from the Measurable Selection Theorem in 
    \cref{appendix:MeasurableSelection}, by setting $(\Omega, \mathcal{F}) = 
    (\SpaceH_t,\mathcal{B}_{\SpaceH_t})$ and with finite action
    space $\SpaceA$, that there exists a maximizer $\pi^*_t:\SpaceH_t\to
    \SpaceA$ of $g$, which satisfies $\pi^*(h_t)\in\argmax_{a\in\SpaceA} 
    g(h_t,a)$ for each $h_t\in\SpaceH_t$, and that the map
    \[
        h_t \mapsto g^*(h_t) := \max_{a\in\SpaceA} g(h_t,a) = g(h_t,\pi^*(h_t))
    \]
    is $\mathcal{B}_{\SpaceH_t}$-measurable.
\end{proof}

Under certain conditions, one could in principle find optimal value
functions and optimal policies explicitly via dynamic programming. Assume 
that, in addition to $\SpaceA$, the observation space $\SpaceO$ is finite, and 
that the transition kernel $P_t$ and the reward function $R_t$ are known at 
each step $t$. This allows one to solve the optimal value iteration
\eqref{eq:HDPOptimalValueIteration} through evaluating for each history $h_t$ 
the term in the curly brackets, and subsequently selection an action that 
maximizes it, which gives the optimal value function at step $t$. By keeping a 
record of the set of maximizing actions at each step $t\in[T]$, and for all 
histories $h_t$, optimal policies can be derived by choosing from those sets. 
However, evaluating \eqref{eq:HDPOptimalValueIteration}, would mean to compute 
the expectation
\[
    \sum_{o\in\SpaceO} P_t(o|h_t,a) V^*_{t+1}(h_t,a,o)
\]
for each possible history at each step. With $\abs{\SpaceA}=n\in\SetN$ and 
$\abs{\SpaceO} = m\in\SetN$, the number of possible histories over the whole 
episode is $\abs{\SpaceH_T} = (mn)^Tm$. Thus, computing the expectation at time
$t\in[T-1]$ requires 
$
    \abs{\SpaceH_t}\abs{\SpaceO}\abs{\SpaceA} 
    = (\abs{\SpaceO}\abs{\SpaceA})^t \abs{\SpaceO}\abs{\SpaceO}\abs{\SpaceA} 
    = m^{t+2}\,n^{t+1}
$ multiplications, and the whole dynamic programming procedure amounts to 
\[
    \sum_{t=0}^{T-1} \abs{\SpaceH_t} \abs{\SpaceO} \abs{\SpaceA}
    = \sum_{t=0}^{T-1} m^{t+2}\,n^{t+1}
\]
operations, rendering this approach infeasible even for small action and
observation spaces, and short episodes length $T$.

 
%%%%%%%%%%%%%%%%%%%%%%%%%%%%%%%%%%%%%%%%%%%%%%%%%%%%%%%%%%%%%%%%%%%%%%%%%%%%%%%
\chapter{Signature-Q-Learning} 
\label{chapter:SignatureQLearning}
%%%%%%%%%%%%%%%%%%%%%%%%%%%%%%%%%%%%%%%%%%%%%%%%%%%%%%%%%%%%%%%%%%%%%%%%%%%%%%%

We now turn to the question of how to find a maximizing policy 
\eqref{eq:HDPMaximizerPolicy} in a reinforcement learning setting, in which 
transition probabilities and reward functions are unknown. This will again 
require Q-functions, which in the history-based setting, will be defined as 
functions of histories, and we provide an explicit method, 
\textit{Signature-Q-Learning}, to learn estimates approximating them.

%------------------------------------------------------------------------------
\section{Q-Functions in History-Based Decision Processes}
\label{chapter:SignatureQLearning:section:QFunctionsInHDP}
%------------------------------------------------------------------------------

Consider a HDP given by \cref{eq:HDPTuple}. Based on the optimal value 
functions \eqref{eq:HDPOptimalValueFunction} we define optimal 
Q-functions for history-action pairs $(h_t,a)\in\SpaceH_t \times \SpaceA$.

\begin{definition}[Optimal Q-functions in HDPs]
\label{def:HDPQFunctions}
    For an HDP given by \cref{eq:HDPTuple}, the \textit{optimal Q-function} 
    $Q^*_t:\SpaceH_T\times\SpaceA\to\SetR$ at time $t\in[T-1]$ is defined as 
    \begin{equation}
    \label{eq:HDPOptimalQFunction}
        Q_t^*(h_t,a) 
        := R_t(h_t,a) + \gamma\, \E_{P_t(\dd{o}|h_t,a)}\bigl[
            V^*_{t+1}(h_t, a, O_{t+1})
        \bigr], 
    \end{equation}
    for any $(h_t,a)\in\SpaceH_t\times\SpaceA$, and for $t=T$ as $Q_T(h_T,a)
    := R_T(h_T)$ for any $(h_T,a)\in\SpaceH_T\times\SpaceA$.
\end{definition}

The optimal Q-value $Q^*_t(h_t,a)$, gives the maximal attainable value over
the remaining episode, if action $a$ is chosen in history $h_t$ at time $t$.
This value is decomposed into the reward generated from the pair $(h_t,a)$ at
time $t$, and the expected optimal value from following an optimal policy from 
$t+1$ onwards. The optimal value functions and optimal Q-functions are 
connected in the following way,
\begin{equation} 
\label{eq:HDPValueQConnection}
    V^*_t(h_t) = \max_{a\in\SpaceA} Q^*_t(h_t, a),
\end{equation}
and by the definition of $Q^*_t$ and the characterization as maximizing policy 
in \cref{thm:HDPOptimalValueIteration}, an optimal policy $\pi^*$ satisfies
\begin{equation}
\label{eq:HDPGreedyPolicy}
    \pi^*_t(h_t) \in \argmax_{a\in\SpaceA} Q^*_t(h_t,a),
    \quad \text{for } h_t\in\SpaceH_t,
\end{equation}
and each $t\in[T-1]$. A policy that satisfies \eqref{eq:HDPGreedyPolicy} for 
all $h_t\in\SpaceH_t$ at each $t\in[T-1]$, is called \textit{greedy} with 
respect to the collection $(Q^*_t)_{t\in[T]}$. Analogue to the MDP case, 
\eqref{eq:HDPGreedyPolicy} implies that knowledge of $Q^*_t$ is sufficient for 
acting optimally in the underlying HDP, by following a policy that is greedy 
w.r.t $Q^*$. Substituting \eqref{eq:HDPValueQConnection} for $V^*_{t+1}$ in 
\cref{eq:HDPOptimalQFunction} yields the \textit{optimal Q-iteration in
HDPs},
\begin{equation}
\label{eq:HDPOptimalQIteration}
    Q^*_t(h_t,a) 
    = R_t(h_t,a) + \gamma \, \E_{P_t(\dd{o}|h_t,a)}\Bigl[
        \max_{a'\in\SpaceA} Q^*_{t+1}((h_t, a, O_{t+1}), a')
    \Bigr],
\end{equation}
for each $(h_t,a)\in\SpaceH_t\times\SpaceA$, with the terminal condition
$Q^*_T(h_T,a) = R_T(h_T)$ for each $(h_T,a) \in\SpaceH_T \times \SpaceA$.
In contrast to the optimal value iteration
\eqref{eq:HDPOptimalValueIteration}, the maximum in the optimal Q-iteration is 
taken inside the expectation. This enables us to sample 
\cref{eq:HDPOptimalQIteration} through sampling histories $h_T$ by interaction 
with the environment, and utilize those samples to learn estimates of 
$(Q^*_t)_{t\in[T]}$ exploiting their recursive structure. If the learned 
estimates are close to the true optimal Q-functions, a policy that is greedy 
with respect to the estimates is close to optimal. The optimal Q-iteration 
will thus be the basis for updating the estimates in the proposed method.

%------------------------------------------------------------------------------
\section{The Signature-Q-Learning Algorithm}
\label{chapter:SignatureQLearning:section:Algorithm}
%------------------------------------------------------------------------------

This section introduces the reinforcement learning method 
\textit{Signature-Q-Learning}, details the method's approximation argument and provides 
an algorithm for the procedure. Analogous to Q-Learning in MDPs, 
Signature-Q-Learning learns an estimate of the optimal Q-functions 
$(Q_t^*)_{t\in[T]}$, albeit in history-dependent domains modeled through 
HDPs. In this setting function approximation is crucial, since an enumeration 
of the high-dimensional inputs is not feasible, even if the observation space 
was finite. Given the signature's properties, detailed in 
\cref{chapter:Signature:section:BoundedVariation}, it is a natural 
choice to act as feature map on history spaces. More precisely, our aim is to 
approximate the optimal Q-functions at all time steps $t\in[T]$ by a single 
estimate, which is parametrized as a linear functional $\ell\in T^{(N)}
(\SetR^d)$, acting on the time-augmented histories' signatures. We justify this approximation by the universal 
approximation property from \cref{thm:UniversalApproximation}, which holds 
uniformly in time. The functional's parameters\footnote{
    We identify $\SetR^{d_N}$ with the truncated tensor algebra 
    $T^{(N)}(\SetR^d)$ up to level $N$ of dimension $d_N$ without specifying an explicit isomorphism.
}  
are updated utilizing the optimal Q-iteration \eqref{eq:HDPOptimalQIteration} 
with histories sampled through interaction with the environment in an online 
fashion. 

To apply \cref{thm:UniversalApproximation} as the algorithm's approximation 
argument, we will interpret the optimal Q-functions $(Q_t^*)_{t\in[T]}$ as one 
functional acting on the space of linearly interpolated, time-augmented 
histories. Recall the definitions of the metric spaces $(\Omega_t^0, 
d_{1,[0,t]})$, for $t\in[0,T]$, and $(\Lambda_T, d_{\Lambda_T})$ from 
\cref{chapter:Signature:section:BoundedVariation}, where here we will consider 
paths $X\in\mathcal{V}^0([0,T],\SetR^{d+2})$.  We make the following 
assumption, with the objective of removing the translation invariance of the 
1-variation norm \eqref{eq:1-VariationNorm} of an interpolated path by fixing 
the histories' starting value. A discussion on this assumption in 
practice is provided below in \cref{rem:SpacesContainOrigin}.

\begin{assumption}
\label{ass:SpacesContainOrigin}
    The sets $\SpaceO$ and $\SpaceA$ contain the origin in $\SetR^d$ and 
    $\SetR$, respectively, and for the remaining section we consider histories 
    starting at $o_0=0\in\SetR^d$ and $a_0=0\in\SetR$. To emphasise this, for 
    $t\in[T]$, let $\hat\SpaceH^0_t = \{ \hat h_t\in\hat\SpaceH_t \mid 
    \hat h_0 = (0,o_0,a_0) = 0 \in\SetR^{d+2}\}$ be the space of all time-augmented, 
    time-$t$ histories starting at the origin.
\end{assumption}
Fix $t\in[T]$. For a time-augmented history $\hat{h}_t \in\hat{\SpaceH}^0_t$, 
recall that its concatenation with an action $a\in\SpaceA$ is given as the 
stream of data $\hat{h}_t^{a} = (u, o_u, a_u)_{u\in[t]} \in \hat\SpaceH^0_t 
\times\SpaceA$, where we set $a_t := a$, and with components $(u,o_u,a_u) \in 
\SetR^{d+2}$ for $u\in[t]$. Its linear interpolation $\hat H^{a}:[0,t] \to 
\SetR^{d+2}$ is a continuous, piecewise linear path of finite variation with 
values in $\SetR^{d+2}$, starting at the origin. Define the 
space of all trajectories on $[0,t]$ of such linearly interpolated, 
time-augmented histories of length $t\in[T]$ as
\begin{equation}
    \mathcal{K}^0_t 
    := \{\hat H^{a}_{[0,t]} 
    \mid \hat H^{a} \text{ is linear interpolation of } 
    \hat h^{a}_t\in\hat\SpaceH^0_t\times\SpaceA\}.
\end{equation}
Further, the union of all linearly interpolated, continuous-time
histories of arbitrary length is defined as
\begin{equation}
\label{eq:UnionLinearInterpolations}
    \Lambda_{[T]} := \bigcup_{t\in[T]} \mathcal{K}^0_t.
\end{equation}
Due to \cref{ass:SpacesContainOrigin} we have 
$\mathcal{K}^0_t\subset\Omega_t^0$ for all $t\in[T]$, and thus 
$\Lambda_{[T]}\subset\Lambda_T$. Additionally, the following lemma holds.
\begin{lemma}
    For each $t\in[T]$, the set $\mathcal{K}^0_t\subset\Omega_t^0$ is a compact with respect to the topology induced by $d_{1,[0,t]}$ on $\Omega_t^0$. Moreover, the set $\Lambda_{[T]}\subset\Lambda_T$ is compact with respect to the topology induced by $d_{\Lambda_T}$ on $\Lambda_T$.
\end{lemma}
\begin{proof}
    For any $m\in\SetN$ we consider $\SetR^m$ endowed with the 1-norm 
    $\norm{\cdot}_{1,\SetR^m}$.\footnote{
        Note that for finite products $C := C_1\times\cdots\times C_n$, with 
        $C_i\subset\SetR^{m_i}$ and $n,m_i\in\SetN$, the product topology on 
        $C$ coincides with the subspace topology from $\SetR^{m_1\cdots m_n}$.
    }
    We use the notation
    \begin{equation*}
    \begin{split}
        h^a_t = (o_0,a_0,\dots ,o_t,a_t), \quad \text{and} \quad
        \mfh_t^\mfa &= (\mfo_0, \mfa_0, \dots, \mfo_t,\mfa_t),        
    \end{split}
    \end{equation*}
    to differentiate between two (not necessarily) distinct histories in $\SpaceH_t\times\SpaceA$ for any $t\in[T]$, and analogously 
    their time-augmented versions $\hat h^a_t$ and $\hat \mfh_t^\mfa$.

    For the first claim, recall that $\SpaceO\subset\SetR^d$ and 
    $\SpaceA\subset\SetR$ are by assumption compact, which implies the finite product 
    $\SpaceH_{t-1}\times\SpaceA\subset\SetR^{d+1}$ is compact. Since 
    $\hat\SpaceH_t^0\times\SpaceA$ is canonically homeomorphic to 
    $\SpaceH_{t-1}\times\SpaceA$ it is also compact. If we show that the 
    interpolation map $\Phi_t:\hat\SpaceH_t^0\times\SpaceA\to\Omega_t^0$, 
    which sends a history $\hat{h}_t^a$ to its linear interpolation 
    $\Phi(h^a_t) = \hat H^a_{[0,t]}$, is continuous, then 
    $\Phi_t(\hat\SpaceH_t^0\times\SpaceA)=\mathcal{K}^0_t$ is compact in 
    $\Omega_t^0$ as the continuous image of a compact set. For a 
    piecewise linear path with values in $\SetR^{d+2}$, its total variation is 
    simply the sum of the norms of its successive linear increments, and 
    the difference of two piecewise linear paths is again piecewise linear. It follows by a simple calculation that the 1-variation distance between the linear interpolations for any two histories $\hat h^a_t,\hat \mfh^\mfa_t \in 
    \SpaceH_t^0\times\SpaceA$ is
    \begin{multline*}
    %\begin{split}
        d_{1,[0,t]}\bigl(\Phi_t(\hat h^{a}_t), \Phi_t(\hat\mfh_t^\mfa)\bigr)
        = \bigl\Vert
            \Phi_t(\hat h^{a}_t) - \Phi_t(\hat\mfh_t^\mfa)
        \bigr\Vert _{1,[0,t]} \\
        = \sum_{u=1}^t \Bigl(
            \bigl\Vert(o_{u}-o_{u-1})-(\mfo_u-\mfo_{u-1})\bigr\Vert_{1,\SetR^d}
            + \bigl|(a_u-a_{u-1})- (\mfa_u-\mfa_{u-1})\bigr|
        \Bigr).
    %\end{split}
    \end{multline*}
    The right hand side is a finite sum of continuous functions of the 
    coordinates of $\hat h^a_t$ and $\hat\mfh^\mfa_t$, where the time component vanishes since it is the same in both histories, hence $\Phi_t$ is 
    continuous and the first claim follows.

    For the second claim, note that for each $t\in[T]$ the canonical 
    embedding $\Omega_t^0\hookrightarrow\Lambda_T$ is continuous with 
    respect to $d_{\Lambda_T}$, which implies $\mathcal{K}^0_t$ remains compact when 
    viewed as subset of $\Lambda_T$. Finally, $\Lambda_{[T]}$ is defined as a finite union of compact sets \eqref{eq:UnionLinearInterpolations}, and hence 
    itself compact in the metric space $(\Lambda_T, d_{\Lambda_T})$.
\end{proof}

\begin{remarks}
\label{rem:SpacesContainOrigin}
    \begin{enumerate}
        \item Without fixing the starting value of histories in 
        \cref{ass:SpacesContainOrigin}, the 1-variation distance $d_{1,[0,t]}$ 
        would be translation invariant in the observation and action component 
        and thus only defines a pseudometric on $\mathcal{K}^0_t$ rather than a 
        true metric. In the case $(0,0) \notin \SpaceO\times\SpaceA$, the 
        lemma above remains valid, if we assume the first point of each 
        history is fixed at an arbitrary $(o,a)\in\SpaceO\times\SpaceA$ and 
        translate all streams by $-(o,a)$. The resulting space of translated 
        linear interpolations is compact, and the original space is obtained 
        from this by a continuous translation. Moreover, the requirement 
        $a_0=\alpha$ for some fixed $\alpha\in\SpaceA$ can be omitted if 
        $\Omega_t^0$ is redefined to contain all paths whose $d$ components, 
        which correspond to the observations in a linearly interpolated path, 
        start at the origin, while the component corresponding to the actions 
        may start at an arbitrary value. Fixing the initial observation already 
        removes the translation invariance of the 1-variation distance.
    
        \item In applications one can always add a fixed artificial basepoint 
        $(o_{-1}, a_{-1})$ prior to any actually observed history $\hat h_t$ 
        to make sure that \cref{ass:SpacesContainOrigin} is satisfied. Adding 
        a basepoint furthermore helps with the calculation of the signature at 
        the initial step $t=0$ which is further discussed on 
        \cref{chapter:SignatureQLearning:section:QImplementation}. Moreover, 
        we empirically evaluate the effect of using different basepoint values 
        in the experiment in \cref{chapter:Experiments:section:MountainCar}.
    \end{enumerate}
\end{remarks}

As next step, we define the path functional $Q^*:\Lambda_{[T]} \to \SetR$ such that for any 
linearly interpolated history $\hat{H}^{a}_{[0,t]}\in\Lambda_{[T]}$, it 
recovers the respective time-$t$ optimal $Q$-function with input $(h_t,a) 
\in \SpaceH_t\times\SpaceA$, from which the interpolation was derived, i.e.
\begin{equation}
\label{eq:HDPOptimalQWithTime}
    Q^*\Bigl( \hat{H}^{a}_{[0,t]} \Bigr)
    := Q^*(t,h_t, a) = Q^*_t(h_t,a).
\end{equation}
Note that the dependence of $Q^*$ on the time component in $\hat{H}^a_{[0,t]}$
is only through its last point $(t,o_t,a)$. The functional $Q^*$ is going to be learned in a RL setting, justified via the universal approximation 
property for signatures from \cref{thm:UniversalApproximation}. 
Assume $Q^*$ is continuous on the subset of all linearly interpolated paths 
$\Lambda_{[T]}\subset \Lambda_T$. Since $\Lambda_{[T]}$ is compact it is also 
closed, and by Tietze's extension theorem 
~\cite[Satz~12.10.3]{kaltenbaeck2014_fa}, there exists a continuous extension 
of the functional $Q^*$ to all of $\Lambda_T$ which agrees with the original $Q^*$ on $\Lambda_{[T]}$.%
\footnote{
    Specifically, there exists a continuous map $F:\Lambda_T\to\SetR$ with 
    ${F|}_{\Lambda_{[T]}} = Q^*$ and 
    $\sup\{|F(\hat X_{[0,t]})| : \hat X_{[0,t]}\in\Lambda_T \} 
    = \sup\{|Q^*(\hat H^a_{[0,t]})|: \hat H^a_{[0,t]} \in \Lambda_{[T]}\}$. 
    }
With a slight abuse of notation we denote this extension also by $Q^*$. Since 
$\mathcal{K}^0_T\subset\Omega^0_T$ is compact it follows by 
\cref{thm:UniversalApproximation} that for any $\epsilon>0$, there exists a 
linear functional $\ell\in T(\SetR^{d+2})$, such that
\begin{equation}
\label{eq:HDPQApproxContinuous}
    \sup_{\hat H^a\in\mathcal{K}^0_T} \sup_{t\in[0,T]} \,
    \abs{
        Q^*\bigl(\hat H^a_{[0,t]}\bigr) 
        - \bigl\langle \ell, S_t(\hat H^a) \bigr\rangle
    } < \epsilon.
\end{equation}
Taking the supremum over $\mathcal{K}^0_T$ is equivalent to taking the 
supremum over histories in $\SpaceH^0_T\times\SpaceA$, and then passing their 
time-augmented linear interpolations as inputs in \eqref{eq:HDPQApproxContinuous}. Since 
$[T]\subset[0,T]$, and by the construction of $Q^*$, such that it agrees with 
the true optimal Q-functions $(Q^*_t)_{t\in[T]}$ on $\Lambda_{[T]}$, we 
obtain
\begin{equation}
\label{eq:HDPQApproxDiscrete}
    \sup_{h^{a}_T\in\SpaceH^0_T\times\SpaceA} \sup_{t\in[T]} \,
    \abs{ 
        Q_t^*\bigl(h_{t}, a_t\bigr)
        - \bigl\langle \ell, S_t\bigl(\hat h^{a}_T\bigr) \bigr\rangle
    } < \epsilon,
\end{equation}
where the time-$t$ signature of the discrete-time history $\hat h^a_T$ is 
defined as the signature of its linear interpolation $S_t(\hat h^{a}_T) = 
S_t(\hat H^a)$ as in \cref{def:SignatureOfStream}, and the input $( 
h_t,a_t) = h^{a_t}_{t}$ to the true time-$t$ optimal Q-function $Q_t^*$ 
is understood as the truncation of the time-$T$ history $h^{a}_T$ at step 
$t\in[T]$. \Cref{eq:HDPQApproxDiscrete} constitutes the main approximation 
argument of the proposed Signature-Q-Learning method. We highlight that this 
approximation to the true optimal Q-functions is realized uniformly in $t\in[T]$ via a single linear 
functional $\ell\in T(\SetR^{d+2})$.

\begin{remark}
    For the applicability of \cref{thm:UniversalApproximation} the continuity 
    of $Q^*$ w.r.t to the metric $d_{\Lambda_{T}}$ on $\Lambda_{[T]}$ is 
    required, which is implicitly assumed for the experiments in 
    \cref{chapter:Experiments}. To ensure this assumption analytically, 
    additional smoothness assumptions on the observation kernels $\KernelP$ 
    and reward functions $\KernelR$ of the HDP model \eqref{eq:HDPTuple} are 
    necessary. An alternative to constructing $Q^*$ and estimating it via 
    \cref{thm:UniversalApproximation} would be to approximate $Q^*_t$ at each 
    $t\in[T]$ with a separate linear functional $\ell^t \in T^{(N)}
    (\SetR^{d+2})$. This approach relies on a similar approximation theorem 
    for paths of fixed lengths \cite[Theorem~1]{kiraly2019_ksod}. It does not 
    require the continuity of $Q^*$ in the time component, but rather 
    continuity of $Q^*_t$ in $(h_t,a)$ for each $t\in[T]$, which is a weaker assumption. The elegance of Signature-Q-Learning, however, is its uniform approximation over time, which is furthermore advantageous from a computational perspective, as training separate linear functionals $\ell_k^t$ at each step requires substantially less sample efficient.    
\end{remark}


In the implementation we have to truncated signatures at 
some level, and the Signature-Q-Learning algorithm aims at learning a 
truncated version of a linear functional $\ell$ that satisfies \eqref{eq:HDPQApproxDiscrete} for some $\epsilon>0$. The 
procedure is as follows. Fix a truncation order $N\in\SetN$. For any linear 
functional $\ell\in T^{(N)}(\SetR^{d+2})$, we define the \textit{approximate 
Q-function corresponding to $\ell$}, acting on the space of all time-augmented 
histories of arbitrary length, as
\begin{equation}
    Q(\,\cdot\,|\,\ell): 
        \bigcup_{t\in[T]} (\hat\SpaceH^0_t\times\SpaceA) \to \SetR 
    \quad \text{ by } \quad
    Q(\hat h^a_t|\,\ell) := \bigl\langle     
        \ell, S^{(N)}_t\bigl(\hat{h}^a_t\bigr) 
    \bigr\rangle,    
\end{equation}
where the truncated signature up to level $N$ of the discrete time-history 
$\hat h^a_t$ is defined as the truncated signature of its linear interpolation 
$\hat H^a_{[0,t]}$. The algorithm starts with some initial 
linear functional $\ell_0 \in T^{(N)}(\SetR^{d+2})$, where the parameters of 
$\ell_0$ are initialized randomly, and its corresponding approximate 
Q-function $Q(\,\cdot\,|\,\ell_0)$. Let $\ell_k \in T^{(N)}(\SetR^{d+2})$ be 
the linear functional after the $k$-th update step, such that $Q(\,\cdot\,|\,\ell_k)$ is the algorithms 
approximation to $Q^*$ after update step $k=m t$ corresponding to step 
$t\in[T]$ in episode $m\in\SetN$. By following some exploration policy, for 
example an $\varepsilon$-greedy policy with respect to the current approximate 
Q-function $Q(\,\cdot\,|\,\ell_k)$, see \cref{rem:SignatureQAlgoRemarks}, we 
sample histories and rewards. At time $t\in[T-1]$ in the current episode, the 
exploration policy selects action $a_t\in\SpaceA$, based on the 
history $h_t\in\SpaceH_t$ observed so far. Subsequently, we receive the reward 
$r_t = R_t(h_t,a_t)$, and the environment transitions to the next observation 
$o_{t+1}$. This results in a transition sample 
$(h_t,a_t,r_t,o_{t+1})$, from which we create the \textit{Q-target value} $y_t$ at 
step $t$ for update step $k+1$ as
\begin{equation}
\label{eq:HDPQTargetValue}
    y_t := r_t + \gamma\, \max_{a\in\SpaceA} 
        Q\bigl( \hat{h}^a_{t+1}\mid\ell_k \bigr),
\end{equation}
based on the current $Q^*$ estimate provided by $\ell_k$. The Q-target value $y_t$ constitutes a sample of the right-hand side of the
optimal Q-iteration \eqref{eq:HDPOptimalQIteration}. It is a supposedly 
better estimate of the true value $Q^*_t(h_t,a_t)$ than our current estimate 
$Q(\hat h_t^{a_t}\,|\,\ell_k)$, due to the actual reward $r_t$ received from the 
environment, and the iterative nature of the optimal Q-functions in \eqref{eq:HDPOptimalQIteration}. 
Thus, the linear functional $\ell_{k+1}$, which provides an updated approximate Q-function, is selected such that it minimizes the error between the Q-target value $y_t$ and the estimate for the true optimal Q-value of the tuple $(h_t,a_t)$, measured by some loss functional $\mathcal{L}:
\SetR^2\to\SetR$, concretely,
\[
    \ell_{k+1} \in \argmin_{\ell}\mathcal{L}\bigl(y_t, Q(\hat h_t^{a_t} \,|\,\ell)\bigr).
\]
In Signature-Q-Learning the parameters of the functional $\ell_k$ are updated towards 
this minimum by one-step-gradient descent. If $\mathcal{L}$ is 
given as the squared error loss $\mathcal{L}(x,y)=\tfrac{1}{2}(x-y)^2$, for 
$x,y\in\SetR$, the update amounts to
\begin{equation}
\label{eq:HDPQApproxUpdate}
    \ell_{k+1} \gets \ell_k + \alpha_t \,
    \Big(y_t - Q\bigl(\hat{h}^{a_t}_t \,|\, \ell_k\bigr)\Big)
    \nabla_{\ell}\, Q\bigl(\hat{h}^{a_t}_t \,|\, \ell\bigr)\Big|_{\ell=\ell_k},
\end{equation}
with learning rate $\alpha_k \geq 0$ at update step $k$. After a total number of 
$M\in\SetN$ episodes and $K=MT$ update steps the algorithms has learned the 
approximate Q-function $Q(\cdot\,|\,\ell_K)$ as final estimate for the true 
optimal Q-functions $(Q^*_t)_{t\in[T]}$. We can then follow a policy 
$\pi^K\in\Pi^\text{HD}$ that is greedy with respect to $Q(\cdot\,|\,\ell_K)$, 
by selecting the action at step $t\in[T-1]$ as 
\begin{equation}
\label{eq:GreedyApproxQFunction}
    a_t = \pi^K(h_t)\in\argmax_{a\in\SpaceA} Q(\hat h^a_t\,|\,\ell_K).    
\end{equation}
for any $h_t\in\SpaceH_t$. If the learned function $Q(\cdot\,|\,\ell_K)$ is a 
sufficiently good approximation of the true optimal Q-functions, then $\pi^K$ 
is approximately optimal. Note, that the policy $\pi^K$ can be considered a 
stationary in the sense that it applies the time-invariant decision rule 
$\argmax_{a\in\SpaceA} \langle \ell, \,\cdot\, \rangle$ to any input element
from $\bigcup_{t\in T}S^{(N)}_t(\hat \SpaceH^0_T\times\SpaceA) \subset T^{(N)}
(\SetR^{d+2})$. The full Signature-Q-Learning procedure is provided in 
\cref{alg:SignatureQLearning}, and \cref{rem:SignatureQAlgoRemarks} comment 
on some aspects of it.

%------------------------------------------------------------------------------
% Signature-Q-learning algorithm WITH actions

%\algrenewcommand\algorithmicrequire{\textbf{Input:}}
\begin{algorithm}
	\caption{Signature-Q-learning}
	\label{alg:SignatureQLearning}
	\begin{algorithmic}[1] 
    % optional parameter [1] means line numbering for every line 
		\Require Number of episodes $M\in\SetN$, signature truncation order $N\in\SetN$, initial linear functional 
            $\ell_0\in T^{(N)}(\SetR^{d+2})$, loss function $\mathcal{L}$, learning rates 
            $\alpha_{n,t}$, exploration rates $\varepsilon_{n,t}$, for
            $t\in[T-1]$ and $n\in\SetN$.
            
		%\State Initialize parameter $\ell \gets \ell_0$.
		
		% for loop, one for each episode
		\For{episode $n = 1,2,\dots,M$}
            \State Observe $o_0$ and set $\hat{h}_0 = (0,o_0)$.
            \If{basepoint = True}
                \State Append basepoint to $\hat{h}_0$ at time step $t=-1$ (details in 
                \cref{chapter:SignatureQLearning:section:QImplementation}). 
            \EndIf
            \State Calculate $S_0^{(N)}(\hat{h}^a_0)$ for each $a\in\SpaceA$.
            \For{time step $t=0,1,\dots,T-1$}
                \State \label{alg:SignatureQLearning:EpsGreedy} 
                Sample action $a_t$ according to $\varepsilon$-greedy policy based on current estimate $\ell_k$:
                \[
                    a_t = 
                    \begin{cases}
                    \text{selected from }
                    \argmax_{a\in\SpaceA} \bigl\langle
                            \ell_k, S_t^{(N)}(\hat{h}^a_t) 
                    \bigr\rangle, 
                            & \mbox{with probability } 1-\varepsilon_{n,t}, \\
                        \text{sampled uniformly from } \SpaceA, 
                            & \mbox{with probability } \varepsilon_{n,t}. \\
                    \end{cases}
                \]
			    \State Take action $a_t$, receive reward $r_t$ and 
                    observe $o_{t+1}$.
                \If{$t=T-1$}
                    \State Receive terminal reward $r_T$.
                \EndIf
                \State Update the time-augmented history 
                    $\hat{h}_{t+1} = (\hat{h}_t, a_t, t+1, o_{t+1})$.
                \State Update the respective signature
                    using Chen's identity,
                    \begin{displaymath}
                        S_{t+1}^{(N)}(\hat{h}^a_{t+1}) 
                        = S_t^{(N)}(\hat{h}^{a_t}_{t}) 
                            \otimes S_1^{(N)}(\hat{h}^a_{t:t+1}),
                            \quad \text{for each } a\in\SpaceA.
                    \end{displaymath} 
                    \label{alg:SignatureQLearning:Chen}
                \State 
                    Create the $Q$-target value as
                    \begin{displaymath}
                        y_t = \begin{cases}
                            r_t + r_T & \mbox{if } t = T-1, \\
                            r_t + \max_{a\in\SpaceA} \bigl\langle 
                                \ell_k, S_{t+1}(\hat{h}^a_{t+1}) 
                            \bigr\rangle, & \mbox{else.}
                        \end{cases}
                    \end{displaymath}
                    \label{alg:SignatureQLearning:TargetValue}
                    
                \State Update the linear functional parameters via stochastic
                    gradient descent with learning rate $\alpha_{n,t}$ and loss
                    function $\mathcal{L}$,
                    \[
                        \ell_{k+1} \leftarrow \ell_{k} + \alpha_{n,t}\, \nabla_{\ell}\, \mathcal{L}\Bigl(y_t, \bigl\langle \ell, S_t^{(N)}(\hat{h}^{a_t}_t) \bigr\rangle\Bigr)\Big|_{\ell=\ell_k}
                    \]                    
            \EndFor
		\EndFor
	\end{algorithmic}
\end{algorithm}

\begin{remarks}
\label{rem:SignatureQAlgoRemarks}
    \begin{enumerate}
        \item The update of the estimate $Q(\cdot \mid \ell_k)$ towards the
        target value $y_t$ given by \cref{eq:HDPQApproxUpdate}, corresponds to
        the usage of the squared error loss $\mathcal{L}(x,y) = \frac{1}{2}
        (x-y)^2$. Other loss functions may be used,
        such the absolute error loss $\mathcal{L}(x,y) = \abs{x-y}$ or the 
        smooth-$L_1$ loss,
        \begin{equation}
        \label{eq:SmoothL1Loss}
            \mathcal{L}(x,y)
            = \begin{cases}
                \frac{1}{2} \big(x-y\big)^2, & 
                    \mbox{if } \abs{x-y} < 1, \\
                \abs{x-y} - \frac{1}{2}, & 
                    \mbox{else},
            \end{cases}                     
        \end{equation}
        which combines the absolute loss with the squared loss in a continuously
        differentiable way. In the numerical experiments presented in
        \cref{chapter:Experiments}, the smooth-$L_1$ is used.

        \item At time $T-1$ of each episode, the Q-target value in
        \cref{alg:SignatureQLearning:TargetValue} of
        \cref{alg:SignatureQLearning} is provided as the sum $r_{T-1} + r_T$.
        Inline with the handling of terminal rewards in episodic reinforcement learning tasks~\cite{pardo2018_tlrl}, no
        bootstrapping takes place at time $T$, since there are no more rewards
        available afterwards, and as such the last target value in each
        episode is given as $y_{T-1} = r_{T-1} + r_T$.
        
        \item \Cref{alg:SignatureQLearning} applies  
        $\varepsilon$-greedy as exploration policy in
        \cref{alg:SignatureQLearning:EpsGreedy}. This exploration scheme is
        often used in online Q-Learning for MDPs and derives its name from the
        parameters $\varepsilon_{n,t} \in [0,1]$, that govern the
        exploration-exploitation trade-off. At each step $t$, action $a_t$ is
        sampled uniformly from $\SpaceA$ with probability $\varepsilon_{n,t}$,
        and otherwise selected greedily based on the current estimate $Q(\cdot\,|\,
        \ell_k)$ as in \cref{eq:GreedyApproxQFunction}. Setting 
        $\varepsilon_{n,t} =\varepsilon \in [0,1]$
        for all $n\in\SetN$ and $t\in[T]$ amounts to a constant exploration rate
        and ensures that each action is sampled infinitely often if the number
        of episodes $n$ tends to infinity. On the other hand, the exploration rates $\varepsilon_{n,t}$ may also be defined as a decreasing series approaching some value $\varepsilon_\infty \geq 0$
        for $n\to\infty$ and $t\to T$, to ensure high exploration early during training, and a focus on exploitation later on when the learned approximate Q-functions are already more accurate.
        Possible decay schemes include exponential and linear decays, with
        examples displayed in \cref{fig:EpsilonDecaySchemes}. The
        decay may take place over update steps or episode-wise where $\varepsilon_{n,t} = \varepsilon_{n}$ for all $t\in[T]$ and episode $n\in\SetN$. Note, that the decision about
        whether action $a_t$ is sampled uniformly or not is made \textit{for
        each} action. Since the objective of Signature-Q-Learning is to select
        actions which maximize future expected values \textit{based on
        histories} instead of single observations or states, an alternative
        exploration approach could be to decide at the beginning of each
        episode if it includes uniformly sampled actions or not. This would
        provide the algorithm with a mix of histories including exploration,
        alternating with histories for which actions are exclusively selected
        greedily w.r.t. to the current estimate alone.

        % epsilon decay schemes
        \begin{figure}
            \centering
            \begin{tikzpicture}
                \begin{axis}[
                    axis x line = bottom,
                    axis y line = left,
                    axis line style={|-stealth},
                    axis line shift = 5pt,
                    xlabel = episode $n$,
                    ylabel = exploration parameter $\varepsilon_n$,
                    width = 10cm,
                    height = 7cm
                ]
                    \addplot [
                        domain=0:200, 
                        samples=100, 
                        color=blue,
                        line width=1pt,
                    ]
                    {
                        0.97^x
                    };
                    \addlegendentry{
                        exponential decay
                    }
        
                    \addplot [
                        domain=0:200,
                        samples=100,
                        color=red,
                        line width=1pt,
                    ]
                    { 
                        max((100 - x)/100, 0.0)
                        %(200 - x)/200
                    };
                    \addlegendentry{
                        linear decay
                    }
                \end{axis}
            \end{tikzpicture}
            \caption[Illustration of exploration parameter decay schemes]{Two
                exemplary decay schemes for the parameters 
                $\varepsilon_{n,t}$ of the $\varepsilon$-greedy exploration
                policy from \cref{alg:SignatureQLearning} are displayed.
                Parameter values are assumed to be fixed during episodes, i.e.
                $\varepsilon_n := \varepsilon_{n,0} = \varepsilon_{n,t}$ for
                all steps $t\in[T]$. For the starting value $\varepsilon_{0} =
                1$, the linear decay scheme reduces the value linearly to zero
                over 100 episodes, and leaves it fixed for the remaining 100
                episodes. The exponential decay reduces the initial value by
                the constant factor 0.97 over 200 episodes.}
            \label{fig:EpsilonDecaySchemes}
        \end{figure}        
    
        \item Instead of an $\varepsilon$-greedy exploration policy in
        \cref{alg:SignatureQLearning}, one might use a softmax exploration
        policy $\pi_\tau$, based on the current optimal Q-function estimate.
        Specifically, with current linear functional $\ell_k$,
        and given history $h_t$, the probability of sampling action
        $a\in\SpaceA$ under $\pi_\tau$ is given by
        \begin{equation}
            \label{eq:SoftmaxPolicy}
            	\PM^{\pi_\tau}[A_t=a \mid H_t = h_t] =  
                \frac{
            		\exp(\tau^{-1} \Bigl\langle 
                        \ell_k, S_t^{(N)}\bigl(\hat{h}^a_t\bigr) 
                    \Bigr\rangle)
            		}{
            		\sum_{\alpha \in \SpaceA} 
            			\exp(\tau^{-1} \Bigl\langle 
                            \ell_k, S_t^{(N)}\bigl(\hat{h}^\alpha_t\bigr)
                        \Bigr\rangle)
            	},
            \end{equation}
        where $\tau>0$ is the exploration parameter, called temperature, that determines the exploration behaviour. For $\tau
        \to 0+$, the policy $\pi_\tau$ converges to the uniform distribution
        on $\SpaceA$, while for $\tau\to\infty$, the sampling policy becomes
        greedy w.r.t the current estimate and the action with the highest
        estimated expected reward is selected. Analogue to 
        $\varepsilon$-greedy, decay schemes such as the ones displayed in \cref{fig:EpsilonDecaySchemes} may be employed to the temperature parameter $\tau$ to steer exploration behaviour during training.

        \item \Cref{alg:SignatureQLearning} constitutes an online
        algorithm, where the estimates $Q(\cdot\,|\,\ell_k)$ are updated at each step while
        interacting with the environment, and, with $\varepsilon$-greedy as exploration policy, the updated estimate determines
        the action selection at the subsequent step. It is also possible to
        design the algorithm in an offline fashion, where first, a number of full histories $\{h_T^i\}_{i=0}^n, n\in\SetN$ are sampled by following some fixed exploration policy $\pi$, and only afterwards the linear functional's parameters are learned, with updates determined by the sampled histories. The key difference in the offline version, is that the exploration policy $\pi$ is not informed by any already learned estimates and as such, there is no exploitation component in the action selection when collecting histories to learn from at a later step. Depending on the environment, it may be hard to collect histories with high value only though random action selection. An example is the environment used in the experiment in \cref{chapter:Experiments:section:MountainCar}, in which the algorithm must learn to reach a predetermined goal in as few steps as possible. If actions are selected, for example, solely by uniformly sampling the action space, the probability of reaching the goal is close to zero. Thus, the collected sample of histories does not contain any information about the goal, and in the following training phase, the algorithm likely fails to learn an approximate Q-function which is sufficiently close to the true optimal Q-functions.
        \end{enumerate} 
\end{remarks}   

%------------------------------------------------------------------------------
\section{Implementation of Approximate Q-Functions}
\label{chapter:SignatureQLearning:section:QImplementation}
%------------------------------------------------------------------------------

Recall that $\SpaceO\subset\SetR^d$ and $\SpaceA\subset\SetR$ with 
$\abs{\SpaceA} = m < \infty$, and a time-augmented, action-concatenated history $\hat h_t^a\in\SpaceH\times\SpaceA$ is the stream of data
$\hat{h}^a_t = (i, o_i, a_i)_{i=0}^t$ is a stream of data with sample points
in $\SetR^{d+2}$ for any $a_t \in \SpaceA$ and $\hat{h}_t\in\SpaceH_t$.
As already mentioned in \cref{chapter:Signature:section:StreamComputation}, the
computation of the signature of linearly interpolated streams of data may be
handled with various packages in Python
\cite{kidger2021_sdcsltb, lyons2002_esig, reizenstein2020_a1ileciss}. To
implement the linear functional computationally on top of the signature
computation, one may again exploit existing libraries, that provide machine
learning functionalities. Popular Python libraries include the packages
\texttt{PyTorch} \cite{paszke2017_adp} and \texttt{TensorFlow}
\cite{martinabadi2015_tlmlhs}. The advantage of the usage of such packages is,
that it readily provides access to back-propagation functionalities, 
gradient-based optimization algorithms such as ADAM~\cite{kingma2015_amso} and
methods to include the decay of hyperparameters in training as displayed in
\cref{fig:EpsilonDecaySchemes}. By identifying a linear
functional $\ell\in T^{(N)}(\SetR^{d+2})$ as a parameter vector $\theta_\ell
\in \SetR^{d_N}$, with $d_N$ the dimension of $T^{(N)}(\SetR^{d+2})$, it can
be implemented as single linear neural network layer, which takes as input the
time-augmented histories' signature, again identified with a vector in 
$\SetR^{d_N}$, and outputs approximate Q-values for each action. Given 
$\hat{h}_t\in\SpaceH_t$, the approximate Q-values for $(\hat h_t^a)_{a \in 
\SpaceA}$ are then obtained by looping through actions and evaluating the 
vector product
\[
    Q(\hat{h}^a_t\mid \ell) = \theta_\ell^\top \cdot S^{(N)}(\hat{h}^a_t).
\]
With this approach there is a single parameter vector $\theta_\ell$ for all 
actions that is updated at each step, irrespective of the action that was 
chosen. Further, this requires the application of Chen's identity to compute 
the updated signature $S^{(N)}(\hat{h}^a_{t+1})$ in 
\cref{alg:SignatureQLearning:Chen}, $\abs{\SpaceA}$-times, before selecting 
the maximal action value.

A second approach is to use $\hat{h}_t$ as input to the linear functional, 
without the concatenation with next possible action at the end of the 
history, and have the linear functional output multiple approximate Q-values, 
each corresponding to one action. Since $\hat{h}_t$ contains an uneven number 
of actions and observations, by concatenating a dummy action\footnote{
    Often $\alpha\in\SpaceA$ is possible, for example many environments 
    permit a ``do-nothing'' action, as it is the case in the experiments in 
    \cref{chapter:Experiments}, where the agent simply observes the 
    environment's transition and which constitutes a generic choice for 
    $\alpha$.
} 
$\alpha\in\SetR$ at the beginning of $\hat h_t$ before $(0, o_0)$, we obtain the stream of data $\leftindex[][h]^{\alpha}
{\hat{h}}_t := (a_i,i,o_i)_{i=0}^t$, where $a_0 := \alpha$, again with
sample points in $\SetR^{d+2}$. The signature of $\leftindex[][h]^{\alpha}
{\hat{h}}_t$ can readily be computed, and the linear functional $\ell$ is then
implemented via the parameters $\Theta = (b,W)$, with $b\in\SetR^m$ and
$W\in\SetR^{m\times d_N}$, corresponding to the usual bias vector and weight
matrix of a linear neural network layer, which outputs the vector of approximate 
Q-values by
\begin{equation}
\label{eq:QImplementationApproach2}
    \bigl(
        Q\bigl(\hat{h}^{a_1}_t\mid \ell\bigr),
        \dots, Q\bigl(\hat{h}^{a_m}_t\mid \ell\bigr)
    \bigr)
    = \Theta \, S^{(N)}\bigl(\leftindex[][h]^{\alpha}{\hat{h}}_t\bigr)^\top
    = b + W \bigl(
        S^{1}\bigl(\leftindex[][h]^{\alpha}{\hat{h}}_t\bigr),
        \dots,
        S^{N}\bigl(\leftindex[][h]^{\alpha}{\hat{h}}_t\bigr)
    \bigr)^\top.
\end{equation}
This approach avoids looping over actions and requires only one evaluation of
the linear layer \eqref{eq:QImplementationApproach2} to provide the Q-values
for all actions. Note, that this implementation effectively amounts to one linear
functional per action, which for action $a_i\in\SpaceA$ is given by the $i$-th
bias vector component $b_i$ and $i$-th weight matrix row $(W_{i,j})_{j=1,\dots,
d_N}$.

The third approach, is analogue to the second, with the difference, that
histories contain only observations and actions are discarded. This is
possible since the action space at each step does not depend on the
observation the environment is in and is inspired by the approximation of
optimal Q-functions in MDPs by neural network, where often the input is the
current state and the network outputs all action-values at once. A 
time-augmented history containing only observations has components $(t, o_t) 
\in \SetR^{d+1}$ at each step and its truncated signature of order
$N\in\SetN$ is an element in $T(\SetR^{d+1})$. The implementation is then 
given by \eqref{eq:QImplementationApproach2} with the only difference that 
the weight matrix's number of rows $d_N$ is the dimension of 
$T(\SetR^{d+1})$. 

Finally, we comment on the handling of the signature computation at the
beginning of an episode. At time $t=0$ there is only one tuple available,
depending on the approach it is given as either $\hat{h}^a_0 = (0, o_0, a)$
for the first, $\leftindex[][h]^{\alpha}{\hat{h}}_0 = (\alpha, 0, o_0)$ for
the second, and $(0,o_0)$ for the third approach. Since analytically the
signature of a continuous path at a single time point evaluates to zero,
compare \cref{def:Signature} with $s=t$, one can simply set $S^{N}(
\hat{h}^a_0) = \bm{1}$, with $\bm{1}$ the unit element in $T^{N}(\SetR^{d+2}
)$, and analogue for the other two initial tuples. In the first approach, this
results in the approximate Q-value for each $a\in\SpaceA$ being given by the
first component of $\theta_\ell$ which implies that the algorithm is not able
to differentiate between action at the beginning of episodes. For the other two
approaches, the vector of approximate Q-values
\eqref{eq:QImplementationApproach2} at $t=0$ is then given by the bias vector
$b$. As such the algorithm can learn to differentiate actions at the start of
an episode, this differentiation however, is static and can not depend on the
information provided by the initial observation $o_0$. To be able to
incorporate the information provided by $o_0$ at the start of an episode, we
propose the addition of a basepoint before the initial tuple. In approach
three, this basepoint would be given as tuple $(0, o_{-1})$ for some
artificial observation $o_{-1}\in\SetR^{d}$, such that $t=0$, there is already
the stream $((0,o_{-1}), (0, o_0))$ available for the computation of the
signature. A simple solution would be to set $o_{-1} = 0$, such that the 
Q-function approximation at $t=0$ depends on the value of $o_0$ and the bias
vector $b$. However, any other value is possible and the choice can be
tailored to the environment, as in the experiment in
\cref{chapter:Experiments:section:MountainCar}. For the other two approaches a
basepoint is provided analogously as $(0, o_{-1}, a_{-1})$ with some 
artificial action $a\in\SetN$ for the first approach, and$(\alpha, 0, o_{-1})$ 
for the second approach, respectively.

% screenshot and comment on exclusion of reward in history
\begin{comment}
\begin{figure}[htpb!]
    \centering
    \includegraphics[width=0.7\textwidth]{graphics/inspiration2.png}
    \caption{Screenshot from Neural Fitted Q Iteration, Martin Riedmiller. Could include similar info on why I do not use rewards in transition tuples for SigQ.}
    \label{fig:inspiration2}
\end{figure}
\end{comment}

\begin{comment}

% this comment section includes the initially planned policy gradient method
% and some results for its implementation in the mountain car environment

%%%%%%%%%%%%%%%%%%%%%%%%%%%%%%%%%%%%%%%%%%%%%%%%%%%%%%%%%%%%%%%%%%%%%%%%%%%%%%%
\newpage
\section{Signature REINFORCE} 
\label{section:SignatureReinforce}
%%%%%%%%%%%%%%%%%%%%%%%%%%%%%%%%%%%%%%%%%%%%%%%%%%%%%%%%%%%%%%%%%%%%%%%%%%%%%%%

In order to show that our approach is promising, we adapt the REINFORCE 
policy gradient algorithm such that incorporates a policy parametrized as 
a linear functional of the process' history's signature. Assuming a 
Markovian policy (after a possible enlargement of the observation space)
the gradient with respect to the parameter vector has the same form as in 
the MDP setting. We derive an explicit expression for the gradient and plug
it into the algorithm.

In an episodic POMDP setting given as 
$(\SpaceS, \SpaceA, P, r, \SpaceO, Q, \gamma)$, assume the following:

\begin{itemize}
	\item $\SpaceS \subseteq \SetR^{d_s}, \SpaceO \subset \SetR^{d_o}$ 
		with $d_s,d_o \in \SetN$,
	\item $\SpaceA = \{a_1, a_2, \dots, a_n \}$ for some $n \in \SetN$ 
		and $a_i \in \SetR^{d_a}$ for $i=1,\dots,n$,
	\item $N \in \SetN$ is the number of iterations (episodes),
	\item $T < \infty$ is the length for each trajectory.
\end{itemize} 

Let $\SpaceH_t \subseteq (\SpaceO \times \SpaceA)^{t-1} \times \SpaceO$
be the space of all observation-action histories of length $t \leq T_{max}$ 
where $h_t \in \SpaceH_t$ is given by
\[
	h_t = (o_0, a_0, o_1, a_1, \dots, a_{t-1}, o_t)
\]
and $\SpaceH := \bigcup_{t=0}^{T-1} \SpaceH_t$  the space of all finite 
histories of the underlying process. For simplicity we here consider 
observation-action histories without rewards. Note that we exclude 
histories of length $T$ since the episodes terminate at this point and
no further step is taken. For $h\in\mathcal{H}$ we shall denote its 
concatenation with action $a\in\SpaceA$ as $h^a$, specifically, 
for time $t$ we have
\[
	h^a_t 
	= (o_0, a_0, o_1, a_1, \dots, a_{t-1}, o_t, a) \in (\SetR^{d_o\times d_a})^t
%	= (o_i, a_i)_{i=0}^t \in (\SetR^{d_o\times d_a})^t
\]
This is notational convenience, since we can readily compute the signature 
of $h^a_t$, but not of $h_t$ due to the different number of actions and 
observations in the sequence.
%
The signature transform for history-action pairs is given by the mapping
$\phi: \SpaceH \times \SpaceA \to T(\SetR^{d_o \times d_a})$. At any time 
$t<T$ it takes as input a history-action pair $(h_t,a)$ (suitably embedded
into a continuous path, see \cite{levin2013_lfppsfles}) and returns its
signature on the interval $[0,t]$ as element of the extended tensor algebra
of $\SetR^{d_o \times d_a}$. We denote it as $\mathds{H}^a_t := \phi(h_t,a)$
and by a slight abuse of notation we also write $\phi(h^a_t)$ instead of
$\phi(h_t,a)$. For some arbitrary $h\in\SpaceH$ we omit the time subscript 
and only write $\mathds{H}^a$.
%
Similarly, when we consider a truncated history 
$h_{s:t} = (o_s,a_s,\dots,a_{t-1},o_t)$ with $s<t$, we write $h^a_{s:t}$ 
for the concatenated version and $\mathds{H}^a_{s:t}$ for its signature
on the interval $[s,t]$.

For an adaptation of the REINFORCE Monte Carlo policy gradient algorithm,
we parametrize the policy as linear functional of the process' history's 
signature. That is, for parameter vector $\theta$ the policy is given by 
\[
	\pi^{sig}(h,a \mid \theta) 
	 := f(\langle \theta, \phi(h,a) \rangle)
	 = f(\langle \theta, \mathds{H}^a \rangle),
\]
where $\langle \theta,\cdot \rangle$ is a linear functional with coefficients
$\theta$ acting on the signature (add details how) and $f$ is a function such
that $\pi(h, \cdot \mid \theta) \in \text{Prob}(\SpaceA)$. A common choice
for $f$ in the literature is the softmax function
\begin{equation}
\label{eq:SoftmaxPolicy}
	f(\langle \theta, \mathds{H}^a \rangle) 
	:= \frac{
		\exp( \langle \theta, \mathds{H}^a \rangle)
		}{
		\sum_{\alpha \in \SpaceA} 
			\exp( \langle \theta, \mathds{H}^{\alpha} \rangle)
	}.
\end{equation}
The policy provides current history-action preferences, that is 
$\pi(h,a \mid \theta)$ is the probability of the agent taking action $a$, 
given the observed history $h$ under the current policy parametrized by
$\theta$. 

Using the policy parametrization from \eqref{eq:SoftmaxPolicy},
\cref{alg:SigReinforce} gives an adaption of the classic REINFORCE policy
gradient method. Note that since $h^a_0 = (o_0, a)$ contains only one 
observation and action respectively, its signature is $0$ for each 
$a\in\SpaceA$ and thus, it follows from \eqref{eq:SoftmaxPolicy} that the
first step's distribution over actions is uniform.

%------------------------------------------------------------------------------
% Signature REINFORCE algorithm

\algrenewcommand\algorithmicrequire{\textbf{Input:}}
%\algnewcommand\Input{\textbf{Input:}}
\begin{algorithm}
	\caption{Signature REINFORCE}
	\label{alg:SigReinforce}
	\begin{algorithmic}[1] % optional parameter [1] means line numbering for every line 
		\Require number of iterations $N$, 
			learning rate $\beta$, 
			length of trajectories $T$, 
			signature policy parametrization $\pi^{sig}(h,a\mid\theta)$
		\State Initialize paremeter vector $\theta \gets \theta_0$.
		
		% for loop, one for each episode
		\For{$n=0,\dots,N-1$}
			% initialize values, take first action, receive first reward
			\State Initialize $\mathds{H}^a_0 = 1$ for each $a\in\SpaceA$
			\State Sample $a_0 \sim \pi^{sig}(h_0,a\mid\theta)$ 
				acc. to \eqref{eq:SoftmaxPolicy}
				\Comment{Uniform distribution over $\SpaceA$}
			\State Receive $r_0$

			% for loop to create a sample trajectory
			\For{$t=1,\dots,T$}
				\Comment {Sample trajectory}
				\State Observe $o_t$
				\State \label{alg:stp:SignatureUpdate}
                    Update the signature for each $a\in\SpaceA$ 
					using Chen's identity 
					\[
						\mathds{H}^a_t 
						\gets \mathds{H}^{a_{t-1}}_{t-1} 
						\otimes \mathds{H}^a_{t-1:t}
					\]
				\State Sample $a_t \sim \pi^{sig}(h_t,a\mid\theta)$
					acc. to \eqref{eq:SoftmaxPolicy}
				\State Receive $r_t$
			\EndFor
			\For{$t=0,\dots,T$}
				\State Calculate return $G_t = \sum_{s=t}^{T-1} \gamma^{s-t}r_s$
				\State \label{alg:stp:ParameterUpdate}
                    Update the policy parameter
					\begin{equation} \label{eq:ReinforceUpdate}
						\theta \gets \theta 
						+ \beta \gamma^t G_t 
						\nabla_\theta \log \pi^{sig}(h_t, a_t \mid \theta)
					\end{equation}
			\EndFor
		\EndFor
	\end{algorithmic}
\end{algorithm}

We used the expected cumulative discounted rewards 
$J_t(\theta)=\E_{h\sim\pi^{sig}_\theta}[\sum_{s=t}^{T-1} \gamma^{s-t} r_s]$ 
as objective function and a simple Monte Carlo based estimate for its 
gradient w.r.t $\theta$. (results in high variance, approach detailed 
in \cite{azizzadenesheli2020_pgpoeac}, check is policy gradient formula 
still correct). Given our policy parametrization \eqref{eq:SoftmaxPolicy},
the gradient in \eqref{eq:ReinforceUpdate} has the explicit form
\begin{equation} 
    \label{eq:GradientSoftmax}
    \nabla_{\theta} \log \pi^{sig}(h_t,a_t\mid \theta)
    = \mathds{H}_t^{a_t} - \sum_{
        \alpha\in\SpaceA} {
            \left( \pi^{sig}(h_t,\alpha \mid \theta)\, \mathds{H}_t^\alpha \right),
        }
\end{equation}
which can be derived using the definitions and elementary calculus. This
expression enables us to implement the update in \cref{alg:SigReinforce} 
directly. 

One drawback of using the softmax function on the discrete actionspace 
in our new signature policy approach is, that for the direct gradient
calculation at each step $t$ of our trajectory we again need the signature 
of $h_t^a$ for all $a\in\SpaceA$. The question that arises here is whether
the signatures should be stored when first calculated in \cref{alg:stp:SignatureUpdate} 
or be computed again at each update step in \cref{alg:stp:ParameterUpdate}. 
Depending on the truncation order as well as trajectory length in 
implementation, the first approach would imply storing very large arrays in
each episode, while the latter approach means additional computation time, for 
something that has already been computed at an earlier step in the algorithm.
For computational efficiency it might be of advantage to use a continuous 
action space. A possible choice is a Gaussian policy, for which actions are
drawn from a parametrized (truncated) normal distribution,
\begin{equation}
    \pi^{sig}(h,a \mid \theta) 
    = \frac{ 1 }{ \sqrt{2\pi}\sigma_\theta} 
    \exp(-\frac{(a-\mu_\theta)^2}{2\sigma_\theta^2}),
\end{equation}
where in our case the mean is represented by 
$\mu_\theta := \langle \theta,\mathds{H}^a \rangle$ and 
$\sigma_\theta$ is the parametrized standard deviation, see for instance 
\cite{williams1992_ssgacrl} and . In this continuous case, the gradient 
w.r.t. $\theta$  takes the form
\begin{equation}
    \label{eq:GradientGaussian}
    \nabla_\theta \pi^{sig}(h,a \mid \theta) 
    = \frac{
        (a - \mu_\theta)\, \mathds{H}^a
    }{
        \sigma_\theta^2
    }.
\end{equation}
%
Using \eqref{eq:GradientGaussian} instead of the discrete softmax 
gradient \eqref{eq:GradientSoftmax} in the parameter update step of
\cref{alg:SigReinforce}, one can see that at every step $t$ we only
need the signature of $h_t^{a_t}$, where $a_t$ is the one action that
was actually taken. Whether these values are stored or calculated a 
second time, this should result in faster computation for the algorithm,
compared to using the softmax gradient in connection with signatures.

A possibly alternative approach would be to use of automated 
differentiation and leveraging existing gradient optimization methods
for neural networks. However, further considerations have to be carried
out on how these methods feasibly connect with the sequential nature of
signature calculations.

\newpage
%-------------------------------------------------------------------------
\subsection{Some notation} 
\label{subsection:SomeNotation}
%-------------------------------------------------------------------------

The notation established in this section is the same as outlined in
\cref{section:SignatureReinforce}, enumerated for quick reference.
Let $0\leq s < t < T$.
\begin{itemize}
	\item	Process' observation-action history $h_t$ up to $t$ is given by
		\[
			h_t = (o_0,a_0,\dots,a_{t-1},o_t).
		\]
	
	\item	$h^a_t$ denotes the history up to $t$ 
		concatenated with some $a\in\SpaceA$, specifically
		\[
			h^a_t = (o_0,a_0,\dots,a_{t-1},o_t, a) = h_t \sqcup (a).
		\]
	
	\item $h_{s:t}$ is the history starting at time $s$ and running until 
        time $t$, analogously $h^a_{s:t}$.
	
	\item $\SpaceH_t$ is the space of all possible histories of length $t$, 
		$\SpaceH_t := \bigcup_{t=0}^{T-1} \SpaceH_t$ the space of all finite 
        histories and $\SpaceH_{s:t}$ contains all histories starting from 
        $s$ and running until $t$.
	
	\item $\phi: \SpaceH\times\SpaceA\to T((\SetR^{d_o \times d_a})), (h,a)\mapsto\phi(h,a)$ 
        denotes the signature transform for history-action pairs. We write 
        $\mathds{H}^a:=\phi(h,a)$ when an arbitrary pair $(h,a)\in\SpaceH\times\SpaceA$ 
        is given and with a slight abuse of notation also $H$ as well as 
        $\mathds{H}^a_t$ and $\mathds{H}^a_{s:t}$ for the histories $h_t$ 
        and $h_{s:t}$ respectively.

\end{itemize}

\end{comment}




%%%%%%%%%%%%%%%%%%%%%%%%%%%%%%%%%%%%%%%%%%%%%%%%%%%%%%%%%%%%%%%%%%%%%%%%%%%%%%%
\chapter{Numerical Experiments and Results} 
\label{chapter:Experiments}
%%%%%%%%%%%%%%%%%%%%%%%%%%%%%%%%%%%%%%%%%%%%%%%%%%%%%%%%%%%%%%%%%%%%%%%%%%%%%%%

%------------------------------------------------------------------------------
\section{Github Repositories}
\label{chapter:Experiments:section:Github}
%------------------------------------------------------------------------------

The code for the experiments in this chapter can be found on GitHub: The first experiment in \cref{chapter:Experiments:section:MountainCar} is implemented in repository 
\href{https://github.com/phaelicks/SignatureQLearning-MountainCar}{
    \texttt{SignatureQLearning-MountainCar}}\footnote{
    \url{https://github.com/phaelicks/SignatureQLearning-MountainCar}
} 
and the second experiment in \cref{chapter:Experiments:section:Execution} in the repository 
\href{https://github.com/phaelicks/SignatureQLearning}{
    \texttt{SignatureQLearning-Execution}}\footnote{
    \url{https://github.com/phaelicks/SignatureQLearning}
}.
Since the second experiment's implementation was more involved due to the
integration of the market simulator ABIDES, detailed in
\cref{chapter:Experiments:section:Execution:sub:ABIDES}, the experiments are
provided as two distinct repositories. Both experiments are coded in Python
3.8. The main third-party packages used are \texttt{Signatory}
\cite{kidger2021_sdcsltb} for the signature computations, the machine learning
library \texttt{PyTorch} \cite{paszke2017_adp}, the reinforcement learning API
\texttt{Gymnasium} \cite{towers2023_gv0} and its predecessor OpenAI's
\texttt{Gym} \cite{brockman2016_og}, and the market simulator ABIDES
\cite{byrd2020_athmms}, which is described in more detail in
\cref{chapter:Experiments:section:Execution:sub:ABIDES}. Python's version 3.8,
is the latest version which is compatible with the last stable release of the
package \texttt{Signatory} for Apple's Mac. The full list of package
requirements and their versions can be found in the \texttt{requirements.txt}
file in the respective repository. All training was performed on a Mac Book
Pro with Intel(R) Core(TM) i7-6700HQ CPU @ 2.60GHz processor. 

%------------------------------------------------------------------------------
\section{Mountain Car Experiment}
\label{chapter:Experiments:section:MountainCar}
%------------------------------------------------------------------------------

In this section we test the Signature-Q-learning method given by 
\cref{alg:SignatureQLearning} in an environment modeled as partially
observable Markov Decision Process. The environment is an adaptation of the
Markovian reinforcement learning problem ``Mountain Car'', provided as part of
the Python package \texttt{Gymnasium} \cite{towers2023_gv0}. The package is an
open source reinforcement learning library that provides ready to use
environments and a standard API for the communication with custom learning
algorithm. Since in the adaptation the agent has only access to a partial
observation, it has to use the history of observations to infer the hidden true
state of the system and learn an approximate optimal Q-function.

%------------------------------------------------------------------------------
\subsection{The Environment} 
\label{chapter:Experiments:section:MountainCar:sub:Environment}
%------------------------------------------------------------------------------

The ``Mountain Car'' problem consists of a car placed at the bottom
of a 2-dimensional valley, with the starting position chosen uniformly over
some area in the valley. The agent can accelerate the car in either direction
and the goal is to reach the top of the hill located to the right hand side of
the valley. An illustration of the setting is given in \cref{fig:MountainCar}.
The agent's challenge is to learn a policy that drives the car up the left side
of the valley (potentially multiple times) first, in order to gain enough
momentum while coming down from the left, to subsequently being able to reach
the top right. 

The dynamics in the original environment provided by the package
\texttt{Gymnasium} follow a deterministic MDP. The state $s_t=(p_t,v_t)$ is
given as the car's position along the x-axis $p_t$ and its current velocity
$v_t$, which takes values in the state space $\SpaceS = [-1.2, 0.6] \times
[-0.07, 0.07]$. The initial position $p_0$ is sampled from $\mathcal{U}([-0.6,
-0.4])$ while the initial velocity is set to $v_0=0$. The action space is 
$\SpaceA=\{0, 1, 2\}$ where actions correspond in increasing order to
accelerating the car to the left, not accelerating and accelerating it to the
right, respectively. Following the selection of action $a\in\SpaceA$, the
environment transitions to the new state $s_{t+1} = (p_{t+1}, v_{t+1})$
according to the deterministic, Markovian dynamics,  
\begin{align}
\label{eq:MountainCarDynamics}
    v_{t+1} &= v_{t} + (a-1)f -\cos(3p_t)g, \\
    p_{t+1} &= p_t + v_{t+1},
\end{align}
with acceleration force $f\equiv0.001$ and gravity on slopes $g\equiv 0.0025$.
The maximum episode length is $T=200$ and an episode is solved successfully if
$p_t \geq 0.5$ for some $t\leq 200$, which corresponds to the position of the
flag in \cref{fig:MountainCar}. 

\begin{figure}[htbp!]
    \centering
    \includegraphics[width=0.4\textwidth]{graphics/mountaincar/mountainCar.png}
    \caption[Illustration of the Mountain Car environment]{%
        An illustration of Gymnasium's Mountain Car environment.\footnote{
        The illustration is taken from the documentation of the Gymnasium
        library \cite{towers2023_gv0}, which can be found at 
        \url{https://gymnasium.farama.org/}.
    } The goal is to strategically accelerate the car such that it reaches the
    flag on the mountain top to its right. Since its power is not sufficient
    to go up by just accelerating right, the agent has to learn a policy that
    first moves the car up the opposite hill and uses its momentum when going
    down again.
    }
    \label{fig:MountainCar}
\end{figure}

%------------------------------------------------------------------------------
\paragraph{Observation and Action Space}
\label{chapter:Experiments:section:MountainCar:sub:Environment:par:Spaces}

We modify the environment such that it becomes partially observable one, which
justifies the need of using the environment's history for learning. Instead of
having access to the full state $s_t$, we define an observation as the car's
position $o_t = p_t$, taking values in the observation space $\SpaceO = [-1.2,
0.6]$, and the agent receives the time-augmented observation
\[
    \hat{o}_t = \Bigl(1-\frac{t}{T}, p_t\Bigr) 
    \in [0,1]\times\SpaceO,
\]
where the time component is given as the normalized remaining episode time.
The action space remains as in the original environment as $\SpaceA=\{0,1,2\}$
for the possibility to accelerate right, not accelerate and accelerate left.
The environment still progresses according to the dynamics in
\cref{eq:MountainCarDynamics}, however since the agent only observes 
$\hat{o}_t$, the problem now a (deterministic) partially observable Markov
decision process. In order to infer the hidden velocity and  an
(approximately) optimal policy, observation histories $\hat{h}_t = (\hat{o}_0, 
\dots, \hat{o}_{t})$ may be used.

%------------------------------------------------------------------------------
\paragraph{Rewards}
\label{chapter:Experiments:section:MountainCar:sub:Environment:par:Reward}

The objective is to reach the top right hill as quickly as possible. For the
agent to learn a policy that does so, keeping an episode short must result in
a higher reward. The original environment achieves this by assigning the
reward $R_t(s_t, a_t) = -\mathds{1}_{[-1.2, 0.5)}(p_t)$ for any $t\in[T]$,
which simply penalizes the agent for each position different to a goal
position and serves as incentive to keep the length of an episode short. With
Q-learning methods, however, it can be difficult to solve the problem with
this specific reward structure. The agent has no knowledge about the existence
of observations or trajectories thereof, such that the total reward
per episode is greater than $-200$,\footnote{
    This corresponds to the undiscounted reward per episode if the goal has not been reached. Depending on the discount factor, the total reward the agent receives may be different.
} before a position $p_t \geq 0.5$ is reached
for the first time. Until then, the randomly initialized Q-functions
approximation, contains no information about histories with higher values.
Since the exploration behaviour is based on the current approximation, it
is extremely difficult to reach a goal position by accident. For this reason,
we adapted the reward functions to incorporate some guidance for the agent, to
\begin{equation}
\label{eq:MountainCarReward}
    R_t(\hat{h}_t) = R_t(\hat{o}_t)
    := \mathds{1}_{\{x<0.5\}}(p_t) \frac{p_t - 0.5}{T} 
        + \mathds{1}_{\{x\geq0.5\}}(p_t)\,C\,\bigl(1-\frac{t}{T}\bigr)^2,
\end{equation}
for $t\in[200]$ and
with $C\geq 0$. Note, that the reward function depends on $\hat{h}_t$ only
through the current time-augmented observation $\hat{o_t}$. The reward
obtained is the sum of two components. The first is a running reward the agent
receives as long as the goal has not been reached. It provides some guidance
as to where the goal is located and at the same time serves as incentive to
keep episode short, since it only takes negative values. At the same time, the
agent must learn to accept negative rewards of higher absolute value when
moving the car to the left, in order to increase the overall episode reward by
managing to reach the goal earlier. It is normalized for stability. The second
component is a terminal reward, which provides a positive signal if the goal
is reached at time $t\leq T$, and which increases with smaller $t$. We include
the terminal reward in the definition of the reward functions $R_t$ since it
depends on the variable time step at which the episode ends. In the
experiments $C=0.01$ was chosen.

As an alternative to setting the reward functions to
\eqref{eq:MountainCarReward}, one could use the environment's original reward
structure in conjunction with action repetition to increase to probability of
the agent reaching the goal position for a first time. Instead of selecting an
action at each step, fix $k\in\SetN$ (w.l.o.g. assume that $T\mod k = 0$) and
select actions only at steps $0, k-1, 2k-1, \dots, T-k-1,T-1$. In between
those times the action selected last is repeated, that is if $\tilde a
\in\SpaceA$ is selected at time $t=nk-1$ for some $n\in \{1,\dots,T/k\}$, we
have $a_{t+i} = \tilde a$ for $i=0,1,\dots,k-1$.

%------------------------------------------------------------------------------
\subsection{Training Details} 
\label{chapter:Experiments:section:MountainCar:sub:TrainingDetails}
%------------------------------------------------------------------------------

We perform training in three different set-ups. Set-up one is a single
hyperparameter configuration, to emphasize that Signature-Q-Learning is able
to successfully solve the problem. Set-up two studies the effect of different
signature truncation orders, and set-up three the behaviour of the algorithm
with different basepoint positions $p_1$. 

Before providing details on the different set-ups in the paragraphs below, we 
describe the training configuration's parts that are mutual to all three.
First, we consider only the time-augmented observation history. That is, the
input to the approximate Q-function at time $t$ is the stream
$(\hat{o}_i)_{i=0}^t = (1 - i/T, p_i)_{i=0}^t$. The approximate Q-function is
implemented as single linear layer, which outputs the Q-value for each one of
the three actions simultaneously. Details on this implementation are given in 
\cref{chapter:SignatureQLearning:section:QImplementation}. For the signature
truncated at order $N\in\SetN$, this results in $d_N$ trainable parameters for
each action, with $d_N$ the dimension of $T(\SetR^2)$, which consist of a weight
matrix $W \in \SetR^{3\times (d_N-1)}$ and a bias vector $b\in\SetR^3$, cf.
\eqref{eq:QImplementationApproach2}. The entries in $W$ are initialized
randomly from $\mathcal{U}(-\sqrt{k},\sqrt{k})$ with $k=6/(d_N^{-1}+3)$, which
corresponds to \texttt{PyTorch}'s uniform Xavier initialization
\cite{glorot2010_udtdfnn}, and the bias vector's components are set to
$b_i=0.1, i=1,2,3$. Using histories containing only observations proved to be
more stable in training, than the other two implementation approaches described
in \cref{chapter:SignatureQLearning:section:QImplementation}. 

Further, the environment exhibits two possible terminations. The first one is
due to reaching the maximum episode length $T=200$, and the second due to
reaching a goal position $p_t\geq 0.5$. Both are treated as environmental
terminations and no bootstrapping is performed at either one, as proposed in
\cite{pardo2018_tlrl}. The Q-target values at $k$-th update step are thus
given as
\[
    y_t = \begin{cases}
        r_t, & \mbox{if } t=200 \mbox{ or } p_t\geq 0.5, \\
        r_t + \argmax_{a\in\SpaceA} Q_t(\hat{h}^a_{t+1}\mid \ell_k), 
            & \mbox{otherwise,}
    \end{cases}
\]
where $r_t$ is the sampled reward at time $t$ according to the reward function 
\eqref{eq:MountainCarReward}. W.l.o.g. we assume there are no ties between
actions such that the argmax indeed returns a single action. Otherwise, ties
are broken up by sampling uniformly from the actions that are part of the tie. 

\paragraph{Set-Up One: Fixed parameter configuration} 
\label{chapter:Experiments:section:MountainCar:sub:TrainingDetails:par:SetUpOne}

The first training set-up consists of a fixed hyperparameter configuration,
which is detailed in \cref{tab:MountainCarTrainingParameters}. In total, 10
training runs are performed in set-up one. The hyperparameter configuration
was found through an informal search. The discount factor of $\gamma=0.99$
enabled stable learning behaviour in reducing the Q-function's horizon to
roughly $1/(1-\gamma) = 100$ steps. Regarding the handling of the signature
computation at $t=0$, see
\cref{chapter:SignatureQLearning:section:QImplementation}, the basepoint 
$\hat{o}_{-1} = (1, p_{-1}) = (1, -0.5)$ was used. This amounts to an
artificial position $p_{-1}=-0.5$ before $t=0$, exactly in the middle of the
interval $[-0.6, -0.4]$ in which the car is placed randomly at the beginning
of each episode. As detailed in
\cref{chapter:SignatureQLearning:section:QImplementation}, this choice should
enable the algorithm to differentiate between initial observations $o_0$, or
more precisely initial positions $p_0$, beyond the values provided by the bias
vector. For the initial position $p_0 = -0.5$, the initial observation 
$\hat{o}_0$ is equal to the basepoint $\hat{o}_{-1}$ and the truncated
signature at $t=0$ evaluates to $\bm 1 \in T^{(N)}(\SetR^{d+1})$. In this
case, the selection of the first action $a_0$ solely depends on the values in
the bias vector $b\in\SetR^3$. For any other $p_0\in[-0.6, -0.4] \setminus \{-0.5\}$,
the truncated signature is distinct from $\bm 1$ and as such contains
information about $p_0$. Other values for $p_{-1}$ are considered in training
set-up three, \cref{chapter:Experiments:section:MountainCar:sub:%
TrainingDetails:par:SetUpThree}

The exploration policy used is $\varepsilon$-greedy. The learning rates and
the exploration rates are decreased over training, with the decay schemes
detailed in \cref{tab:MountainCarDecaySchemes}. Since the environment exhibits
a clear goal for each episode, both parameter values are decreased only after
successful episodes to facilitate learning and exploration when the number of
successful episodes is still low. During the episode both values are held
fixed. Using the example of the exploration rate, this means that 
$\varepsilon_{n} := \varepsilon_{n,0} = \varepsilon_{n,t}$ for all $t\in[T]$,
and this value is only decreased to the next smaller value 
$\varepsilon_{n+1}$, if in the current episode the goal was reached, else it
is left the same for the subsequent episode. 

\begin{table}[htpb!]
    \centering
    \caption[%
        Training parameter values in set-up one in the Mountain 
        Car experiment
    ]{%
        The hyperparameter values used in training set-up one for the
        Mountain Car experiment. The parameters correspond to number of
        episodes, maximum episode length, discount factor, loss function,
        gradient descent optimizer, signature truncation order, and 
        basepoint, in the order of columns. The smooth-$L_1$ loss function is
        provided by \eqref{eq:SmoothL1Loss}, the gradient descent optimizer Adam is introduced in \cite{kingma2015_amso}. 
    }
    \label{tab:MountainCarTrainingParameters}        
    \begin{tabular*}{\textwidth}{l @{\extracolsep{\fill}} *{8}{c} @{}}
        \toprule
        \textbf{Parameter} & 
            $M_\text{train}$ & $T$ & $\gamma$ & $\mathcal{L}$ 
                & optimizer & $N$ & $\hat{o}_{-1}$ \\
        \midrule
        \textbf{Value} & 
            4000 & 200 & 0.99 & smooth-$L_1$ & Adam & 5 & $(1, -0.5)$ \\
        \bottomrule
    \end{tabular*}
\end{table}

\begin{table}[htpb!]
    \centering
    \caption[Parameter decay schemes in the Mountain Car experiment]{Details on the decay schemes of the learning rates and $\varepsilon$-greedy exploration rate in the Mountain Car experiment.}
    \label{tab:MountainCarDecaySchemes}    
    \begin{tabular*}{\textwidth}{l @{\extracolsep{\fill}} *{4}{c}}
        \toprule
        \textbf{Parameter} & \textbf{Start value} & \textbf{End value} &
                \textbf{Decay epochs} & \textbf{Decay mode} \\
        \midrule
        Learning rates $\bm\alpha_n$ &
            0.01 & $10^{-5}$ & 3000 successes & Linear \\          
        Exploration rates $\bm\varepsilon_n$ &
            0.3 & 0.05 & 2000 successes & Linear \\          
        \bottomrule
    \end{tabular*}
\end{table}

\paragraph{Set-Up Two: Truncation Orders} 
\label{chapter:Experiments:section:MountainCar:sub:TrainingDetails:par:SetUpTwo}

The second training set-up investigates different signature truncation orders
$N\in\SetN$, a hyperparameter distinct to Signature-Q-Learning. For each
truncation order four training runs are performed, respectively. Apart from
the truncation orders, other parameter values are the same as in set-up one and
detailed in \cref{tab:MountainCarTrainingParameters}. The $\varepsilon$-greedy
exploration rates $\varepsilon_n$ are again decayed over successful episodes
with the decay schemes detailed in \cref{tab:MountainCarDecaySchemes}.
Distinct to set-up one, in set-up two the learning rate is held fixed at 
$\alpha = \alpha_n =0.005$ for all episodes $m\in[M]$ in each training run.
Since the observation history is only two dimensional, it is possible to
truncate at a relatively high order without running into computational
restrictions due to the exponential scaling of the truncated signature's
dimension $d_N$. We compare truncation orders
\[
    N\in\{2,3,4,5,6,7,8\}.
\]
We expect that the algorithm performs better with a higher truncation order in
terms of higher rewards and number of successfully solved episodes.

\paragraph{Set-Up Three: Basepoints}
\label{chapter:Experiments:section:MountainCar:sub:TrainingDetails:par:SetUpThree}

The third set-up investigates different basepoint configurations for handling
the signature computation at $t=0$, as detailed in
\cref{chapter:SignatureQLearning:section:QImplementation}. Specifically, we
perform training runs with three different basepoints and without basepoint.
For each configuration four training runs are performed, respectively. 
Each basepoint is given as a tuple $\hat{o}_{-1} := (1, p_{-1})$, with
$p_{-1}\in\SetR$ an artificial position prior to $t=0$. This tuple is
added to the observation history, such that at $t=0$, the agent receives the 
stream $(\hat{o}_{-1}, \hat{o}_0)$ to compute the signature. The value for
$p_{-1}$ is based on the interval $[-0.6, -0.4]$ from which $p_0$ is sampled
uniformly in each episode, and we compare the values
\[
    p_{-1} \in \{-0.65, -0.5, -0.35\}.
\]
In the case of no basepoint, at $t=0$, the truncated signature is set to 
$\bm 1\in T^{(N)}(\SetR^{d+1})$. Thus, its components apart form the first
entry 1 are zero and only take non-zero values at $t=1$ when there are two
observations available. This results in the selection of action $a_0$ based
solely on the bias vector $b\in\SetR$, since no information about the initial
observation $\hat{o}$ is available to the agent through the signature. 

Apart from the basepoint configuration, the remaining hyperparameters are
again provided in \cref{tab:MountainCarTrainingParameters}. The 
$\varepsilon$-greedy exploration rates $\varepsilon_n$ are again decayed over
successful episodes with the decay schemes detailed in
\cref{tab:MountainCarDecaySchemes} and the learning rate held fixed at
$\alpha=0.005$ for all basepoint configurations and runs.

%------------------------------------------------------------------------------
\subsection{Testing Details} 
\label{chapter:Experiments:section:MountainCar:sub:TestingDetails}
%------------------------------------------------------------------------------

For each training run, we asses the learned approximate Q-function by
performing a test run consisting of $M_\text{test} = 500$ episodes. Let
$Q(\cdot \,|\,\ell_K)$ be the approximate Q-function from a training run, with 
$\ell_K$ the final linear functional of the signature approximation after a
total of $K=T M_\text{train}$ update steps during training. In the
following test run, actions are selected greedily with respect to $Q(\cdot \,|
\, \ell_K)$. That is, the policy followed in testing at $t\in[T-1$] is
given as
\begin{equation}
\label{eq:MountainCarTestPolicy}
    \pi^{\text{test}}_t(h_t) 
    := \argmax_{a\in\SpaceA} Q(\hat{h}_t^a\,|\,\ell_K),
\end{equation}
for any $h_t\in\SpaceH_t$, where we assume that there are no ties between
actions, such that the action that maximizes the right hand side of 
\eqref{eq:MountainCarTestPolicy} is unique\footnote{
    If there happens to be a tie between action values in the experiment, an
    action may be selected by sampling uniformly from the set of actions 
    involved in the tie.
}

For each test run, we report the test statistics in
\cref{tab:MountainCarTestStatistics}. We provide a brief description of their 
calculation. For test episode $n\in\{1,\dots,M_\text{test}\}$,
let $r^{(n)}$ be the sum of undiscounted rewards obtained at
each step in the episode. The mean cumulative, undiscounted episode reward 
$\overline{r}$ and the episode reward standard deviation $\sigma_r$ are calculated
as the sample mean and standard deviation of the sample $\{r^{(n)}\}_{n\in
\{1,\dots,M_\text{test}\}}$. Analogously, let $m^{(n)}:=\inf\{t\in[T]\mid p^{(n)}_{t} 
\geq 0.5\}$ be the length of episode $n$, with $p_t^{(n)}$ the car's position
at step $t$ in episode $n$, and the convention $\inf \varnothing := T$. The 
mean episode length $\bar{m}$ and the episode length standard deviation 
$\sigma_m$ are calculated as the sample mean and sample standard deviation of
the sample $\{m^{(n)}\}_{ n\in\{1,\dots,M_\text{test}\}}$. Further, we have 
$\min_m := \min\{m^{(n)} : n\in\{1,\dots,M_\text{test}\}\}$ and $\max_m := 
\max\{m^{(n)}: n\in\{1,\dots,M_\text{test}\}\}$. For the mean first observation 
value $\bar v^*_0$, recall that in each episode the car's starting position $p_0$,
which is the agent's initial observation $o_0$, is sampled from $\mathcal{U}
[-0.6, -0.4]$. Let $o^{(n)} = p_0^{(n)} \in [-0.6, -0.4]$ be the initial 
position in test episode $n\in\{1,\dots,M_\text{test}\}$ of the current test 
run and $\hat{o}^{(n)}_0 = (1,p^{(n)}_0)$. We define $v_0^{*,n}$ as the 
value of the first observation $\hat{o}^{(n)}_0$ as provided by the final 
approximate Q-function $Q(\cdot\,|\, \ell_K)$, and $\bar{v}_0^*$ as the 
average over test episodes,
\begin{equation}
\label{eq:MountainCarFirstObservationValue}
    v_0^{*,n} := \max_{a\in\SpaceA} Q(\hat{o}^{(n)}_0, a\,|\,\ell_K),
    \qquad \text{and} \qquad
    \bar v^*_0 := \frac{1}{M_\text{test}} \sum_{n=1}^{M_\text{test}} v_0^{*,n}.
\end{equation}
Finally, the mean cumulative, discounted episode reward $\overline{r}_\gamma$, 
with discount factor $\gamma=0.99$, is calculated as the sample mean over the 
sample $\{r_\gamma^{(n)}\}_{n\in\{1,\dots,M_\text{test}\}}$, 
where for each episode $n$, the cumulative, discounted episode reward 
$r_\gamma^{(n)}$ is given as as the sum
\begin{equation}
\label{eq:MountainCarDiscountedReward} 
    r_\gamma^{(n)} = \sum_{t=0}^{m^{(n)}} \gamma^t r_t^{(n)},
\end{equation}
with $r_t^{(n)}$ the reward obtained at step $t$ in test episode $n$.

If the algorithm converged to the true optimal Q-functions, then 
$Q(\hat{o}_0, a\,| \,\ell_K)$ is an estimate of the true Q-value $Q_0^*(o_0,
a)$ for any $(o_0, a)\in\SpaceO\times\SpaceA$. This implies that due to 
\cref{eq:HDPValueQConnection}, the first observation value $v_0^{*,n}$ is an 
estimate of the true optimal value $V_0^*(o^{(n)}_0)$, and by its definition 
as sample mean in \eqref{eq:MountainCarFirstObservationValue}, 
$\bar v_0^*$ is an estimate of the expected optimal value at $t=0$,
\[
    \bar v_0^* \approx \E[V_0^*(O_0)], 
    \qquad \text{with } 
    O_0\sim\mathcal{U}[-0.6, -0.4].
\]
In absence of an optimal policy we can use for comparison, however, we do not 
know whether the algorithm has converged only by looking at the value of $\bar
v_0^*$. Still, it may offer an indication about convergence at least to 
\textit{some} solution, if we compare it to $\overline{r}_\gamma$. Recall 
\cref{def:HDPValueFunctions} for the value functions $(V^\pi_t)_{t\in[T]}$ 
corresponding to a given policy $\pi\in\Pi^\text{HR}$. In the test run, 
actions are selected according to $\pi^\text{test}$, thus, $r^{(n)}_\gamma$ 
is a sample of $V_0^{\pi^\text{test}}(o^{(n)}_0)$, and by its definition as 
sample mean, $\overline{r}_\gamma$ is an estimate of the expected initial 
value of policy $\pi^\text{test}$,
\[
    \overline{r}_\gamma \approx \E[V^{\pi^\text{test}}_0(O_0)]
    \qquad \text{with } 
    O_0\sim\mathcal{U}[-0.6, -0.4].    
\]
If we obtain $\overline{r}_\gamma \approx \bar v_0^*$, we may conclude, that in
this test run the algorithm has converged to some Q-functions. More 
specifically, the approximate Q-function correctly returns the action-%
values of the policy $\pi^\text{test}$.

\begin{table}[htpb!]
    \caption[
        Description of reported test statistics in the Mountain Car Experiment
    ]{
        Description of the reported test statistics in the Mountain Car
        Experiment. Statistics are calculated for each performed test run and
        reported in
        \cref{chapter:Experiments:section:MountainCar:sub:Results}.
    }
    \label{tab:MountainCarTestStatistics}    
    
    \centering    
    \begin{tabularx}{\textwidth}{>{\centering}m{4cm}X}
        \toprule
         ${\overline r}$ & Mean undiscounted reward per episode \\
         ${\sigma_r}$ & 
            Standard deviation of undiscounted reward per episode \\
         ${\overline{r}_\gamma}$ & Mean discounted reward per episode \\
         {succ.} & Percentage of successful episodes in test run \\
         ${\bar m}$ & Mean episode length until termination \\
         ${\sigma_m}$ & 
            Standard deviation of episode length until termination \\
         ${\min_m}$ & Minimum episode length until termination \\
         ${\max_m}$ & Maximum episode length until termination \\
         ${\bar v^*_0}$ & Mean first observation value
            \eqref{eq:MountainCarFirstObservationValue} \\
        \bottomrule
    \end{tabularx}
\end{table}


%------------------------------------------------------------------------------
\subsection{Results} 
\label{chapter:Experiments:section:MountainCar:sub:Results}
%------------------------------------------------------------------------------

The following paragraphs present the results for the three training set-ups
described in \cref{chapter:Experiments:section:MountainCar:sub:TrainingDetails}.
Results are compared based on the statistics in
\cref{tab:MountainCarTestStatistics} calculated from performed test runs.

%------------------------------------------------------------------------------
\paragraph{Set-Up One: Fixed parameter configuration}
\label{chapter:Experiments:section:MountainCar:sub:Results:par:SetUpOne}

\Cref{tab:MountainCarResultsStatisticsSetUpOne} provides a comparison of the different
test runs performed for each training run in set-up with a fixed hyperparameter configuration. We note that, Signature-Q-Learning algorithm
was able to consistently solve the problem in all runs, showing that the
approach proposed in this work indeed works for a history-dependent problem
given as the partially observable MDP at hand. In all 10 runs, the 
algorithm was able to
learn a Q-function approximation that reaches the goal position in 100\% of
test episodes. The difference between the lowest and highest mean undiscounted 
episode reward reward $\overline r$ is 3.86\% of the highest absolute value, 
indicating that for this specific hyperparameter configuration the algorithm is 
not sensitive to parameter initialization and the seed used to simulate the 
environment. This is also apparent from the mean \textit{discounted} episode 
rewards $\overline r_\gamma$, which is the measure the algorithm optimizes for. 
Its sampled values are practically the same, agreeing up to a prevision of 
$10^{2}$. At the same time the standard deviation of rewards in each run is 
relatively high, with all values in the interval $[-0.145, -0.165]$. This
variability can be explained by large differences in episodes lengths. While
the mean episode length $\bar{m}$ is very similar for all runs, the
difference between the minimum $\min_m$ and maximum $\max_m$ episode lengths is 
more than 100 steps in all runs but one. This has a direct consequence on
rewards, due to the reward structure \eqref{eq:MountainCarReward}, that is 
revealed by \cref{fig:MountainCarScatterplots}. The figure shows episode 
lengths and obtained undiscounted episode rewards as a function of the starting 
position $p_0$ for test run 8. The plots show a clear split around 
$p_0\approx -0.45$ with vastly different behaviour for smaller and larger
starting positions for both statistics, which results in a distribution of 
episode rewards that is heavily clustered around two values.
Although this causes high variance in both statistics, it is also an indication that the algorithm learned to
differentiate starting positions and follows different action sequences
depending on them, to increase rewards. This is in particular visible when
comparing the result in set-up one, to the ones in the following set-up two
with the same truncation order $N=5$, where a more homogeneous policy was
learned that however delivers smaller overall episode rewards.

% test result statistics table Mountain Car set-up
\begin{table}[htpb!]
    \centering
    \caption[%
        Test statistics in the Mountain Car experiment for training set-up one
    ]{%
        The table displays test statistics for the Mountain Car experiment
        \textbf{set-up one} in
        \cref{chapter:Experiments:section:MountainCar:sub:%
        TrainingDetails:par:SetUpOne}, when following the obtained approximate
        Q-function from each training run for 500 test episodes. Bold values
        highlight the best result over runs: the highest reward, the lowest
        standard deviation, or the shortest episode length. We refrained from
        highlighting the column containing success rates, since for each run
        it is 100\%, and also the column with initial optimal values, because
        the obtained values are not meaningful due to the overestimation of 
        Q-values. $\bar v^*_0$ is calculated as in
        \cref{eq:MountainCarFirstObservationValue} with the basepoint 
        $\hat{o}_{-1} = (1, p_{-1}) = (1, -0.5)$ for the computation of the
        signature at time $t=0$.
    }
    \label{tab:MountainCarResultsStatisticsSetUpOne}        
    \begin{tabular*}{\textwidth}{%
        >{\centering}m{0.7cm} | @{\extracolsep{\fill}}  *{9}{c}}
    %\begin{tabularx}{\textwidth}{ >{\centering}m{0.7cm} | *{9}{C} }
        \toprule
        Run
            & ${\overline r}$ & ${\overline{r}_\gamma}$ & ${\sigma_r}$
                & {succ.} & ${\bar m}$ & ${\sigma_m}$ 
                    & ${\min_m}$ & ${\max_m}$ & ${v^*_0}$ \\
        \midrule
        \midrule

1 & $-0.6475$ & $-0.3538$ & 0.1606 & 100\% & 144.79 & 29.98 & $\bm{86}$ & 199 & 267.67 \\
2 & $-0.6638$ & $-0.3590$ & $\bm{0.1456}$ & 100\% & 147.44 & 27.10 & 87 & 197 & 263.70\\
3 & $-0.6546$ & $-0.3563$  & 0.1537 & 100\% & 145.72 & 28.45 & $\bm{86}$ & $\bm{185}$ & 263.12\\
4 & $-0.6591$ & $-0.3570$  & 0.1474 & 100\% & 147.04 & 28.00 & 87 & 195 & 268.42\\
5 & $-0.6514$ & $-0.3558$  & 0.1522 & 100\% & 144.61 & 27.13 & 87 & 193 & 269.61\\
6 & $-0.6580$ & $-0.3562$  & 0.1523 & 100\% & 147.29 & 28.98 & $\bm{86}$ & 199 & 272.03\\
7 & $-0.6603$ & $-0.3577$ & 0.1465 & 100\% & 146.88 & $\bm{27.03}$ & 87 & 196 & 271.70 \\
8 & $-0.6431$ & $-0.3523$ & 0.1627 & 100\% & 144.28 & 30.23 & $\bm{86}$ & 199 & 267.23 \\
9 & $\bm{-0.6391}$ & $\bm{-0.352}$ & 0.1644 & 100\% & $\bm{142.54}$ & 29.99 & $\bm{86}$ & 196 & 275.27\\
10 & $-0.6476$ & $-0.3534$ & 0.1635 & 100\% & 145.22 & 31.76 & $\bm{86}$ & 199 & 260.34 \\
        \bottomrule
    \end{tabular*}
\end{table}

% scatterplots Mountain Car set-up one
\begin{figure}[htbp!]
    \centering
    \begin{subfigure}{0.49\textwidth}
        \includegraphics[width=\textwidth]{%
            graphics/mountaincar/20240905_A_start_pos_vs_steps.png
        }
        \caption{Start position vs. steps till termination.}
        \label{fig:MountainCarScatterplots:sub:start_pos_vs_steps}
    \end{subfigure}
    \hfill
    \begin{subfigure}{0.49\textwidth}
        \includegraphics[width=\textwidth]{%
            graphics/mountaincar/20240905_A_start_pos_vs_reward.png
        }
        \caption{Start position vs. loss per episode.}
        \label{fig:MountainCarScatterplots:sub:start_pos_vs_reward}
    \end{subfigure}            
    \caption[Test results scatterplots in the Mountain Car experiment
    ]{%
        The figures displays scatter plots for test run 9 from set-up one in
        the Mountain Car experiment. It shows that the learned Q-function
        approximation differentiates between the car's starting positions
        $p_0\in[-0.6, -0.4]$ and follows different action sequences based on
        them.
        Plot (\subref{fig:MountainCarScatterplots:sub:start_pos_vs_steps}) displays
        the number of steps until termination as a function of the starting
        position $p_0$ for each episode from test run 9.
        Plot (\subref{fig:MountainCarScatterplots:sub:start_pos_vs_reward}) displays
        the total undiscounted reward obtained as a function of the starting
        position $p_0$ for each episode from test run 9.
    }
    \label{fig:MountainCarScatterplots}
\end{figure}

\Cref{fig:MountainCarTrainingResults} shows training performance in set-up 
one, averaged over the 10 performed training runs. It again highlights the
variability in episode rewards
(\subref{fig:MountainCarTrainingResults:sub:Reward}), and in steps until
termination (\subref{fig:MountainCarTrainingResults:sub:Steps}), due to the
heterogeneous actions sequences followed based on starting positions. Overall
however, all runs exhibit a similar training behaviour and convergence to the
same Q-functions approximation. The algorithms success of learning an
approximation that solves the problem is apparent from
\cref{fig:MountainCarTrainingResults:sub:Steps}, which visualized the mean end
position for each training episode. It shows that after around 2000 episodes,
all runs managed to reach a goal position of $p_t \geq 0.5$ in almost all
episodes, and towards the end of the training runs in 100\% of episodes.

% averaged training result plots
\begin{figure}[htbp!]
    \centering
    \begin{subfigure}{0.49\textwidth}
        \includegraphics[width=\textwidth]{graphics/mountaincar/20240905_A_mean_reward.png}
        \caption{Average reward per episode.}
        \label{fig:MountainCarTrainingResults:sub:Reward}
    \end{subfigure}
    \hfill
    \begin{subfigure}{0.49\textwidth}
        \includegraphics[width=\textwidth]{graphics/mountaincar/20240905_A_mean_loss.png}
        \caption{Average loss per episode.}
        \label{fig:MountainCarTrainingResults:sub:Loss}
    \end{subfigure}
    \hfill
    \begin{subfigure}{0.49\textwidth}
        \includegraphics[width=\textwidth]
            {graphics/mountaincar/20240905_A_mean_end_position.png}
        \subcaption{Average end position per episode.}
        \label{fig:MountainCarTrainingResults:sub:EndPosition}
    \end{subfigure}
    \hfill
    \begin{subfigure}{0.49\textwidth}
        \includegraphics[width=\textwidth]
            {graphics/mountaincar/20240905_A_mean_steps_to_termination.png}
        \caption{Average steps till termination per episode.}
        \label{fig:MountainCarTrainingResults:sub:Steps}
    \end{subfigure}
            
    \caption[Averaged training results of set-up one in the Mountain Car experiment
    ]{%
        Averaged training results of the Signature-Q-Learning algorithm in the
        Mountain Car experiment set-up one \cref{chapter:Experiments:%
        section:MountainCar:sub:TrainingDetails:par:SetUpOne}. 10 training runs
        were performed in total with the hyperparameter set-up detailed in
        \cref{tab:MountainCarTrainingParameters}. The sub-figures show the
        respective statistic as the mean over training runs for each episode.
        (\subref{fig:MountainCarTrainingResults:sub:Reward}) shows the average 
        cumulative, undiscounted reward obtained per episode.
        (\subref{fig:MountainCarTrainingResults:sub:Loss}) the average total 
        loss obtained per episode.
        (\subref{fig:MountainCarTrainingResults:sub:EndPosition}) shows the 
        average of the car's end position at the episode termination.
        (\subref{fig:MountainCarTrainingResults:sub:Steps}) shows the average 
        number of steps until the episode terminated, either because the car 
        reached a goal position or the maximum number of steps $T=200$.     
    }
    \label{fig:MountainCarTrainingResults}
\end{figure}

One issue Signature-Q-Learning suffers from in this particular problem, is
the overestimation of Q-values. Although the algorithm was able to consistently
reach the goal position, the approximate Q-values for pairs $(h_t,a)$ did not
converge to their true values provided by the true optimal Q-functions $Q_t^*,
t\in[T]$. This is apparent from the high values obtained for the mean first 
observation value $\bar v_0^*$. Recall, its definition from 
\cref{eq:MountainCarFirstObservationValue}, and that it should approximately
match the cumulative, discounted reward $\overline r_\gamma$ for
any run in 
which the approximate Q-function has converged. Clearly, for any test run we 
have $\bar v_0^* \not\approx \overline r_\gamma$ in
\cref{tab:MountainCarTestStatistics}. Additionally, due to the reward
structure \eqref{eq:MountainCarReward}, the true expected first observation
value $\E[V_0^*(O_0)]$ must take its value in the interval $[-0.735, 0]$. As 
such any sufficiently accurate estimate must fall into the same interval, 
however, clearly $\bar v_0^* \notin [-0.735, 0]$ for all test runs. The 
overestimation further  results in high episode losses during training, 
visible in \cref{fig:MountainCarTrainingResults:sub:Loss}. While the loss
stabilized around the value 3000, its variability over episodes remains high,
which may be due to the combination of overestimated Q-values, large
differences in episodes lengths, and the fact that one set of parameters is
used to approximation optimal Q-functions at all time-steps. As an heuristic
example, consider an episode in which at $t=86$ a goal position
$p_{86}\geq 0.5$ is reached, which corresponds to the minimum episode
length for most runs in \cref{tab:MountainCarResultsStatisticsSetUpOne}. For
simplicity, assume the loss function $\mathcal{L}$ in 
\cref{alg:SignatureQLearning} is the absolute loss $\mathcal{L}(x,y) = 
\abs{x-y}$ for any $x,y\in\SetR$. The following reasoning, however, also 
applies to the smooth-$L_1$ loss \eqref{eq:SmoothL1Loss} used in the 
experiment. Recall the definition of the Q-target value $y_t$ from 
\cref{eq:HDPQTargetValue} used to update the Q-functions approximation. As 
described in the training details 
\cref{chapter:Experiments:section:MountainCar:sub:TrainingDetails}
no bootstrapping is performed at terminations and the target value at $t-1=85$
is given as $y_{85} = R_{85}(\hat{o}_{85})=0.016245$, for the reward function
\eqref{eq:MountainCarReward}. In any other episode that does not terminate at
$t=86$, the target value additionally includes the maximum action value 
$\max_{a\in\SpaceA} Q(\hat{h}_{86}^a\,|\,\ell)$ at the next step $t=86$, which
presents an estimate of the maximal attainable cumulative discounted reward
over potentially 100 or more remaining steps. An overestimation of this value
results in vastly different errors for step $t=85$, which in both cases is
calculated as $\bigl| y_{85} - Q(\hat{h}_{85}^{a_{85}}\,|\,\ell) \bigr|$, to
update the set same set of parameters $\ell$. If during training short
episodes alternate with significantly longer episodes, which is the case in
this experiment as apparent from
\cref{fig:MountainCarScatterplots:sub:start_pos_vs_steps}, the overestimation
issue results in large errors at a number of steps.

Overestimation of Q-values is a known issue for Q-Learning in MDPs
\cite{hasselt2010_dq,hasselt2016_drlwdq}, especially in combination with
function approximation as in Deep Q-Networks \cite{mnih2015_hctdrl}. In
\cite{hasselt2016_drlwdq}, the authors argue that estimation errors, that 
emerge from the use of function approximation or environmental noise,
introduce an upwards bias in Q-function approximations, and they provide a
lower bound for the error (Theorem 1 \cite{hasselt2016_drlwdq}). Their 
statement readily transfers to the history-based setting considered in the 
present work. Assume $\abs{\SpaceA} >1$, and consider a time $t\in[T]$ and 
history $h_t\in\SpaceH_t$, for which the true optimal Q-values are equal at 
$Q_t^*(h_t^a) = V^*_t(h_t)$ for all $a\in\SpaceA$, with $V_t^*$ the true 
optimal value function \eqref{eq:HDPOptimalValueFunction} at $t$. Further 
consider some Q-function approximation provided by the linear functional 
$\ell$, such that
%\footnote{As mentioned before, in the Mountain Car experiment only observations are used in time-augmented histories, that serve as input to the Q-function approximation. However for ease of exposition and consistency with prior sections, we also write $Q(\hat{h}^a_{t}\,|\,\ell)$ for the Q-function approximation of the pair $(\hat{h}_t,a)$, with the understanding that only the time-augmented observation history serves as input to the approximation and its output corresponding to the action value $a$ is chosen.}
\[
    \sum_{a\in\SpaceA} Q(\hat{h}^a_t\,|\,\ell) -  V_t^*(h_t) = 0,
    \quad \text{and} \quad
    \abs{\SpaceA}^{-1} \sum_{a\in\SpaceA} 
        \bigl(Q(\hat{h}^a_t\,|\,\ell) -  V_t^*(h_t)\bigr)^2 = C
\]
for some $C>0$. Then, we have the lower bound
\[
    \max_{a\in\SpaceA} Q(\hat{h}^a_t\,|\,\ell) 
    \geq V_t^*(h_t) + \sqrt{\frac{C}{\abs{\SpaceA}-1}}.
\]
While this lower bound decreases with the number of actions $\abs{\SpaceA}$,
the authors in \cite{hasselt2016_drlwdq} show that the actual overestimation
typically increases with the size of $\SpaceA$. The usage of the maximum 
action value $\max_{a\in\SpaceA} Q(h_t,a\,|\,\ell)$ to approximate 
\textit{expected} maximum action value from the optimal Q-iteration 
\eqref{eq:HDPOptimalQIteration} in the bootstrapping procedure 
\eqref{eq:HDPQApproxUpdate}, further reinforces the overestimation and
propagates the error through estimates, see also \cite{hasselt2010_dq}. An
overestimation however, does not have to result in a poor policy, as argued in
\cite{hasselt2016_drlwdq}, and which is also apparent in the experiment at
hand. If the overestimation occurs uniformly over state-action pairs, in our
case history-action pairs, the relative action preferences, and in this way
maximizing actions, are maintained. 

To overcome the issue of overestimation of Q-values in MDPs,
\cite{hasselt2010_dq} proposed the \textit{Double Q-Learning} algorithm in a
tabular setting, which was later extended to the approximation of Q-functions
with deep neural networks, resulting in Double Deep-Q-Networks
\cite{hasselt2016_drlwdq}. The principle behind Double Q-learning is to
disentangle the agent's action selection from the estimation of the selected
action's value through the target value $y_t$. Two distinct Q-function
approximation either in tabular or parametrized form are used, one for each
of the task, and alternately updated with the encountered environmental
transitions. In MDPs, this proofed to avoid the overestimation of Q-values and
improve performance, especially in conjunction with function approximation
\cite{hasselt2016_drlwdq}. The same principle could be used here, to
investigate whether it mitigate the overestimation of Q-values, by using one
linear functional with parameters $\ell_k^{(1)}$ to select the action at each
step, and another functional with parameters  $\ell_k^{(2)}$ to create the
target value $y_t$.

%------------------------------------------------------------------------------
\paragraph{Set-Up Two: Truncation Orders}
\label{chapter:Experiments:section:MountainCar:sub:Results:par:SetUpTwo}

\Cref{tab:MountainCarResultsStatisticsSetUpTwo} shows the test results for
training set-up two with different truncation orders. As expected, the
performance increased with higher truncation order $N$, however only up to a
certain order. For $N=2$, Signature-Q-Learning did not learn a Q-function
approximation that solves the problem consistently in any run and obtained
rewards are generally lower than with higher truncation orders. The results
for $N=3$ are similarly poor, with the exception of run 4, which performed
best in a number of criteria, indicated by values highlighted in bold. It
obtained the highest mean reward over truncation orders not only in run 4 but
over all runs, as well as the overall shortest mean episode length,
corresponding shortest standard deviation and smallest maximum episode length.
However, the outperformance compared to successful runs of higher order is
very narrow. Further, considering that only one out of four runs was able to
achieve 100\% success rate, indicates that the order $N=3$ is still too low
for consistent successful training.

For truncation orders $N\geq 4$ Signature-Q-Learning achieved a 100\% success
rate in all test runs. While performance still increases from truncation order
$N=4$, with $N=5$ exhibiting slightly higher mean reward and shorter mean
episode length, higher orders yielded no further improvement. For $N\geq5$, no
truncation order emerges as clearly best performing, with all reported
statistics taking very similar values regardless the value of $N$. The only
exception to this is the mean first observation value $v_0^*$, which increases
with the truncation order, indicating a growing overestimation of Q-values with
larger $N\in\SetN$. This effect may caused by the larger number of parameters
with higher truncation order together with the fixed learning rate 
$\alpha=0.005$. Overall, the results show, that in the Mountain Car
environment, a low truncation order of $N=5$, which results is sufficient.
Note, that for $N=5$, the Q-function approximation, implemented as described
in the training details in
\cref{chapter:Experiments:section:MountainCar:sub:TrainingDetails}, contains
in total $\abs{\SpaceA}\cdot d_N = 189$ trainable parameters. 

In terms of the policy derived from the approximate Q-function, all
successful truncation orders $N\geq5$ learned the same solution. This is
apparent from the remarkably similar values for mean episode length and
corresponding standard deviation as well as minimum and maximum episode
length. It is evident, that for this specific hyperparameter set-up and 
basepoint configuration, \cref{alg:SignatureQLearning} is robust in Signature
truncation order $N\in\SetN$.

For truncation order $N=5$ in set-up two, the only difference to set-up one in
the training configuration is the the use of a fixed rate $\alpha=0.005$ in
set-up two. This resulted in a different Q-function approximation and thus,
also in a different policy, which is apparent from comparing the statistics
reported for $N=5$ in \cref{tab:MountainCarResultsStatisticsSetUpTwo} for 
set-up two, to the ones in \cref{tab:MountainCarResultsStatisticsSetUpOne} for
set-up one. While the success rate for all involved runs is 100\%, there is a
clear difference in obtained reward and episode characteristics. Notable, the
policy followed in set-up two, results in a much smaller range of episode
lengths, with the difference between minimum and maximum length roughly 40
steps, compared to over 100 steps in set-up one. This more homogeneous policy
w.r.t to episode lengths however, comes at the cost of a lower reward. \vskip 2em

\begin{table}[htpb!]
    \centering
    \caption[%
        Test statistics in the Mountain Car experiment for training set-up two
    ]{%
        The table displays test statistics for the Mountain Car experiment
        \textbf{set-up two} in
        \cref{chapter:Experiments:section:MountainCar:sub:Results:par:SetUpTwo},
        when following the obtained approximate Q-function from each training
        run for 500 test episodes. Bold values highlight the best result over 
        truncation orders $N$ for each run: highest reward, lowest
        standard deviation, highest successes rate or shortest episode length.
        Note, that truncation order $N=3$ resulted in the lowest standard 
        deviation $\sigma_r$ for Run 1, however, this configuration is excluded
        from the ranking since the goal position was not reached in any episode.
        Further, we refrained 
        from highlighting the column containing mean first observation values 
        $\bar v^*_0$, due to the overestimation of those values, and from 
        reporting the cumulative, discounted reward $\overline r_\gamma$, since its 
        comparison with $\bar v_0^*$ would be superfluous. $\bar v^*_0$ is 
        calculated as in \cref{eq:MountainCarFirstObservationValue} with the
        basepoint $\hat{o}_{-1} = (1, p_{-1}) = (1, -0.5)$ for the
        computation of the signature at time $t=0$.
    }
    \label{tab:MountainCarResultsStatisticsSetUpTwo}    
    
    \begin{tabular*}{0.98\textwidth}{%
        >{\centering}m{0.4cm} c | @{\extracolsep{\fill}}  *{8}{c} 
    }
        \toprule
        ${N}$ & {Run} & 
            ${\overline{r}}$ & ${\sigma_r}$ & {succ.} & ${\bar{m}}$ & 
                ${\sigma_m}$ & ${\min_m}$ & ${\max_m}$ & 
                    ${\bar v^*_0}$ \\    
        \midrule
        \midrule

% truncation order 2
\multirow{4}{*}{${2}$} 
& 1 & $-1.0515$ & 0.0677 & 8.0\% & 197.41 & $\bm{8.86}$ & 164 & 200 & 12.82 \\
& 2 & $-0.6981$ & 0.1213 & 91.6\% & 151.03 & 24.79 & $\bm{86}$ & 200 & 10.00 \\
& 3 & $-0.7930$ & 0.0480 & 71.6\% & 173.17 & 20.33 & 148 & 200 & 7.01 \\
& 4 & $-0.9572$ & 0.0736 & 74.8\% & 176.74 & 14.91 & 163 & 200 & 8.51 \\

        \midrule

% truncation order 3
\multirow{4}{*}{${3}$} 
& 1 & $-0.9472$ & 0.0197 & 0.0\% & 200.00 & 0.00 & 200 & 200 & 1.23 \\
& 2 & $-0.7457$ & 0.1239 & 82.4\% & 152.62 & 28.35 & 132 & 200 & 2.96 \\
& 3 & $-0.8994$ & 0.0925 & 43.8\% & 183.03 & 19.28 & 157 & 200 & 0.49 \\
& 4 & $\bm{-0.6801}$ & 0.0239 & $\bm{100.0\%}$ & $\bm{145.29}$ & $\bm{8.63}$ & 147.24 & $\bm{164}$ & 139.97 \\       
    
        \midrule

% truncation order 4
\multirow{4}{*}{${4}$} 
& 1 & $-0.6948$ & 0.0224 & $\bm{100.0\%}$ & 148.98 & 9.35 & 135 & 171 & 251.46 \\
& 2 & $-0.6926$ & 0.0214 & $\bm{100.0\%}$ & 148.40 & $\bm{9.01}$ & 135 & 170 & 248.97 \\
& 3 & $-0.6918$ & 0.0225 & $\bm{100.0\%}$ & 148.24 & 9.39 & 133 & 170 & 255.19 \\
& 4 & $-0.6925$ & $\bm{0.0221}$ & $\bm{100.0\%}$ & 148.24 & 9.26 & 135 & 170 & 240.24 \\ 

        \midrule
        
% truncation order 5
\multirow{4}{*}{${5}$} 
& 1 & $-0.6907$ & 0.0239 & $\bm{100.0\%}$ & 147.64 & 9.89 & 130 & 171 & 307.49 \\
& 2 & $-0.6894$ & 0.0217 & $\bm{100.0\%}$ & 147.21 & 9.22 & 133 & 170 & 283.08 \\
& 3 & $\bm{-0.6841}$ & 0.0244 & $\bm{100.0\%}$ & $\bm{145.46}$ & $\bm{9.21}$ & $\bm{130}$ & 170 & 302.91 \\
& 4 & $-0.6896$ & $\bm{0.0221}$ & $\bm{100.0\%}$ & 147.74 & 9.60 & 133 & 168 & 285.18 \\

        \midrule
        
% truncation order 6
\multirow{4}{*}{${6}$} 
& 1 & $-0.6887$ & 0.0237 & $\bm{100.0\%}$ & $\bm{146.98}$ & 9.66 & $\bm{129}$ & 170 & 317.27 \\
& 2 & $\bm{-0.6872}$ & $\bm{0.0205}$ & $\bm{100.0\%}$ & $\bm{146.90}$ & $\bm{9.10}$ & 132 & $\bm{166}$ & 296.99 \\
& 3 & $-0.6884$ & 0.0226 & $\bm{100.0\%}$ & 147.62 & 9.98 & 131 & $\bm{168}$ & 291.95 \\
& 4 & $-0.6899$ & 0.0229 & $\bm{100.0\%}$ & 147.43 & 9.63 & 131 & 170 & 322.98 \\

        \midrule

% truncation order 7
\multirow{4}{*}{${7}$} 
& 1 & $-0.6884$ & 0.0230 & $\bm{100.0\%}$ & 147.02 & 9.78 & 131 & 169 & 329.42 \\
& 2 & $-0.6895$ & 0.0231 & $\bm{100.0\%}$ & 147.67 & 9.86 & 131 & 168 & 312.88 \\
& 3 & $-0.6878$ & $\bm{0.0211}$ & $\bm{100.0\%}$ & 147.02 & 9.25 & 132 & $\bm{168}$ & 313.78 \\
& 4 & $-0.6875$ & 0.0233 & $\bm{100.0\%}$ & 146.70 & 10.12 & $\bm{129}$ & 169 & 319.66 \\   
        
        \midrule
        
% truncation order 8
\multirow{4}{*}{${8}$} 
& 1 & $\bm{-0.6879}$ & $\bm{0.0210}$ & $\bm{100.0\%}$ & 147.42 & 9.63 & 130 & $\bm{167}$ & 310.34 \\
& 2 & $-0.6890$ & 0.0225 & $\bm{100.0\%}$ & 147.32 & 9.34 & 130 & 169 & 323.90 \\
& 3 & $-0.6891$ & 0.0219 & $\bm{100.0\%}$ & 147.55 & 9.58 & 131 & $\bm{168}$ & 313.16 \\
& 4 & $-0.6847$ & 0.0242 & $\bm{100.0\%}$ & 145.72 & 9.60 & 130 & 169 & 321.06 \\

        \bottomrule
    \end{tabular*}
\end{table} 

%------------------------------------------------------------------------------
\paragraph{Set-Up Three: basepoints}
\label{chapter:Experiments:section:MountainCar:sub:Results:par:SetUpThree}

\Cref{tab:MountainCarResultsStatisticsSetUpThree} reports the test 
results for training set-up three. It shows that the different basepoint 
configurations have a large effect on Signature-Q-Learning performance 
and the policy derived from the learned Q-function approximations. In 
contrast to the truncation orders in set-up two, there usually is a clear 
basepoint configuration which performed best w.r.t to a certain test 
statistic.

In terms of obtained reward, a basepoint with artificial position $p_{-1}
=-0.35$ performed best in all four runs by a significant margin. On 
average, the mean reward obtained is 8.75\% higher than for the 
artificial position $p_{-1}=-0.5$, which achieved the second highest 
reward in all but one runs. Notably, the reward with $p_{-1}=-0.35$ is 
also larger than in set-up one, with the tuned hyperparameter setting, 
showing the effect the basepoint configuration has on performance. 
Further, this configuration also exhibits the shortest mean episode 
length in all but one run, and the shortest minimum episode length, with 
the algorithm reaching a goal position in 84 steps in the fastest case. 
At the same time however, the learned Q-function estimates do not solve 
the problem consistently, with success rates up to two percentage points 
lower than 100\%. The failure to solve all test episodes, yields the 
highest variability in obtained rewards and episode lengths of any
configuration, visible from the reported corresponding standard 
deviations. This indicates that the policy followed, strongly 
differentiates between starting positions, managing to reach the goal 
very quickly for some, and not at all for others. This is clearly visible 
in \cref{fig:MountainCarBasePointEpisodeLengths:sub:-0.35}, where there 
is a clear cut in episode length as a function of the starting position 
for the episodes in test run 1 at around $p_0\approx -0.4875$. 
Interestingly, this cut is closer to the middle of the starting interval 
$[-0.6, -0.4]$, than the one in 
\cref{fig:MountainCarScatterplots:sub:start_pos_vs_steps} from set-up one,
although there a basepoint position of $-0.5$ was used.

The basepoint position $p_{-1}$ appeared as the least performing with the
reported statistics over the performed runs showing no consistent 
picture. While two runs managed to achieve a 100\% success rate, the 
corresponding rewards are lower than with positions $p_{-1}\in\{-0.35, 
-0.5\}$. Further, in the successful runs, this configuration did not 
differentiate between starting positions, but followed a policy that 
resulted in a linear decrease in the number of steps until termination, 
with larger $p_0\in[-0.6, -0.4]$, visible in 
\cref{fig:MountainCarBasePointEpisodeLengths:sub:-0.65}. This linear 
relationship is reversed for rewards, with a linear decrease in obtained
rewards the larger $p_0$.

For basepoint position $p_{-1}=-0.5$ and the configuration with no 
basepoint, denoted as \textit{None} in
\cref{tab:MountainCarResultsStatisticsSetUpThree} the algorithm solved the
problem consistently, achieving 100\% success rate in all four runs,
respectively. Setting $p_{-1}=-0.5$ yields the same training set-up as in 
set-up two and the corresponding values in
\cref{tab:MountainCarResultsStatisticsSetUpThree} are consistent with the 
ones reported in \cref{tab:MountainCarResultsStatisticsSetUpTwo} for the 
previous set-up. Compared to the other basepoint configurations, it the 
obtained rewards are the second highest in all but run 3, and it 
delivered a success rate of 100\% in all four runs. The learned policy is 
qualitatively the same as with the previous basepoint position, apparent 
from the similarity of
\cref{fig:MountainCarBasePointEpisodeLengths:sub:-0.5} to the one for 
$p_{-1} = -0.65$.

Using no basepoint for the computation of the signature at $t=0$, 
yielded a perceptibly different policy to the other configurations. While 
this configuration yielded the lowest rewards of all basepoint 
configurations, it was also the most consistent, with the the value of 
$\overline r$ falling in the interval $[-0.7, -0.71]$ in each run, and an 
absolute difference of $0.006$ between the largest and smallest mean 
cumulative reward. Furthermore, it yielded the lowest standard deviation 
of mean rewards, being below $0.0155$ in each run. The learned policies 
follow the same action sequence, regardless of the starting position 
$p_0$. Achieving a 100\% success rate in all runs, the largest difference 
between minimum and maximum episode length is 7 steps, resulting in a 
standard deviation of less than 1.3 steps in each run. This effect is 
consistent expectations, due to the translation invariance of the 
signature if histories without a fixed basepoint are used as input. The 
homogeneous distribution of episode lengths is also visible in
\cref{fig:MountainCarBasePointEpisodeLengths:sub:None}, which shows 
distinctly different picture than for the other configurations. Using no 
basepoint further resulted in the lowest maximum episode length of all 
configurations, where in the longest episode the car reached a goal 
position at $t=153$.

In summary, the results from set-up three show that the basepoint
configuration can have a large effect on the algorithms performance and
learned Q-function approximation, with each configuration resulting in
different policies. Except for $p_{-1}$, the results within a configuration
are robust, converging to the same Q-function approximation in all four runs.
Further, the findings are consistent with the theory w.r.t incorporation
information about the initial observation. For a fixed basepoint, the 
Q-function approximation contains this information and is able to
differentiate between starting positions, resulting in varying action
sequences depending on $p_0$. Without a basepoint, this information is lost
and the translation invariance of the signature maintained in this case,
yields a policy that does not differentiate between start, but follows the
same action sequence regardless of the value $p_0$. Finally, the results in
\cref{tab:MountainCarResultsStatisticsSetUpThree} indicate some trade-off
between higher reward and a consistent success rate. The high rewards with
$p_{-1} =-0.35$ come at the cost of high variance and a below 100\% success 
rate, while for the homogeneous and low variance policy obtained with no 
basepoint, one has to accept the lowest rewards of all tested configuration,
and with the performance of basepoint position $p_{-1}=-0.5$ falling
somewhere in the middle of the other two. Depending on the specific criteria
one wants to optimize in a problem, a certain basepoint configuration can be
selected.

% test statistics set-up three basepoints
\begin{table}[htpb!]
    \centering
    \caption[%
        Test statistics in the Mountain Car experiment for training set-up
        three
    ]{%
        The table displays test statistics for the Mountain Car experiment 
        \textbf{set-up three} in
        \cref{chapter:Experiments:section:MountainCar:sub:Results:%
        par:SetUpThree}, when following the obtained approximate Q-function
        from each training run for 500 test episodes. Bold values highlight the
        best result over basepoint positions $p_{-1}$ for each run: highest
        reward, lowest standard deviation, most successes or shortest episode
        length. We refrained from highlighting the column containing mean 
        first observation values $\bar v^*_0$, due to the overestimation of 
        those values, and from reporting the cumulative, discounted reward 
        $\overline r_\gamma$, since its comparison with $\bar v_0^*$ would be 
        superfluous. $\bar v^*_0$ is calculated as in 
        \cref{eq:MountainCarFirstObservationValue} with basepoints 
        $\hat{o}_{-1} = (1, p_{-1})$ where $p_{-1}$ is set to the respective
        value in the first column.
    }
    \label{tab:MountainCarResultsStatisticsSetUpThree}
    
    \begin{tabular*}{\textwidth}{cc| @{\extracolsep{\fill}}  *{8}{c} }
        \toprule
        ${p_{-1}}$ & {Run} & 
            ${\overline{r}}$ & ${\sigma_r}$ & {succ.} & ${\bar{m}}$ & 
                ${\sigma_m}$ & ${\min_m}$ & ${\max_m}$ & 
                    ${\bar v^*_0}$ \\    
        \midrule
        \midrule
        
% basepoint position -0.65
\multirow{4}{*}{${-0.65}$} 
& 1 & $-0.6934$ & 0.0340 & $\bm{100.0\%}$ & 145.51 & 10.79 & 127 & 168 & 246.53 \\
& 2 & $-0.8456$ & 0.0495 & 65.8\% & 184.16 & 11.67 & 172 & 200 & 44.74 \\
& 3 & $-0.6554$ & 0.1116 & 98.6\% & $\bm{139.86}$ & 22.51 & 90 & 200 & 237.59 \\
& 4 & $-0.7017$ & 0.0378 & $\bm{100.0\%}$ & 145.02 & 10.57 & 127 & 170 & 272.35 \\     

        \midrule
        
% basepoint position -0.5
\multirow{4}{*}{${-0.5}$} 
& 1 & $-0.6905$ & 0.0232 & $\bm{100.0\%}$ & 147.36 & 9.65 & 130 & 171 & 307.97 \\
& 2 & $-0.6844$ & 0.0262 & $\bm{100.0\%}$ & 145.64 & 9.52 & 130 & 164 & 305.54 \\
& 3 & $-0.6873$ & 0.0239 & $\bm{100.0\%}$ & 146.80 & 9.12 & 129 & 163 & 296.40 \\
& 4 & $-0.6888$ & 0.0234 & $\bm{100.0\%}$ & 147.00 & 9.12 & 130 & 171 & 301.21 \\

        \midrule
        
% basepoint position -0.35
\multirow{4}{*}{${-0.35}$}
& 1 & $\bm{-0.6163}$ & 0.2203 & 98.0\% & $\bm{138.16}$ & 41.13 & $\bm{85}$ & 200 & 314.15 \\
& 2 & $\bm{-0.6343}$ & 0.2210 & 99.4\% & $\bm{142.62}$ & 40.47 & $\bm{86}$ & 200 & 310.75 \\
& 3 & $\bm{-0.6269}$ & 0.2220 & 98.4\% & 141.79 & 41.03 & $\bm{86}$ & 200 & 326.89 \\
& 4 & $\bm{-0.6327}$ & 0.2189 & 98.2\% & $\bm{141.46}$ & 41.47 & $\bm{84}$ & 200 & 330.34 \\

        \midrule
        
% basepoint position None
\multirow{4}{*}{{None}}
& 1 & $-0.7019$ & $\bm{0.0122}$ & $\bm{100.0\%}$ & 146.76 & $\bm{1.26}$ & 145 & $\bm{150}$ & 335.42 \\
& 2 & $-0.7039$ & $\bm{0.0144}$ & $\bm{100.0\%}$ & 147.50 & $\bm{1.23}$ & 145 & $\bm{151}$ & 305.56 \\
& 3 & $-0.7072$ & $\bm{0.0153}$ & $\bm{100.0\%}$ & 148.81 & $\bm{1.27}$ & 146 & $\bm{153}$ & 307.00 \\
& 4 & $-0.7079$ & $\bm{0.0154}$ & $\bm{100.0\%}$ & 149.03 & $\bm{1.26}$ & 147 & $\bm{152}$ & 310.69 \\

        \bottomrule   
    \end{tabular*}
\end{table}


\begin{figure}[htbp!]
    \centering
    \begin{subfigure}{0.49\textwidth}
        \includegraphics[width=\textwidth]{%
            graphics/mountaincar/20240903_A_start_vs_steps_bp_position_-0.65.png
        }
        \caption{$p_{-1} = -0.65$}
        \label{fig:MountainCarBasePointEpisodeLengths:sub:-0.65}
    \end{subfigure}
    \hfill
    \begin{subfigure}{0.49\textwidth}
        \includegraphics[width=\textwidth]{%
            graphics/mountaincar/20240903_A_start_vs_steps_bp_position_-0.5.png
        }
        \caption{$p_{-1} = -0.5$}        
        \label{fig:MountainCarBasePointEpisodeLengths:sub:-0.5}
    \end{subfigure}
    \hfill
    \begin{subfigure}{0.49\textwidth}
        \includegraphics[width=\textwidth]{%
            graphics/mountaincar/20240903_A_start_vs_steps_bp_position_-0.35.png
        }
        \caption{$p_{-1} = -0.35$}        
        \label{fig:MountainCarBasePointEpisodeLengths:sub:-0.35}
    \end{subfigure}
    \hfill
    \begin{subfigure}{0.49\textwidth}
        \includegraphics[width=\textwidth]{%
            graphics/mountaincar/20240903_A_start_vs_steps_bp_position_None.png
        }
        \caption{No basepoint}        
        \label{fig:MountainCarBasePointEpisodeLengths:sub:None}
    \end{subfigure}
    \caption[%
        Effect of different basepoint configurations in the Mountain Car
        experiment.
    ]{%
        The figure illustrates the effect of different basepoint
        configurations on the learned policy in the Mountain Car experiment.
        Each subplot shows the number of steps until termination as a function
        of the starting position for all episodes in test run 1, for the
        specific basepoint position $p_{-1}$ and for no basepoint.
        Plot (\subref{fig:MountainCarBasePointEpisodeLengths:sub:-0.65}) shows the
        results for the basepoint position $p_{-1}=-0.65$. Note, that test
        run 1 is one of two test runs that achieved a 100\% success rate for
        this configuration.
        Figure (\subref{fig:MountainCarBasePointEpisodeLengths:sub:-0.5}) shows the
        results for basepoint position $p_{-1}=-0.5$, with a 100\% success
        rate in test run 1.
        Figure (\subref{fig:MountainCarBasePointEpisodeLengths:sub:-0.35}) shows the
        results for basepoint position $p_{-1}=-0.35$, with a 98\% success
        rate in test run 1.
        Plot (\subref{fig:MountainCarBasePointEpisodeLengths:sub:None}) shows the
        results without a basepoint, with a 100\% success rate in test run 1. 
    }
    \label{fig:MountainCarBasePointEpisodeLengths}    
\end{figure}


%\cleardoublepage
%------------------------------------------------------------------------------
\section{Simple Execution Experiment}
\label{chapter:Experiments:section:Execution}
%------------------------------------------------------------------------------

In the second experiment we test the Signature-Q-Learning algorithm in an
environment, in which the agent has to solve a simple execution problem, by
trading in a single stock, simulated financial market. The environment 
provides an (approximately) optimal policy as benchmark, which albeit not 
analytically available, offer the possibility to assess the algorithm's 
learning performance. Before detailing the environment and its components, we 
introduce the agent-based market simulator, which is used to
generate the underlying stock market to trade in.

%------------------------------------------------------------------------------
\subsection{The Market Simulator ABIDES}
\label{chapter:Experiments:section:Execution:sub:ABIDES}
%------------------------------------------------------------------------------

The simple execution problem's environment relies on the discrete event
simulator ABIDES \cite{byrd2020_athmms}, and more specifically on it's
extension ABIDES-Gym \cite{amrouni2022_aemdesafm}, to simulate a stock market.
We describe the simulator and its extension, and provide details on the
generation of realistic financial market environments. 

ABIDES stands for \textit{Agent-Based Interactive Discrete Event Simulation}.
It is an open-source market simulator developed to support research in agent-%
based modeling and simulation within financial markets. At its core, ABIDES is
a discrete event simulation engine, which manages the timing of events within
the market. Instead of the market's state being updated at fixed steps, time
progresses by processing those events sequentially, and the simulation moves
forward by executing events in order of their specific timestamps. This
approach is particularly efficient for market simulations, where events like order
submissions, trades, and price changes happen at irregular intervals. The
simulated market consists of a set of trading agents and an exchange agent.%
\footnote{
    In general, those agents are not reinforcement learning agents in the 
    sense of \cref{chapter:RLAndMDP:section:RL}
}
Trading agents are autonomous entities that represent market participants, such
as traders or institutional investors. Each trading agent has its own strategy
and decision-making process, which can range from simple rule-based logic (e.g.
trend-following or market-making with a fixed spread) to more complex
algorithms. The exchange agent functions as the central market infrastructure,
handling the matching of orders, keeping an order book, and executing trades.
It mimics the behavior of real-world exchanges, following protocols like
NASDAQ's ITCH and OUCH, used for order book management and trade execution,
giving rise to a realistic order processing. In particular, no assumptions
about order queue positions for limit orders and order cancellations have to be
taken, as opposed to running simulations based on historical data, compare e.g.
\cite{spooner2018_mmrl}. Another of ABIDES’ key features is its ability to
simulate network latency between agents to reflect real-world conditions,
especially in high-frequency scenarios.

ABIDES-Gym is an extension of ABIDES, designed to integrate its simulation
capabilities with the modular OpenAI Gym framework. It provides a standardized
interface to apply reinforcement learning algorithms to different market
scenarios, enabling the training and evaluation of learning agents, in a
simulated environment that mimics the complexities of real-world financial
markets. Each simulation within ABIDES-Gym is represented as Gym environment,
where reinforcement learning agents interact with the market by sending actions
(e.g. placing or canceling orders) and receiving observations (e.g. market
prices, order book states) along with corresponding rewards in a continuous
feedback loop. The design separates the underlying market simulation, driven by
ABIDES, from the agent's learning algorithm, enabling an easy use of custom
methods or third-party reinforcement learning libraries, as demonstrated in
\cite{amrouni2022_aemdesafm}. This provides a controlled environment for
training RL agents and evaluating their performance in executing trades,
managing risk, and maximizing profits.

%------------------------------------------------------------------------------
\paragraph{Trading Agents}
\label{chapter:Experiments:section:Execution:sub:ABIDES:par:TradingAgents}

The ABIDES-Gym framework provides pre-specified market configurations that set
up and start the desired market environment with exchange agent settings, and a
population of trading agents. To any selected configuration a RL agent is added
through the Gym interface, employing for example Signature-Q-Learning or any
other method to learn a trading strategy. The configuration may be adapted by
changing the trading agent population or the exchange settings, to generate
different market conditions. At the same time, more than one learning agent
could be added through the Gym interface to create a multi-agent or adversarial
RL environment. The following types of trading agents may constitute a set of
market participants:
\begin{itemize}
    \item \textit{Noise agents}, also referred to as \textit{Zero Intelligence}
        agents \cite{byrd2020_athmms}, wake up once during the trading day and
        place one limit buy or sell order in the market. Wake up time, order
        type are determined randomly, as is the order size, drawn from some
        distribution.
        
    \item \textit{Value agents}, also referred to as \textit{Heuristic Belief
        Learning} agents \cite{byrd2020_athmms}, maintain a believe of the
        stock's fundamental value, and trade when they think the stock is
        mispriced, placing a limit buy order if the mid-price is below, and a
        limit sell order if it is above their believe, respectively. The order
        size is determined stochastically. An agent's belief is updated
        in a Bayesian manner, based on noisy observations of some extrinsic
        fundamental time series representing the stock's fundamental value
        (described in more detail below in this section).
    
    \item \textit{Momentum agents} wake up at fixed time intervals throughout
        the trading day and follow a simple momentum strategy. At each wake up
        they compare the average of the past 50 mid-price observations, 
        $\bar{p}_{50}$, with the average of the past 20 observations 
        $\bar{p}_{20}$. If $\bar{p}_{20} \geq \bar{p}_{50}$, an agent assumes
        the stock's price will increase, and it places a limit buy order,
        otherwise a limit sell order, assuming the price will fall.
    
    \item \textit{POV Market Maker agents} provide liquidity to the market,
        through placing limit order pairs at fixed distances to, and outside
        some window around the current mid-price. The window and distance
        between price levels is determined at each wake up and based on the
        observed spread in the market. POV stands for \textit{percentage of
        volume} and refers to the orders' size which is given as a percentage
        of the overall volume traded since the last wake up time.
\end{itemize}

%------------------------------------------------------------------------------
\paragraph{Fundamental Value Time Series}
\label{chapter:Experiments:section:Execution:sub:ABIDES:par:FundamentalValue}

An important feature of the simulation is the utilization of an external time 
series, that defines the fundamental value of the stock. It can be thought of
as a consensus valuation derived from all available information at a given time,
reflecting the ``true'' stock price. No market participant (with the exception
of the exchange agent) has access to this true fundamental value. Instead,
a trading agent is, at best, only able to receive noisy observations of it. An
example is the above introduced Value agent, whose fundamental value belief is
updated after it obtained such a noisy observation. Crucially, there is only
one fundamental value time series per stock in the simulation, and it is not
affected by actions taken by trading agents. In the ABIDES simulation, the
fundamental value time series is modeled as a mean-reverting stochastic 
process $X=(X_t)_{t \in [0,T]}$, with $T\in\SetR$ the simulation's terminal
time. Its dynamics are given by
\begin{equation}
\label{eq:ABIDESFundamentalValue}
    \dd{X}_t = -\gamma (X_t - \mu) \dd{t} + \sigma \dd{W_t} + \dd{Y}_t,
    \quad X_0 = x_0,
\end{equation}
where $W=(W_t)_{t\in[0,T]}$ is a standard Brownian motion, $Y = (Y_t)_{t \in
[0,T]}$ a compound Poisson process independent of $W$, $x_0\in\SetR$ the
initial value, and $\mu\in\SetR$, $\sigma>0$ and $\gamma\in (0,1]$ are
constants. Over time, the process tends to its long term mean $\mu$ at the
constant mean reversion rate $\gamma$. The volatility $\sigma$ determines how 
strong the process fluctuates around its tendency to the mean. Without the 
$\dd{Y_t}$ term, \cref{eq:ABIDESFundamentalValue} would be standard Ornstein-%
Uhlenbeck process. Its addition introduces jumps to the trajectory. The compound Poisson $Y$ with rate $\lambda>0$ 
given as
\begin{equation}
\label{eq:CompoundPoisson}
    Y_t = \sum_{i=1}^{N_t} J_n,
\end{equation}
where, $(N_t)_{t\in[T]}$ is a homogeneous Poisson counting process with rate 
$\lambda$, and the jumps $(J_n)_{n\in\SetN}$ are i.i.d. random variables. The
distribution of $J_1$ is the mixture of the two normal distributions with means
$\mu_J\in\SetR$ and $-\mu_J$, common variance $\sigma_J^2 >0$, and mixing
parameter $p=1/2$. Importantly, the parameters satisfy $\abs{\mu_J} \gg 0$ and 
$\sigma_J^2 \gg \sigma^2$, resulting in a bimodal distribution with zero mean,
which introduces large jumps relative to the continuous fluctuations from the
Ornstein-Uhlenbeck component. In a financial context, the Brownian motion term 
$\sigma\dd{W_t}$ models continuous, small movements around the fundamental time
series' mean due to high-frequency market interactions, while the compound
Poisson term $\dd{Y_t}$ resembles mega-shocks to the stock price, caused by
infrequent arrival of extrinsic information. With a variance $\sigma_J^2$
significantly larger than $\sigma^2$, these shocks are able to sustainably
shifts the consensus valuation of the stock. Note, that while the jump 
component $Y$ is part of the fundamental value time series as implemented in 
the ABIDES simulator, in the pre-specific market configuration used in the 
experiment and detailed in 
\cref{chapter:Experiments:section:Execution:sub:Environment}, the value for 
the jump intensity $\lambda$ is set so small such that practically there 
are no jumps the trading window.  

In general, the simulation advances at the nanoseconds time-scale, where
events can occur at any nanosecond time-step during the simulation time.
However, activity of single agents is set-up such that it rather occurs at the 
order of seconds, resembling a more retail-like behaviour and resulting in a 
sparser distribution of events during the simulation. Due to the continuous 
time process \eqref{eq:ABIDESFundamentalValue} used as fundamental value time 
series and the simulation progressing event-wise instead of in fixed 
time-intervals, arbitrary time intervals with no activity may be skipped 
without any computational cost. Further details on background agents to 
populate a simulation, and the construction and necessity of a exogenous 
fundamental value time series are provided in \cite{byrd2020_athmms, 
byrd2019_eafms}, and in the ABIDES github repository's 
\href{https://github.com/abides-sim/abides/wiki}{wiki}.\footnote{
    \url{https://github.com/abides-sim/abides/wiki}
}
Details on the specific market configuration used in the simple execution
environment are given in the following section.

%------------------------------------------------------------------------------
\subsection{The Environment}
\label{chapter:Experiments:section:Execution:sub:Environment}
%------------------------------------------------------------------------------

The problem set-up is as follows: The agent faces a simple execution problem,
wherein it starts a trading session with positive inventory of the stock 
traded in the market. Its objective is simply to sell all its holdings until 
the end of the session. The observations the agent receives only contain information 
about is current inventory and the remaining time. The transition dynamics 
however, depend on the mid-price and order flow, both not available to the 
agent, which makes the environment partially observable. The agent can trade 
by posting limit orders in both directions, and additionally it may choose not 
to trade at all. The reward it receives will be a quadratic function of its 
current inventory, with the magnitude determined by the change in mid-price. 
The rewards functions are designed in such a way that it penalizes
the agent for holding (large) absolute inventory values. To maximize 
cumulative expected rewards, the agent has to learn to reduce its inventory
as quickly as possible towards zero, and subsequently stop trading. This 
provides a clear optimal baseline policy for the agent, through which 
Signature-Q-Learning's performance may be assessed. Following the framework of 
ABIDES-Gym \cite{amrouni2022_aemdesafm}, the environment is implemented as a 
custom Gym environment, on top of the underlying market simulation. In the 
rest of this section details on the simulated market configuration and the 
environment's components are given.

%------------------------------------------------------------------------------
\paragraph{Market Configuration}
\label{chapter:Experiments:section:Execution:sub:Environment:par:MarketConfig}

The market configuration used in the simple execution environment is provided 
by the \texttt{rmsc04} configuration file from the
\href{https://github.com/jpmorganchase/abides-jpmc-public}
{ABIDES-Gym github repository}.%
\footnote{
    \url{https://github.com/jpmorganchase/abides-jpmc-public}. The
    configuration file \texttt{rmsc04} is located in the repository 
    sub-directory \texttt{/abides-markets/abides-market/configs}.
}
It consists of one Exchange Agent, two POV Market Maker agents, 102 Value
agents, 12 Momentum agents, and 1000 Noise agents. Together, this population of
trading agents ensures rich market dynamics, mimicking real world
environments. All agents trade one single stock, whose fundamental value time
series is modeled by the stochastic process \eqref{eq:ABIDESFundamentalValue},
with its parameter values given in \cref{tab:FundamentalValueParams}. Note, 
that while the fundamental value time series' jump component $Y$ is present in 
the implementation, due to the small value for the jump intensity $\lambda$ 
there are practically no jumps in the trading window\footnote{
    The probability of at least one jump within the 35 minutes simulated 
    trading interval is $1-\exp(-\lambda N)\approx \lambda N 
    \approx 9.72\times 10^{-8}$, with $N=3.5\times 10^{10}$ nanoseconds in the 
    window. Even if a full trading day of 8 hours is considered, this 
    probability only increases to about $2.34\times10^{-6}$.
}, and the process \eqref{eq:ABIDESFundamentalValue} is effectively a standard 
Ornstein-Uhlenbeck process.
Note,
that the fundamental value time series is given in USD cents, such that 
$\mu=10^5$ corresponds to a mean value of USD 1000. An episode consists of
the simulation of a 35 minutes of trading. To sufficiently populate the order 
book, the reinforcement learning agent's first wake up is five minutes into 
the session, which is defined as step $t=0$ at which it receives its first 
observation $o_0$ and the Signature-Q-Learning algorithm is invoked.\footnote{
    The environment is implemented such that one can define an additional 
    interval at the beginning of an episode, in which the agent merely 
    receives observations, but is not allowed to trade. This may be useful if 
    for example price or order-book features are included in the observations, 
    such that the first trade decision is informed by their trajectories in 
    this observation interval.
}
During the remaining 30 minutes of an episode, the agent wakes up every ten 
seconds of the simulation to receive a new observation and select an action, 
resulting in $T=180$ steps until termination. No transaction costs are 
incurred in trading. Order priority and queue position as well as market 
impact and time delays are handled automatically through the ABIDES simulation.

\begin{table}
    \caption[%
        Parameter values for ABIDES' fundamental value time series
    ]{%
        The parameter values for the ABIDES fundamental value time series
        \eqref{eq:ABIDESFundamentalValue} in Simple Execution experiment. The
        values are taken from the configuration file \texttt{rmsc04} used in
        the experiment to set-up the market, which can be found in the 
        ABIDES-Gym github repository. Note, that the simulation advances on 
        the nanosecond time-scale, and that due to the small jump intensity 
        $\lambda$ effectively there are no jumps in the trading window.
    }
    \label{tab:FundamentalValueParams}
    \centering
    \begin{tabular*}{\textwidth}{ c @{\extracolsep{\fill}} *{6}{c} }
        \toprule
        $\bm\gamma$ & $\bm\mu$ & $\bm{\sigma^2}$ & $\bm{x_0}$ & 
            $\bm\lambda$ & $\bm{\mu_J}$ & $\bm{\sigma_J^2}$ \\
        \midrule
        $1.67\cdot10^{-16}$ & $10^{5}$ & $2.5\cdot10^{-10}$ & $10^{5}$ &
            $2.77778\cdot 10^{-18}$ & $10^{3}$ & $5\cdot10^{4}$ \\
        \bottomrule        
    \end{tabular*}
\end{table}

%------------------------------------------------------------------------------
\paragraph{Observation Space}
\label{chapter:Experiments:section:Execution:sub:Environment:%
    par:ObservationSpace}

For $t\in[T]$, let $i_t \in \SetZ$ denote the amount of stock the
agent currently holds. An inventory $i_t<0$ means that the agent is short 
$\abs{i_t}$ units of the stock. Furthermore, we restrict the maximal inventory
the agent is allowed to hold in either direction to $\abs{i_t} \leq i_{\max}$
for any $t \in [T]$, where $i_{\max}\in\SetN$ will be a multiple of the fixed
order size the agent trades, see below in the description of the action space. 
Observations are defined as the current inventory as percentage of the maximal 
inventory $o_t := i_t/i_{\max}$ taking values in $\SpaceO := [-1,1]$. At each 
$t\in[T]$, the agent receives the time-augmented observation
\[
    \hat{o}_t := \biggl(1- \frac{t}{T}, o_t\biggr) 
    = \biggl(1-\frac{t}{T}, \frac{i_t}{i_{\max}}\biggr) 
    \in [0,1] \times \SpaceO,
\]
where the time-component is the normalized, remaining time in the episode. 
Each episode starts at some initial observation $\hat{o}_0 = (1, o_0) = 
(1, i_0/i_{\max})$ with the initial inventory $i_0 \in [-i_{\max}, i_{\max}] 
\cap\SetZ$. In the specific environment configuration given in 
\cref{tab:ExecutionEnvironmentParameters} we set $i_0$ at a fixed positive 
value, such that the agent has to learn to liquidate its position over the 
trading window. Note, however, that this choice is arbitrary and that $i_0$ 
could also be sampled according to some initial distribution $\mu$ on 
$[-i_{\max}, i_{\max}]\cap\SetZ$. Note further, that although the agent does 
not observe any other market variable, such as the mid-price or order book 
statistics, the observation transitions still depend on them, making the 
environment partially observable. More specifically, assume the agent holds 
$i_t$ units of stock and wants to reduce its holdings to $i < i_t$ by posing a 
limit order of size $s=i_t - i$. The probability that at $t+1$ the agent 
indeed observes $i_{t+1} = i$ depends on two things: First, the best bid price 
in the order book has to match or exceed the agent's limit sell price during 
the time interval $[t,t+1]$. And second, the order has be entirely executed, 
since a partial fill would yield a new inventory $i_{t+1} > i_t - i$.

%------------------------------------------------------------------------------
\paragraph{Action Space}
\label{chapter:Experiments:section:Execution:sub:Environment:par:ActionSpace}

The action space $\SpaceA$ is provided in \cref{tab:ActionsExecution}, and it
consists of posting a limit sell order, posting a limit buy order, and not
trading. The limit prices $p_t^{bid}$ and $p_t^{ask}$ are the best bid and 
best ask price available at time $t\in[T]$, which corresponds to the first 
price level in the order book. This guarantees a high probability that the 
order will at least partially be filled. If one side of the order book is 
empty, the respective limit price is set to the price of the last transaction 
in the market. The agent is restricted to one single order of fixed size $\rho 
\in \SetN$ at each step, which satisfies $\rho\leq i_{\max}$ and 
$(i_{\max}\mod\rho)=0$. Any existing orders of the agent in the order book are 
cancelled before a new order is submitted to the market. To satisfy the 
constraint $\abs{i_t}\leq i_{\max}$ for all $t\in[T]$, whenever $\abs{i_t} +
\rho \geq i_{\max}$, and the agent selects the action that would further 
increase $\abs{i_t}$, Action 2 (Do nothing) is selected instead.

\begin{table}[htbp!]
    \caption[Action space in the Simple Execution experiment]{%
        The action space in the simple execution the execution environment.
        The prices $p_t^{bid}$ and $p_t^{ask}$ correspond to the best bid and 
        best ask price available in the order book at time $t$, respectively. 
    }
    \label{tab:ActionsExecution}    
    \centering
    \begin{tabular*}{\textwidth}
        {>{\centering}m{2cm} >{\centering}m{3cm} l @{\extracolsep{\fill}} l}
    \toprule
        & {Action 0} 
            & Post limit buy order at price $p_t^{bid}$ of size $\rho$. &\\
        & {Action 1} 
            & Post limit sell order at price $p_t^{ask}$ of size $\rho$. &\\
        & {Action 2} & Do nothing & \\         
    \bottomrule
    \end{tabular*}
\end{table}


% alternative description of action space
\begin{comment}
Let $b_t^{bid,1}$ and $b_t^{ask,1}$ be the best bid and best ask price
available at time step $t\in[T]$, which corresponds to the first price level in
the order book. The mid-price $m_t$ and spread $s_t$ are then given as
\[
    m_t = \frac{b_t^{bid,1} + b_t^{ask,1}}{2}, 
    \quad 
    s_t = b_t^{ask,1} - b_t^{bid,1},
\]
and the limit prices at which the agent posts its orders are calculated as 
\begin{equation}
\label{eq:SimpleExecutionLimitPrices}
    p_t^{bid} = m_t - \frac{s_t}{2} 
    \quad \text{and} \quad
    p_t^{ask} = m_t + \frac{s_t}{2}.
\end{equation}
with $p_t^{bid}$ the limit buy, and $p_t^{ask}$ the limit sell price,
respectively. Defined like this, the prices $p_t^{bid/ask}$ are simply the
prices at the best bid and ask levels in the order book at time $t$ with $m_t$
denoting the mid price and $s_t$ the spread at time $t$, respectively.
\end{comment}

%------------------------------------------------------------------------------
\paragraph{Rewards}
\label{chapter:Experiments:section:Execution:sub:Environment:par:Rewards}

The policy we want the agent to learn is to reduce its absolute inventory
towards zero as quickly as possible, then refrain from trading until the end 
of the episode. Thus, the reward functions should discourage the agent from 
holding large absolute inventory values and provide an incentive for keeping 
its holdings close to zero, or more specifically, at a value that satisfies 
$\abs{i_t}\leq \rho/2$ with $\rho\in\SetN$ the agent's fixed order size. A 
detailed reasoning for this is provided in the description of the baseline 
policy in \cref{chapter:Experiments:section:Execution:sub:BaselinePolicy}.

For any $t\in[T]$, let $m_t:=(p_t^{bid}+p_t^{ask})/2$ be the mid-price at step 
$t$, with $p_t^{bid}$ and $p_t^{ask}$ the best bid and ask price in the order 
book at step $t$, respectively, and define the absolute mid-price change at $t\in[T]\setminus\{0\}$ as $\Delta m_t := \abs{m_{t} - m_{t-1}}$.\footnote{
    All prices here are assumed to be in USD terms. Recall, that prices in the 
    simulation are given as integers representing USD cents, thus a simple 
    division by 100 does the trick, compare 
    \cref{chapter:Experiments:section:Execution:sub:ABIDES:%
    par:FundamentalValue}
}
Then, for all $t\in[T-1]$ and triplets $(h_t,a_t,o_{t+1}) \in \SpaceH_t \times 
\SpaceA \times \SpaceO$ the reward function is defined as 
\begin{equation}
    \label{eq:RewardExecution}
    R_t\bigl(h_t, a_t, o_{t+1}\bigr)
    := \Delta m_{t+1} \biggl[
            \Bigl(\frac{\rho/2}{i_{\max}}\Bigr)^2
            %- \Bigl(\frac{i_{t+1}}{i_{\max}}\Bigr)^2
            - o_{t+1}^2
        \biggr],
    %\quad t\in[T-1].
\end{equation}
with $o_{t+1} = i_{t+1}/i_{\max}$.\footnote{
    Note, that $R_t$ is defined as a function of $o_{t+1} = i_{t+1}/i_{\max}$ 
    only (compare also \cref{rem:HDPRewardsRemark}). While technically it is 
    also a function of $\Delta m_{t+1}$, we omit this dependence since the 
    mid-price information is hidden from the agent, highlighting again the partial observability it faces.
}
The reward signal's qualitative behaviour is controlled by the inventory the 
agent holds at $t+1$, which resulted from its trading decision at $t$ and the 
market events in between, and its magnitude determined by the absolute 
mid-price change $\Delta m_{t+1}$. Whenever $\abs{i_{t+1}}>\rho$, the signal 
is negative, and it decays quadratically to zero for $\abs{i_{t+1}} \to 
\rho/2$. For values $\abs{i_{t+1}} > \rho/2$, the reward signal is positive 
with its maximum at for $i_{t+1}=0$. \Cref{fig:RewardFunctionExecution} 
illustrates the reward functions behaviour for inventory values around the 
interval $[-\rho/2, \rho/2]$, and this specific design will become more 
apparent with the optimal baseline policy provided in 
\cref{chapter:Experiments:section:Execution:sub:BaselinePolicy}

To provide an additional inventive to end an episode with $i_T\in [-\rho, 
\rho]$, the terminal reward function is defined similarly as
\begin{equation}
\label{eq:TerminalRewardExecution}
    R_T\bigl(h_T\bigr) 
    := C \, \biggl[
            \Bigl(\frac{\rho/2}{i_{\max}}\Bigr)^2
            - o_T^2,
            %- \Bigl(\frac{i_{t+1}}{i_{\max}}\Bigr)^2 
        \biggr],
\end{equation}
for some fixed scaling factor $C\geq 0$ and with terminal observation $o_T = 
i_T/i_{\max}$ in the  history $h_T$.

% pgfplot of reward functions
\begin{figure}
    \centering
    \begin{tikzpicture}
            \begin{axis}[
                axis lines = left,
                axis line style={|-stealth},
                axis line shift = 5pt,
                width = 11cm,
                height = 7cm,                
                % x axis
                xmin = -0.12, xmax = 0.12,                
                xtick = {-0.12, -0.05, 0.0, 0.05},
                xticklabels={$-i_{\max}$, $-\rho/2$, $0$, $\rho/2$},
                axis x discontinuity=crunch,                
                xlabel = $i_{t+1}$,  
                % y axis
                ymin = -0.01, ymax = 0.0035,    
                restrict y to domain = -0.01:0.0035,
                ytick = {-0.01, 0, 0.0025},
                yticklabels = {
                    $\left(\frac{\rho/2}{i_{\max}}\right)^2-1$, $0$, 
                    $\left(\frac{\rho/2}{i_{\max}}\right)^2$
                },
                scaled y ticks={real:1.0},
                ytick scale label code/.code={$\times \Delta m_{t+1}$}, 
                axis y discontinuity=crunch,                                
                ylabel = $R_t(i_{t+1}/i_{\max})$,
            ]
                \addplot [
                    domain=-0.105:0.105, 
                    samples=101, 
                    color=blue,
                    line width=2pt,
                ]
                {
                    -(x^2 - 0.0025)
                };
                % horizontal lines
                \addplot[dashed]coordinates{(-0.12, 0.0) (0.12,0.0)};
                \addplot[dashed]coordinates{(-0.12, 0.0025) (0.12,0.0025)};
                %vertical lines
                \addplot[dashed]coordinates{(-0.05,-0.01) (-0.05,0.0025)};
                \addplot[dashed]coordinates{(0.05,-0.01) (0.05,0.0025)};       
            \end{axis}
    \end{tikzpicture}
    \caption[Reward functions in the Simple Execution experiment]{
        The figure shows the behaviour of the reward functions 
        $R_t,t\in[T-1]$ from \cref{eq:RewardExecution}, for inventory values 
        symmetrically around zero. Specifically, the agent receives a positive 
        reward for $i_{t+1}\in[-\rho/2, \rho/2]$, with $\rho$ the fixed order 
        size it can trade, and a negative reward otherwise. The maximum value 
        is always attained at $i_{t+1}=0$, and the functions decay 
        quadratically in both directions to its minima at $\abs{i_{+1}}
        =i_{\max}$. Note, that the functions' magnitude is determined by the 
        absolute absolute mid-price change $\Delta m_{t+1}$, as indicated by 
        the $y$-axis scaling.
    }
    \label{fig:RewardFunctionExecution}
\end{figure}        

%------------------------------------------------------------------------------
\paragraph{Environment Parameters}
\label{chapter:Experiments:section:Execution:sub:Environment:par:EnvSetUp}

The parameter values provided in \cref{tab:ExecutionEnvironmentParameters}
define the specific instance of the environment used in the simple execution
experiment. The agent starts an episode with $i_0=700$, which 
corresponds to an inventory of 70\% of the maximal allowed value $i_{\max}
=1000$, thus $\hat{o}_0 = (1,0.7)$. With the fixed order size $\rho=50$, the 
agent would need 14 orders in total to reduce its holdings to zero if each 
order was fully filled. Note, that this is seldom the case in the market 
simulation. The scaling factor for the terminal reward function is set to 
$C=0.7$, which was found to yield balanced training.

\begin{table}[htpb!]
    \caption[Environment parameter values in the Simple Execution experiment]
    {The environment parameter values that define the specific instances of
    the observation space, the action space, and the reward functions used in
    the Simple Execution experiment.}
    \label{tab:ExecutionEnvironmentParameters}
    \centering
    \begin{tabularx}{\textwidth}{CCCC}
        \toprule
         ${i_0}$ & ${i_{\max}}$ & $\rho$ & ${C}$ \\
         \midrule
         700 & 1000 & 50 & 0.7 \\
         \bottomrule
    \end{tabularx}    
\end{table}

%------------------------------------------------------------------------------
\subsection{Baseline Policy}
\label{chapter:Experiments:section:Execution:sub:BaselinePolicy}
%------------------------------------------------------------------------------

Albeit an optimal policy is not analytically available, the Simple Execution 
environment offers a deterministic, (possibly approximately) optimal policy, 
which serves as baseline. At any $t\in[T-1]$, this \textit{optimal 
baseline policy} $\pi^b$ is defined as
\begin{equation}
\label{eq:ExecutionBaselinePolicy}
    \pi_t^b(h_t) = \pi^b_t(o_t)
    := \begin{cases}
        \text{Action 0 – limit buy},  & \text{if } i_t < \rho/2, \\
        \text{Action 1 – limit sell}, & \text{if } i_t > \rho/2, \\
        \text{Action 2 – do nothing}, & \text{else,}
    \end{cases}
\end{equation}
for any history $\hat{h}_t \in \hat{\SpaceH}_t$ with observation $\hat o_t = 
(1-t/T, i_t/i_{\max})$ at step $t$ and arbitrary starting inventory $i_0 \in
[-i_{\max}, i_{\max}]\cap\SetZ$. When following $\pi^b$, the agent reduces its 
absolute inventory value $\abs{i_t}$ as quickly as possible with a series of 
limit buy orders (action 0), respectively, limit sell orders (action 1), 
depending on the sign of the starting inventory. As soon as its absolute 
inventory is at or below 50\% of the fixed order size $\rho$, the agent stops 
trading until the end of the episode (action 2). Note, that due to the fixed 
order size $\rho$, when following $\pi^b$ with arbitrary starting inventory 
$i_0\in[-i_{\max}, i_{\max}]\cap\SetZ$, there almost\footnote{
    Here, \textit{almost} means, save the extreme cases where there is less 
    than $i_0 - \rho/2$ units of the agent's total posted volume executed 
    during the whole trading window. A highly unrealistic scenario where $i_0$ 
    must amount to a large fraction of the total traded volume in the market 
    over the whole trading window, and negligible for the market configuration 
    detailed in 
    \cref{chapter:Experiments:section:Execution:sub:Environment}.
} always exists a $t\in[T]$, such that 
\[
    i_t \in I_{\rho/2} := \Bigl[-\frac{\rho}{2}, \frac{\rho}{2}\Bigr],
\]
which is considered an optimal inventory value, in the sense that an 
additional limit 
order of size $\rho$ in either direction would result in an inventory 
$i_{t+1}$ with $\abs{i_{t+1}}\geq\rho/2 \geq \abs{i_t}$, if the order is fully 
filled. As such, the agent has no further inventive to continue trading as 
soon as $i_t \in I_{\rho/2}$. Note, that we do not claim $\pi^b$ to 
be strictly optimal in the sense of \cref{def:HDPOptimalPolicy}, since in 
absence of observation transition probabilities and specifically, order 
execution dynamics, we do not know whether a different behaviour around 
absolute inventory values close to $\rho/2$ is better.\footnote{
    Consider the example, of the probability of the event, that 
    less than half the volume $\rho$ posted in an order is filled between 
    steps $t$ and $t+1$, being consistently higher than for its complement at 
    any time step $t$. In this case, a policy that continues to trade when 
    $\abs{i_t}$ is slightly below $\rho/2$, would on average yield a inventory 
    values closer to zero than $\pi^b$, and thus, overall higher rewards.
} 
Nevertheless, since $\pi^b$ is very close to optimal, we use it as a proxy for 
a truly optimal policy, and as benchmark to asses the performance of any 
policy learned by Signature-Q-Learning.

The specific environment configuration used in the experiment, given in 
\cref{tab:ExecutionEnvironmentParameters}, starts with a positive initial 
inventory $i_0=700$. Thus, the optimal baseline policy amounts to posting 
limit sell orders (action 1) until $i_{t} \leq \rho/2$ for some $t\in[T]$, and 
subsequently doing nothing (action 2) until the episode terminates. 
\Cref{fig:BaselineTrajectories} shows two typical action sequences and the 
resulting inventory trajectories when following $\pi^b$ in the environment 
specified in \cref{tab:ExecutionEnvironmentParameters}. It is 
clearly visible that the policy only selects action 1, until the inventory 
falls within the interval $I_{\rho/2}$, indicated by the dotted lines, after 
which it switches to action 2 and stops trading. The specific time step 
$t\in[T]$ at which $i_t \in I_{\rho/2}$ depends on how quickly the limit 
sell orders are filled in the respective episode, which is apparent by the two 
different inventory trajectories.

The policy $\pi^b$ will be run for $M_b=2000$ baseline episodes. Denote the 
average, cumulative reward obtained per episode as $\overline r_b$. Since 
$\pi^b$ is (approximately) optimal, the value $\overline r_b$ constitutes an 
empirical estimate of the true optimal value $V_0^*(o_0)$ for the initial 
observation $o_0 = i_0/i_{\max}$. To assess the convergence of 
Signature-Q-Learning, for any learned approximate Q-function $Q(\cdot|
\ell_K)$, we can compare the estimate $\overline r_b$ to: (a) the average, 
cumulative reward per episode when following the greedy-policy with respect to 
$Q(\cdot|\ell_K)$ for a number of test episodes, and (b) the estimate for 
$V_0^*(o_0)$ provided by $\max_{a\in\SpaceA}Q(\hat o_0, a|\ell_K)$.

\begin{figure}[htbp!]
    \centering
    \begin{subfigure}{0.46\textwidth}
        \includegraphics[width=\textwidth]{%
            graphics/execution/baseline_actions.png
        }
        \caption{Baseline policy action sequences.}
        \label{fig:BaselineTrajectories:sub:Actions}
    \end{subfigure}
    \hfill
    \begin{subfigure}{0.46\textwidth}
        \includegraphics[width=\textwidth]{%
            graphics/execution/baseline_inventories.png
        }
        \caption{Baseline policy inventory trajectories.}
        \label{fig:BaselineTrajectories:sub:Inventories}
    \end{subfigure}            
    \caption[Action and inventory trajectories from baseline policy]{
        Display of typical (\subref{fig:BaselineTrajectories:sub:Actions})
        action sequences, and (\subref{fig:BaselineTrajectories:sub:Inventories})
        inventory trajectories for two episodes when following the optimal
        baseline policy $\pi^b$ from \cref{eq:ExecutionBaselinePolicy} in the
        Simple Execution experiment with configuration given in
        \cref{tab:ExecutionEnvironmentParameters}. The dotted lines in 
        (\subref{fig:BaselineTrajectories:sub:Inventories}) correspond to the 
        boundaries of interval $ I_{\rho/2} = [-\rho/2, \rho/2]$ for fixed order size $\rho=50$.
    }
    \label{fig:BaselineTrajectories}
\end{figure}

%------------------------------------------------------------------------------
\subsection{Training Details}
\label{chapter:Experiments:section:Execution:sub:TrainingDetails}
%------------------------------------------------------------------------------

\Cref{tab:ExecutionTrainingParameters} provides the Signature-Q-Learning 
hyperparameter configuration used for training in the simple execution 
environment, with the values found through an informal search. In total, 10 
training runs are performed. As in the Mountain Car experiment, we use only 
the time-augmented observation history as input to the approximate Q-function, 
which at time $t$ is the stream $(\hat{o}_i)_{i=0}^t = (1 - i/T, i_t/
i_{\max})_{i=0}^t$. The approximate Q-function is implemented as single linear 
neural network layer, with the identity map as activation function, which 
outputs the Q-value for each of the three actions simultaneously. This amounts 
to three linear functionals of the signature, one for each of the actions to 
estimate their action-values given any history. For truncation
order $N=7$, we obtain $d_N=255$ trainable parameters for each action,
with $d_N$ the dimension of $T^{(N)}(\SetR^2)$, which are collected in a weight
matrix $W \in \SetR^{3\times (d_N-1)}$ and a bias vector $b\in\SetR^3$, compare
\cref{eq:QImplementationApproach2} and 
\cref{chapter:SignatureQLearning:section:QImplementation} for details on this 
implementation. 

The entries in $W$ are initialized
randomly, where each entry is sampled uniformly from the interval $(-\sqrt{k}, 
\sqrt{k})$ with $k=6/(d_N+3)$, which corresponds to \texttt{PyTorch}'s uniform 
Xavier initialization \cite{glorot2010_udtdfnn}. The bias vector's entries are 
initialized as $b_i=0.01$ for $i=1,2,3$. A basepoint of $o_{-1} = (0,0)$ was 
used for computing the truncated signature at $t=0$, compare 
\cref{chapter:SignatureQLearning:section:QImplementation}. Note, that since  
the initial observation is deterministic, the truncated signature's value at 
$t=0$ is the same in each episode. Nevertheless, since we include a basepoint, 
all of the truncated signature's components contribute to the selection of the 
first action $a_0$, as opposed to using no basepoint, in which case only the 
bias vector would contain non-zero entries that determine the choice of $a_0$.
%Using histories containing only observations proved to be more stable
%in training, than the other two implementation approaches described in
%\cref{chapter:SignatureQLearning:section:QImplementation}.

\begin{table}[htpb!]
    \centering
    \caption[Training parameter values in the Simple Execution experiment]
    {%
        The training hyperparameter values in the Simple Execution experiment. 
        In column order the parameters correspond to number of episodes, 
        episode length, discount factor, loss function, gradient descent 
        optimizer, signature truncation order, and basepoint.
        The smooth-$L_1$ loss function is provided by \eqref{eq:SmoothL1Loss},
        the gradient descent optimizer Adam is introduced in
        \cite{kingma2015_amso}. 
    }
    \label{tab:ExecutionTrainingParameters}        
    \begin{tabular*}{\textwidth}{>{\centering}m{2.5cm} @{\extracolsep{\fill}} *{7}{c} @{}}
        \toprule
        $M_\text{train}$ & $T$ & $\gamma$ & $\mathcal{L}$ 
            & optimizer & $N$ & $\hat{o}_{-1}$ \\
        \midrule 
        3000 & 180 & 1 & smooth-$L_1$ & Adam & 7 & $(0, 0)$ \\
        \bottomrule
    \end{tabular*}
\end{table}

The exploration policy used is $\varepsilon$-greedy. Exploration rates and 
learning rates are decayed during training, with the decay schemes
detailed in \cref{tab:MountainCarDecaySchemes}. Contrary to the Mountain Car
experiment, decay happens after each single episode. During an episode, the 
values of both parameter values are held fixed, that is $\varepsilon_{n} :=
\varepsilon_{n,0} = \varepsilon_{n,t}$ and $\alpha_{n} := \alpha_{n,0} =
\alpha_{n,t}$ for all $t\in[T-1]$, and episode $n\in\{1,\dots,M_{\text{train}}\}$.
We do not discount rewards in the simple execution environment and set 
$\gamma=1$. This results in the Q-functions horizon being the full episode 
length and sampled reward values are propagated throughout all steps. 

\begin{table}[htpb!]
    \centering
    \caption[Parameter decay schemes in the Simple Execution experiment]
    {Details on the decay schemes of the learning rates and $\varepsilon$-greedy exploration rate in the Simple Execution experiment. End values are rounded and correspond to start value multiplied by the decay factor $M_\text{train}=3000$ times.}
    \label{tab:ExecutionDecaySchemes}    
    \begin{tabular*}{\textwidth}{l @{\extracolsep{\fill}} *{4}{c}}
        \toprule
        \textbf{Parameter} & \textbf{Start value} & \textbf{Decay factor} &
            \textbf{End value} & \textbf{Decay mode} \\
        \midrule
        Learning rates ${\alpha_n}$ & $5\times10^{-5}$ & $0.999$ &
            $2.4856\times10^{-6}$ & Exponential \\          
        Exploration rates ${\varepsilon_n}$ & $1$ & $0.997$ & 
            $1.2175\times 10^{-4}$ & Exponential \\          
        \bottomrule
    \end{tabular*}
\end{table}

%------------------------------------------------------------------------------
\subsection{Testing Details} 
\label{chapter:Experiments:section:Execution:sub:TestingDetails}
%------------------------------------------------------------------------------

As a benchmark, we run the optimal baseline policy $\pi^b$ defined in 
\cref{eq:ExecutionBaselinePolicy} for $M_b=2000$ episodes. Then, for each 
training run, we test the learned approximate Q-function by performing a test 
run of $M_\text{test} = 2000$ episodes. Let $Q(\cdot \,|\,\ell_K)$ be the 
approximate Q-function obtained from a training run, with $\ell_K$ the final 
linear functional after a total of $K=T\times M_\text{train}$ update steps 
during training. In the following test run, actions are selected greedily with 
respect to $Q(\cdot\,|\,\ell_K)$, such that for any $t\in[T-1$], the policy 
followed in testing is given as
\begin{equation}
\label{eq:ExecutionTestPolicy}
    \pi^{\text{test}}_t(h_t) 
    := \argmax_{a\in\SpaceA} Q(\hat{h}_t^a\,|\,\ell_K),
\end{equation}
for any $h_t\in\SpaceH_t$, where we assume there are no ties between
actions, such that the $\argmax$ is well-defined.\footnote{
    If there happens to be a tie between actions in the experiment, we select 
    by sampling uniformly from the set of actions involved in the tie.
} 
For each test run, we report the test statistics in
\cref{tab:ExecutionTestStatistics}. Note, that there is no discounting of
rewards since $\gamma=1$. For test episode $n \in \{1,\dots, M_\text{test}\}$, let 
$r^{(n)}$ be the sum of (undiscounted) sampled rewards obtained over the 
episode. The mean cumulative episode reward $\overline{r}$ and the episode reward 
standard deviation $\sigma_r$ are calculated as sample mean and standard 
deviation over $\{r^{(n)}\}_{n \in \{1,\dots,M_\text{test}\}}$. Analogously, let 
$i^{(n)}_T\in[-i_{\max}, i_{\max}]\cap\SetZ$ be the inventory the agent 
holds at terminal time $T$ in test episode $n$. The mean terminal inventory 
$\overline{i_T}$ and the terminal inventory standard deviation $\sigma_{i_T}$ are 
calculated as sample mean and sample standard deviation over 
$\{i^{(n)}_T\}_{n \in \{1,\dots,M_\text{test}\}}$, with $\min_{i_T}$ and $\max_{i_T}$ 
the minimal and maximal value in the sample, respectively. The first 
observation value $v^*_0$ is the value of the initial observation $\hat{o}_0 = 
(1,o_0) = (1, 0.7)$ as provided by the final estimate $Q(\cdot\,|\, \ell_K)$, 
calculated as
\begin{equation}
\label{eq:ExecutionFirstObservationValue}
    v^*_0
    %:= \frac{1}{5}\sum_{i=1}^5 V^{*,i}(\hat{o}_0) 
    = \max_{a\in\SpaceA} Q(\hat{o}_0, a\,|\,\ell_K).
\end{equation}
Since in the simple execution environment the initial observation is 
deterministic, $v_0^*$ takes the same value in each test episode, as opposed 
to the initial values $v_0^{*,n}$ in 
\eqref{eq:MountainCarFirstObservationValue} from the Mountain Car experiment. 
Statistics for the optimal baseline policy are calculated analogously, with 
$\overline r_b$ and $\sigma_{r_b}$ referring to the episode rewards' sample mean 
and standard deviation in the baseline run, respectively.

For each run, the value $v_0^*$ offers an indication about whether the 
algorithm has converged, by comparing $v_0^*$ to the average episode reward 
$\overline r$ from following the test policy $\pi^\text{test}$, and to the average 
episode reward $\overline r_b$ obtained from the optimal baseline policy 
\eqref{eq:ExecutionBaselinePolicy}. Recall
\cref{def:HDPValueFunctions,def:HDPOptimalValueFunctions} for 
the value functions $(V^\pi_t)_{t\in[T]}$ corresponding to a given policy 
$\pi\in\Pi$, and the optimal value functions $(V_t^*)_{t\in[T]}$.
%
Since $\pi^b$ is (approximately) optimal, $\overline{r}_b$ is a sample-based
estimate of $V_0^*(o_0)$ with $o_0 = 0.7$, that is the optimal value over an 
entire episode for initial observation $o_0$. In a test run actions are 
selected according to $\pi^\text{test}$, thus, $\overline{r}$ is an estimate of 
$V^{\pi^\text{test}}_0(o_0)$. 
%
Assume that the algorithm has perfectly converged to the true optimal 
Q-functions. Then $\pi^\text{test}$ is optimal and we have
\[
    \overline r_b \approx V_0^*(o_0) = V_0^{\pi^\text{test}}(o_0) \approx \overline r.
\]
In general we expect $\overline{r}<\overline{r}_b$ in any test run, with $\overline{r}_b - 
\overline{r} \approx 0$.
%
Furthermore, for any $(o_0, a)\in\SpaceO\times\SpaceA$, we have the equality 
$Q(\hat{o}_0, a\,|\,\ell_K) = Q_0^*(o_0, a)$, and by 
\cref{eq:HDPValueQConnection,eq:ExecutionFirstObservationValue} we obtain
\[
    \overline r_b \approx V_0^*(o_0)
    = \max_{a\in\SpaceA} Q_0^*(o_0, a)
    = \max_{a\in\SpaceA} Q(\hat{o}_0, a\,|\, \ell_K) 
    = v_0^*
\]
This implies, that for any test run the fact that $\overline r_b \approx 
v_0^*\approx \overline r$ is an indication that the algorithm has converged to the 
optimal Q-functions.

In the case that $Q(\cdot|\ell_K)$ is not optimal, then we distinguish two 
cases. In the first case we have we get $v_0^*\approx \overline r_b$ but $\overline r <
\overline r_b$ and $\overline r < \bar v_0^*$. This implies that $Q(\cdot|\ell_K)$ 
correctly represents the true Q-function $Q_0^*$ at $t=0$, and from this also 
the true value $V_0^*$, however there is some deviation at later steps $t>0$, 
resulting in a lower average reward $\overline r$ from following $\pi^\text{test}$ 
than from $\pi^b$.
%
In the second case, we get $v_0^*\approx \overline r$ but $\overline r < \overline r_b$ and 
$v_0^* < \overline r_b$. This implies, that while the algorithm was not able to 
learn the true optimal Q-functions, it still converged to some sub-optimal 
Q-functions, and $Q(\cdot\,|\,\ell_K)$ correctly representing the action
-values of the policy $\pi^\text{test}$. That $v_0^*\approx\overline{r}$ need not 
be the case in general is clear from the results of the experiment in 
\cref{chapter:Experiments:section:MountainCar}, where the approximate 
Q-function vastly overestimates Q-values, however it produces action sequences 
that solve the problem.

\begin{table}[htpb!]
    \centering
    \caption[%
        Description of reported test statistics in the Simple Execution
        Experiment
    ]{
        Description of the reported statistics in the Simple Execution
        experiment for the test and baseline runs. The statistics are 
        calculated for each performed test or baseline run with the values 
        are reported in 
        \cref{chapter:Experiments:section:Execution:sub:Results}. The 
        subscript $b$ indicates that the respective statistic corresponds 
        to the baseline runs. Recall that $\rho=50$ corresponds to the minimal
        order size.
    }
    \label{tab:ExecutionTestStatistics}        
    \begin{tabularx}{\textwidth}{>{\centering}m{3.5cm}X}
        \toprule
        ${\overline r}$ / ${\overline r_b}$ & Mean reward obtained per episode
            for test / baseline runs \\
        ${\sigma_r}$ / ${\sigma_{r_b}}$ & Standard deviation of reward per
            episode for test / baseline runs \\
        $ m$ / ${m_b}$ & Median reward obtained per episode for 
            test / baseline runs\\
        ${\overline{i_T}}$ & Mean terminal inventory per episode \\
        ${\sigma_{i_T}}$ & 
            Standard deviation of terminal inventory per episode \\
        ${\min_{i_T}}$ & Minimum terminal inventory in test run \\
        ${\max_{i_T}}$ & Maximum terminal inventory in test run \\
        ${v^*_0}$ & First observation value
            \eqref{eq:ExecutionFirstObservationValue} (only for test runs) \\
        $\%I_\rho$ & Percentage of episodes in a test run with 
            $i_T\in I_\rho = (-\rho, \rho]$ \\
        $\%I_{\rho/2}$ & Percentage of episodes in a test run with 
            $i_T\in I_{\rho/2} = (-\rho/2, \rho/2]$ \\
        \bottomrule
    \end{tabularx}
\end{table}

\begin{comment}
We have three quantities we can compare:
\begin{itemize}
    \item Let $\overline{r}_b$ be the average reward obtained per episode when
    following the optimal baseline policy $\pi^{*,b}$ for $M_\text{test}$. From the definition of $\pi^{*,b}$, it follows that $\overline{r}_b$ is a sample-based estimate of $V_0^*(o_0)$, i.e. the optimal value over an entire episode if we start in observation $o_0=0.7$. 
    
    \item In a test run actions are selected according to $\pi^\text{test}$, thus, $\overline{r}$ is an estimate of $V^{\pi^\text{test}}_0(\hat{o_0})$. If $A(\cdot\,|\,\ell_k)$ has converged to the true Q-functions, and returns the true Q-value for history-action pairs at least for $t=0$, then we must have $\overline{r}\approx\overline{r}_b$. Due to errors, and imperfect convergence (finite number of episodes) we expect $\overline{r}<\overline{r}_b$, but hope for
    $\abs{\overline{r}-\overline{r}_b}\approx 0$.

    \item Since $Q(\hat{o}_0,\cdot\,|\,\ell_k)$ is supposed to be an estimate of the true Q-value $Q_0^*(\hat{o}, \cdot)$, and due to \cref{eq:HDPValueQConnection}, $v_0^*$ is supposed to be an estimate of $V_0^*(o_0)$. Again, if $Q(\hat{o}_0,\cdot\,|\,\ell_k)$ converged to the true Q-functions, then we should have $v_0^*\approx \overline{r}_b$.

    \item There is the possibility that neither $v_0^*$ nor $\overline{r}$ approximate the reward obtained when following the optimal baseline policy. More specifically, we would have $\overline{r} > \overline{r}_b$ since the policy derived from the learned approximate Q-functions is sub-optimal. Nevertheless, in this case the two estimators may still agree, i.e. $v_0^*\approx\overline{r}$. While it is clear that the algorithm has not learned to approximate the true optimal Q-functions, it still converged to some Q-functions and that $Q(\cdot\,|\,\ell_k)$ correctly represents 
    That this need not be the case in general is clear from the results of the experiment in \cref{chapter:Experiments:section:MountainCar}, where the learned approximate Q-function vastly overestimates any possible Q-values, however it still manages to solve the problem.
\end{itemize}
\end{comment}
    

%------------------------------------------------------------------------------
\subsection{Results}
\label{chapter:Experiments:section:Execution:sub:Results}
%------------------------------------------------------------------------------

First, the results from running the optimal baseline policy are given, then  
we describe the Signature-Q-Learning algorithm's training behaviour, and 
subsequently results from testing the obtained approximate Q-functions are 
presented.

%------------------------------------------------------------------------------
\paragraph{Baseline Results}

\Cref{tab:ExecutionBaselineStatistics} displays the results from running the 
optimal baseline policy $\pi^b$ for $M_b=2000$ episodes. By design, the mean 
terminal inventory $\overline{i_T}$ is approximately zero and in each episode 
$i_T$ falls within the interval $(-\rho/2, \rho/2]$ with minimal order size
$\rho=50$, as indicated by the minimum and maximum terminal inventory values 
over episodes. The histogram and boxplot in 
\cref{fig:BaselineRewardDistribution} show a unimodal and left-skewed 
distribution of the episode rewards, which explains the median episode reward 
being larger than the mean. The sample's Fisher-Pearson coefficient of skewness 
is $-2.32$, indicating moderate skewness. The mean episode reward $\overline 
r_b = -0.0884$ will serve as benchmark for policies learned by Signature-Q-%
Learning. The mean episode reward's $95\%$ confidence interval is $(-0.08763, 
-0.08117)$, which corresponds to a deviation of $\pm 3.83\%$ from the sample 
mean $\overline r_b$. Its calculation is based on the central limit theorem 
with bounds given by
\[
    \overline r_b \pm t_{M_b-1, 0.975} \frac{\sigma_b}{\sqrt{M_b}},
\]
with an estimated standard error of the sample mean $\sigma_b/\sqrt{M} = 
0.00165$, and where $t_{M_b-1,0.975}$ is the $0.975$-quantile of the Student's 
$t$ distribution with $M_b-1=1999$ degrees of freedom. Due to the large sample 
size and moderate skewness, this confidence interval calculation is deemed 
appropriate. Additionally, a $95\%$ confidence interval based on bootstrapping 
with replacement was computed, with the interval's bounds agreeing with the 
ones presented above up to an precision of $10^{-4}$. Values are rounded to 5 
decimal points.

% baseline result statistics table execution
\begin{table}[htpb!]
    \centering
    \caption[%
        Baseline policy results in the Simple Execution experiment.
    ]{%
        Reward and inventory statistics received when following the optimal
        baseline policy $\pi^b$ from \cref{eq:ExecutionBaselinePolicy} in the 
        Simple Execution environment for $M_b = 2000$ episodes. Values are rounded to five and two decimal points, respectively.
    }
    \label{tab:ExecutionBaselineStatistics}        
    \begin{tabular*}{\textwidth}{c @{\extracolsep{\fill}} *{6}{c}}
        \toprule
        ${\overline r_b}$ & ${\sigma_{r_b}}$ & ${m_b}$ &
            ${\overline{i_T}}$ & ${\sigma_{i_T}}$ & ${\min_{i_T}}$ & 
                ${\max_{i_T}}$\\
        \midrule
    $-0.08440$ & $0.07359$ & $-0.06188$ & $-0.22$ & $13.10$ & $-24$ & $25$ \\
        \bottomrule
    \end{tabular*}
\end{table}

\begin{figure}[htbp!]
    \centering
    \includegraphics[width=0.85\linewidth]{graphics//execution/baseline_reward_distribution.png}
    \caption[Distribution of episode rewards from baseline policy]{Histogram (left) and boxplot (right) of episode rewards when following the optimal baseline policy $\pi^b$ for $M_b=2000$ episodes.}
    \label{fig:BaselineRewardDistribution}
\end{figure}


%------------------------------------------------------------------------------
\paragraph{Training Results}

Analysing Signature-Q-Learning's during training shows robust learning, with 
all ten runs exhibiting a similar and stable training behaviour. This is 
evident from \cref{fig:ExecutionTrainingResults}, which displays the average 
cumulative reward, the average cumulative loss and the average terminal 
inventory $i_T$ per episode during training. Each statistic is averaged over 
the ten training runs for each training episode. The standard deviation, 
represented by the shaded areas around the mean values, becomes small over the 
training phase, indicating that the ten runs converged similarity. 
Especially the average loss shows no variation at end of training, meaning 
that each run converged to a stable approximate Q-function. Rewards increase 
quickly in the first third of the training phase and tend towards a value 
slightly below zero. Moreover, 
\cref{fig:ExecutionTrainingResults:sub:TerminalInventory} shows that on 
average, the agent learns to end the trading window with a terminal inventory 
$i_T$ close to zero, a characteristic that is expected from an (approximately) 
optimal policy. 

Further evidence for the algorithm's convergence is provided by the evolution 
of the initial observation value $v_0^*$ during training. Recall the 
definition of $v_0^*$ from \cref{eq:ExecutionFirstObservationValue}, as the 
value of the initial observation $\hat{o}_0 = (1, 0.7)$ provided by the
learned approximate Q-function at the end of training. Analogously, for each 
training episode $n\in\{1,\dots,3000\}$, we can calculate $v_0^*(n)$ as the value of 
$\hat{o}_0 = (1,0.7)$ with the \textit{current} estimate $Q(\cdot |
\ell_{k(n)})$, where $\ell_{k(n)}$ is the corresponding linear functional at 
the end of training episode $n$, after update step $k(n)$. 
\Cref{fig:ExecutionFirstObsValue:sub:AllRuns} shows 
the evolution of $v_0^*(n)$ during training for each run, together with their 
sample mean and standard deviation calculated over runs for each episode.
Clearly, $v_0^*(n)$ converges towards the same value for each run, and 
importantly, the limit value corresponds to the average episode reward 
$\overline r^* = -0.0844$ from following the optimal baseline policy. Since 
$\overline r^*$ is an estimate of the true optimal value $V_0^*$, convergence 
of $v_0^*(n)$ provides us with numerical evidence, that Signature-Q-Learning 
indeed successfully learned the true optimal Q-function $Q_0^*$. A closer look 
at the last third of the training episodes in 
\cref{fig:ExecutionFirstObsValue:sub:RewardvsFirstObs} shows the average 
$v_0^*(n)$ values tending to $\overline r_b = -0.0844$ with decreasing 
standard deviation towards the end. Additionally, also the average episode 
reward roughly settles around $\overline r_b$ indicating convergence, albeit 
it fluctuating with high variability due to the rewards depending on the 
mid-price trajectory in each episode.

\cref{fig:ExecutionObservationActionTrajectories} offers a more qualitative 
insight, which again highlights the successful training. It shows the 
trajectories of selected actions $(a_t)_{t\in[T]}$ and resulting inventories 
$(i_t)_{t\in[T]}$ for all runs in a selected number of training episodes. 
Starting from a random behaviour in the first training episode due to the 
exploration parameter's value $\varepsilon_0 = 1$, in all runs 
Signature-Q-Learning is able to improve it's action selection as training 
progresses, until in the last training episode 
the trajectories clearly resemble the ones generated from following the 
optimal baseline policy, displayed in \cref{fig:BaselineTrajectories}. This
shows that in all but one training run the baseline strategy is learned in the 
end. Only the action Run 5 deviates slightly from an optimal action trajectory 
by selecting action 0 (limit buy) around step 20. Note, however, that this may 
be due to the $\varepsilon$-greedy policy followed, since even in the last 
training episode the probability of selecting an action randomly is positive. 

Overall, Signature-Q-Learning demonstrated robust training and exhibited 
convergence to the truly optimal Q-functions represented by the optimal 
baseline policy $\pi_b$. Recall that in the Mountain Car environment in 
\cref{chapter:Experiments:section:MountainCar} episode losses and first 
observation values diverged in training, despite the algorithm successfully 
solving the problem, and it was hypothesis, that this is a consequence of the 
varying episode lengths. The difference to training 
behaviour observed for the Simple Execution environment which features a fixed 
episode length, might indeed support this hypothesis, however more analysis is needed.

% averaged training result plots
\begin{figure}[htpb!]
    \centering
    \begin{subfigure}{0.45\textwidth}
        \includegraphics[width=\textwidth]{graphics/execution/custom_execution_mean_reward_20250127_A.png}
        \caption{Average reward per episode.}
        \label{fig:ExecutionTrainingResults:sub:Reward}
    \end{subfigure}
    \hfill
    \begin{subfigure}{0.45\textwidth}
        \includegraphics[width=\textwidth]{graphics/execution/custom_execution_mean_loss_20250127_A.png}
        \caption{Average loss per episode.}
        \label{fig:ExecutionTrainingResults:sub:Loss}
    \end{subfigure}
    \hfill
    \begin{subfigure}{0.9\textwidth}
        \centering
        \includegraphics[width=0.5\textwidth]{graphics/execution/custom_execution_mean_terminal_inventory_20250127_A.png}
        \caption{Average terminal inventory $i_T$ per episode.}
        \label{fig:ExecutionTrainingResults:sub:TerminalInventory}
    \end{subfigure}
            
    \caption[Averaged training results in the Simple Execution experiment.
    ]{%
        Averaged training results of the Signature-Q-Learning algorithm in the
        Simple Execution experiment. 10 training runs were performed in total, each with 3000 episodes and
        with the hyperparameter set-up detailed in
        \cref{tab:ExecutionTrainingParameters}. The sub-figures show the
        respective statistic per episode averaged over training runs.
        (\subref{fig:ExecutionTrainingResults:sub:Reward}) shows the cumulative,
        undiscounted reward obtained per episode.
        (\subref{fig:ExecutionTrainingResults:sub:Loss}) shows the average total loss 
        obtained per episode.
        (\subref{fig:ExecutionTrainingResults:sub:TerminalInventory}) shows the 
        average inventory held at terminal time $T=180$ per episode.
    }
    \label{fig:ExecutionTrainingResults}
\end{figure}

% first observation values
\begin{figure}
    \centering
    \begin{subfigure}{0.93\textwidth}
        \includegraphics[width=\linewidth]{graphics/execution/first_observation_values_20250127_A.png} 
        \caption{
            First observation values $v_0^*$ during training for each run.
            Sample mean and standard deviation is calculated at each episode 
            over all runs. The dotted line corresponds to the average baseline 
            episode reward $\overline r_b = -0.0844$.
        }
        \label{fig:ExecutionFirstObsValue:sub:AllRuns}
    \end{subfigure}
    \hfill
    \begin{subfigure}{0.93\textwidth}
        \includegraphics[width=\textwidth]{graphics/execution/reward_vs_first_observation_value_20250127_A.png}
        \caption{
            Comparison of average reward (\textbf{left}) and average first 
            observation value $v_0^*$ (\textbf{right}) in the last 1000 
            training episodes. The dotted black line in each subplot 
            corresponds to the mean baseline episode reward $\overline r_b = 
            -0.0844$.
        }
        \label{fig:ExecutionFirstObsValue:sub:RewardvsFirstObs}
    \end{subfigure}     
    \caption[First observation value during training in the Simple Execution environment]{The figures illustrate the development of the first observation value $v_0^*(n)$ as a function of the training episode $n$, and compares it to the average episode reward.}
    \label{fig:ExecutionFirstObsValue}
\end{figure}

% training behaviour
\begin{figure}[htpb!]
    \centering
    \includegraphics[width=0.9\textwidth]{graphics/execution/custom_execution_observation_action_histories_20250127_A.png}
    \caption[Inventory and action trajectories in training in the Simple Execution experiment]
    {
        The figure shows trajectories of observations $o_t=i_t/i_{\max}$ 
        (\textbf{left column}) and actions $a_t$ (\textbf{right column}) for 
        $t\in[T]$, with the number of episode steps $T=180$, and for all performed runs in a number of 
        training episodes: At the start of training in episode 1 (first row), 
        in episode 1000 (second row), in episode 2000 (third row), and at the 
        end of training in episode 3000 (fourth row). The labels of the $x$-axis of the graphs in the bottom row apply analogously to the graphs in the rows above, and correspond to the time step $t\in[T]$ of an episode.
    }
    \label{fig:ExecutionObservationActionTrajectories}
\end{figure}


%%%%%%%%%%%%%%%%%%%%%%%%%%%%%%%%%%%%%%%%%%%%%%%%%%%%%%%%%%%%%%%%%%%%%%%%%%%%%%%
\clearpage
\paragraph{Testing Results}

Each of the ten obtained approximate Q-function from training is tested for 
$M_\text{test}=2000$ episodes. Results regarding rewards obtained in the test 
runs are displayed in \cref{tab:ExecutionTestResultsRewardsStatistics}. The 
first thing to note is, that all runs performed similarity in testing in terms 
of mean episode reward, indicating that for this specific hyperparameter 
configuration the algorithm is not sensitive to parameter initialization and 
the seed used to simulate the environment. The values obtained for each 
statistic are of the same order of magnitude for all test runs as well as for 
the baseline run. Test run 4 stands out as best performing run in terms of the 
episode rewards' mean and median and variation, measured by theirs standard 
deviation. The largest absolute difference in mean reward between any two runs 
is $8.68\times 10 ^{-3}$, which corresponds to roughly $10.5\%$ of the best 
run's mean value. Notably, despite the apparent convergent behaviour seen in 
training, the mean episode reward is consistently lower than the mean baseline 
episode reward of $\overline r_b=-8.44 \times 10^{-2}$. The minimal absolute 
difference is obtained in run 4, and the maximal absolute difference in run 7, 
with values provided in \cref{tab:ExecutionTestEstimateDifferences}. Furthermore, the 
variability in episode rewards is high, with the standard deviations of the 
same order of magnitude as the means. This is consistent with the values 
obtained for the baseline policy from \cref{tab:ExecutionBaselineStatistics}. 
The boxplots of episode rewards in \cref{fig:ExecutionTestingRewardsBoxplot} 
show a left-skewed distribution for each test run, similar to the one obtained 
for the baseline policy, albeit with heavier left tails, indicated by larger 
interquartile ranges and lower outlier values for most runs. The seed to 
initialized the environment used in the baseline run was different to any test 
run to avoid correlation in results due to the same pseudo-random sampling 
behaviour.

% test result statistics table Mountain Car set-up
\begin{table}[htpb!]
    \centering
    \caption[%
        Test statistics for rewards in the Simple Execution experiment.
    ]{%
        Reward statistics obtained from testing the approximate Q-function from 
        each training run for $M_b=2000$ test episodes in the Simple Execution 
        experiment. Bold values highlight the best result over test runs: 
        highest reward, lowest standard deviation, lowest $v_0^*$ value. 
        $v^*_0$ is calculated as in \cref{eq:ExecutionFirstObservationValue} 
        for the initial observation $\hat{o}_0 = (1, 0.7)$ and with the 
        basepoint $\hat{o}_{-1} = (0, 0)$ for the truncated signature's 
        computation at time $t=0$. Values are rounded to five decimal points. 
    }
    \label{tab:ExecutionTestResultsRewardsStatistics}        
    \begin{tabularx}{\textwidth}{ >{\centering}m{1.5cm} *{4}{C} }
        \toprule
        Run & ${\overline r}$ & ${\sigma_r}$ & 
            ${m}$ & ${v^*_0}$ \\
        \midrule
        \midrule

        $ 1$ & $-9.290\times10^{-2}$ & $8.177\times10^{-2}$ 
            & $-6.769\times10^{-2}$ & $-8.817\times10^{-2}$ \\
            
        $ 2$ & $-9.323\times10^{-2}$ & $9.122\times10^{-2}$ 
            & $-6.439\times10^{-2}$ & $-8.194\times10^{-2}$ \\

        $ 3$ & $-9.392\times10^{-2}$ & $8.452\times10^{-2}$ 
            & $-6.847\times10^{-2}$ & $-8.479\times10^{-2}$ \\

        $ 4$ & $\bm{-8.628\times10^{-2}}$ & $\bm{7.814\times10^{-2}}$ 
            & $\bm{-6.092\times10^{-2}}$ & $-8.777\times10^{-2}$ \\

        $ 5$ & $-9.237\times10^{-2}$ & $8.296\times10^{-2}$ 
            & $-6.766\times10^{-2}$ & $-8.242\times10^{-2}$ \\

        $ 6$ & $-9.220\times10^{-2}$ & $8.553\times10^{-2}$ 
            & $-6.430\times10^{-2}$ & $-8.474\times10^{-2}$ \\

        $ 7$ & $-9.496\times10^{-2}$ & $8.317\times10^{-2}$ 
        & $-6.797\times10^{-2}$ & $\bm{-8.163\times10^{-2}}$ \\

        $ 8$ & $-9.418\times10^{-2}$ & $7.865\times10^{-2}$ 
            & $-7.004\times10^{-2}$ & $-8.403\times10^{-2}$ \\

        $ 9$ & $-9.385\times10^{-2}$ & $8.327\times10^{-2}$ 
            & $-6.841\times10^{-2}$ & $-7.945\times10^{-2}$ \\

        ${10}$ & $-9.456\times10^{-2}$ & $8.177\times10^{-2}$ 
            & $-6.848\times10^{-2}$ & $-8.528\times10^{-2}$ \\
    
        \bottomrule
    \end{tabularx}
\end{table}

The first observation value $v_0^*$ provides another estimate of the true 
optimal $V_0^*(o_0)$, and a means of judging the algorithms performance. 
With the exception of run 4, $v_0^*$ is consistently larger than $\overline r$, 
estimating the expected cumulative episode reward $V_0^*(o_0)$ higher, than 
what was actually obtained from following the learned approximate Q-function in 
a test run. The range of absolute differences between the two sampled estimates 
is given in \cref{tab:ExecutionTestEstimateDifferences}, with the minimal value again 
in run 4. The training results already indicated that the values $v_0^*$ 
approached the sampled average baseline episode reward $\overline r_b =
-8.44\times10^{-2}$, which is visible \Cref{fig:ExecutionFirstObsValue}, and 
this is confirmed by the testing results. Again with the exception of run 4, 
$v_0^*$ is consistently closer to the average baseline episode reward 
$\overline r_b$ than $\overline r$ from the respective test run. In particular, 
the minimal and maximal absolute difference $\abs{v_0^* - \overline r_b}$ over 
runs, are an order of magnitude smaller than the difference between remaining 
estimate combinations, respectively, which is apparent from 
\cref{tab:ExecutionTestEstimateDifferences}. As already mentioned in 
\cref{chapter:Experiments:section:Execution:sub:ABIDES:par:TradingAgents}, the 
fact that $v_0^*$ provides a better approximation to the sampled $\overline 
r_b$ than $\overline r$ in each run, is a sign the algorithm learned an 
approximate Q-function $Q(\cdot\,|\,\ell_K)$ which estimates the true optimal 
Q-function $Q_0^*$ at $t=0$, and as such $V_0^* \approx r_b$, sufficiently 
well, however there is still exist discrepancies to the true Q-functions 
$Q_t^*$ at some later time steps $t>0$. Furthermore, for seven out of the then 
test runs, $v_0^*$ falls within the $95\%$-confidence interval for the mean 
baseline episode reward, given by $(-8.763, -8.117)\times 10^{-2}$. 

\begin{table}[tbp!]
    \caption[Differences between $V_0^*(o_0)$ estimates in the Simple Execution environment]{Absolute differences between the various $V_0^*(o_0)$ estimates obtained in training and testing in the Simple Execution experiment.}
    \label{tab:ExecutionTestEstimateDifferences}
    \centering
    \begin{tabularx}{\textwidth}{C CCCC}
        \toprule
         \multirow{2}*{Difference} & \multicolumn{2}{c}{Minimum} 
            & \multicolumn{2}{c}{Maximum} \\
        \cmidrule{2-5}
        %\cline{2-5}
             & Run & Value & Run & Value \\
        \midrule
        ${|\overline r - \overline r_b|}$ 
           & 4 & $1.880\times10^{-3}$ & 7 & $1.055\times10^{-2}$ \\
        ${|v_0^* - \overline r|}$ 
           & 4 & $1.489\times10^{-3}$ & 9 & $1.440\times10^{-2}$ \\
        ${|v_0^* - \overline r_b|}$ 
           & 6 & $3.334\times10^{-4}$ & 9 & $4.949\times10^{-3}$ \\
        \bottomrule
    \end{tabularx}
\end{table}

\Cref{tab:ExecutionTestResultsInventoryStatistics} displays statistics about 
terminal inventories from the ten test runs. On average, the learned 
approximate Q-functions resulted in policies that ended episodes with an 
inventory close to zero: 
For half of the test runs $\overline{i_T}$ lies within $I_{\rho/2}$, meaning 
that, on average, the terminal inventory is optimal in the sense that a 
further, fully filled trade in either direction would yield an inventory value 
further away from zero, as explained in
\cref{chapter:Experiments:section:Execution:sub:BaselinePolicy}. For the other 
half, $\overline{i_T}$ still lies within $I_\rho$, indicating sub-optimality, 
in the sense that the inventory could be brought closer to zero with further 
trades. 
%
Further, terminal inventory values exhibit some variability, which is apparent 
from their sample standard deviation $\sigma_{i_T}$, and the wide range 
between $\min_{i_T}$ and $\max_{i_T}$. Since a larger absolute terminal 
inventory translates to a lower reward, this explains the variability in 
episode rewards, reported in \cref{tab:ExecutionTestResultsRewardsStatistics}, 
as well as the lower mean episode reward obtained in test runs, when compared 
to the one obtained with the baseline policy, reported in 
\cref{tab:ExecutionTestEstimateDifferences}.
%
Notably, the best run in terms of rewards, test run 4, does not have the 
lowest mean terminal inventory, but ends episodes most consistently with $i_T$ 
in the intervals $I_{\rho/2}$ and $I_\rho$, respectively. In general, the 
percentages of episodes with $i_T\in I_{\rho/2}$ are lower, with more than 
50\% of episode in three test runs. The share of episodes with $i_T\in I_\rho$ 
however, increases greatly for all runs with all but one above $70\%$. Still, 
the number of episodes with sub-optimal $i_T$ larger than the fixed order size 
$\rho=50$, is above 11\%, even for runs 4 and 6, which are best in terms of 
highest mean episode reward $\overline r$ and lowest mean terminal inventory 
$\overline {i_T}$.

\Cref{fig:ExecutionTestRunsInventoryTrajectories} displays the action 
trajectories for all episodes for each test run, and 
\cref{fig:ExecutionTestRunsInventoryTrajectories} shows the corresponding 
inventory trajectories. Both reveal mixed qualitative behaviour in the test 
runs. Most notably, in most runs there is still a deviation from the optimal 
action trajectory in \cref{fig:BaselineTrajectories:sub:Actions} provided by 
the baseline policy. Only the action sequences of runs 7 and 10 are ... however, in both runs, the policy tends to oversell ending with negative inventories smaller than $-\rho$, visible in the corresponding inventory trajectory plots. Action sequences of runs 5 and 6 are close to the baseline trajectories, but in both the policy occasionally chooses to post a limit buy (action 0) before switching to do nothing. Nevertheless, they perform better in terms of mean rewards and percentage of runs in $I_\rho$ and $I_{\rho/2}$, respectively, and run 6 even exhibits the smallest standard deviation of terminal inventories $\sigma_{i_T}=24.8$, which is visible in inventory trajectories too. In Test run 4 the policy overshoots $-\rho$ by some extend and subsequently starts to post limit buy orders (action  0) to bring the inventory back closer to zero, a behaviour which resulted in the highest mean episode reward. The test run with the lowest mean terminal inventory, run 9, exhibits a similar behaviour, however, instead of switching to limit buy orders in the middle of the episode it does so at the end in a try to increase its inventory closer to zero. In general, we observe some erratic behaviour towards the end of the episodes in half of the runs, an effect likely due to the terminal inventory reward, which encourages the policy to bring the inventory closer to zero before receiving the reward.

In summary, upon observing the test runs' action trajectories the overall 
policies appear not as robust as the training behaviour has indicated. 
Notably, in training runs, the learned initial values $v_0^*$ indeed 
approached the overall baseline episode reward given by $\overline r_b$, 
indicating their convergence to the initial values provided by the true 
optimal Q-function $Q^*_0$. Test results, however, hint at the fact that 
Q-values later during the episode might not have converged sufficiently. This 
is visible from the fact that mean episode rewards $\overline r$ obtained in 
test runs are below mean episode rewards $\overline r_b$ from baseline runs, 
and from the action selection in test runs, which shows the desired behaviour  
to some degree, however, still deviates from $\pi^b$. Approaches to improve 
the learned Q-functions, and bring action selection closer to the optimal 
baseline policy, include longer training, a systematic search for suitable 
training hyperparameters, and the tuning of the reward structure by balancing 
running and terminal rewards through parameter $C$ in 
\eqref{eq:TerminalRewardExecution}. 

% test result statistics table Mountain Car set-up
\begin{table}[htpb!]
    \centering
    \caption[%
        Test statistics for terminal inventories in the Simple Execution experiment
    ]{%
        Terminal inventory statistics obtained from testing the approximate 
        Q-function from each training run for $M_b=2000$ test episodes in the 
        Simple Execution experiment. Bold values highlight the best result over 
        test runs: smallest absolute inventory values, lowest standard 
        deviation, highest percentage points. Values are rounded to one decimal 
        point. 
    }
    \label{tab:ExecutionTestResultsInventoryStatistics}        
    \begin{tabularx}{\textwidth}{ *{7}{C} }
        \toprule
        Run & ${\overline{i_T}}$ & ${\sigma_{i_T}}$ 
            & ${\min_{i_T}}$ & ${\max_{i_T}}$ & ${\%I_{\rho/2}}$ 
                & ${\%I_\rho}$ \\
        \midrule
        \midrule
        $1$  & $37.5$ & $43.4$ & $-76$ & $501$ & $34.9$ & $75.6$\\
        $2$ & $13.0$ & $44.0$ & $-257$ & $110$ & $41.6$ & $80.5$ \\
        $3$ & $37.9$ & $38.5$ & $-51$ & $164$ & $41.8$ & $68.2$ \\
        $4$ & $-13.9$ & $29.4$ & $-123$ & $132$ & $\bm{65.4}$ 
            & $\bm{88.8}$ \\
        $5$ & $14.4$ & $38.8$ & $-62$ & $188$ & $52.2$ & $83.2$ \\
        $6$ & $28.7$ & $\bm{24.8}$ & $\bm{-33}$ & $\bm{120}$ 
            & $48.4$ & $82.3$ \\
        $7$ & $34.9$ & $33.3$ & $-201$ & $169$ & $33.1$ & $72.4$ \\
        $8$ & $34.5$ & $42.4$ & $-50$ & $202$ & $43.9$ & $70.4$ \\
        $9$ & $\bm{5.4}$ & $39.8$ & $-279$ & $132$ & $53.2$ & $85.1$ \\
        $10$  & $22.8$ & $36.8$ & $-185$ & $136$ & $39.3$ & $78.8$ \\
        \bottomrule
    \end{tabularx}
\end{table}


\begin{figure}[htbp!]
    \centering
    \includegraphics[width=1\linewidth]{graphics//execution/testing_boxplot_all_rewards_20250127_A_2000_episodes.png}
    \caption[Boxplots of episode rewards from testing in the Simple Execution experiment]{%
        Boxplot of the episode rewards in the Simple Execution environment 
        obtained in the ten test runs and the baseline run with $M_\text{test} 
        = M_b = 2000$ episodes.
    }
    \label{fig:ExecutionTestingRewardsBoxplot}
\end{figure}

\begin{figure}[htbp!]
    \centering
    \includegraphics[width=0.94\linewidth, height=0.9\textheight, keepaspectratio]
        {graphics/execution/test_runs_action_trajectories.png}
    \caption[Inventory trajectories in test runs 4 and 6 in the Simple Execution experiment]{%
        Inventory trajectories in test runs. The dotted lines 
        correspond to the inventory values $\{-50, 0, 50\}$, respectively, with 
        $\rho=50$ the fixed order size in the environment. Subfigures in a column share their $x$-axis, and plots in a row share their $y$-axis.
    } 
    \label{fig:ExecutionTestRunsActionTrajectories}
\end{figure}

\begin{figure}[htbp!]
    \centering
    \includegraphics[width=0.94\linewidth, height=0.9\textheight, keepaspectratio]
        {graphics/execution/test_runs_inventory_trajectories.png}
    \caption[Inventory trajectories in test runs 4 and 6 in the Simple Execution experiment]{%
        Inventory trajectories in test runs. The dotted lines 
        correspond to the inventory values $\{-50, 0, 50\}$, respectively, with 
        $\rho=50$ the fixed order size in the environment. Subfigures in a column share their $x$-axis, and plots in a row share their $y$-axis.
    } 
    \label{fig:ExecutionTestRunsInventoryTrajectories}
\end{figure}



%%%%%%%%%%%%%%%%%%%%%%%%%%%%%%%%%%%%%%%%%%%%%%%%%%%%%%%%%%%%%%%%%%%%%%%%%%%%%%%
% Appendix
\appendix 
%%%%%%%%%%%%%%%%%%%%%%%%%%%%%%%%%%%%%%%%%%%%%%%%%%%%%%%%%%%%%%%%%%%%%%%%%%%%%%%

%%%%%%%%%%%%%%%%%%%%%%%%%%%%%%%%%%%%%%%%%%%%%%%%%%%%%%%%%%%%%%%%%%%%%%%%%%%%%%%
\chapter{Appendix} 
\section{Stochastic Kernels} 
\label{chapter:PreliminariesStochasticKernels}
%%%%%%%%%%%%%%%%%%%%%%%%%%%%%%%%%%%%%%%%%%%%%%%%%%%%%%%%%%%%%%%%%%%%%%%%%%%%%%%

Some results on stochastic kernels are collected. They are used to 
construct the components of the decision processes treated in this thesis. 
The contents in this section mainly follow the exposition in
\cite{klenke2020_ptcc}. Exact references are given with the respective statements in the following.

\begin{definition}[Definition~8.25 in \cite{klenke2020_ptcc}]
\label{def:StochasticKernel}
    Let $(\Omega_1, \mathcal{F}_1)$ and $(\Omega_2, \mathcal{F}_2)$ be two
    measurable spaces. A \textit{stochastic kernel} $\kappa$ from $(\Omega_1,
    \mathcal{F}_1)$ to $(\Omega_2, \mathcal{F}_2)$ is defined as the map
    $\kappa: \Omega_1 \times \mathcal{F}_2  \to [0,1]$
    %\[
    %    \kappa: \mathcal{F}_2 \times \Omega_1 \to [0,1],
    %    \quad (A,\omega) \mapsto \kappa(A, \omega),
    %\]
    such that
    \begin{enumerate}
        \item for any fixed $A\in\mathcal{F}_2$, the map $\omega \mapsto 
            \kappa(\omega, A)$ is $\mathcal{F}_1$-measurable,
            
        \item \label{def:StochasticKernel:item2}
            for any fixed $\omega\in\Omega_1$, the map $A \mapsto 
            \kappa(\omega, A)$ is a probability measure on $\mathcal{F}_2$.
            %which will be denoted as $\kappa(\,\cdot\,|\,x)$.
    \end{enumerate}
    The probability measure defined by some $\omega$ in 
    \ref{def:StochasticKernel:item2} will be denoted as $\kappa(\cdot \,|\,
    \omega) := \kappa(\omega, \cdot)$ and we also write $\kappa(\dd y\,|\,
    \omega)$, such that $\kappa(A\,|\,\omega) = \kappa(\omega, A)$ is the 
    probability of event $A$ under this measure. If it is clear from the 
    context, we will say $\kappa$ is a stochastic kernel from $\Omega_1$ 
    to $\Omega_2$ without explicitly mentioning the corresponding $\sigma$-
    algebras.
\end{definition}

In the literature, a kernel $\kappa$ is often referred to as \textit{Markov} 
kernel or \textit{transition} kernel. We do not use the term \textit{Markov} 
here to avoid ambiguities when defining the more general history-based 
decision processes and we will use the term transition kernel explicitly for 
the more specific kernels in the definition of decision processes.
        
Similar to product measures we can define the product of stochastic kernels.

\begin{theorem}[Theorem~14.25 in \cite{klenke2020_ptcc}]
\label{thm:ProductKernel}
    Let $(\Omega_i, \mathcal{F}_i), i=0,1,2$ be measurable spaces, $\kappa_0$
    a stochastic kernel from $(\Omega_0, \mathcal{F}_0)$ to $(\Omega_1,
    \mathcal{F}_1)$ and $\kappa_1$ a stochastic kernel from $(\Omega_0\times
    \Omega_1, \mathcal{F}_0\otimes\mathcal{F}_1)$ to $(\Omega_2,
    \mathcal{F}_2)$. Then, the \textbf{product} of
    $\kappa_0$ and $\kappa_1$, defined as the map
    \begin{equation}
    \label{eq:KernelProduct}
        \begin{split}
            \kappa_0 \otimes \kappa_1:
                \Omega_0 \times (\mathcal{F}_1\otimes\mathcal{F}_2) 
                    &\to 
                        [0,\infty), \\
            (\omega, A) 
                &\mapsto 
                    \int_{\Omega_1} \int_{\Omega_2} {
                        \mathds{1}_A(\omega_1, \omega_2) 
                        \,\kappa_1((\omega_0, \omega_1), \dd{\omega_2})
                        \,\kappa_0(\omega_0, \dd{\omega_1})
                    },
        \end{split}
    \end{equation}
    is well-defined and a stochastic kernel from $(\Omega_0, \mathcal{F}_0)$
    to $(\Omega_1\times\Omega_2, \mathcal{F}_1\otimes\mathcal{F}_2)$.
\end{theorem}

\begin{corollary}
\label{cor:MeasureKernelProduct}
    Let $(\Omega_1, \mathcal{F}_1, \mu)$ be a probability space and $(\Omega_2,
    \mathcal{F}_2)$ a measurable space. Let $\kappa$ be a stochastic kernel
    from $\Omega_1$ to $\Omega_2$. Then there exists a unique probability 
    measure $\mu \otimes \kappa$ on $(\Omega_1\times\Omega_2, \mathcal{F}_1
    \otimes \mathcal{F}_2)$ defined on measurable rectangles as
    \begin{equation}
        \mu\otimes\kappa\, (A_1\times A_2) := 
            \int_{A_1} \kappa(\omega_1, A_2) \,\mu(\dd{\omega_1}).
    \end{equation}
\end{corollary}
\begin{proof}
    This follows directly from \cref{thm:ProductKernel} if we consider $\mu$ 
    as a stochastic kernel not depending on the $\Omega_0$-coordinate and set 
    $\kappa_0(\omega_0, \cdot\,) = \mu$ and $\kappa_1 = \kappa$.
\end{proof}

As for ordinary product measures, there is a analogue of Fubini's theorem for 
stochastic kernels. 
\begin{theorem}[Theorem~14.32 in \cite{klenke2020_ptcc}]
\label{thm:KernelFubini}
    Assume the setting of \cref{cor:MeasureKernelProduct} and let the function 
    $f:\Omega_1 \times \Omega_2 \to \overline{\SetR}$ be measurable with 
    respect to the product $\sigma$-algebra $\mathcal{F}_1 \otimes 
    \mathcal{F}_2$. If both or either of the conditions $\int f^+\dd{\PM\otimes\kappa}<+\infty$ and $\int f^-\dd{\PM\otimes\kappa}<+\infty$ hold, then
    %alternatively $f\leq 0$ and
    %f\geq 0$ or $\int\abs{f}\dd{\PM\otimes\kappa}<+\infty$,
    \[
        \int_{\Omega_1\times\Omega_2} 
            f \dd{(\mu\otimes\kappa)}
        = \int_{\Omega_1} \int_{\Omega_2}
            f(\omega_1, \omega_2) 
            \kappa(\omega_1, \dd{\omega_2})
            \mu(\dd{\omega_1}).
    \]
\end{theorem}

To ease notation in the following, for measurable spaces $(\Omega_k, \mathcal{F}_k), k\in\SetN$ we set
\begin{equation}
\label{eq:ProductMeasurableSpaces}
    (\Omega^{(i,j)}, \mathcal{F}^{(i,j)}) := 
        \Bigl(
            \prod_{k=i}^{j}\Omega_k, \bigotimes_{k=i}^{j}\mathcal{F}_k
        \Bigr),
        \quad \text{for } i,j\in\SetN \text{ with } i\leq j,
\end{equation}
and with elements $\omega_{i:j}:=(\omega_i,\omega_{i+1},\dots,
\omega_j)\in\Omega^{(i,j)}$.

\begin{corollary}
\label{cor:NKernelProduct}
    For $n\in\SetN$, let $(\Omega_i, \mathcal{F}_i), i=0,1\dots,n$ be 
    measurable spaces and let $\kappa_i$ be a stochastic kernel from 
    $(\Omega^{(0,i)}, \mathcal{F}^{(0,i)})$ to $(\Omega_{i+1},
    \mathcal{F}_{i+1})$ for $i=0,\dots,n-1$. Then, the recursion 
    %$
    %\bigotimes_{j=0}^i \kappa_j := 
    %    \bigl( \bigotimes_{j=0}^{i-1} \kappa_j \bigr) \otimes \kappa_i
    %$
    $\kappa^{(0,i)} := \kappa^{(0,i-1)}\otimes\kappa_i$
    defines a stochastic kernel from $(\Omega_0, \mathcal{F}_0)$ to 
    $(\Omega^{(1,i+1)}, \mathcal{F}^{(1,i+1)})$ for each $i=0,\dots,n-1$.
    %
    Furthermore, if $\mu$ is a probability measure on $(\Omega_0, 
    \mathcal{F}_0)$ then, $\PM_i:= \mu\otimes\kappa^{(0,i-1)}$,
    defines a probability measure on $(\Omega^{(0,i)}, \mathcal{F}^{(0,i)})$,
    for each $i=0,1,\dots,n$, with the convention $\PM_0 \equiv \mu$.
\end{corollary}
\begin{proof}
    This follows by inductively applying \cref{thm:ProductKernel}
    as well as from \cref{cor:MeasureKernelProduct}.
\end{proof}

The family of probability measures ${(\PM_i)}^n_{i=0}$ defined in
\cref{cor:NKernelProduct} are by construction projective, by which we mean 
that for any $i,j,k\in\{1,\dots, n\}$ with $i,j\geq k$, and for any $A\in
\mathcal{F}^{(0,k)}$, we have
\[
    \PM_i[A\times\Omega_{k+1}\times\dots\times\Omega_i] 
    = 
    \PM_j[A\times\Omega_{k+1}\times\dots\times\Omega_j].
\]
%
In the setting of \cref{cor:NKernelProduct}, we define the \textit{canonical 
probability space} given the measurable spaces $(\Omega_i,\mathcal{F}_i)$, the 
stochastic kernels $\kappa_i$ and initial distribution $\mu$, as the product space
\begin{equation}
\label{eq:CanonicalProbSpace}
    (\Omega, \mathcal{F}, \PM) 
    := 
    \bigl(\Omega^{(0,n)}, \mathcal{F}^{(0,n)}, \PM_n\bigr).
    %= (
    %    \prod_{i=0}^n \Omega_i, 
    %    \bigotimes_{i=0}^n \mathcal{F}_i, 
    %    \PM_0\otimes\bigotimes_{i=1}^{n}\kappa_i
    %)
\end{equation}
The measure $\PM$ is the unique probability measure on the probability space 
\eqref{eq:CanonicalProbSpace} that satisfies 
\[
    \PM[A\times \Omega_{k+1}\times\dots\times\Omega_n] = \PM_k[A],
\]
for any $k\in\{0,1,\dots n\}$ and $A\in\mathcal{F}^{(0,k)}$. For $i\in\{0,1,
\dots,n\}$, let the $i$-th coordinate map be denoted as $X_i(\omega) := 
\omega_i$ for $\omega\in\Omega$. The \textit{canonical process} 
$X:= (X_i)_{i\in\{0,1,\dots,n\}}$ is a stochastic process on $(\Omega,
\mathcal{F})$, and by construction, the finite dimensional distributions of 
$X$ under $\PM$ are determined by the kernels and the initial
probability measure,
\begin{equation}
    \PM_{X_0} = \mu, \quad \text{and} \quad
    \PM_{X_0, X_1, \dots, X_i} = \mu\otimes\kappa^{(0,i-1)}
    \quad \text{for } i \in\{1,\dots,n\},    
\end{equation}
where $\PM_{X_0, X_1, \dots, X_j}$ is the push-forward measure of $\PM$ under
$X_0,X_1,\dots, X_j$. 

Stochastic kernels define regular conditional distributions 
\cite[Definition~8.28]{klenke2020_ptcc}. With the stochastic process $X$ 
constructed as above, the kernel $\kappa_i$ from 
\cref{cor:MeasureKernelProduct} corresponds to a regular conditional 
distribution of $X_{i+1}$, given $X_{0:i} = (X_0,X_1,\dots,X_i)$. This mean 
that it determines the transition dynamics of $X$ at point $i$, and for 
$\omega_{0:i}\in \Omega^{(0,i)}$ and $F\in\mathcal{F}_{i+1}$,
\begin{equation}
    \PM[X_{i+1}\in F \,|\, X_{0:i} 
    = \omega_{0:i}] = \kappa_{i}(\omega_{0:i}, F).
\end{equation}
This allows us to express conditional expectations of functions of $X_{i+1}$, 
given $X_{0:i}=\omega_{0:i}$, as integrals with respect to the kernel 
$\kappa_i$. Fix $i\in\{0,1,\dots,n-1\}$ and consider an 
$\mathcal{F}_{i+1}$-measurable function $f:\Omega_{i+1}\to\SetR$
which satisfies $\E[\abs{f(X_{i+1})}]<\infty$. Then, the conditional 
expectation of $f(X_{i+1})$ provided that $X_{0:i}=\omega_{0:i}$ can be 
expressed as
\begin{equation}
\label{eq:ConditionalExpectationSingle}    
    \E_{\kappa_i(\dd{x} |\omega_{0:i})}\bigl[f(X_{i+1})\bigr] 
    := \E\bigl[f(X_{i+1}) \,\big|\, X_{0:i}=\omega_{0:i} \bigr]
    = \int_{\Omega_{i+1}} f(x) \,\kappa_i(\dd{x} | \, \omega_{0:i}).
\end{equation}
The notations $\E_{\kappa_i(\dd{x} |
\omega_{0:i})}$ or equivalently $\E_{\kappa_i(\cdot |
\omega_{0:i})}$, if the integration variable is clear from the context, to indicate that the conditional expectation is determined by 
the kernel $\kappa_i$, as integral w.r.t. to the measure $\kappa(\cdot|
\omega_{0:i})$ for inputs $\omega_{0:i}$. 
%
If the function $f$ is instead given as an $\mathcal{F}^{(0,i+1)}$-measurable 
function $f:\Omega^{(0,i)}\times\Omega_{i+1}\to\SetR$ that satisfies 
$\E[\abs{f(X_{0:i},X_{i+1})}]<\infty$, then given $\omega_{0:i} \in 
\Omega^{(0,i)}$, we define again
\begin{equation}
\label{eq:ConditionalExpectationMultipleOne}
    \E_{\kappa_i(\dd{x} |\omega_{0:i})}\bigl[f(\omega_{0:i}, X_{i+1})\bigr] 
    := \E\bigl[f(X_{0:i}, X_{i+1}) \,\big|\, X_{0:i}=\omega_{0:i}\bigr]
    = \int_{\Omega_{i+1}} f(\omega_{0:i}, x) \,\kappa_i(\dd{x}|\,\omega_{0:i}).
\end{equation}
Furthermore, the map $\kappa^{(i,n-1)}$ defined recursively by the product of kernels
\begin{equation}
\label{eq:KernelRecursionBackwards}
    \kappa^{(i,n-1)} :=  \kappa_i \otimes \kappa^{(i+1,n-1)}
    \quad \text{for } i= n-1, n-2, \dots,1,0,
\end{equation}
is a stochastic kernel from $(\Omega^{(0,i)}, \mathcal{F}^{(0,i)})$ to 
$(\Omega^{(i+1,n)}, \mathcal{F}^{(i+1,n)})$ for each $i\in\{0,1,\dots,n-1\}$. 
Note that $\kappa^{(0,n-1)}$ obtained through the recursion 
\eqref{eq:KernelRecursionBackwards} coincides with the one obtained from the 
recursion in \cref{cor:NKernelProduct}. For each $i\in\{0,1,2,\dots,n-1\}$, 
the kernel $\kappa^{(i,n)}$ is a regular conditional distribution of 
$X_{i+1:n}$ given $X_{0:i}$. For a $\mathcal{F}$-measurable function 
$f:\Omega\to\SetR$, which satisfies $\E[\abs{f(X_{0:n})}] < \infty$,
analogue to \eqref{eq:ConditionalExpectationSingle}, we can express 
conditional expectations given $X_{0:i}=\omega_{0:i}$ as integral w.r.t. to 
$\kappa^{(i,n)}$ as
\begin{equation}
\label{eq:ConditionalExpectationMultipleTwo}
    \E[
        f(X_{0:n}) \, | \, X_{0:i} = \omega_{0:i}
    ]
    = \int_{\Omega^{(i+1,n)}} 
        f( 
            \omega_{0:i}, x_{i+1:n}
        ) \, 
        \kappa^{(i,n)} \bigl( 
            \dd{x_{i+1:n}} | \, \omega_{0:i} \bigr
        ) 
\end{equation}
for any $\omega_{0:i}\in\Omega^{(0,i)}$, and where $x_{i+1:n} = (x_{i+1},
\dots,x_n$) are the integration variables and $\omega_{0:i}$ is the input to 
the kernel $\kappa^{(i,n)}$. Since $\kappa^{(i,n)} = \kappa_i \otimes 
\kappa^{(i+1,n)}$, application of \cref{thm:KernelFubini} with the measure 
$\mu = \kappa_i(\cdot|\omega_{0:i})$ and the stochastic kernel $\kappa$ 
defined as
\begin{align*}
    \kappa: 
    \Omega_i\times\mathcal{F}^{(i,n)} &\to [0,1] \\
    (\omega_i, F) &\mapsto \kappa(\omega_i, F) 
        := \kappa^{(i,n)}(F\,|\,\omega_{0:i},\omega_i)
\end{align*}
yield the equivalence of the integral of 
\cref{eq:ConditionalExpectationMultipleTwo} to
\begin{equation}
\label{eq:ConditionalExpectationMultipleEquivalence}
\begin{aligned}
    \int_{\Omega_{i+1}} \int_{\Omega^{(i+2,n)}}
        &
        f( 
            \omega_{0:i}, x_{i+1}, x_{i+2:n}
        ) \,
        \kappa^{(i+1:n)}\bigl(
            \dd{x_{i+2:n}} | \, \omega_{0:i}, x_{i+1} \bigr            
        ) \,
        \kappa_i(
            \dd{x_{i+1}} | \, \omega_{0:i}
        )   \\    
    &= \int_{\Omega_{i+1}}
        \E\bigl[
            f(X_{0:n}) \,\big|\, X_{0:i} = \omega_{0:i}, X_{i+1} = x_{i+1}
        \bigr] \,
    \kappa_i(
        \dd{x_{i+1}} | \, \omega_{0:i}
    ), 
\end{aligned}
\end{equation}
where the inner integral is simply the conditional expectation of 
$f(X_{0:n})$ given $X_{0:i}=\omega_{0:i}$ and $ X_{i+1}=x_{i+1}$ as in 
\eqref{eq:ConditionalExpectationMultipleTwo}. Property 
\Cref{eq:ConditionalExpectationMultipleEquivalence} is of particular importance to 
show recursions for value functions in decision processes, which are defined 
as conditional expectations.

The exposition in this section shows that, given an initial probability space
and a finite sequence of stochastic kernels with associated measurable spaces,
there exists a canonical probability space and a discrete-time stochastic 
process on that space, whose transition dynamics are given by the a priori 
fixed stochastic kernels.\footnote{
    This proposition also holds true for countably many stochastic kernels
    according to the \textit{Ionesca Tulca} theorem 
    \cite[Theorem~14.35]{klenke2020_ptcc}.
}
This ensures the existence of the decision processes that model reinforcement 
learning environments in 
\cref{chapter:RLAndMDP:section:MDP,chapter:HistoryRLEnvironment:section:HDP}.

We conclude this section with two lemmas.

\begin{lemma}
\label{lem:KernelsAuxiliaryMeasurableExpectation}
    For measurable spaces $(\Omega_i, \mathcal{A}_i), i=1,2$ let $\kappa$ be a
    stochastic kernel from $\Omega_1$ to $\Omega_2$ and let $g:\Omega_1 \times
    \Omega_2 \to \overline{\SetR}$ be measurable w.r.t. $\mathcal{A}_1 \otimes
    \mathcal{A}_2$. Then, the function
    $\lambda:\Omega_1 \to \overline{\SetR}$ given by 
    \begin{equation}
    \label{eq:KernelIntFunction}
        \lambda(\omega_1) 
        := \int_{\Omega_2} 
            g(\omega_1, \omega_2) \, \kappa(\dd{\omega_2}|\omega_1)
    \end{equation}
    is well-defined and $\mathcal{F}_1$-measurable if $f\geq 0$, or if for
    each $\omega_1\in\Omega_1$ both or either of the conditions
    \begin{equation}     
    \label{eq:Case2Conditions}
        \int_{\Omega_2} g^+(\omega_1,\omega_2) \, 
            \kappa(\dd{\omega_2|\omega_1}) < \infty
        \quad \text{and} \quad
        \int_{\Omega_2} g^-(\omega_1,\omega_2) \,
            \kappa(\dd{\omega_2|\omega_1}) < \infty,
    \end{equation}
    holds. Specifically, maps of the form 
    \begin{equation}
    \label{eq:MeasurabilityCondExpectation}
        \omega_{0:i} \mapsto \E[f(X_{0:n}) | X_{0:i} = \omega_{0:i}],
    \end{equation}
    for conditional expectations as in 
    \cref{eq:ConditionalExpectationMultipleTwo}, are 
    $\mathcal{F}^{(0:i)}$-measurable, with $\mathcal{F}^{(0:i)}$ defined in 
    \cref{eq:ProductMeasurableSpaces}.
\end{lemma}
\begin{proof}
    A proof for the first case $f\geq 0$ is provided in 
    \cite[Lemma~14.23]{klenke2020_ptcc}. For general $f$ note that $f^+$ and
    $f^-$ are non-negative, measurable functions and as such it follows form
    the first case that the integrals in \eqref{eq:Case2Conditions} are well-
    defined. The claim follows then from $f=f^+ - f^-$ and the assumption that
    at least one of the integral \eqref{eq:Case2Conditions} is finite. 
    Finally, the measurability of the map 
    \eqref{eq:MeasurabilityCondExpectation} follows from its representation as 
    integral w.r.t. $\kappa^{(i,n)}(\cdot|\omega_{0:i})$ in 
    \cref{eq:ConditionalExpectationMultipleTwo}.
\end{proof}

\begin{comment}
\begin{lemma}
    \label{lem:KernelsAuxiliaryContinuousExpectation}
    If $\kappa$ and $f$ are continuous, then $\lambda$ is continuous
\end{lemma}
\begin{proof}
    For (i) see \cite[Lemma 14.23]{klenke2020_ptcc} or \cite[Proposition~7.29]{bertsekas1996_socdc}. For (ii) see \cite[Proposition 7.30]{bertsekas1996_socdc}.    
\end{proof}
\end{comment}


\begin{lemma}
    For $n\in\SetN$, let $(\Omega_i, \mathcal{F}_i), i=1,\dots,n$ be
    measurable spaces, and consider the stochastic kernels $\kappa_i$ for 
    $i=0,1\dots,n-1$ and the probability measures $\PM_i$ for $i=0,1,\dots,n$ 
    from \cref{cor:NKernelProduct}. Let $f:\Omega\to\overline{\SetR}$ be a 
    measurable function on the associated canonical probability space 
    \eqref{eq:CanonicalProbSpace}, with both or either of the conditions $\int_\Omega f^+ \dd{\PM}<+\infty$ or $\int_\Omega f^- \dd{\PM}<+\infty$ satisfied. Then,
    \begin{equation}
    \begin{split}
        \int_\Omega f \dd{\PM} 
        = \int_{\Omega_0} \int_{\Omega_1} \dots \int_{\Omega_n}
                f(\omega_0, \omega_1, \dots, \omega_n) \,
            &\kappa_{n-1}(\dd{\omega_n} |\, \omega_{0:n-1}) \\ 
            \dots \,
            &
            \kappa_1(\dd{\omega_2} |\, \omega_{0:1}) \,
            \kappa_0(\dd{\omega_1} |\, \omega_0) \,
            \mu(\dd{\omega_0})
    \end{split}
    \end{equation}
\end{lemma}
\begin{proof}
    The statement follows from inductively applying \cref{thm:KernelFubini}
    to $\PM = \mu\otimes\kappa^{(0,n-1)}$. Since $\mu$ is a probability 
    measure on $(\Omega_0,\mathcal{F}_0)$ and $\kappa^{(0,n-1)}$ is a 
    stochastic kernel from $\Omega_0$ to $\Omega^{(1,n)}$, it follows from 
    \cref{thm:KernelFubini} that
    \[
        \int_{\Omega} f \dd{\PM}
        = \int_{\Omega_0} \int_{\Omega^{(1,n)}}
            f(\omega_0, \omega_{1:n}) \,
        \kappa^{(0,n-1)} (\dd{\omega_{1:n}} | \, \omega_0 ) \,
        \mu(\dd{\omega_0}).
    \]
    Now assume that for $i = 1,\dots,n-2 $ it holds that
    \begin{multline}
    \label{eq:Hypothesis}
        \int_{\Omega} f \dd{\PM}
        = \int_{\Omega_0} \dots \int_{\Omega_i} \int_{\Omega^{(i+1,n)}}
            f(\omega_0, \dots, \omega_i, \omega_{i+1:n}) \,
        \kappa^{(i,n-1)} (\dd{\omega_{i+1:n}} | \, \omega_{0:i} ) \, \\
        \times
        \kappa_{i-1}(\dd{\omega_i} |\, \omega_{0:i-1}) \, 
        \dots \,
        \kappa_0(\dd{\omega_1} |\, \omega_0) \,
        \mu(\dd{\omega_0})
    \end{multline}
    By definition $\kappa^{(i,n-1)}=\kappa_i\otimes\kappa^{(i+1,n-1)}$, and 
    anew application of \cref{thm:KernelFubini} with the stochastic kernel 
    $\kappa^{(i+1,n-1)}$ and the probability measure $\kappa_i(\cdot | 
    \omega_{0:i})$ with fixed input $\omega_{0:i}\in\Omega^{0,i}$, yields for 
    the innermost integral in
    \eqref{eq:Hypothesis},
    \begin{multline}
    \label{eq:Step}
        \int_{\Omega^{(i+1,n)}}
            f(\omega_0, \dots, \omega_i, \omega_{i+1:n}) \,
        \kappa^{(i,n-1)} (\dd{\omega_{i+1:n}} | \, \omega_{0:i} ) \\
        = \int_{\Omega_{i+1}} \int_{\Omega^{(i+2,n)}}
            f(\omega_0, \dots, \omega_i, \omega_{i+1}, \omega_{i+2:n}) \,
        \kappa^{(i+1,n-1)} (\dd{\omega_{i+2:n}} | \, \omega_{0:i}, 
            \omega_{i+1} )
        \kappa_{i}(\dd{\omega_{i+1}} |\, \omega_{0:i}).
    \end{multline}
    Plugging \eqref{eq:Step} into \eqref{eq:Hypothesis} proves the induction 
    step. A similar statement has been proved in \cite[Proposition~7.28]{bertsekas1996_socdc}.
\end{proof}

%%%%%%%%%%%%%%%%%%%%%%%%%%%%%%%%%%%%%%%%%%%%%%%%%%%%%%%%%%%%%%%%%%%%%%%%%%%%%%%
\section{Measurable Selection}
\label{appendix:MeasurableSelection}
%------------------------------------------------------------------------------

A measurable selection theorem is presented, based on \cite[Lemma~16.1.15]{hinderer2016_do}. Let $(\Omega,\mathcal{F})$ be some measurable space and 
$\SpaceA \subset \mathbb{Z}$ with $\abs{\SpaceA}<\infty$. Given a function 
$g:\Omega \times \SpaceA \to\SetR$, we call the function $\pi:\Omega \to
\SpaceA$ a maximizer of $g$, if
\[
    \pi(\omega) \in \argmax_{a\in\SpaceA} g(\omega,a),
    \quad \text{for all } \omega\in\Omega.
\]

\begin{theorem}
\label{thm:SelectionTheorem}
    Let $g:\Omega\times\SpaceA\to\SetR$ be $\mathcal{F} \otimes
    \mathcal{B}_{\SpaceA}$-measurable. Then, for each $\omega\in\Omega$, there
    exists a smallest and a largest maximizer of $g$ and the map $\omega\mapsto
    g^*(\omega) := \max_{a\in\SpaceA} g(\omega,a)$ is $\mathcal{F}$-measurable.
\end{theorem}
\begin{proof}
    This follows literally by setting $(S,\mathfrak{S})=(\Omega,\mathcal{F})$ as
    well as $D(s)=A=\SpaceA$ for each $s\in S$ in 
    \cite[Lemma~16.1.15]{hinderer2016_do}.
\end{proof}




%%%%%%%%%%%%%%%%%%%%%%%%%%%%%%%%%%%%%%%%%%%%%%%%%%%%%%%%%%%%%%%%%%%%%%%%%%%%%%%
% List of Figures
%\chapter*{} % List of figures
\listoffigures
%\cleardoublepage
%%%%%%%%%%%%%%%%%%%%%%%%%%%%%%%%%%%%%%%%%%%%%%%%%%%%%%%%%%%%%%%%%%%%%%%%%%%%%%%

%%%%%%%%%%%%%%%%%%%%%%%%%%%%%%%%%%%%%%%%%%%%%%%%%%%%%%%%%%%%%%%%%%%%%%%%%%%%%%%
% List of tables
%\chapter*{} 
\listoftables
%%%%%%%%%%%%%%%%%%%%%%%%%%%%%%%%%%%%%%%%%%%%%%%%%%%%%%%%%%%%%%%%%%%%%%%%%%%%%%%

%%%%%%%%%%%%%%%%%%%%%%%%%%%%%%%%%%%%%%%%%%%%%%%%%%%%%%%%%%%%%%%%%%%%%%%%%%%%%%%
% bibliography %\newpage
%\chapter*{} 
\setstretch{1.001} % or 1.005 for tiny change
\printbibliography
%\newpage
%%%%%%%%%%%%%%%%%%%%%%%%%%%%%%%%%%%%%%%%%%%%%%%%%%%%%%%%%%%%%%%%%%%%%%%%%%%%%%%


%%%%%%%%%%%%%%%%%%%%%%%%%%%%%%%%%%%%%%%%%%%%%%%%%%%%%%%%%%%%%%%%%%%%%%%%%%%%%%%
%%%%%%%%%%%%%%%%%%%%%%%%%%%%%%%%%%%%%%%%%%%%%%%%%%%%%%%%%%%%%%%%%%%%%%%%%%%%%%%
\end{document}
%%%%%%%%%%%%%%%%%%%%%%%%%%%%%%%%%%%%%%%%%%%%%%%%%%%%%%%%%%%%%%%%%%%%%%%%%%%%%%%
%%%%%%%%%%%%%%%%%%%%%%%%%%%%%%%%%%%%%%%%%%%%%%%%%%%%%%%%%%%%%%%%%%%%%%%%%%%%%%%